\documentclass[11pt]{article}
\usepackage[a4paper,margin=25mm]{geometry}
\usepackage[T1]{fontenc}
\usepackage{lmodern}
\pdfextension mapfile{+lm.map}
\pdfextension mapfile{+cm.map}
\pdfextension mapfile{+symbols.map}
\pdfextension mapfile{+euler.map}
\pdfextension mapfile{+rsfs.map}
\usepackage{amsmath,amssymb,amsthm,mathtools,mathrsfs}
\usepackage{booktabs,longtable,array,enumitem,microtype,xurl}
\usepackage[numbers,sort&compress]{natbib}
\usepackage[unicode,colorlinks=true,linkcolor=blue,citecolor=blue,urlcolor=blue]{hyperref}
\hypersetup{pdftitle={ANABELIAN SUCCESSOR: point sections and reconstruction from etale topoi},
pdfauthor={GPT-6 Astra (OpenAI, ChatGPT Pro mode)},
pdfsubject={Incomplete research return, 16 November 2026}}
\makeatletter
\newcommand{\claimspace}{\par\if@nobreak\else\begingroup
  \@tempdima=5\baselineskip
  \vskip\z@\@plus\@tempdima\penalty-100
  \vskip\z@\@plus-\@tempdima
  \vskip\@tempdima\penalty10000\vskip-\@tempdima
\endgroup\fi}
\newenvironment{claim}[2]{\claimspace\medskip\phantomsection\def\@currentlabel{#1}\noindent\textbf{#1.}\quad\textsc{#2}\par\nobreak\smallskip\nobreak}{\par\medskip}
\makeatother
\newcommand{\F}{\mathbf F}
\newcommand{\Z}{\mathbf Z}
\newcommand{\Q}{\mathbf Q}
\newcommand{\C}{\mathbf C}
\newcommand{\Gal}{\operatorname{Gal}}
\newcommand{\Spec}{\operatorname{Spec}}
\newcommand{\Hom}{\operatorname{Hom}}
\newcommand{\Aut}{\operatorname{Aut}}
\newcommand{\im}{\operatorname{im}}
\newcommand{\Pic}{\operatorname{Pic}}
\newcommand{\et}{\mathrm{\acute et}}
\newcommand{\pp}{p'}
\newcommand{\rank}{\operatorname{rank}}
\newcommand{\coker}{\operatorname{coker}}
\newcommand{\Tr}{\operatorname{Tr}}
\setlength{\emergencystretch}{3em}
\setlength{\parskip}{0.2em}
\title{ANABELIAN\_SUCCESSOR\\Point sections in characteristic $p$\\and reconstruction from \'etale topoi}
\author{AI-produced by GPT-6 Astra\\(OpenAI, ChatGPT Pro mode)}
\date{16 November 2026}
\begin{document}
\maketitle
\begin{abstract}
\textbf{INCOMPLETE research return.}
We prove injectivity of the point-section map for every smooth proper
curve of genus at least two over every finite field, retaining the full
arithmetic quotient and the whole geometric prime-to-$p$ group.
The proof combines fixed-prime cuspidal configuration exactness,
diagonal boundary transport, and a point-recognition and Kummer
evaluation argument. It also proves the corresponding injectivity for
hyperbolic curves over fields finitely generated over $\mathbf F_p$.
We prove the rational relative-class-number criterion, its exact character
formula, its reduction to prime cyclic covers, and a controlled comparison
for nonconstant deck groups. We compute all base-field prime-character
tests for genus-two curves over fields of sizes $4,8,16$, give rational
witnesses for every original pointed genus-two pair over $\mathbf F_4$,
and compute every base-field odd-prime character test on pointed plane
quartics over $\mathbf F_3$ and $\mathbf F_5$. We determine all-extension least primes in
the explicitly certified cases. We audit the geometric
criteria, the exceptional-set formulation, and the norm module. An explicit
supersingular genus-two pencil refutes universal existence of nontrivial
odd cyclic covers with primitive two-rank zero; allowing the identity
cover leaves a separate positive-two-rank question. A cyclic-cover
Kodaira fibration and strong approximation also prove that the full
prime-to-$p$ configuration fiber map is not injective for any prime $p$,
ruling out that sufficient route without producing a collision. In the
topos problem we refute two broad formulations, prove a normal-extension
reconstruction theorem, prove reconstruction for subfields of
$\overline{\mathbf Q}_p$ of finite transcendence degree with bounded
$p$-primary ramification, prove reconstruction for finite extensions of the
fixed finite-class nilpotent extensions of the strict henselization of
$\mathbf Q$ at $p$ and for their subfields, and distinguish
precise model descent from a
bare pro-category assertion. The separate cover-existence, full-census,
abstract norm-module, and field-classification requests remain explicitly
open where they have not been settled. No claim of priority or human
authorship is made.
\end{abstract}
\tableofcontents
\section{Part A. Point sections over finite fields}
% Source: sections/A_foundations.tex
\subsection{The exact question and the specialization comparison}

We retain the entire arithmetic quotient throughout.  If $k$ is a finite
field of characteristic $p$, we put
\[
 \Pi_C=\pi_1(C_{\bar k})^{(p')},\qquad
 1\longrightarrow\Pi_C\longrightarrow E_C^{(p')}
       \longrightarrow G_k\longrightarrow1.
\]
The superscript means a quotient of the \emph{geometric} kernel, not the
maximal prime-to-$p$ quotient of the arithmetic group.  A chosen geometric
point above $P\in C(k)$ gives a representative $s_P$; changing its path
changes it by conjugation in $\Pi_C$.  We call distinct points a collision
only when their section classes agree in this whole extension.

\begin{claim}{A.1}{Proved; comparison used from R1}
For a smooth proper geometrically connected curve $C/k$ of genus at least
two and $P,Q\in C(k)$, the following are equivalent:
\begin{enumerate}
\item $s_P$ and $s_Q$ are conjugate by $\Pi_C$.
\item For every finite $L/k$ and every connected prime-to-$p$ finite
\'etale Galois cover $D\to C_L$ with constant deck group and a specified
$P'\in D(L)$ above $P$, there is a $Q'\in D(L)$ above $Q$ such that
$[Q'-P']\in J_D(L)[p^\infty]$.
\end{enumerate}
The point $Q'$ in (2) is unique.  One may equivalently allow covers whose
geometric Galois closure has prime-to-$p$ degree.  Thus the covering
condition of cote~149, scan~36, is an exact reformulation of a collision.
Its final implication $P=Q$ is proved below in A.3.
\end{claim}
\begin{proof}
The forward implication and uniqueness are checked in A.4.  For a cover
that is only geometrically Galois, use the same fixed-fibre argument:
the open arithmetic subgroup containing the chosen representative of
$s_P(G_L)$ contains the equal representative of $s_Q(G_L)$, which fixes a
point of the fibre above $Q$.  The induced topology on its geometric
group is the full prime-to-$p$ topology.  Indeed a core inside a
prime-to-$p$ Galois closure embeds in a product of prime-to-$p$ groups;
the finite quotient and its extension still have order prime to $p$.
The lifted sections therefore agree in the required full extension.
Prime-to-$p$ Kummer theory gives the asserted torsion class.

For the converse, use $P$ to write $E_C^{(p')}=\Pi_C\rtimes G_k$.
The group $\Pi_C$ is topologically finitely generated by the surface
presentation in SGA~1, XIII, Corollary~2.12 \cite{AR_SGA1}.
It consequently has cofinally many open characteristic subgroups $N$.
The open subgroup $N\rtimes G_k$ defines a geometrically connected
finite \'etale cover $D_N\to C$ with a distinguished rational point
$P_N$ and geometric deck group $\Pi_C/N$.  The continuous arithmetic
action on this finite group becomes trivial over some finite $L/k$.
Hypothesis (2) then gives $Q_N\in D_N(L)$ with $p$-primary difference.
Uniqueness holds over $\bar k$, by A.4's Hilbert~90 argument.  Every
element of $G_k$ fixes $P_N$ and preserves the torsion condition, so
$Q_N$ is in fact rational over $k$.

For $N'\subset N$ the finite morphism $D_{N'}\to D_N$ sends a point
difference to its image under the norm homomorphism of Jacobians.  The
image of $Q_{N'}$ therefore has the required torsion difference and is
$Q_N$ by uniqueness.  The compatible rational points form a $G_k$-fixed
point in the inverse limit of fibres above $Q$.  This limit is the
prime-to-$p$ \'etale path torsor from $P$ to $Q$.  Its fixed point is
precisely a geometric conjugator between the two sections.  This proves
the converse in the whole profinite group, rather than only at each
uncoordinated finite quotient.  Allowing the additional covers does not
alter the equivalence because the forward implication applies to them.
\end{proof}

\begin{claim}{A.2}{Proved; the finitely generated field corollary}
Universal absence of collisions for proper curves of genus at least two
over finite fields is equivalent to injectivity of the geometrically
prime-to-$p$ point-section map for every hyperbolic curve over every
field finitely generated over $\mathbf F_p$.  Both assertions hold,
by A.3 and the comparison proved here.  Independently of A.3, this
injectivity holds in genus zero and for affine curves of genus one.
For a proper curve over a finite field, $J_C(k)[p]=0$ already suffices.
\end{claim}
\begin{proof}
We verify the comparison in AP-A.8, since a field extension or
specialization cannot silently preserve every rational torsion test.
Over a finitely generated field $K/\mathbf F_p$, the kernel of the
prime-to-$p$ Kummer map of a semiabelian variety is its $p$-primary
torsion.  For its abelian quotient this follows from Lang--N\'eron and
finite generation over a finite constant field.  For a split torus it
follows by taking divisors on a normal proper model: $K^\times$ is a
finite constant group extended by a free abelian group.  Pass to a
finite splitting field for a general torus.  For the extension by an
abelian variety, lifts of finitely many generators leave a free toric
summand, so the same intersection-of-multiples assertion holds.  This
checks the Kummer step of AP-A.1 used here.

The generalized Jacobian of a punctured genus-zero curve is a torus.
Its geometric $p$-primary torsion point group is zero; after splitting
the punctures its Abel map contains a separating cross-ratio
coordinate.  Equality of point sections is therefore impossible for
distinct points.  For an affine genus-one curve $E\setminus B$, choose
the alleged colliding point $P$ as origin.  Its difference from $Q$
has order $p^a$.  Choose $m$ prime to $p$ with $m^2>p^a$ and use the
finite \'etale cover
\[
 E\setminus[m]^{-1}B\ \xrightarrow{[m]}\ E\setminus B.
\]
After a finite separable extension defining its deck group and boundary,
a collision lifts to $0,Q'$.  The generalized-Jacobian difference has
order $p^a$: its elliptic projection has that order and its toric kernel
has no $p$-primary geometric points.  There is a rational function with
divisor $p^a(Q'-0)$ and value $1$ at every boundary point.  Subtracting
$1$ gives at least $m^2$ distinct zeros to a function of degree $p^a$,
a contradiction.  These arguments include $p=2,3$.

For genus at least two, passage from an affine curve to its smooth
proper compactification induces a continuous quotient of arithmetic
extensions over $G_K$ and maps geometric conjugators to geometric
conjugators.  It remains to specialize a proper collision.  Spread the
smooth proper curve and its two distinct sections over a smooth
connected finite-type scheme $S/\mathbf F_q$ and shrink so that the
sections are disjoint.  At a closed point $s$, put
$R=\mathcal O_{S,s}^{h}$, $L=\operatorname{Frac}(R)$ and
$k_s=\kappa(s)$.  The map $G_L\twoheadrightarrow\pi_1(R)=G_{k_s}$ is
surjective because a connected finite \'etale cover of the normal
integral henselian local scheme remains connected generically.
SGA~1, XIII, Proposition~4.3, Examples~4.4 and 4.5.3 give the
specialization comparison for the smooth proper family with section;
its geometric prime-to-$p$ group is constant under this comparison.
Thus there is a continuous homomorphism
\[
 E_{C_L}^{(p')}\longrightarrow E_{C_s}^{(p')}
\]
which is an isomorphism on geometric kernels, induces the displayed
surjection on arithmetic quotients, and sends the section of each
specialized rational point to its section after that quotient map.
An equality of section representatives upstairs consequently gives an
equality downstairs.  The sections remain distinct as points.  The
residue field is finite, which proves the required reduction.  The
reverse universal implication includes finite fields among finitely
generated fields.  Finally, for a finite field the Kummer kernel is
$\bigcap_{(n,p)=1}nJ_C(k)=J_C(k)[p^\infty]$; the Abel embedding
settles the last assertion when this group is zero.
\end{proof}

\begin{claim}{A.3}{Proved; universal injectivity of proper point sections}
\label{clm:ARuniversalpointsections}
Every smooth proper geometrically connected curve of genus at least two
over every finite field has an injective geometrically prime-to-$p$
point-section map.  Equivalently, the last implication in A.1 always
holds, or the universal conclusion of A.2 holds.
The arithmetic quotient is the whole $G_k$, and conjugation is by the
whole geometric prime-to-$p$ group.  In particular the assertion includes
characteristics two and three and every finite constant extension.
\end{claim}
\begin{proof}
We assemble the three group-level comparisons proved in A.88--A.90.
Suppose $P,Q\in C(k)$ have geometrically conjugate sections, and choose
any prime $r\ne p$.  A.89 proves exactness for the configuration fiber
after retaining its entire proper geometric prime-to-$p$ quotient and
replacing only the cuspidal kernel by its maximal pro-$r$ quotient.
Its arithmetic comparison retains $G_k$ and identifies each punctured
curve group with the actual inverse image of its point section under
the second projection of the configuration group.

A.90 identifies the puncture decomposition subgroup with the
intersection with the global diagonal boundary group.  A geometric
conjugator between $s_P(G_k)$ and $s_Q(G_k)$ lifts to that boundary
group.  Conjugation by this lift, followed by a geometric inner
correction in the target fiber, gives a continuous isomorphism
between the maximal cuspidally pro-$r$ groups of $C-P$ and $C-Q$.
It is over the identity of the full proper arithmetic quotient and
carries the actual cusp decomposition group to a conjugate of the
actual target cusp group.  A.88 proves that such an isomorphism
forces $P=Q$.  This contradicts the existence of a collision.

For clarity, the implication in A.88 uses the entire family of finite
covers visible in this cuspidal quotient.  Their full same-prime
compact cohomology and Frobenius are preserved.  Their zeta functions
recognize noncuspidal point sections by a compatible inverse-limit
argument; the degree-one cuspidal central extension then turns this
recognition into a Kummer evaluation contradiction.  No finite-level
test is substituted for whole-group conjugacy.

The full prime-to-$p$ configuration fiber map is itself noninjective,
as proved in A.87.  The fixed-$r$ cuspidal exactness used here is a
different statement with its own finite-quotient proof.  The proof
does not use the disputed full configuration assertion in the
cuspidalization literature.  Together with A.1 and A.2, it supports
the injectivity contemplated in cote~149, scans~35--36, including
the question ``Il se pourrait qu'on puisse en d\'eduire $P=Q$''.
\end{proof}

% Source: agent_criteria/criteria_section.tex
% Candidate section owned by finite_field_criteria. Root renumbers AC.*.
\subsection{Rational torsion and the relative class number}

We use ``cover'' below to mean a connected, geometrically connected finite
\'etale cover.  A witness in criterion (R) is required to have degree greater
than one.  This requirement is missing from the wording of L1 and A3: with
$G=1$, the numerical condition is automatic, while the proposed contradiction
has no nonidentity deck transformation to use.

\begin{claim}{A.4}{Proved}
Let $L$ be a finite field of characteristic $p$, let $C/L$ be smooth proper
of genus $g\ge2$, and let $P,Q\in C(L)$ be distinct points with conjugate
sections in the geometrically prime-to-$p$ arithmetic fundamental group.
Let $\pi:D\to C$ be an \'etale Galois cover with constant deck group $G$,
$p\nmid |G|$, and let $P'\in D(L)$ lie over $P$.  There is a unique
$Q'\in D(L)$ above $Q$ with
\[
                 T=[Q'-P']\in J_D(L)[p^\infty].
\]
It is nonzero.  If $|G|$ is odd, every $1\ne\delta\in G$ satisfies
$\delta T\ne T$.  For any homomorphism of abelian varieties
$h:J_D\to B$ over $L$ with $B(L)[p]=0$, we have $h(T)=0$.
\end{claim}
\begin{proof}
We check precisely the part of AP-A.3 used here.  With the rational point
$P$ as base point, $\pi$ corresponds to an open subgroup of the arithmetic
extension containing the chosen representative of $s_P(G_L)$.  Conjugating
the representative of $s_Q$ by a geometric element so that it equals $s_P$
exhibits a fixed element of the finite \'etale fibre over $Q$.  This gives
$Q'\in D(L)$, and the lifted point sections are equal.  The prime-to-$p$
Kummer maps of $J_D(L)$ then kill $[Q'-P']$.  Since $J_D(L)$ is finite,
\[
            \bigcap_{(n,p)=1} nJ_D(L)=J_D(L)[p^\infty].
\]
This proves existence without enlarging $L$.

For uniqueness, we check the proper-curve instance of AP-A.2.  If two such
lifts differ by a deck transformation $a$ of order $m>1$, then
$n(aR-R)$ is principal for some power $n$ of $p$.  Normalize a function
$f$ with this divisor so that its norm for the cyclic action of $a$ is
one.  Hilbert 90 writes $f=a(u)/u$.  Hence
$\operatorname{div}(u)-nR$ is an invariant divisor of degree $-n$.
Every orbit has length $m$, since the action is free on the proper curve.
Consequently $m\mid n$, a contradiction.  The Abel map of a curve of
positive genus is injective on geometric points: a principal divisor
$Q'-P'$ for distinct points would give a degree-one map to $\mathbf P^1$.
Thus $T\ne0$.

We also check AP-A.C2.7 in the full stated characteristic range.  If
$\delta T=T$, then the distinct effective divisors
$\delta Q'+P'$ and $Q'+\delta P'$ are linearly equivalent.  They have no
common term, because their terms above $P$ and $Q$ are disjoint and a
nonidentity deck transformation has no fixed point.  They give a degree-two
map to $\mathbf P^1$, so $D$ is hyperelliptic.  Its unique hyperelliptic
involution is central in its automorphism group.  Choose an odd prime-order
element $a\in G$.  Its induced action on the genus-zero quotient has the
same order and has a geometric fixed point.  The fibre above that point
contains at most two geometric points.  An odd prime-order permutation of
this fibre is trivial, so $a$ fixes a point of $D$, contradicting \'etaleness.
This argument works also in characteristics two and three.  Finally $h(T)$
is an $L$-rational point of $p$-power order, and any nonzero such point has a
nonzero multiple of order $p$.  The hypothesis on $B(L)$ therefore kills it.
\end{proof}

\begin{claim}{A.5}{Proved}
In the situation of A.4, assume $|G|>1$ is odd.  Put
$h_L(X)=\#J_X(L)$.  If
\[
                       p\nmid h_L(D)/h_L(C),
\]
then $P$ collides with no distinct point of $C(L)$, and hence with no
point rational over any specified subfield of $L$.  The quotient in this
display is a positive integer.  More generally the same conclusion holds
for any prime-to-$p$ group $G$ containing a nontrivial odd-order element,
provided that element acts freely on $D$ as it does for an \'etale cover.
\end{claim}
\begin{proof}
Let $B=(\ker\pi_*)^0$ be the connected kernel of the norm
$\pi_*:J_D\to J_C$.  Both this morphism and the pullback
$\pi^*:J_C\to J_D$ are homomorphisms over $L$ and satisfy
\[
 \pi_*\pi^*=[d],\qquad \pi^*\pi_*=\sum_{\gamma\in G}\gamma,
 \qquad d=|G|.
\]
Thus $J_C\times B\to J_D$, $(x,b)\mapsto\pi^*x+b$, is an isogeny.
Point counts of abelian varieties over a finite field are isogeny
invariants, so $h_L(D)/h_L(C)=\#B(L)$ is an integer.  On the finite
$p$-primary group $M=J_D(L)[p^\infty]$, division by $d$ is unambiguous.
The endomorphisms
\[
             e=d^{-1}\sum_{\gamma\in G}\gamma,
             \qquad 1-e
\]
are complementary idempotents.  Their images are respectively the invariant
part and the part mapping into $B[p^\infty]$.  For the second assertion,
$\pi_* (1-e)=0$; any possible component of $\ker\pi_*$ has prime-to-$p$
order, since the isogeny above has kernel killed by a power of $d$.
Alternatively $(1-e)M$ is the image on $p$-primary points of
$[d]-\pi^*\pi_*$, whose connected image is contained in $B$.
The numerical hypothesis gives $B(L)[p^\infty]=0$, so $M=eM$.
A hypothetical collision's $T$ would be fixed by every deck transformation,
contradicting A.4 for any nonidentity odd-order element.

The rigid target here is the new abelian factor $B$ with
$B(L)[p]=0$.  The proof uses rational torsion over this exact $L$;
it makes no assertion that $B$ has geometric $p$-rank zero.
\end{proof}

\begin{claim}{A.6}{Proved}
For a fixed $C/L$ and $P\in C(L)$, the following are equivalent:
\begin{enumerate}
\item There is a nontrivial constant odd-order prime-to-$p$ Galois
$P$-split \'etale cover satisfying criterion (R).
\item There is such a cover with cyclic group of odd prime order.
\end{enumerate}
The extension field $L$ need not be changed.  Consequently nonabelian
odd-order groups and noncyclic abelian groups cannot enlarge the set of
$(C,P,L)$ reached by the whole-relative-class-number criterion L1.
\end{claim}
\begin{proof}
The odd-order theorem of Feit--Thompson, the unnumbered main theorem in
Chapter I, \S1, printed p.~775, states: ``All finite groups of odd order are
solvable'' \cite{AC_FeitThompson1963}.  We use only its stated hypothesis
that $G$ is finite of odd order.  A nontrivial finite solvable group has a
nontrivial abelianization, hence a quotient $G\twoheadrightarrow
\mathbf Z/\ell$ for an odd prime $\ell\ne p$.  If $N$ is its kernel, let
$E=D/N$.  The intermediate cover $E\to C$ and its deck group are defined
over $L$, and the image of $P'$ is rational, so $P$ splits completely.
For any finite separable map of smooth proper curves over $L$, pullback and
norm give an isogeny decomposition of the source Jacobian with the target
Jacobian as a factor.  Therefore both ratios in
\[
        \frac{h_L(D)}{h_L(C)}=
        \frac{h_L(D)}{h_L(E)}\frac{h_L(E)}{h_L(C)}
\]
are positive integers.  If their product is prime to $p$, so is the second
one.  This proves (1)$\Rightarrow$(2); the reverse implication is immediate.
The argument says nothing about a criterion applied only to a selected
primitive or nonabelian representation factor.
\end{proof}

\subsection{The exact character test and the base-point twist}

\begin{claim}{A.7}{Proved}
Let $C/L$ be as above, $q=\#L$, $P\in C(L)$, and let
$a:J_C(L)\twoheadrightarrow A$ be a nontrivial finite abelian quotient
of order prime to $p$.  For a nontrivial character
$\psi:A\to\overline{\mathbf Q}^{\times}$ define
\[
 c_d(\psi)=\sum_{E\in\operatorname{Div}^{+}_d(C)(L)}
           \psi\bigl(a([E-dP])\bigr),
 \qquad \kappa=[K_C-(2g-2)P].
\]
Here an effective divisor is a Galois-invariant effective geometric
divisor; it need not be a sum of rational points.  The corresponding
$P$-split Lang cover has twisted polynomial
\[
 L_P(\psi,T)=\sum_{d=0}^{2g-2}c_d(\psi)T^d,
 \qquad c_{2g-2}(\psi)=q^{g-1}\psi(a(\kappa)).
\]
For $g-1\le d\le2g-2$ one has
\[
 c_d(\psi)=q^{d+1-g}\psi(a(\kappa))
                   c_{2g-2-d}(\psi^{-1}).                 \tag{A.4}
\]
In particular, modulo any prime above $p$,
\[
               L_P(\psi,1)\equiv
                  \sum_{d=0}^{g-1}c_d(\psi).             \tag{A.5}
\]
For genus two this is exactly
$L_P(\psi,1)\equiv1+\sum_{R\in C(L)}\psi(a([R-P]))$.
\end{claim}
\begin{proof}
The Euler product for the unramified rank-one character, with Frobenius
at a closed point $x$ sent to $a([x-(\deg x)P])$, expands as the displayed
effective-divisor series.  Its coefficient in each degree is $c_d$.
All degree-$d$ line bundles defined over $L$ are parametrized by
$J_C(L)$ after subtracting $dP$.  A line bundle $M$ with
$h^0(M)=r$ has exactly $(q^r-1)/(q-1)$ effective divisors over $L$.
This follows by dividing its nonzero rational sections by $L^\times$;
$\operatorname{Br}(L)=0$ ensures that rational Picard classes are represented
by line bundles.  For $d>2g-2$, Riemann--Roch makes $r=d+1-g$
independent of the class, and summing the nontrivial character over
$J_C(L)$ gives zero.  The series is therefore a polynomial of degree at
most $2g-2$.

For $d\ge g-1$, Riemann--Roch and the change of variables
$M\mapsto\omega_C\otimes M^{-1}$ give
\[
 \frac{q^{h^0(M)}-1}{q-1}
 =q^{d+1-g}\frac{q^{h^0(\omega_C\otimes M^{-1})}-1}{q-1}
   +\frac{q^{d+1-g}-1}{q-1}.
\]
The second term has character sum zero.  The first gives (A.4), including
its degree-zero coefficient and hence the leading coefficient.  Every
coefficient of degree $d\ge g$ is divisible by $q$, proving (A.5).
Thus no cohomological interpretation is needed for the genus-two numerical
formula; it is an exact Riemann--Roch calculation.
\end{proof}

\begin{claim}{A.8}{Proved}
With $A$ as in A.7 put
\[
 \Theta_{C,P,A}=\sum_{d=0}^{g-1}
   \sum_{E\in\operatorname{Div}^{+}_d(C)(L)}[a([E-dP])]
                      \quad\hbox{in }\mathbf F_p[A].
\]
The full $A$-cover satisfies criterion (R) if and only if
\[
     \bar\psi(\Theta_{C,P,A})\ne0
       \quad\text{for every nontrivial }
          \bar\psi:A\to\overline{\mathbf F}_p^{\times}.    \tag{A.6}
\]
One can equivalently sum through $2g-2$ in the definition of $\Theta$.
For $A=\mathbf Z/\ell$, with $\ell$ an odd prime different from $p$,
\[
 \frac{h_L(D)}{h_L(C)}
       =\prod_{i=1}^{\ell-1}L_P(\psi^i,1)
       =N_{\mathbf Q(\zeta_\ell)/\mathbf Q}L_P(\psi,1).
\]
If $p$ generates $(\mathbf Z/\ell)^\times$, failure of (R) is equivalent
to all the coefficients of $\Theta$ in $\mathbf F_p[\mathbf Z/\ell]$
being equal.
\end{claim}
\begin{proof}
On a closed fibre of an abelian Galois cover, multiplication by its
Frobenius element has eigenvalues $\psi(a_x)$ in the regular
representation.  Factoring its permutation determinant into these
characters and multiplying over closed points gives
\[
                 Z(D,T)=\prod_{\psi\in\widehat A}L_P(\psi,T).
\]
The trivial factor is $Z(C,T)$.  The common genus-zero denominator of the
two zeta functions cancels, and evaluation of their numerator polynomials
at $1$ gives the relative class-number formula.  By A.7, reduction of each
factor is its character evaluation on $\Theta$.  Because $p\nmid|A|$,
all characters reduce to distinct characters over
$\overline{\mathbf F}_p$ and exhaust those characters.  The product is
nonzero modulo $p$ exactly when all nontrivial factors are nonzero, which
proves (A.6).

For prime cyclic $A$, the factors are all Galois conjugates of a fixed
primitive character, giving the norm.  If $p$ is primitive modulo $\ell$,
$1+X+\cdots+X^{\ell-1}$ is irreducible over $\mathbf F_p$.
A polynomial of degree at most $\ell-1$ vanishes at one nontrivial
$\ell$th root exactly when it is a scalar multiple of this polynomial.
The same irreducibility shows that vanishing at one character is equivalent
to vanishing at all nontrivial characters.  This proves the last assertion.
\end{proof}

\begin{claim}{A.9}{Proved}
For the same fixed character $\psi$ and rational base points $P,P_0$,
\[
 L_{P_0}(\psi,T)
       =L_P\bigl(\psi,\psi(a([P-P_0]))T\bigr).             \tag{A.7}
\]
The associated covers differ by the constant torsor with arithmetic
Frobenius value $a([P-P_0])$ (up to reversing all Frobenius conventions
consistently).  Thus changing the base point is not innocuous for (R).
For a finite extension $L'/L$, the base change of a $P$-split Lang cover
corresponds to $a\circ N_{L'/L}:J_C(L')\to A$ and remains $P$-split.
If a fixed cover fails (R) over $L$, it fails over every extension $L'$.
If it satisfies (R), it need not do so over every extension.
\end{claim}
\begin{proof}
For every effective divisor of degree $d$,
$[E-dP_0]=[E-dP]+d[P-P_0]$.  Substitution in each coefficient proves
(A.7).  At closed points this is precisely multiplication of Frobenius
characters by the indicated constant degree character.  Class field
functoriality under constant extension follows directly from the norm of
divisors and the product of Frobenius elements over a closed point; it gives
the asserted norm character.  The cover stays geometrically connected and
has its rational lift, so the resulting quotient is still onto.

For the last assertions use $h(D)/h(C)=\#B(L)$ from A.5.  A nonzero
$L$-rational $p$-point of $B$ remains rational after every extension.
Conversely a variety with positive geometric $p$-rank acquires nonzero
$p$-torsion over some finite extension even when it has none over $L$.
More precisely, if $F$ is arithmetic Frobenius on $B[p](\overline L)$,
then the criterion survives a degree-$m$ extension exactly when
$\ker(F^m-1)=0$.  This identifies the field dependence rather than
replacing rational torsion by geometric torsion.
\end{proof}

\subsection{The genus-two count and its limits}

\begin{claim}{A.10}{Proved}
Suppose $g=2$ and $\ell\ne p$ is an odd prime for which $p$ is a primitive
root.  Put $V=J_C(L)/\ell J_C(L)$ and
$N_b=\#\{R\in C(L):[R-P]\bmod\ell J_C(L)=b\}$.
If $P$ collides with a distinct point, then
\[
                      N_b\equiv N_0+1\pmod p
                       \quad(0\ne b\in V).               \tag{A.8}
\]
If $|V|>\#C(L)$, this implies
$N_0\equiv-1\pmod p$ and $N_b\equiv0\pmod p$ for $b\ne0$.
For one surjection $a:J_C(L)\twoheadrightarrow\mathbf Z/\ell$,
let $H=\ker a$, $N=\#C(L)$.  If
\[
                           N\ge2|H|,
                           \qquad \ell>N,                \tag{A.9}
\]
then its cover satisfies (R).
\end{claim}
\begin{proof}
Every nontrivial character of $V$ factors through a quotient of order
$\ell$.  Under a collision, A.5 forces (R) to fail for each such quotient;
A.8 and the primitive-root hypothesis force every corresponding
nontrivial character evaluation of
$[0]+\sum_R[[R-P]]$ to vanish.  Fourier inversion is valid because
$p\nmid|V|$.  Therefore this group-ring element is a scalar multiple of
$\sum_{b\in V}[b]$, giving (A.8).  Since $P$ occupies the zero coset,
$|V|>N$ leaves an empty nonzero coset, proving the second assertion.

For the last assertion we verify the geometric input from AP-A.C2.21.
For each nonzero class $h$ in the Jacobian of a genus-two curve, at most
two ordered pairs of distinct geometric points have difference $h$.
Two distinct representations $R-S=R'-S'$ give a degree-two linear
system, hence the unique hyperelliptic pencil, and force
$(R',S')=(\iota S,\iota R)$.  This includes ramification points and
all characteristics.  Thus at most $2(|H|-1)$ rational points share an
$a$-coset with another rational point: each such point occurs as the first
member of an ordered pair with nonzero difference in $H$.
Under (A.9), at least two points occupy singleton cosets; at most one is
in the zero coset.  There is therefore a nonzero singleton coset.  If (R)
failed, A.8 applied to this one cyclic quotient, together with an empty
coset (since $\ell>N$), would make each nonzero occupancy divisible by
$p$.  A singleton contradicts this.  Notice that the primitive-root
hypothesis is essential to this argument and is retained explicitly.
\end{proof}

\begin{claim}{A.11}{Proved}
For any genus $g\ge2$, any finite quotient $A$ of $J_C(L)$, and
$N=\#C(L)$, let $N_b$ count the Abel-image points in the $A$-cosets.
The number of ordered distinct pairs in the same coset satisfies
\[
 \sum_b N_b(N_b-1)
 \le\frac{N^2}{|A|}
       +\frac{|A|-1}{|A|}(2g-2)^2q-N.                     \tag{A.10}
\]
If the nonzero occupancies are multiples of $p$, the same quantity is
at least $(p-1)(N-N_0)$.  These two estimates alone give a contradiction
only if their displayed numerical bounds conflict; $|A|>N$ by itself
and the inequality $p>(2g-2)^2$ are not the full hypotheses of such a
finite-field contradiction.
\end{claim}
\begin{proof}
Katz, \emph{Twisted L-Functions and Monodromy}, Proposition 9.1.3,
printed p.~185, applies to a smooth proper geometrically connected curve
$C/L$, a point $P\in C(L)$ and any nontrivial character $\rho$ of
$J_C(L)$.  In its $d=1$ case it bounds the absolute value of
$\sum_R\rho([R-P])$ by $(2g-2)\sqrt q$ \cite{AC_Katz2002}.
Its proof checks that the Abel map preserves the geometric abelianized
fundamental group, so the pulled-back nontrivial Lang local system is
geometrically nontrivial; this verifies the needed hypothesis for all
characters factoring through $A$.
Character orthogonality gives exactly
\[
   \sum_bN_b^2=\frac1{|A|}\left(N^2+
             \sum_{\rho\ne1}|\sum_R\rho([R-P])|^2\right).
\]
Apply the bound and subtract $N$ to obtain (A.10).  For each $b\ne0$
with $N_b>0$ a multiple of $p$, one has
$N_b(N_b-1)\ge(p-1)N_b$.  Sum these inequalities.

In particular the $-N$ in (A.10) cannot simply be changed to $-q$.
The weaker bound stated in L3 with coefficient $(2g-2)^2-1$ follows
from this argument only if
$N+(2g-2)^2q/|A|\ge q$.  A large-field asymptotic argument also has to
control $N^2/|A|$ and $N_0$ and provide the required quotients over those
fields.  More fundamentally, for $g>2$ the premise on the occupancies
has not followed from a collision: the correct character test A.5 has
additional effective-divisor terms.  Thus no pair $(p,g)$ is settled
universally by the displayed L3 counting argument without an additional
existence theorem.  A.10 is the rigorous genus-two sufficient statement.
\end{proof}

\begin{claim}{A.12}{Refuted: the point-only extension to genus three}
The genus-two expression $1+\sum_{R\in C(L)}\psi([R-P])$ does not test
relative $p$-divisibility in genus three, even for a prime cyclic cover
with $p$ a primitive root.  An explicit example is
\[
        C/\mathbf F_2:\quad y^2+y=x^7+x^3,\qquad P=\infty.
\]
There is a surjection $\chi:J_C(\mathbf F_2)\to\mathbf Z/3$ whose
point-coset counts are $(3,1,1)$ and whose twisted polynomial is
\[
                 L_P(\chi,T)=1+2T+3T^2+4T^3+4T^4.
\]
Its relative class number is $196$, while the point-only expression
reduces to $1$ modulo $2$.  This is a counterexample to exporting that
numerical test; it is not a collision or a counterexample to A1.
\end{claim}
\begin{proof}
The affine curve is smooth because the derivative in $y$ is one.
The pole of the Artin--Schreier function at infinity has order seven,
so the smooth proper normalization has one rational point there and
genus three.  Rational reduced Mumford divisors $(u,v)$ are given by
$u$ monic, $\deg u\le3$, $\deg v<\deg u$ and
$u\mid v^2+v+x^7+x^3$.  Their complete enumeration has $27$ elements.
The class $a=[(0,0)-\infty]$ generates the group of order $27$;
$[(1,0)-\infty]=6a$ and the conjugates have classes $-a$ and $-6a$.
Define $\chi(a)=1\bmod3$.  This specifies the quotient completely.

The five rational points have character exponents $0,1,2,0,0$.
There are no degree-two closed points.  The fifteen effective degree-two
divisors obtained from unordered rational-point pairs have coset counts
$(7,4,4)$.  Consequently, with $\zeta^2+\zeta+1=0$,
\[
 c_1=3+\zeta+\zeta^2=2,
       \qquad c_2=7+4\zeta+4\zeta^2=3.
\]
Since $K_C\sim4\infty$, formula (A.4) gives $c_3=4$ and $c_4=4$.
Thus $L_P(\chi,1)=14$ and its quadratic norm is $196$.
The point-only expression is $1+c_1=3$, but the correct reduction is
$1+c_1+c_2=6\equiv0$.

The attached program \path{genus3_countercheck.py} enumerates all reduced
divisors, computes the full quotient table, checks closure, inverses,
commutativity and all $27^3=19683$ associativity identities for its exact
Cantor arithmetic, and prints all coefficients.  Its output is
\path{genus3_countercheck_output.json}.  The finite computation is used
for the explicit example, while A.7 proves the general replacement.
\end{proof}

\begin{claim}{A.13}{Proved: a finite-prime obstruction for the rational method}
Let $\overline C/\overline{\mathbf F}_p$ have genus at least two and let
$S$ be a finite set of odd primes different from $p$.  Suppose that for
every connected cyclic \'etale cover of degree $\ell\in S$, its new
Jacobian factor has positive geometric $p$-rank.  There is a finite-field
model $C/k$ with $P\in C(k)$ such that no cyclic $P$-split cover of degree
in $S$ satisfies (R), over any finite extension of $k$.
\end{claim}
\begin{proof}
There are finitely many geometric cyclic covers of each fixed degree:
they are determined by the finitely many nonzero homomorphisms from the
finite group $\pi_1(\overline C)^{\mathrm{ab}}/\ell$ to $\mathbf Z/\ell$,
up to changing the deck generator.  Choose a geometric point $P$ and a
lift to each cover.  Descend all these curves, lifts, maps and deck actions
to one finite field $k$.  Enlarge $k$ finitely so that the new factor of
each cover has a nonzero rational point of order $p$; such a point exists
by the rank hypothesis, and there are only finitely many factors.

After any extension $L/k$, a constant cyclic torsor with one of the
specified geometric covers differs from its chosen model by a constant
$\mathbf Z/\ell$-torsor.  Requiring a rational point above $P$ makes this
twist trivial: evaluate the difference character on the section of $P$.
It is therefore the base change of that chosen $P$-split model (up to
changing the deck generator).  The already rational $p$-point persists,
so the relative class number is divisible by $p$.
\end{proof}

\subsection{The primitive rational factor}

\begin{claim}{A.14}{Proved}
Let $C/L$ be smooth proper of genus $g\ge2$ and not geometrically a
separable double cover of a genus-one curve.  Let $D\to C$ be a
connected \'etale cyclic cover of odd degree $n>1$ prime to $p$, with
constant group $\langle\sigma\rangle$ and $P'\in D(L)$ over
$P\in C(L)$.  Define the homomorphism over $L$
\[
       A=\prod_{\ell\mid n}(1-\sigma^{n/\ell}):J_D\to J_D,
       \qquad P_n=\operatorname{im}(A),
\]
where the product is over distinct prime factors.  If
$p\nmid\#P_n(L)$, then $P$ has no collision.  Equivalently, a collision
forces $p\mid\#P_n(L)$ on every such cover.  This proves L4 with all
field and nontriviality hypotheses explicit.
\end{claim}
\begin{proof}
A.4 gives a nonzero rational $p$-primary $T=[Q'-P']$ under a collision.
The numerical hypothesis gives $A T=0$.  We check the precise geometric
endgame of AP-A.C2.20 rather than imposing zero geometric $p$-rank.
Put $r=\omega(n)$ and $d=2^r$.  Expanding $A$ gives a principal divisor
with $d$ distinct simple zeros and $d$ distinct simple poles, all supported
on the two free disjoint deck orbits of $P'$ and $Q'$.  The corresponding
function $f$ has separable degree $d$.  The subset products of the elements
$\sigma^{n/\ell}$ are distinct because their orders are distinct primes.
One has $n>(d-1)^2$: for $r=1,2$ this follows from $n\ge3,15$; for
$r\ge3$ use $n\ge3\cdot5^{r-1}>4^r$.

Over $\overline L$ let $M=\overline L(\gamma f:\gamma\in G)$ and
$e=[\overline L(D):M]$.  Poonen's Proposition 2.1, Lemma 2.2 and
Corollary 2.3 in \emph{Gonality of modular curves in characteristic $p$},
author p.~3, apply over the perfect field $\overline L$ and give
\cite{AC_Poonen2007}
\[
                     e\mid d,\qquad g(M)\le(d/e-1)^2.
\]
For completeness, the induction behind the assertion is this.  If an
intermediate field has index $d/s$ and genus at most $(s-1)^2$, adjoining
one more degree-$d$ function with relative degree $a$ gives, by
Castelnuovo--Severi, genus at most
\[
        a(s-1)^2+(a-1)(as-1)
          =(as-1)^2-as(a-1)(s-1)\le(as-1)^2.
\]
This also handles a compositum with more than two generators.
The alternative $e=1$ contradicts
$g(D)=n(g-1)+1\ge n+1>(d-1)^2$.  Thus $e>1$ is a power of two.
Let $Z$ be the smooth curve with field $M$.  Its induced $G$-action is
faithful because the kernel order divides both $e$ and $n$.  It is free:
a prime-order stabilizer of odd order would partition the fibre divisor
of $D\to Z$ into orbits of that prime length, forcing that prime to divide
$e$.  Riemann--Hurwitz gives
$g(Z)-1=n(g(Z/G)-1)$, and the genus bound forces $g(Z)=1$.

Write $f=h\circ v$, $v:D\to Z$ of degree $e$.  Over a zero of $h$, the
fibre of $v$ is reduced, since the zero divisor of $f$ is reduced.  It
lies in the two original deck orbits, each of which maps injectively to
$Z$ by the free action there.  Hence $e\le2$, and so $e=2$.
The separable quadratic involution of $D\to Z$ commutes with $G$, because
$G$ preserves its fixed field and a group of order two has no nontrivial
automorphism.  Taking the quotient by $G$ gives a separable degree-two
map $C\to Z/G$, whose target has genus one.  This contradicts the
hypothesis on $C$.

Every vanishing used before this last geometric contradiction holds over
the original finite field $L$.  The target is the rationally
$p$-torsion-free primitive factor $P_n$, and the remaining argument is
geometric.  Thus the ``controlled-field endgame'' set aside in R8 is not
an additional gap in this conditional rational criterion; the outstanding
issue is constructing a suitable pointed cover over a field on which the
required rational torsion vanishes.
\end{proof}

\begin{claim}{A.15}{Open}
The nontrivial-cover existence assertion A3 has not been proved or refuted
by A.4--A.14.  In its precise reduced form, it asks whether for every
$C/k$ of genus at least two over a finite field and every $P\in C(k)$
there are a finite extension $L/k$, an odd prime $\ell\ne p$ and a
surjection $J_C(L)\to\mathbf Z/\ell$ such that all its nontrivial
twisted values $L_P(\chi^i,1)$ are $p$-adic units.  This is equivalent to
the version allowing every nontrivial odd-order Galois group by A.6.
\end{claim}
\begin{proof}[Record of the unresolved step]
A.7--A.8 provide an exact test and A.10 supplies a genus-two sufficient
numerical condition.  They do not construct primes with the required
quotient size and primitive-root property, or a substitute for those
conditions.  In higher genus, replacing the effective-divisor test by
point counts is invalid by A.12.  The random-character estimate in L2 is
not a theorem: character values for a fixed curve are constrained,
and no independence or equidistribution statement with those parameters
has been established.  A.6 rules out the suggested enlargement to
nonabelian odd-order groups for this particular criterion.  A.14 gives
a weaker primitive-factor target but does not construct it.  These are
failures of the attempted methods, not counterexamples to the existence
assertion or to A1.
\end{proof}

\begin{claim}{A.16}{Proved}
Fix any prime $p$, any genus $g\ge2$, and any finite set $S$ of odd
primes different from $p$.  There are a finite field $k$ of characteristic
$p$, a smooth proper geometrically connected ordinary genus-$g$ curve
$C/k$, and $P\in C(k)$, such that no $P$-split cyclic cover of degree
in $S$ satisfies (R), over any finite extension $L/k$.  By A.6 the same
holds for all odd-order Galois covers whose group has prime divisors only
in $S$.  In particular there is no bound depending only on $(p,g)$ for
the least prime of a rational L1 witness, even if every $(C,P)$ eventually
has some witness.
\end{claim}
\begin{proof}
Ozman--Pries, Theorem 1.1, printed p.~456, assumes $\ell\ne p$ prime,
$g\ge2$ and $0\le f\le g$, with $f\ne0$ when $g=2$.
For every irreducible component $S_f$ of the $p$-rank-$f$ stratum,
its inverse image in the finite cover-moduli space is irreducible and
has generically ordinary Prym \cite{AC_OzmanPries2019}.  Take $f=g$;
all hypotheses then hold.  Ordinariness is an open condition, and the
forgetful map for covers is finite, as stated in their \S2.2.
The image of the closed nonordinary-Prym locus is therefore closed and
omits the generic point.  Thus on a dense open of the ordinary
curve-moduli component, every degree-$\ell$ cyclic cover has ordinary
Prym, of positive dimension $(\ell-1)(g-1)$.

Intersect these dense opens for the finite set $S$.  The intersection
is nonempty.  An \'etale scheme atlas has a closed point over a finite
extension of $\mathbf F_p$ in that open; after a further finite extension
it supplies an actual curve and a rational marked point.  This gives the
geometric hypothesis of A.13, which then gives the desired finite-field
model and excludes those primes over every extension.  An odd-order
Galois witness supported in $S$ would have a prime-degree witness with
prime in $S$ by A.6, and is therefore excluded too.
\end{proof}

\subsection{A semilinear obstruction and the least prime for the two hard classes}

\begin{claim}{A.17}{Proved}
Let $k$ be a finite field of characteristic two, $C/k$ a genus-two curve,
and $P\in C(k)$.  Let $\ell$ be an odd prime with
$\operatorname{ord}_{\ell}(2)=\ell-1$.  Suppose that every geometric
cyclic degree-$\ell$ \'etale cover has ordinary new Jacobian factor.
Suppose further that, for some $d>0$, the action of arithmetic Frobenius
$F^d$ on $J_C[\ell](\overline k)$ is scalar $2^j$, with
$\gcd(j,\ell-1)=1$, and that every eigenvalue of $F$ on this space has
exact order $d(\ell-1)$.  Then no constant-deck cyclic degree-$\ell$
$P$-split cover satisfies (R) over any finite extension of $k$.
\end{claim}
\begin{proof}
Fix a geometric cyclic cover, equivalently a quotient of the geometric
abelianized fundamental group modulo $\ell$.  Its kernel is stable under
$F^d$, because that action is scalar.  Using the section of the actual
point $P$, the open subgroup generated by this kernel and $s_P(G_{k_d})$
constructs a pointed model $D_0\to C_{k_d}$ over the degree-$d$ extension
$k_d/k$.  Its deck group need not yet be constant; Frobenius acts on it
by the scalar $2^j$.  This construction is legitimate before the deck
group is rational.  The $P$-point fixes the twist.

Let $V$ be the new summand of $J_{D_0}[2](\overline k)$, cut out by the
odd-order averaging idempotent.  Ordinariness gives
$\dim_{\mathbf F_2}V=\ell-1$.  Each nontrivial geometric deck character
occurs once, by the decomposition of
$H^1(D_0,\mathcal O_{D_0})$ into degree-zero character lines and their
one-dimensional first cohomology spaces; the Artin--Schreier exact
sequence transfers this to the ordinary stable part.  Since two is
primitive modulo $\ell$, $V$ is the single nontrivial irreducible
$\mathbf F_2[\mathbf Z/\ell]$-module.  Identify it with
$E=\mathbf F_{2^{\ell-1}}$, on which a deck generator acts by multiplication
by a primitive $\ell$th root $\zeta$.

Let $\Phi$ denote arithmetic Frobenius over $k_d$ on $V$.  The relation
$\Phi\zeta\Phi^{-1}=\zeta^{2^j}$ makes it a semilinear map
\[
                            x\longmapsto a x^{2^j}
                  \qquad(a\in E^\times),
\]
possibly with $j$ replaced by $-j$ under the opposite action convention.
Either exponent generates $\operatorname{Gal}(E/\mathbf F_2)$.
Therefore
\[
                 \Phi^{\ell-1}(x)
                   =N_{E/\mathbf F_2}(a)x=x.
\]
All new two-torsion is rational over $k_{d(\ell-1)}$.

A constant cyclic quotient over $k_m$ requires a nonzero quotient of
$J_C[\ell]$ on which $F^m$ is identity.  The exact eigenvalue-order
hypothesis forces $d(\ell-1)\mid m$.  Its $P$-pointed form is the base
change of the constructed pointed model: two abelian torsor forms with
the same geometric kernel differ by a constant character, and evaluation
on $s_P$ makes that difference zero.  It consequently has nonzero rational
new two-torsion over $k_m$.  A.5 identifies this with even relative class
number, proving the assertion for all finite extensions, not just the
first field of constant deck transformations.
\end{proof}

\begin{claim}{A.18}{Proved; independent small-prime check for AP-A.C2.24}
Let $\alpha^2+\alpha+1=0$, $\beta=\alpha^2$, and
\[
 C_0/\mathbf F_4:\quad v^2+v=z+\frac{\beta(z+1)}{z^2+z+\alpha}.
\]
Its geometric Raynaud theta divisor has no nonzero point of order three
or five.  Every connected geometric cyclic \'etale cover of either
prime degree has ordinary Prym.  The same statements hold for the
coefficientwise-square conjugate curve.
\end{claim}
\begin{proof}
Put $h=z^2+z+\alpha$ and $Y=hv$, so that
\[
             Y^2+hY=f=h\bigl(zh+\beta(z+1)\bigr).
\]
There are three simple Artin--Schreier poles, namely infinity and the
roots of $h$.  Let their points on $C_0$ be $W_\infty,W_1,W_2$.
The different formula gives
\[
     \operatorname{div}(dz)=2W_1+2W_2-2W_\infty,
     \qquad B=W_1+W_2-W_\infty.
\]
The line of locally exact differentials is $\mathcal O(B)$, since its
local transition functions square to those of differentials and $dz$
is exact.  Thus theta consists of the classes $[R-B]$.
The curve is ordinary: equivalently its three simple Artin--Schreier
poles give two-rank two, or the two independent logarithmic
Cartier-fixed differentials at these poles do.  The origin is not on
theta.  At any of the three ramification points $R$, the class $[R-B]$
is a nonzero two-torsion class; hence an odd-order theta point must have
$R$ away from these points.

If $\ell[R-B]=0$ for odd $\ell$, a function of divisor
$\ell R-\ell W_1-\ell W_2+\ell W_\infty$, multiplied by
$h^{(\ell+1)/2}$, is a function $F$ satisfying
\[
            \operatorname{div}(F)
                 =\ell R+W_1+W_2-(\ell+2)W_\infty.        \tag{A.11}
\]
The basis of functions regular away from infinity gives
$F=h c(z)+b(z)Y$.  The factor $h$ in the first coefficient is forced by
vanishing at $W_1,W_2$, where $Y=0$.
Taking the hyperelliptic norm and writing $r=z(R)$ gives, after scaling,
\[
        h(c^2+cb)+b^2\bigl(zh+\beta(z+1)\bigr)
                          =(z+r)^\ell.                   \tag{A.12}
\]
The odd pole order in (A.11) determines the degree and leading coefficient
of $b$, so none of the following normalizations drops a possible solution.

For $\ell=3$ we have $c=a$ constant and $b=1$.
Comparison of the four coefficients in (A.12), with $s=a^2+a$, gives
\[
      r=1+s,\qquad r^2=r,\qquad r^3=\alpha r+1.
\]
If $r=0$ the last equation reads $0=1$; if $r=1$ it reads
$1=\alpha+1$.  Both are impossible.

For $\ell=5$ write $c=az+b_0$ and $b=z+d_0$.
Set $A=a^2+a$, $B_0=ad_0+b_0+d_0^2$ and
$C_0=b_0^2+b_0d_0$.  Coefficient comparison gives
\[
 \begin{gathered}
 r=1+A,\qquad B_0=r,\qquad C_0=\beta r+1,\\
 r^4=r+1+\beta d_0^2,\qquad
 r^5=r+\alpha+\beta d_0^2.
 \end{gathered}
\]
Consequently, writing
\[
 \begin{aligned}
 R_5(r)&=r^5+r^4+\beta,\qquad
 S(r)=\alpha(r^4+r+1),\\
 T(r)&=S(r)+r,\\
 U(r)&=r^2+(r+1)S(r)+S(r)^2+\beta r+1,
 \end{aligned}
\]
we must have
\[
                 R_5(r)=0,\qquad d_0^2=S(r),
                 \qquad d_0T(r)=U(r).
\]
Squaring the last equality forces
$U(r)^2+S(r)T(r)^2=0$.  These last two polynomial conditions have
no common root: exact Euclidean division over $\mathbf F_4$ gives
\[
              \gcd\bigl(R_5,U^2+ST^2\bigr)=1.             \tag{A.13}
\]
Indeed the obstruction reduces modulo $R_5$ to
$R_4=\beta+\beta r+\alpha r^2+r^3+\alpha r^4$, and the exact identity is
\[
 (1+\beta r+\alpha r^2+\beta r^3)R_5
       +(\beta+\alpha r^3+\alpha r^4)R_4=1.
\]
The attached program \path{least_prime_semilinear.py} verifies both the
reduction and this B\'ezout identity, as well as the unreduced identity.

Finally we check the use of R5.  For each nontrivial degree-zero character
line $L^i$ of a cyclic cover, $H^1(C_0,L^i)$ has dimension one.
The locally exact differential sequence identifies the kernel of its
Frobenius arrow to $H^1(C_0,L^{2i})$ with
$H^0(\mathcal O(B)\otimes L^i)$.  The absence of theta points at orders
three and five makes every such arrow bijective, so the whole new
coherent cohomology is Frobenius-stable.  The Artin--Schreier exact
sequence identifies its stable dimension with the new two-rank.
This proves ordinariness of every cover at those orders.  Squaring all
coefficients preserves the proof and gives the conjugate curve.
\end{proof}

\begin{claim}{A.19}{Proved}
For $C_0$ in A.18, its conjugate, and every original point
$P\in C_0(\mathbf F_4)$, no cyclic $P$-split cover of degree three or
five satisfies (R) over any finite extension of $\mathbf F_4$.
Thus whenever a degree-seven witness is exhibited for such a point,
seven is its least possible odd prime, with the quantifier over all
finite extensions included.
\end{claim}
\begin{proof}
Direct point counts give $N_4=7$, $N_{16}=15$, hence the characteristic
polynomial of fourth-power Frobenius is
\[
                  P(X)=X^4+2X^3+X^2+8X+16.
\]
The attached certificate checks these counts by direct finite-field
arithmetic and verifies the following exact polynomial identities:
\[
 \begin{array}{c|c|c|c}
 \ell&d&X^d\bmod(P,\ell)&
      \text{order of every root of }P\bmod\ell\\\hline
 3&5&2&10\\
 5&13&2&52.
 \end{array}
\]
For the last column, it checks $P\mid X^n-1$ and
$\gcd(P,X^m-1)=1$ for every proper divisor $m\mid n$, with
$n=10$ or $52$.  The polynomials are separable, since these $n$ are
prime to the corresponding $\ell$.  Cayley--Hamilton therefore supplies
the scalar Frobenius identity on the actual $\ell$-torsion, without
requiring a cyclic integral lattice.  Two is primitive modulo both three
and five.  The rank hypothesis of A.17 was independently proved in
A.18.  Apply A.17 with $(\ell,d)=(3,5)$ and $(5,13)$.
The same Frobenius polynomial and theta argument apply to the conjugate.
No assumption that $P$ is a Weierstrass point was made.
\end{proof}

\subsection{The twisted Frobenius determinant, with coefficients specified}

\begin{claim}{A.20}{Proved}
Let $D\to C$ be the constant abelian $A$-cover of A.7 over
$L=\mathbf F_q$.  Let $E\subset\overline{\mathbf F}_p$ be a finite field
containing $L$ and all $|A|$th roots of unity, and let
$\bar\psi:A\to E^\times$ be nontrivial.  Deck transformations act on
cohomology by pullback.  Define
\[
 V_{\bar\psi}=
 \left(H^1(D,\mathcal O_D)\otimes_L E\right)_{\bar\psi},
 \qquad
 \Phi=(f\mapsto f^q)^*\otimes_L 1_E.
\]
Here $(f\mapsto f^q)^*$ means the map on coherent cohomology induced by
the $q$th-power map of the structure sheaf over $L$; it is $L$-linear.
Thus $\Phi$ is its \emph{$E$-linear} extension, not the absolute
$q$-semilinear map on coefficients in $E$.
It preserves $V_{\bar\psi}$, which has dimension $g-1$.  With the local
Artin convention of A.7 and reduction $\psi\mapsto\bar\psi$,
\[
           \overline{L_P(\psi,T)}
              =\det_E(1-T\Phi\mid V_{\bar\psi}).        \tag{A.14}
\]
The inverse-character convention replaces $\bar\psi$ on both sides.
In particular this gives the cohomological form of L2 in every genus,
and in genus two
\[
 \Phi\bigm|V_{\bar\psi}
    =-\sum_{R\in C(L)}\bar\psi\bigl(a([R-P])\bigr).
\]
\end{claim}
\begin{proof}
We use the opened theorem of SGA $4\frac12$, \emph{Fonctions $L$ modulo
$\ell^n$ et modulo $p$}, Theorem 2.2(b), with its hypotheses as given
in \S2.1: a separated finite-type scheme over $\mathbf F_q$, a
constructible complex of finite Tor dimension, and a noetherian reduced
coefficient ring of characteristic $p$.  It identifies the Euler product
with the alternating determinant on compactly supported \'etale
cohomology.  Formula (2.4.1) states the consequence for a constructible
sheaf over a coefficient field \cite{AC_Deligne1977}.
Take that field to be $E$, and the sheaf to be the rank-one character
summand of $\pi_*E$.  This meets every hypothesis: $C$ is smooth proper,
the cover is finite \'etale, and the sheaf is locally constant of rank
one.  In the fibre over a closed point $x$, arithmetic Frobenius
translates the torsor by $a([x-(\deg x)P])$.  Its inverse acts on the
space of functions with deck-pullback character $\bar\psi$ by the scalar
$\bar\psi(a([x-(\deg x)P]))$.  Its Euler product is therefore exactly
the reduction of the polynomial in A.7, with no unnoticed character
inversion.

The theorem makes this reduction the determinant of geometric Frobenius
on the $\bar\psi$-summand of $H^1_{\mathrm{\acute et}}
(D_{\overline L},\mathbf F_p)\otimes_{\mathbf F_p}E$.
There is no degree-zero factor because $D$ is geometrically connected
and $\bar\psi$ is nontrivial.  There is no higher factor: the
Artin--Schreier sequence on the proper curve, together with the
surjectivity of $F-1$ on a finite-dimensional semilinear vector space
over an algebraically closed field, gives
$H^i_{\mathrm{\acute et}}(D_{\overline L},\mathbf F_p)=0$ for $i\ge2$.

Katz, SGA 7 II, Expos\'e XXII, Proposition 2.2.5, gives the functorial
isomorphism between the stable part of coherent cohomology (after scalar
extension) and characteristic-$p$ \'etale cohomology.  It identifies the
linear extension of coherent $q$th-power Frobenius with the inverse of
arithmetic Galois Frobenius; Proposition 2.2.4 records equality of their
characteristic polynomials \cite{AC_Katz1973}.  Because the isomorphism
is functorial for deck automorphisms, apply the idempotent
\[
          e_{\bar\psi}=|A|^{-1}
              \sum_{\gamma\in A}\bar\psi(\gamma)^{-1}\gamma^*.
\]
It gives the same comparison on this character summand.  The nilpotent
coherent part contributes the determinant $1$, so the stable part can
be replaced by all of $V_{\bar\psi}$.  This proves (A.14), including its
Frobenius convention.

Finally the character summand of $\pi_*\mathcal O_{D_E}$ is a nontrivial
torsion line bundle of degree zero on $C_E$.  Its zeroth cohomology
vanishes, and Riemann--Roch gives first-cohomology dimension $g-1$.
For $g=2$ compare the coefficient of $T$ in (A.14) with A.7; it is
minus the scalar $\Phi$, which proves the last formula.
\end{proof}

% A.21--19 deliberately record a broader criterion than the constant-deck
% convention of the prompt.  They are placed here only for integration.

\begin{claim}{A.21}{Proved; broader pointed criterion}
Let $L$ be finite of characteristic $p$, let $C/L$ have genus at least
two, and let $P\in C(L)$.  Suppose $\pi:D\to C$ is a finite \'etale
cover, geometrically connected and geometrically Galois, of nontrivial
odd degree prime to $p$.  Assume only that a point $P'\in D(L)$ above
$P$ is specified; the geometric deck group need not be constant over
$L$.  If $p\nmid h_L(D)/h_L(C)$, then $P$ has no collision.

This condition genuinely enlarges the constant-deck criterion over a
fixed field.  The reduction to a prime cyclic quotient in A.6 is not
asserted for this enlargement: a prime cyclic quotient of the geometric
group need not be stable under arithmetic Frobenius.
\end{claim}
\begin{proof}
A hypothetical collision still gives the rational, nonzero, $p$-primary
class $T=[Q'-P']$ over this very $L$.  The lift argument of A.4 uses an
open subgroup containing $s_P(G_L)$ and does not require the arithmetic
cover to be Galois.  Its geometric Galois closure has prime-to-$p$ degree,
as required in R1.  The geometric stabilizer argument in A.4 does not
require deck transformations to be rational either.

Put $d=\deg\pi$ and $B=(\ker\pi_*)^0$.  The isogeny
$J_C\times B\to J_D$ and the averaging projector
\[
                 e=d^{-1}\pi^*\pi_*
\]
on rational $p$-primary points are defined over $L$.
Over $\overline L$, $\pi^*\pi_*=\sum_{\gamma\in G}\gamma$; arithmetic
Frobenius merely permutes the summands.  Thus the same point-count and
projector proof as in A.5 gives $(1-e)T=0$, so every geometric deck
transformation fixes $T$.  This contradicts the geometric odd-order
stabilizer statement.  Extending $L$ to make the deck group constant
would be an invalid shortcut: it can introduce rational $p$-torsion
and destroy the numerical hypothesis.  A.22 proves the claimed strict
same-field improvement by a fully specified example.
\end{proof}

\begin{claim}{A.22}{Proved; explicit nonconstant-deck witness}
Use $\mathbf F_4=\{0,1,\alpha,\alpha+1\}$ with
$\alpha^2+\alpha+1=0$, and let $C_0$ be the curve of A.18.
Its hyperelliptic quadratic twist
\[
 C_0^{\mathrm t}:\quad
 v^2+v=z+\frac{(\alpha+1)(z+1)}{z^2+z+\alpha}+\alpha
\]
has $\#J_{C_0^{\mathrm t}}(\mathbf F_4)=8$.
At its point $P=\infty$ there is a pointed geometrically cyclic
degree-seven \'etale cover over $\mathbf F_4$ with nonconstant deck
group and relative class number $8051=97\cdot83$, hence odd.
No nontrivial constant odd-order $P$-split Galois cover over
$\mathbf F_4$ satisfies the hypotheses of L1, irrespective of its
relative class number.  This last assertion concerns the fixed field;
it makes no exclusion over its extensions.
\end{claim}
\begin{proof}
Write $h=z^2+z+\alpha$, $Y=hv$.  For $C_0$ and $C_0^{\mathrm t}$ the
polynomials on the right of $Y^2+hY=f$ have ascending coefficient lists
\[
 (1,1,0,\alpha,0,1),\qquad
 (0,1,\alpha,\alpha,\alpha,1),
\]
since twisting adds $\alpha h^2$.  Direct counts give $(N_4,N_{16})$
equal to $(7,15)$ and $(3,15)$ respectively.  Thus the twisted
Frobenius characteristic polynomial is
$X^4-2X^3+X^2-8X+16$, whose value at one is eight.
These counts and both Jacobian orders are independently enumerated by
\path{nonconstant_deck_twist.py}.

Define the original degree-seven cover $D\to C_0$ as the quotient of
the Lang cover, based at $P=\infty$, by the unique order-four subgroup
$H=7J_{C_0}(\mathbf F_4)$.  Let $P'=(P,0)$ be its distinguished point.
The hyperelliptic involution $\iota$ of $C_0$ satisfies
$[\iota R-P]=-[R-P]$.  Inversion on the Jacobian commutes with the Lang
map and preserves $H$.  Consequently it gives a lift
\[
                 \widetilde\iota:(R,z)\longmapsto(\iota R,-z)
\]
to $D$, defined over $\mathbf F_4$, of order two, fixing $P'$ and
conjugating every deck transformation to its inverse.
Twist the pair $(D,C_0)$ by this pair of involutions and the nontrivial
quadratic character of $G_{\mathbf F_4}$.  It gives
$D^{\mathrm t}\to C_0^{\mathrm t}$; because $P'$ is fixed by the lift,
it remains rational.  Arithmetic Frobenius acts by inversion on the
geometric deck group $\mathbf Z/7$, so that group is nonconstant.

For completeness we compute the new numerator rather than infer its
parity from an isomorphism after extension.  Direct enumeration of the
original quotient gives rational-point occupancies
$(1,1,0,2,2,0,1)$, with the point at infinity in the zero coset.
For a primitive seventh root $\zeta$, put $t=\zeta+\zeta^{-1}$.
Then
\[
 t^3+t^2-2t-1=0,
 \qquad S=1+\zeta+\zeta^{-1}+2(\zeta^3+\zeta^{-3})
           =3-t-2t^2,
 \qquad L_{\chi}(T)=1+ST+4T^2.
\]
The involution interchanges the $\chi$ and $\chi^{-1}$ cohomology
summands and commutes with the original Frobenius.  On their direct
sum, twisted Frobenius therefore has determinant
\[
       \det(1-T\widetilde\iota F)
              =L_\chi(T)L_\chi(-T).
\]
Taking the product over the three inverse pairs, or the norm from the
real cubic cyclotomic field, gives the relative numerator
\[
 1+10T^2+65T^4+279T^6+1040T^8+2560T^{10}+4096T^{12}.
\]
At one it is $8051$.  Equivalently, the norms of $5+S$ and $5-S$ are
$97$ and $83$.  The attached exact program verifies the occupancy
vector, the cubic identity, the complete relative polynomial and its
value.  Thus A.21 applies over $\mathbf F_4$ itself.

Finally any nontrivial constant odd-order pointed Galois cover would,
by the group-theoretic step of A.6, yield a constant prime cyclic
$P$-split cover.  Class field theory would give a nonzero odd-prime
quotient of $J_{C_0^{\mathrm t}}(\mathbf F_4)$, which has order eight.
This is impossible, proving the last assertion.
\end{proof}

\begin{claim}{A.23}{Proved; simple Frobenius quotient and exact descent power}
Let $D\to C$ be a cover satisfying the broadened hypotheses of A.21
over a finite field $L$, and suppose $p\nmid h_L(D)/h_L(C)$.
There is an intermediate pointed geometrically Galois cover $E\to C$
over the same $L$, with geometric deck group
\[
                    V\simeq(\mathbf Z/\ell)^r,
             \qquad \ell\text{ odd},\quad \ell\ne p,
             \quad 1\le r\le2g,
\]
such that arithmetic Frobenius acts irreducibly on $V$ and
$p\nmid h_L(E)/h_L(C)$.  Let $m$ be the order of this action, so
$m\mid\ell^r-1$.  Write $L_m/L$ for the degree-$m$ extension.  Then
the deck group becomes constant over $L_m$ and
\[
             \frac{h_{L_m}(E)}{h_{L_m}(C)}
                  =\left(\frac{h_L(E)}{h_L(C)}\right)^m. \tag{A.15}
\]
Consequently there is a constant cyclic degree-$\ell$ pointed witness
for L1 over $L_m$.

For a fixed pair $(C/k,P)$, existence of a witness after some finite
extension is therefore equivalent for the broadened criterion A.21
and for the original constant-deck criterion A3.  Nonconstant deck
groups can lower the field of a witness, as in A.22, but do not enlarge
the class of pairs reached when all finite extensions are allowed.
\end{claim}
\begin{proof}
Let $G$ be the geometric deck group.  The rational point $P'$ gives an
actual Frobenius action on $G$ by comparing the distinguished point in
the geometric fibre; equivalently the associated finite arithmetic
quotient is the semidirect product determined by this section.
Feit--Thompson makes $G$ solvable, so its abelianization $A$ is
nontrivial.  Frobenius acts on $A$.  Choose an odd prime $\ell$ dividing
its order and pass to the nonzero quotient of its $\ell$-primary part
modulo $\ell$.  Choose a maximal proper Frobenius-stable subspace.
The resulting quotient $V$ is a nonzero irreducible
$\mathbf F_\ell[F,F^{-1}]$-module.  Its kernel in $G$ is normal and
Frobenius-stable.  Thus the corresponding geometric intermediate cover
descends to an $L$-cover $E\to C$, and the image of $P'$ is still
rational.  Since $V$ is a quotient of the geometric abelianized
fundamental group modulo $\ell$, its dimension is at most $2g$.

The same Jacobian isogeny argument as in A.6 makes
$h_L(D)/h_L(E)$ and $h_L(E)/h_L(C)$ positive integers.  Their product is
prime to $p$, so the second is prime to $p$.  An irreducible finite
$\mathbf F_\ell[F,F^{-1}]$-module is the additive group of
$\mathbf F_{\ell^r}$, with $F$ acting by multiplication by some
$a\in\mathbf F_{\ell^r}^{\times}$.  Its order $m$ divides
$\ell^r-1$.  Every nonzero vector, and every nontrivial character in
the dual module, has orbit of length exactly $m$: if $a^d\ne1$,
multiplication by $a^d-1$ is invertible.  The same conclusion holds for
the inverse or dual action, so no Frobenius convention affects it.

To prove (A.15), choose an auxiliary prime $\nu\ne p$ and a finite
extension $K/\mathbf Q_\nu$ containing the $\ell$th roots of unity.
On the complement of the invariant summand in
$H^1_{\mathrm{\acute et}}(E_{\overline L},K)$, geometric Frobenius
$\mathcal F$ permutes the nontrivial $V$-character summands.  Group them
by Frobenius orbits.  Each orbit has length $m$.  If $W_0$ is one of
its summands and $A_0=\mathcal F^m|W_0$, transport a basis of $W_0$
around the orbit by $\mathcal F$.  In these bases $\mathcal F$ is a
cyclic block matrix with final block $A_0$.  Elementary elimination
gives
\[
 \det(1-T\mathcal F\mid\text{orbit})
       =\det(1-T^m A_0\mid W_0),\qquad
 \det(1-\mathcal F^m\mid\text{orbit})
       =\det(1-A_0\mid W_0)^m.
\]
Multiply over the orbits and evaluate the first identity at $T=1$.
The invariant part is exactly the cohomology of $C$, because averaging
by $|V|$ splits the finite \'etale pullback and trace.  The determinant
of its complement at $1$ is the relative class number, both over $L$
and over $L_m$.  This proves the exact equality (A.15), without any
assumption that $m$ is prime to $p$ and without estimating eigenvalues.

Over $L_m$ the geometric deck group $V$ is constant and the chosen point
remains rational.  Equality (A.15) preserves the prime-to-$p$ condition.
Any quotient $V\twoheadrightarrow\mathbf Z/\ell$ gives a constant
prime cyclic intermediate cover; integer relative class-number ratios
again preserve the condition.  This gives the asserted original L1
witness.  The reverse implication between the two existence assertions
is immediate.  Thus unrestricted enlargement of fields would indeed
have been unsafe in A.21, but this particular extension after the
irreducible-module reduction is justified by (A.15).
\end{proof}

\begin{claim}{A.24}{Proved; exact least-prime-five witness}
For the genus-two curve
\[
                  C/\mathbf F_4:\quad y^2+y=x^5+x
\]
and its hyperelliptic quadratic twist, the least odd prime degree of
a constant pointed L1 witness at $P=\infty$, with all finite extensions
allowed, is five.  A witness is defined over $L=\mathbf F_{256}$.
For the specified witness its four nontrivial twisted polynomials are
all
\[
               L_P(\chi^i,T)=1-32T+256T^2=(1-16T)^2,
                 \qquad 1\le i\le4,
\]
and its relative class number is
$225^4=2\,562\,890\,625=3^8 5^8$, which is odd.
The covering curve is the smooth Fermat quintic.
\end{claim}
\begin{proof}
We use the explicit polynomial basis
$L=\mathbf F_2[t]/(t^8+t^4+t^3+t+1)$; an integer denotes the element
whose binary digits are the coefficients in this basis.  Set
\[
            r=18,\qquad a=254,\qquad b=79,\qquad c=44,
            \qquad F=y+a x^2+b x+c.
\]
The exact certificate \path{supersingular_least5.py} verifies
\[
 \begin{gathered}
 r^{16}+r+1=0,
 \qquad b=1+r^4,\qquad a=b^2,\qquad c^2+c=r^5,\\
 F(F+1)=(x+r)^5,
 \qquad R=(r,42)\in C(L),
 \qquad \operatorname{div}(F)=5R-5P.
 \end{gathered}
\]
For the divisor assertion, $F$ has pole order five at infinity because
$y$ has that order and $a x^2+b x+c$ has pole order at most four.
The norm identity shows that its only zero is at the one point $R$
where $F=0$, with multiplicity five; at its conjugate $F=1$.

Define $D\to C_L$ by $z^5=F$.  It is unramified because every valuation
of $F$ is divisible by five.  It is geometrically connected: otherwise
$F$ would be a fifth power and $R-P$ principal, impossible on a positive
genus curve.  The group $\mu_5$ is constant over $L$ since $5\mid255$.
At $P$, the function $u=x^2/y$ is a uniformizer and the residue of
$u^5F$ is one.  Thus the fibre consists of five rational points, and we
choose the one with residue $uz=1$.  This proves the field and pointing
requirements, rather than obtaining them by an uncontrolled extension.
Putting $w=(x+r)/z$ in the norm identity gives
$w^5+z^5=1$, so the smooth projective covering curve is the Fermat
quintic, of genus six.

The same program enumerates all $193$ points of $C(L)$ and evaluates
the Kummer character using $F(Q)^{51}\in\mu_5$.
At $R$ the leading unit is also one, since
$F/(x+r)^5=(F+1)^{-1}$; its character is therefore zero, as is the
character at $P$.  With $12$ as the primitive fifth root in the chosen
field representation, the five occupancies are exactly
\[
                          (13,45,45,45,45).
\]
For each nontrivial fifth-root character the point sum is $13-45=-32$.
The canonical divisor is $2P$, so A.7 gives precisely the four
polynomials stated in the claim.  Their product at one gives the
relative class number.  The covering curve has $5\cdot13=65$ rational
points, an additional check of the Kummer character evaluation.

We exclude degree three over every extension, not only the fields in
the numerical search.  The locally exact differential line is
$\mathcal O(P)$, because $\operatorname{div}(dx)=2P$.
If its theta divisor contained a nonzero point of order three, there
would be a point $R_0\ne P$ and a function with divisor $3R_0-3P$.
Every function regular off $P$ has the form $A(x)+B(x)y$.
Its two terms have respectively even pole order and pole order
$5+2\deg B$, so a function with pole order at most three lies in
$\langle1,x\rangle$ and cannot have pole order three.
There is therefore no such theta point.  Since
$\kappa_P=[K_C-2P]=0$, A.25 excludes every cubic witness over every
finite extension.  This proves minimality of five.

For the quadratic twist, choose $d\in\mathbf F_{16}\subset L$ with
$d^2+d=\alpha$, where $\alpha$ is the nontrivial Artin--Schreier
twisting constant in $\mathbf F_4$.  The substitution $y\mapsto y+d$
gives the isomorphism over $L$, preserving $P$.  Pulling back the
specified Kummer function supplies the same pointed cover and the same
polynomials.  The geometric theta-three exclusion and $\kappa_P=0$
also persist, proving minimality for the twist.
\end{proof}

\begin{claim}{A.25}{Proved; exact all-extension cubic criterion}
Let $C/k$ be any genus-two curve over a finite field of characteristic
two, let $P\in C(k)$, and put
\[
                       \kappa_P=[K_C-2P]\in J_C(k).
\]
Let $\Theta$ be the geometric theta divisor of the line of locally
exact differentials.  A constant cyclic degree-three $P$-split cover
satisfying L1 exists over some finite extension of $k$ if and only if
at least one of the following holds:
\begin{enumerate}
\item $\Theta$ contains a nonzero point of order three;
\item $\kappa_P\notin3J_C(k)$.
\end{enumerate}
In the second case a witness already exists over $k$.
If neither holds, no extension can produce a cubic witness.

More precisely, for an individual cubic cover over $L\supseteq k$
with character $\chi:J_C(L)\twoheadrightarrow\mathbf Z/3$, criterion
L1 holds exactly when its geometric new two-rank is zero or
$\chi(\kappa_P)\ne0$.
\end{claim}
\begin{proof}
Reduce a primitive character to $\psi:\mathbf Z/3\to\mathbf F_4^\times$
and write
\[
                 S=\sum_{R\in C(L)}\psi(\chi([R-P]))
                         \quad\text{in }\mathbf F_4.
\]
The middle-coefficient Riemann--Roch identity of A.7 is
$c_1(\psi)=\psi(\chi(\kappa_P))c_1(\psi^{-1})$.
Since inversion on $\mu_3$ is squaring and the coefficients of the
point sum are integers, reduction gives the exact equality
\[
                    S=\psi(\chi(\kappa_P))S^2.          \tag{A.16}
\]
Thus either $S=0$, or $S=\psi(\chi(\kappa_P))^{-1}$.
By A.7--A.8 the relative class number is odd precisely when
$1+S\ne0$; the other character has value $(1+S)^2$, so it gives the
same condition.  If $S\ne0$, this condition is exactly
$\chi(\kappa_P)\ne0$.

It remains to identify $S=0$ correctly.  By A.20, the coherent
$\psi$-eigenspace is one-dimensional and its $q$th-power Frobenius
scalar is $S$.  The other scalar is $S^2$.
If they vanish, the new stable cohomology has dimension zero, hence
the geometric Prym has two-rank zero by Artin--Schreier.  If they do
not vanish, its two-rank is two, so the new factor is ordinary.
This proves the individual-cover assertion.

The locally exact differential sequence used in R5 says that the
Frobenius arrow on $H^1(C,L_0)$ vanishes precisely when
$L_0\in\Theta$, for a nontrivial degree-zero character line.
The two nontrivial character lines of a cubic cover are interchanged
by squaring.  Consequently a geometric cubic cover has new two-rank
zero exactly when its cyclic order-three subgroup meets $\Theta$.
If the first condition of the claim holds, descend that cover and all
its deck transformations to a finite extension and then rationalize
a point over $P$.  Its geometric two-rank zero is preserved by every
extension and twist, so this procedure gives a valid L1 witness.

If the first condition fails, every geometric cubic cover is ordinary.
If $\kappa_P\notin3J_C(k)$, the nonzero image of $\kappa_P$ in the
finite $\mathbf F_3$-vector space $J_C(k)/3J_C(k)$ is detected by a
linear functional to $\mathbf F_3$.  Its Lang cover is a witness by
the individual criterion.  If instead $\kappa_P=3x$ with $x\in J_C(k)$,
the same equality holds in $J_C(L)$ for every finite extension $L$.
Every cubic character over every such $L$ kills $\kappa_P$, and none
can be a witness.  This proves both directions with the extension
quantifier included.
\end{proof}

\begin{claim}{A.26}{Proved; canonical-class sufficient condition}
Let $C/k$ have genus two in any positive characteristic $p$, and let
$P\in C(k)$.  Suppose $\ell\ne p$ is an odd prime with
$-1\in\langle p\rangle\subset(\mathbf Z/\ell)^\times$.
If $\kappa_P=[K_C-2P]\notin\ell J_C(k)$, there is a degree-$\ell$
constant $P$-split L1 witness over $k$ itself.
\end{claim}
\begin{proof}
Choose a quotient $\chi:J_C(k)\twoheadrightarrow\mathbf Z/\ell$
which does not kill $\kappa_P$.  For any nontrivial character $\psi$
of this quotient, let $S_\psi$ be the point-character sum reduced
modulo $p$.  Choose $j$ with $p^j\equiv-1\pmod\ell$.
The same middle-coefficient symmetry as above gives
\[
          S_\psi=\psi(\chi(\kappa_P))S_{\psi^{-1}}
                =\psi(\chi(\kappa_P))S_\psi^{p^j}.
\]
If the genus-two test failed at this character, $1+S_\psi=0$,
so $S_\psi=-1$ and this equality would force
$\psi(\chi(\kappa_P))=1$.  That contradicts the choice of $\chi$,
because every nontrivial character of a group of prime order is
faithful.  Thus no nontrivial factor vanishes modulo $p$, and A.8
proves L1.  The assumption is only that inversion is a power of $p$;
it does not require $p$ to generate the whole multiplicative group.
\end{proof}

\begin{claim}{A.27}{Proved; ramified elliptic-factor obstruction}
Let $L$ be a finite field of characteristic two.  Suppose $C/L$ has
genus at least two and a separable degree-two morphism
$f:C\to E$ to an elliptic curve.  Let $P\in C(L)$ be a ramification
point, and use $f(P)$ as the origin of $E$.
Let $\varphi:E'\to E$ be an \'etale isogeny of odd degree $n>1$
over $L$, with no assumption that its kernel is constant, and let
\[
                         D=C\times_E E'.
\]
Then $D\to C$ is a connected \'etale pointed geometrically Galois
cover with odd geometric deck group, but
\[
                            2\mid h_L(D)/h_L(C).
\]
Consequently a prime cyclic character factoring through this elliptic
quotient cannot be an L1 witness.  The assertion persists over every
finite extension.  Neither L1 nor its nonconstant-deck enlargement
A.21 can use this elliptic-isogeny pullback.
\end{claim}
\begin{proof}
The cover $E'\to E$ is geometrically connected of odd degree, whereas
$C\to E$ has degree two.  Their function fields are linearly disjoint,
since their intersection degree divides both degrees.  Thus the fibre
product is geometrically connected; it is smooth because $D\to C$
is \'etale.  Put $G=\ker\varphi(\overline L)$, a group of order $n$.
The fibre over $P$ consists of ramification points $R_g$ of the
quadratic map $\rho:D\to E'$, indexed by $B_g=g\in G$.
The point $R_0$ over the origin of $E'$ is rational; the remaining
points are permuted by Galois.  Set
\[
        t=n\left[\sum_{g\ne0}R_g-(n-1)R_0\right]\in J_D(L).
\]
This class is rational because its divisor is Galois-invariant.
Each $n[B_g-B_0]=0$, and the quadratic pullback sends $B_g$ to
$2R_g$, so $2t=0$.  We prove $t\ne0$, including geometrically.
Otherwise there would be a function $u$ on $D_{\overline L}$ with
$\operatorname{div}(u)=n\sum_{g\ne0}R_g-n(n-1)R_0$.
The quadratic involution $\tau$
fixes each ramification point, so $\tau(u)/u$ is a constant $c$.
Applying $\tau$ again gives $c^2=1$, and characteristic two forces
$c=1$.  Thus $u$ is pulled back from a function on $E'_{\overline L}$.
Its valuations at these ramification points must be even, contradicting
the odd valuation $n$ at any $R_g$ with $g\ne0$.
Hence $t$ is nonzero rational
two-torsion.

The deck average of every difference $R_g-R_0$ is zero, so
\[
                          \sum_{\gamma\in G}\gamma(t)=0.
\]
On two-primary torsion the averaging projector is defined because
$n$ is odd.  It kills $t$, putting $t$ in the new summand, and hence in
the connected norm kernel on two-primary torsion as in A.5.
The projector is defined over $L$ even when its individual summands
are not, as checked in A.21.
This new abelian factor therefore has an even number of rational
points.  Its point count equals $h_L(D)/h_L(C)$.

Finally, a constant odd prime cyclic unramified cover of $E$ pointed
over its origin is an \'etale elliptic isogeny after using the chosen
point as the origin upstairs.  A character of $J_C(L)$ obtained by
composition with $f_*:J_C(L)\to E(L)$ gives its pullback.  This
establishes the last statement for exactly the covers in question;
it does not exclude characters involving more than one elliptic
factor of $J_C$.
\end{proof}

\begin{claim}{A.28}{Proved; a fixed-field obstruction to the enlarged criterion}
Let $c\in\mathbf F_4\setminus\mathbf F_2$ and
\[
           C_c/\mathbf F_4:\quad y^2+y=x+\frac{c}{x^2+x},
                   \qquad P=\infty.
\]
There is no pointed geometrically Galois cover of nontrivial odd
degree over $\mathbf F_4$ satisfying the relative-class-number
criterion, even if its geometric deck group is nonconstant or
nonabelian.  Thus the variant of A3 requiring the original field
while allowing nonconstant odd deck groups is false.
This does not refute A3, which allows finite extensions.
\end{claim}
\begin{proof}
Choose $\alpha\in\mathbf F_4$ with $\alpha^2+\alpha=1$ and put
$\beta=\alpha^2$.  The two involutions
\[
  \tau_1(x,y)=(x+1,y+\alpha),\qquad
  \tau_2(x,y)=(x+1,y+\beta)
\]
are defined over $\mathbf F_4$, fix $P$, and have product the
hyperelliptic involution $\iota$.  Their elliptic quotients have
explicit coordinates
\[
 \begin{array}{c|c|c}
 &u&w\\\hline
 E_1&x^2+x&y+\alpha x\\
 E_2&x^2+x&y+\beta x
 \end{array}
 \qquad
 \begin{aligned}
 E_1:&\quad w^2+w=\beta u+c/u,\\
 E_2:&\quad w^2+w=\alpha u+c/u.
 \end{aligned}
\]
Each map has degree two and ramifies at $P$.
Using the images of $P$ as origins, let
$\Phi:E_1\times E_2\to J_C$ be the sum of their Jacobian pullbacks,
and $\Psi:J_C\to E_1\times E_2$ their pair of norms.
The hyperelliptic involution is $[-1]$ on $J_C$, so
$\tau_2=-\tau_1$ there.  Pullback followed by norm is $[2]$;
the cross terms vanish, since each is both invariant and anti-invariant
under $\tau_1$ and a homomorphism of abelian varieties killed by two
must be zero.  Consequently
\[
                   \Psi\Phi=[2],\qquad \Phi\Psi=[2].
\]
These are isogenies whose kernels are killed by two, giving a direct
sum decomposition on every odd Tate module.

Counting the two elliptic curves over $\mathbf F_4$ gives the
unordered pair of point counts $8,4$, hence the Frobenius polynomials
\[
                       X^2+3X+4,\qquad X^2-X+4.
\]
The attached certificate verifies these counts for both values of $c$.
Modulo every odd prime $\ell$, these polynomials are coprime: their
difference is $4X$, while neither has zero as a root.
Every Frobenius-simple quotient of the odd geometric abelianized
fundamental group must therefore be supported in just one of the two
elliptic factors.  A pointed cover with that quotient is the pullback
of the corresponding pointed unramified cover of that elliptic curve.
Such a genus-one cover, with its chosen point used as origin, is an
\'etale elliptic isogeny over $\mathbf F_4$, with no assumption that its
kernel be constant.  The pointing removes a possible constant torsor
twist, so this comparison is over the same field.

If an odd nonabelian or nonconstant witness existed, A.23 would first
give such a Frobenius-simple elementary abelian quotient witness over
the same field.  A.27 excludes its elliptic-isogeny pullback, since
both elliptic quotient maps ramify at $P$.  This contradiction proves
the assertion for every odd group, not merely the covers in a finite
search or a multiplication-isogeny family.
\end{proof}

\begin{claim}{A.29}{Proved; stronger all-extension small-prime and field obstructions}
For the pointed curves $(C_c,P)$ of A.28:
\begin{enumerate}
\item No cubic or quintic constant L1 witness exists over any finite
extension of $\mathbf F_4$.
\item No nontrivial constant odd-group L1 witness exists over
$\mathbf F_{4^m}$ for any $1\le m\le15$.
\item At $m=16$, the only possible prime of a prime cyclic witness is
$113$. The positive certificate at this degree is supplied in A.31.
\end{enumerate}
These conclusions about constant covers also hold at every rational
hyperelliptic ramification point of $C_c$.
\end{claim}
\begin{proof}
We first check the geometric small-prime assertions independently of
the previous-return certificates.  Write $h=x^2+x$, $Y=hy$, so
$Y^2+hY=h(xh+c)$.  The locally exact line has divisor
$B=W_0+W_1-P$.  An odd-order theta point is of the form $[R-B]$ with
$R$ away from the three ramification points.  If its order divides
$n=3$ or $5$, multiplication of a principal function by
$h^{(n+1)/2}$ gives a function with divisor
$nR+W_0+W_1-(n+2)P$.

For $n=3$, its pole bound and zeros at $W_0,W_1$ force it, after
scaling, to be $ah+Y$.  Taking its norm and dividing by $h$ gives
\[
             h(a^2+a)+xh+c=(x+r)^3.
\]
Coefficient comparison yields $r^2+r=1$ and $r^3=c$.
The first equation gives $r^3=1$, contradicting $c\ne1$.
Thus theta has no nonzero three-torsion.  Since $K_C\sim2P$,
A.25 excludes every cubic witness over every extension.

For $n=5$, the same pole and zero conditions give
$h(ax+b)+(x+d)Y$.  Put $s=d^2$ and compare its norm, divided by $h$,
with $(x+r)^5$.  The resulting necessary conditions are
\[
 \begin{gathered}
 Q=r^4+r+1+c=0,\qquad s=r^5/c,\qquad a^2+a=r+1,\\
 E=ad+b=s+r+1,\qquad B_0=b^2+bd=r+1+c,\\
 dE=U,\qquad U=B_0+E^2+(r+1)s.
 \end{gathered}
\]
Hence $U^2+sE^2=0$.
Using $c^2+c=1$, its remainder modulo $Q$ is
$R_2=(c+1)(r^2+r)$.  There is the explicit identity
\[
                    cQ+(c+1)(r^2+r+1)R_2=1.
\]
Both this remainder and the identity are checked over $\mathbf F_4$
in \path{bielliptic_obstructions.py}.  Thus theta has no nonzero
five-torsion either.  By the checked locally exact differential
sequence, every geometric cyclic quintic cover has ordinary new
factor, of two-rank four.

We must still exclude rational quintic witnesses, rather than
substituting geometric ordinariness for that assertion.  From A.28,
Frobenius on $J_C[5]$ has characteristic polynomial
\[
                   (X-3)^2(X^2-2X-1)\quad\text{over }\mathbf F_5.
\]
On the first primary part it is $3I+N$, with $N^2=0$; on the second
part its cube is $2I$, and every nonzero vector has period twelve.
Apply the same decomposition to a nonzero quotient functional $v$.
Its period $m$ and the action at $d=m/4$ are exhausted by
\[
 \begin{array}{c|c|c|c}
 \text{support of }v&m&d&vF^d/v\\\hline
 \text{first part, }vN=0&4&1&3\\
 \text{first part, }vN\ne0&20&5&3\\
 \text{second part present, }vN=0&12&3&2\\
 \text{second part present, }vN\ne0&60&15&2.
 \end{array}
\]
The computation checks all $624$ nonzero vectors for each of the two
possible Jordan types, in addition to the displayed algebraic proof.
The geometric cover has a pointed model over $\mathbf F_{4^d}$ with
deck Frobenius equal to $2$ or $3$; both generate
$\mathbf F_5^\times$.  On its ordinary new two-torsion, the argument
of A.17 identifies this Frobenius with a semilinear map of the
one-dimensional $\mathbf F_{16}$-space, whose fourth power is the
norm to $\mathbf F_2$ and therefore the identity.  Thus all new
two-torsion is rational at the first constant-deck field
$\mathbf F_{4^m}$ and remains rational over every later eligible
field.  This excludes every quintic witness.

For the field bound, the isogenies in A.28 identify the odd-prime
quotients of $J_C(\mathbf F_{4^m})$ with those of the two elliptic
groups.  If a character is supported in just one factor, A.27
excludes it.  A remaining prime must therefore divide both elliptic
orders.  With traces $t_1=-3,t_2=1$, define
\[
 u_i(0)=2,\quad u_i(1)=t_i,\quad
 u_i(m)=t_i u_i(m-1)-4u_i(m-2),\quad
 N_i(m)=4^m+1-u_i(m).
\]
The exact gcd computation gives, for $1\le m\le15$,
\[
 \gcd(N_1(m),N_2(m))_{\mathrm{odd}}=
 \begin{cases}
 3,&m=4,8,\\
 135,&m=12,\\
 1,&\text{otherwise}.
 \end{cases}
\]
Only three and five occur, and both were excluded over all
extensions.  At $m=16$ the odd gcd is $339=3\cdot113$.
This proves the two field assertions for prime covers and hence,
by A.6, for all constant odd groups.

Finally any two hyperelliptic ramification points differ by a class
killed by two in the Jacobian.  Every odd-order character over every
extension kills their difference.  The exact twist formula A.9
therefore identifies their pointed abelian covers and relative
polynomials.  A.6 reduces the odd-group assertion to these prime
abelian covers, proving the last statement.
\end{proof}


\begin{claim}{A.30}{Proved; exact Cartier--Mumford rational test}
Let $k=\mathbf F_q$, $q=2^d$, and let
\[
 C/k:\quad Y^2+h(x)Y=f(x),\qquad \deg h\le2,
 \quad f\text{ monic of degree }5
\]
be smooth of genus two, based at its unique point $P=\infty$.
Let $n>1$ be odd and let $a\in J_C(\overline k)$ have exact order
$n$ and satisfy $\sigma_q(a)=[q]a$, where $\sigma_q$ is arithmetic
Frobenius. Choose a finite extension $E/k$ containing $a$ and $\mu_n$.
For $0<i<n$, write the unique reduced Mumford representative of $ia$
as $(u_i,v_i)$, with $u_i$ monic of degree $r_i\in\{1,2\}$, so that
$ia=[D_i-r_iP]$.

For $j=2i\bmod n$, compute $c_i$ as follows.
If $r_i=1$, put $c_i=h_1$, the coefficient of $x$ in $h$.
If $r_i=2$ and $\gcd(u_i,h)$ has degree one, put $c_i=h_2$.
Otherwise $r_i=2$ and $\gcd(u_i,h)=1$; let $v$ be the Cantor
doubling interpolation of degree at most three, with
$u_i^2\mid v^2+hv+f$ and $v\equiv v_i\bmod u_i$.
If the reduced double has degree one, put $c_i=1$.
If it has degree two, let $v_3$ be the cubic coefficient of $v$
and $b_j$ the coefficient of $x$ in $u_j$, and put $c_i=v_3b_j$.
No other case occurs for a nonzero point of odd order. Define
\[
 A_i=\prod_{s=0}^{d-1}c_{2^s i}^{\,2^{d-1-s}}\in E.       \tag{A.17}
\]
There is a unique constant cyclic $P$-split degree-$n$ cover over $k$
with geometric character subgroup generated by $a$.
Its linear coherent $q$-Frobenius scalar on the summand for $ia$ is
$\sigma_q^{-1}(A_i)$. Consequently its relative class number is odd
if and only if
\[
                         A_i\ne1\qquad(0<i<n).           \tag{A.18}
\]
For the primitive factor alone use only $\gcd(i,n)=1$.
These formulas refer to the pointing at infinity; another base point
requires the twist of A.9.
\end{claim}
\begin{proof}
Put $L_i=\mathcal O(D_i-r_iP)$ and rigidify it at $P$ using
$t=x^2/Y$. Monicity of $f$ gives leading residues one for $xt^2$
and $Yt^5$. The differential $dx/h$ has divisor $2P$.
At finite unramified points this follows from $x$ being a local
parameter. At a finite ramification point,
$dx/h=dY/(f'+h'Y)$ is nonvanishing by smoothness.
Thus its canonical divisor is supported at $P$ and has degree two.
A basis of the one-dimensional Serre dual space
$H^0(C_E,\Omega_C\otimes L_i^{-1})$ is
\[
                            \omega_i=u_i\,dx/h.
\]
Indeed $\operatorname{div}(u_i)=D_i+\iota D_i-2r_iP$ and $r_i\le2$.
Let $e_i$ be its dual basis in $H^1(C_E,L_i)$.

Let $g_i$ have divisor $2D_i-D_j-(2r_i-r_j)P$, normalized to
leading residue one at $P$ relative to $t$.
Multiplication by $g_i$ identifies $L_i^2$ with $L_j$ respecting
rigidifications. Absolute Frobenius followed by this identification
is square-semilinear. Cartier is its Serre adjoint:
\[
 \langle g_iF(z),\omega\rangle
             =\langle z,\mathcal C(g_i\omega)\rangle^2.
\]
This identity follows directly in the residue pairing: the coefficient
of $t^{-1}dt$ in $z^2\eta$ is the square of the residue of
$z\mathcal C(\eta)$. Therefore
$\mathcal C(g_i\omega_j)=\sqrt{c_i}\omega_i$ is exactly the assertion
$F(e_i)=c_i e_j$.

We verify all four cases.
If $r_i=1$, the point is unramified, since a ramification point would
give a two-torsion class. Then $D_j=2D_i$, $u_j=u_i^2$, and $g_i=1$.
The identity $\mathcal C(dx/h)=\sqrt{h_1}\,dx/h$ gives $c_i=h_1$.

If $r_i=2$ and $\gcd(u_i,h)=x-w$, write $D_i=W+R$ with $W$
ramified and $R$ unramified. Since $2W\sim2P$, we have $D_j=2R$
and $g_i=x-w$. With $u_i=(x-w)(x-r)$, direct extraction of odd
powers gives
\[
 \mathcal C\bigl((x-w)(x-r)^2dx/h\bigr)
             =\sqrt{h_2}(x-w)(x-r)dx/h.
\]
For this calculation write $(x-w)/h=(x-w)h/h^2$ and use $h(w)=0$:
the square root of the odd part is $\sqrt{h_2}(x-w)$.
Hence $c_i=h_2$. Two ramification points in $D_i$ would give a
two-torsion class and are excluded.

In the remaining cases $\gcd(u_i,h)=1$.
For a degree-two reduced double, Cantor reduction gives
\[
 v^2+hv+f=v_3^2u_i^2u_j,\qquad
                     g_i=(Y+v)/(v_3u_j).
\]
This function has the required divisor and normalized leading residue.
The Artin--Schreier relation for $Y/h$ gives
\[
 \mathcal C((Y+v)dx/h)=v_3\sqrt{b_j}\,u_i dx/h.
\]
Indeed the term $(Y/h)^2dx$ has Cartier zero, and the odd powers of
$f+hv$ are exactly those of $v_3^2u_i^2u_j$, since $v^2$ has only
even powers. Dividing by $v_3$ before Cartier gives
$\sqrt{v_3b_j}\omega_i$.
For a degree-one double, $v$ has degree at most two, the norm
$v^2+hv+f=u_i^2u_j$ is monic of degree five, and
$g_i=(Y+v)/u_j$. The same calculation with $u_j$ monic linear
gives Cartier coefficient one.

Iteration gives $F^d(e_i)=A_i e_{qi}$, with $A_i$ as in (A.17).
The hypothesis $\sigma_q(a)=qa$ makes coefficientwise conjugation
of the reduced representative $(u_i,v_i)$ equal to $(u_{qi},v_{qi})$.
Consequently $\sigma_q(\omega_i)=\omega_{qi}$ and
$\sigma_q(e_i)=e_{qi}$, with rigidifications unchanged.
The linearized coherent operator is $\sigma_q^{-1}F^d$,
whose scalar is $\sigma_q^{-1}(A_i)$.
This explains the $q$th-root correction; it cannot change whether
the scalar is one.

Finally, $\sigma_q(a)=qa$ is precisely invariance of the associated
geometric homomorphism to the constant group $\mathbf Z/n$.
The factor $q$ records Frobenius on $\mu_n$ in the Kummer
identification. The section of $P$ gives its unique pointed descent.
All multiplication functions were normalized at $P$, so no constant
character twist is omitted. Applying A.20 to these one-dimensional
character spaces proves (A.18). Selecting only primitive characters
gives the primitive-factor statement.
\end{proof}


\begin{claim}{A.31}{Proved; mixed witnesses and exact prime/field minima}
For the pointed split curves $(C_c,\infty)$ of A.28, the least
odd prime of a constant pointed L1 witness over all finite extensions
is eleven. A degree-eleven witness is defined over $\mathbf F_{4^{60}}$,
and sixty is the least extension degree for that prime.
The least extension degree permitting any constant odd-group witness
is sixteen, where a degree-$113$ witness exists.

For the following eight $\mathbf F_4$ models, write
$c\in\{\alpha,\alpha^2\}$ and $\delta\in\{0,\alpha\}$:
\[
 y^2+y=x+\frac{c}{x^2+x}+\delta,\qquad
 y^2+y=x+\frac{c}{x^2+x+\alpha}+\delta.
\]
Each original rational point has an eleven witness over
$\mathbf F_{4^{60}}$ and a $113$ witness over $\mathbf F_{4^{16}}$.
Whenever the cubic test A.25 fails, eleven is the least odd prime.
In particular this holds for every rational point of the split
models and for infinity on each quadratic-pole model.
\end{claim}
\begin{proof}
We first specify and independently check the positive certificates.
Use the field
\[
 K=\mathbf F_2[t]/(t^{120}+t^4+t^3+t+1)=\mathbf F_{4^{60}},
\]
with binary polynomial coefficients encoded by integers, and set
\[
 \alpha=931750399618260646238930031330849980.
\]
The model is $Y^2+(x^2+x)Y=x^5+x^3+\alpha x^2+\alpha x$.
The following reduced pair defines the degree-eleven character line:
\[
 \begin{aligned}
 u(x)&=x^2+u_1x+u_0,\qquad v(x)=v_1x+v_0,\\
 u_0&=1286551050156223861353970523064590862,\\
 u_1&=855838071379080258553061447214623248,\\
 v_0&=372422475677710747247258963726964873,\\
 v_1&=880561393008999941135507307093448287.
 \end{aligned}
\]
The certificate verifies irreducibility of the field polynomial,
$\alpha^2+\alpha=1$, validity of the reduced pair, and exact order
eleven of its Jacobian class $a$. Since $q=2^{120}\equiv1\bmod11$,
the descent condition $\sigma_q(a)=qa$ holds over $K$.

All ten reduced multiples and Cartier coefficients are recorded in
\path{bielliptic_mixed_ell11_F4pow60.json}.
The independent verifier \path{verify_mixed_certificates.py} uses
pure Python carryless field arithmetic, with neither the searcher's
compiled field helper nor its Cartier-coefficient routine.
It recomputes every multiple and reconstructs doubling interpolation
from $(f+hv+v^2)/u$ and $h^{-1}\bmod u$.
All ten $q$-Frobenius scalars differ from one. Their product
determinant has the especially compact form
\[
          \overline{P_D(T)/P_C(T)}=1+T^4+T^{10}.
\]
Equivalently, the first scalar has reciprocal minimal polynomial
$1+T^2+T^5$ over $\mathbf F_2$, and its five conjugates each occur
twice. The determinant evaluates to one at $T=1$.
A.30 therefore proves the odd relative class number over $K$,
with the cover pointed at infinity. This certificate determines
the factors modulo two; their full characteristic-zero coefficients
are not reconstructed here.

For the smaller field degree, the second certificate uses
\[
 E=\mathbf F_2[t]/(t^{224}+t^{21}+t^7+t+1)
\]
and its subfield $L=\mathbf F_{2^{32}}=\mathbf F_{4^{16}}$.
The complete reduced pair and all $112$ multiples are in
\path{bielliptic_mixed_ell113_F4pow16.json}.
They have exact order $113$ and satisfy $\sigma_{2^{32}}(a)=16a$,
where $16\equiv2^{32}\bmod113$.
Thus the constant pointed cover descends to $L$, although its
displayed character line uses the splitting field $E$.
The independent verifier checks all $112$ coefficients, applies
the actual inverse-$q$ conjugation in A.30, and obtains an
$\mathbf F_2$ determinant with $31$ nonzero coefficients, whose
value at one is one. This proves the claimed witness over $L$.
Together with A.29 it proves exact minimality of the field degree
sixteen.

A.29 excludes three and five over every extension, and A.33
excludes seven. The eleven witness therefore proves exact minimality
of the prime eleven. To minimize its field, the two elliptic factors
have modulo-eleven eigenvalue orders five and sixty respectively:
the first polynomial factors as $(X-3)(X-5)$, and the second has
both roots of exact order sixty. These facts are verified in
\path{bielliptic_obstructions.py}.
An eleven character using only one elliptic factor is excluded by
A.27; using both requires an extension degree divisible by sixty.
Our witness attains it.

We also check transfer to the stated forms and points.
A quadratic-pole model becomes split after $x=z+s$, with
$s^2+s=\alpha$, followed by an Artin--Schreier constant change of
$y$. Both changes are available over $\mathbf F_{4^4}$ and preserve
infinity. The split quadratic twist is handled by the constant
change alone. This field embeds in both witness fields, of degrees
sixty and sixteen, so the transported witnesses retain exactly
those fields. Coefficientwise squaring transfers the certificates
from $c=\alpha$ to $c=\alpha^2$; no assertion that these two
geometric classes are $\mathbf F_4$-linearly isomorphic is needed.

The original Jacobian orders of these eight forms are $32,12,48,8$,
as given by their two point counts or the root's census. They are
prime to eleven and $113$. Hence any original point difference
$[Q-\infty]$ has order prime to the cover degree, also after transport
to the witness field. Every character of that prime degree kills
the difference. A.9 proves that the transported cover is split
over every original $Q$, with the same relative polynomials.

For minimality, a hypothetical three, five or seven witness on a
form can be extended to the composite with $\mathbf F_{4^4}$.
The extra field degree is a power of two, which preserves the
numerical condition by A.32. The isomorphism then transfers it
to the split model. Five and seven are unavailable at all original
points: their differences from infinity have orders dividing
$32,12,48$ or $8$, hence are divisible by those primes in every
relevant rational Jacobian, so their base-point twists are trivial.
For three we use the exact test A.25. The split untwisted
Jacobians have two-power order, while the three points on each
split quadratic twist are hyperelliptic ramification points;
their cubic test fails. Infinity has $\kappa=0$ on every form.
This proves all minimal-prime assertions in the claim.
\end{proof}

\begin{claim}{A.32}{Proved; safe extensions of characteristic-power degree}
Let $B/k$ be an abelian variety over a finite field of characteristic
$p$, and let $k'/k$ have degree $p^r$. Then
\[
 B(k)[p]=0\quad\Longleftrightarrow\quad B(k')[p]=0.
\]
In particular a relative-class-number $p$-unit for a fixed finite
separable cover is preserved by every constant extension of
$p$-power degree.
If two pointed curves become isomorphic over a $p$-power extension,
their existence of prime-degree witnesses over arbitrary extensions,
and the set of primes furnishing such witnesses, agree.
\end{claim}
\begin{proof}
On the finite-dimensional vector space
$V=B[p](\overline k)$ let $F$ be arithmetic Frobenius.
Rational $p$-torsion is its fixed space, and
\[
                     F^{p^r}-1=(F-1)^{p^r}.
\]
A power of an endomorphism of a finite-dimensional vector space
is invertible exactly when that endomorphism is invertible.
Thus the two fixed spaces are simultaneously zero.
Apply this to the connected new Jacobian factor to get the
relative-class-number assertion.
For pointed forms, the composite of any witness field with the
given $p$-power isomorphism field has $p$-power degree over the
witness field. The witness survives this extension and transports
across the pointed isomorphism. Apply the same argument in the
other direction.
\end{proof}

\begin{claim}{A.33}{Proved; global seventh-prime exclusion for the bielliptic classes}
For $(C_c,\infty)$ of A.28, no constant pointed degree-seven
L1 witness exists over any finite extension of $\mathbf F_4$.
The same holds for its coefficient conjugate and every pointed
form isomorphic to it at infinity over a power-of-two extension.
\end{claim}
\begin{proof}
Modulo seven the Frobenius polynomial from A.28 is
\[
                  (X-2)^2(X^2-X+4).
\]
On the first primary part write $F=2I+N$, $N^2=0$.
On the second part every nonzero vector has period twenty-four
and $F^{16}=2I$. Consequently $F^{112}=2I$ on all of $J_C[7]$.
These polynomial identities and both possible Jordan types are
verified by \path{bielliptic_obstructions.py}.

The locally exact theta divisor is Frobenius-stable and symmetric:
if $B^2\simeq\omega_C$, then
$-[R-B]=[\iota R-B]$. A nonzero theta point $a$ of order seven
would therefore force all six nonzero multiples of $a$ into theta,
because $\langle2,-1\rangle=(\mathbf Z/7)^\times$.
This is impossible by the genus-two difference bound.
For a fixed nonzero Jacobian class, there are at most two ordered
representations as $R-S$: equality of two such representations
gives linearly equivalent effective degree-two divisors; if they
are distinct, both are the unique hyperelliptic pencil.
Thus a subset of size $r$ in the intersection of a theta translate
with a cyclic group of order seven satisfies
$r(r-1)\le2(7-1)$. It cannot have $r=6$.
Theta has no nonzero seven-torsion, and the checked locally exact
differential sequence makes every geometric cyclic seventh-degree
Prym ordinary, of two-rank six.

For a nonzero quotient functional $v$ on $J_C[7]$, its period $m$
and a scalar power of Frobenius are exhausted by
\[
 \begin{array}{c|c|c}
 m&d&vF^d/v\\\hline
 3&1&2\\
 21&7&2\\
 24&16&2\\
 168&112&2.
 \end{array}
\]
The first two cases are supported on the $(X-2)$ primary part,
without or with a nonzero Jordan contribution.
The last two include the quadratic primary part.
The certificate checks all $2400$ nonzero vectors for each Jordan
type, as well as the displayed formulas.

The corresponding geometric cover has a pointed model over
$\mathbf F_{4^d}$ with deck Frobenius equal to two.
Its new two-torsion is the sum of the two nontrivial irreducible
$\mathbf F_2[\mathbf Z/7]$-modules, each a one-dimensional
$\mathbf F_8$-space. On each, this Frobenius is semilinear for
squaring, and its cube is the norm to $\mathbf F_2^\times$,
hence the identity. New two-torsion is therefore rational over
$\mathbf F_{4^{3d}}$.
Here $3d=m$ or $2m$. The first field with constant deck group is
$\mathbf F_{4^m}$, and A.32 shows that rational two-torsion cannot
first appear only in its quadratic extension.
Its relative class number is already even over that first field
and remains even over every extension.
The uniqueness of the pointed descent identifies every eligible
form with this base change. A.32 also proves the final statement
for power-of-two pointed twists.
\end{proof}


\begin{claim}{A.34}{Proved counterexample to a scalar-Frobenius reduction}
Let
\[
 C/\mathbf F_{16}:\quad y^2+y=x^5,\qquad P=\infty.
\]
Arithmetic Frobenius on \(J_C\) is the scalar \([-4]\).
Nevertheless there is a constant \(P\)-split cyclic cover of degree
thirteen over \(\mathbf F_{4096}\), its first possible constant-deck
field, whose new factor is ordinary of two-rank twelve and has odd
relative class number. More precisely its twisted polynomials are
\[
 L_P(\chi^i,T)
 =1+\frac{139-13\left(\frac{i}{13}\right)\sqrt{13}}2\,T
       +4096T^2,\qquad 1\le i\le12,
\]
after choosing the sign of \(\sqrt{13}\) consistently with the
quadratic Gauss sum. Consequently
\[
 \frac{P_D(T)}{P_C(T)}
 =\bigl(1+139T+12473T^2+569344T^3+16777216T^4\bigr)^6
\]
and
\[
 \frac{h_{\mathbf F_{4096}}(D)}{h_{\mathbf F_{4096}}(C)}
                         =17359173^6.
\]
Thus scalar Jacobian Frobenius, even at a rational hyperelliptic
ramification point, does not reduce the rational criterion to
vanishing of the geometric new \(p\)-rank.
\end{claim}
\begin{proof}
The exact counts \(\#C(\mathbf F_{16})=33\) and
\(\#C(\mathbf F_{256})=193\) give Frobenius polynomial
\((X+4)^4\). Frobenius acts semisimply, so it equals \([-4]\)
as a Jacobian endomorphism; equivalently apply
\cite[Theorem~2(d), p.~140]{AC_Tate1966}.
The following computation is independent of the assertion that all
thirteenth-degree Pryms of this curve are ordinary: it proves
ordinariness for the displayed cover itself.

Use binary polynomial notation in
\[
 K=\mathbf F_2[t]/(t^{12}+t^3+1),\qquad \#K=4096.
\]
The reduced divisor class
\[
 a=(u,v),\qquad
 u=x^2+3377x+483,\qquad v=3144x+2626
\]
has exact order thirteen. Its sixteenth-power conjugate is
\([-4]a\), and its \(4096\)-power conjugate is \(a\).
These are checked by exact Mumford arithmetic. Since
\((-4)\bmod13\) has order three and \((-4)^2=16\), all nonzero
geometric thirteenth-degree quotient characters have first
constant-deck field \(\mathbf F_{16^3}\). Also
\(4096\equiv1\bmod13\), so the displayed \(a\) has the constant-deck
descent required in A.30.

For an independent direct character count, write \(u^{13}=M\).
Lift \(v\) to the unique solution of \(w^2+w=x^5\bmod M\) with
\(w\equiv v\bmod u\). Exact linear algebra gives polynomials
\(A,B\), of degrees thirteen and ten, with \(A\) monic and
\[
 A^2+AB+B^2x^5=u^{13},\qquad A+Bw\equiv0\bmod u^{13}.
\]
Their coefficient lists, in increasing degree, are
\[
\begin{aligned}
 A={}&(1334,2010,3185,1642,2095,2911,77,\\
     &\hspace{12mm}567,182,1378,3843,297,2491,1),\\
 B={}&(1250,517,799,137,81,3560,2285,3410,2702,2433,1191).
\end{aligned}
\]
Hence \(F=A+By\) has divisor \(13D_a-26P\), where
\(D_a\) is the effective degree-two divisor representing \(a\).
The norm identity and the vanishing condition distinguish its
two zeroes from their hyperelliptic conjugates.
The cover is the normalization of \(z^{13}=F\).
It is unramified because every valuation of \(F\) is divisible
by thirteen. It is connected because \(a\ne0\), and it is
\(P\)-split because the leading local unit of \(F\) at \(P\)
is one. Its deck group is constant since \(\mu_{13}\subset K\).

Let the primitive field generator be the binary element \(3\);
use its associated primitive thirteenth root to label the
character exponents. At a point where \(F\ne0\), the exponent is
\(\log_3F\bmod13\). At a zero \(R=(r,y)\), the norm identity
shows that the local unit has character exponent
\(-\log_3B(r)\bmod13\): the other factor of \(u^{13}\) contributes
a thirteenth power. Infinity contributes exponent zero.
Enumerating all \(4353\) rational points, with these local units
at the two zeroes, gives
\[
 (N_0,\ldots,N_{12})
  =(399,323,336,323,323,336,336,336,336,323,323,336,323).
\]
Thus the nonzero quadratic residues modulo thirteen have count
\(323\), and the nonresidues have count \(336\).
The quadratic-period identity
\(\sum_{r\ {\rm square}}\zeta_{13}^r=(-1+\sqrt{13})/2\)
gives
\[
 \sum_R\chi([R-P])
 =399+323\frac{-1+\sqrt{13}}2
       +336\frac{-1-\sqrt{13}}2
 =\frac{139-13\sqrt{13}}2.
\]
A.7, with \(K_C\sim2P\), now gives every displayed
character polynomial and their product, including the exact
relative class number.

As a second check, all twelve Cartier scalars of A.30 belong to
\(\{72,73\}\subset K\), with six of each value.
Both satisfy \(X^2+X+1=0\). Thus none is zero or one, and
\[
 \overline{P_D(T)/P_C(T)}
   =(1+T+T^2)^6=1+T^2+T^6+T^{10}+T^{12}.
\]
The degree twelve proves geometric new two-rank twelve, while
the value one at \(T=1\) proves the rational two-unit condition.
The program \path{scalar_frobenius_counterexample.py} executes
both the Cartier test and the direct Kummer-point test and
requires their complete polynomials modulo two to agree.

The failure of the proposed general reduction is structural.
On its rank-one \(\mathbf F_{4096}\)-module, the nonconstant
deck Frobenius over \(\mathbf F_{16}\) is semilinear for
\(x\mapsto x^{2^8}\), since \(-4\equiv2^8\bmod13\).
Its cube is a norm to \(\mathbf F_{16}^{\times}\), which need
not be one. Hyperelliptic reflection further restricts the scalar
to \(\mathbf F_4^\times\); it still leaves precisely the two
nontrivial values that occur here. Passing from the witness field
to its cubic constant extension raises these values to one and
destroys the rational criterion, giving an explicit instance of
the unsafe-extension warning in A.9.
\end{proof}

\begin{claim}{A.35}{Proved; a simultaneous genus-two counting criterion}
Let \(C/k\) have genus two, with \(k\) a finite field of
characteristic \(p\), and let \(P\in C(k)\).
Let \(\ell\ne p\) be an odd prime such that \(p\) is a primitive
root modulo \(\ell\). Set
\[
 V=J_C(k)/\ell J_C(k),\qquad H=\ell J_C(k),\qquad N=\#C(k),
\]
and assume \(V\ne0\). If either
\[
                  |V|>N\qquad\hbox{or}\qquad N\equiv-1\bmod p,
\]
and
\[
                         (p-1)(N-1)>2(|H|-1),                 \tag{A.19}
\]
then some surjection \(J_C(k)\to\mathbf Z/\ell\) gives a
\(P\)-split L1 witness over \(k\).

One may replace (A.19) by the sharper occupancy-dependent
condition
\[
 N_0(N_0-1)+(p-1)(N-N_0)>2(|H|-1),\qquad
 N_0=\#\{R\in C(k):[R-P]\in H\}.
\]
The subgroup \(H\) is the image of multiplication by \(\ell\),
not the much larger kernel of one chosen cyclic character.
No universal assertion about the existence of such a prime and
field is implicit in this conditional criterion.
\end{claim}
\begin{proof}
Suppose every prime cyclic quotient fails the rational criterion.
For each quotient, the nontrivial character evaluations of
\(\Theta=[0]+\sum_{R\in C(k)}[[R-P]\bmod H]\) are permuted
transitively by \(p\)-power Frobenius, because \(p\) is primitive
modulo \(\ell\). Failure means at least one is zero; hence all
are zero. Every nontrivial character of \(V\) factors through
one such quotient, so Fourier inversion in \(\mathbf F_p[V]\)
gives
\[
                 \Theta=c\sum_{b\in V}[b],
                 \qquad c=(N+1)/|V|\in\mathbf F_p.
\]
The denominator is invertible since \(\ell\ne p\).
If \(|V|>N\), some nonzero coset is empty and forces \(c=0\).
The alternative congruence \(N\equiv-1\bmod p\) also forces
\(c=0\) directly. Thus
\[
              N_0\equiv-1\bmod p,\qquad
              N_b\equiv0\bmod p\quad(b\ne0).
\]

For every nonzero Jacobian class \(\delta\), there are at most
two ordered pairs of geometric points \((R,S)\) with
\([R-S]=\delta\).
Indeed \(R-S=R'-S'\) implies
\(R+S'\sim R'+S\). If these effective degree-two divisors are
different, both belong to the unique hyperelliptic pencil;
the only possible second representation is
\((R',S')=(\iota S,\iota R)\).
Consequently the number of ordered distinct pairs in the same
\(H\)-coset is bounded above by \(2(|H|-1)\).
The occupancy congruences give the lower bound
\[
 \sum_bN_b(N_b-1)
   \ge N_0(N_0-1)+(p-1)(N-N_0).
\]
Since \(P\) lies in the zero coset, \(N_0\ge1\); its congruence
improves this to \(N_0\ge p-1\). Therefore
\[
\begin{aligned}
 N_0(N_0-1)+(p-1)(N-N_0)
   &=(p-1)(N-1)+(N_0-1)(N_0-(p-1))\\
   &\ge(p-1)(N-1).
\end{aligned}
\]
Either stated strict inequality contradicts the upper bound.
Thus at least one cyclic quotient is a witness.
The primitive-root hypothesis was used at the first step and
cannot be dropped by this argument.
\end{proof}


\begin{claim}{A.36}{Proved; a Frobenius obstruction with a split coefficient algebra}
Let \(C/k\) have genus two in characteristic two, with
\(P\in C(k)\), and let \(\ell\ge5\) be an odd prime.
Write \(F\) for arithmetic Frobenius on \(J_C[\ell](\overline k)\).
Suppose that, for some \(d>0\),
\[
 F^d=sI,\qquad s\in\mathbf F_\ell^\times
       \text{ of order }\ell-1,
\]
and that every nonzero quotient functional has exact Frobenius
period \(d(\ell-1)\).
Then no \(P\)-split constant cyclic cover of degree \(\ell\)
satisfies the rational criterion over any finite extension of \(k\).
No assumption that two be primitive modulo \(\ell\) is needed.

In particular this conclusion holds if the characteristic
polynomial of \(F\) is irreducible over \(\mathbf F_\ell\)
and \(q=\#k\) generates \(\mathbf F_\ell^\times\).
\end{claim}
\begin{proof}
The locally exact theta divisor is defined over \(k\).
A nonzero point \(a\) of order \(\ell\) on it would, by the
scalar Frobenius identity, put every nonzero multiple of \(a\)
on theta. In genus two theta is a translate of the curve,
so the difference bound proved in A.35 would give
\[
                    (\ell-1)(\ell-2)\le2(\ell-1),
\]
which is impossible for \(\ell\ge5\).
Thus theta has no nonzero \(\ell\)-torsion, and every geometric
cyclic \(\ell\)-cover has ordinary new factor, of two-rank
\(\ell-1\).

For any such geometric cover, its kernel is stable under
\(F^d\), and the rational point \(P\) gives its unique pointed
descent to the degree-\(d\) extension \(k_d\).
Its geometric new two-torsion \(W\) is free of rank one over
the finite reduced algebra
\[
              R=\mathbf F_2[X]/\Phi_\ell(X).
\]
Indeed each nontrivial irreducible deck-character component
occurs with multiplicity one when the new factor is ordinary.
The algebra \(R\) may be a product of several fields.
Frobenius over \(k_d\) acts on it by \(X\mapsto X^s\)
(or its inverse with dual character conventions).
This automorphism \(\sigma\) has order \(\ell-1\) and
\[
                              R^\sigma=\mathbf F_2.
\]
To check this last equality, extend scalars to
\(\overline{\mathbf F}_2\), where \(R\) is the algebra of
functions on the nontrivial characters; the primitive scalar
permutes all of them transitively, so only constants remain.
Descent of a constant function to \(R\) requires its value to
lie in \(\mathbf F_2\).

Choose an \(R\)-basis of \(W\). Frobenius is then
\(w\mapsto a\,\sigma(w)\) with \(a\in R^\times\).
Its \((\ell-1)\)-st power is multiplication by
\[
                 a\,\sigma(a)\cdots\sigma^{\ell-2}(a)
                         \in(R^\sigma)^\times=\{1\}.
\]
Therefore all new two-torsion is rational over
\(k_{d(\ell-1)}\). By the exact period hypothesis this is the
first field with constant deck group. The unique pointed
descent identifies every eligible cover over later fields
with its base change, and its rational two-torsion persists.
This proves the exclusion.

For the last assertion, identify the irreducible
four-dimensional Frobenius module with
\(\mathbf F_{\ell^4}\), on which \(F\) multiplies by an
element \(\lambda\).
The symplectic reciprocity involution on its four conjugate
eigenvalues is \(\lambda\mapsto q/\lambda\).
It commutes with their cyclic Galois permutation, is nontrivial,
and has order two. It is therefore the \(\ell^2\)-power map,
giving
\[
                         \lambda^{\ell^2+1}=q.
\]
Since \(q\) generates \(\mathbf F_\ell^\times\),
the cyclic group \(\langle\lambda\rangle\) contains that whole
scalar subgroup. Put \(m=\operatorname{ord}(\lambda)\) and
\(d=m/(\ell-1)\). Then \(\lambda^d\) is a primitive scalar,
while every nonzero vector and quotient functional has exact
period \(m\). The first part applies.
\end{proof}

\paragraph{What the strengthened counts still leave unresolved.}
The preceding statements give a concrete obstruction to treating
the character values in L2 as independent random elements.
Whole primes may fail for representation-theoretic reasons,
before any individual cover is tested. Conversely A.34 shows
that an ordinary new factor can still be a rational witness even
when Jacobian Frobenius is scalar.

There is also a direct arithmetic limitation on the available
counting criteria. If \(k=\mathbf F_{p^e}\), \(p\) is primitive
modulo \(\ell\), and
\(r=\dim_{\mathbf F_\ell}J_C(k)/\ell J_C(k)>g\), then
\[
                         \ell\le e+1.
\]
Indeed the Frobenius-fixed subspace of \(J_C[\ell]\) is isotropic
for the Weil pairing unless \(q\equiv1\bmod\ell\), and an isotropic
subspace has dimension at most \(g\).
Thus \(r>g\) forces \(\ell\mid q-1\); primitivity gives
\(\ell-1\mid e\).
For large \(e\), a rank greater than \(g\) therefore supplies
only polynomially many cosets, whereas the curve and Jacobian
point counts grow exponentially in \(e\).
In genus two, the threshold in A.35 consequently still demands
large prime divisors in the lower-rank cases.
No existence theorem for those divisors with the required
residue order has been established here.
The exact Cartier test of A.30 gives a finite decision for
each proposed character; it does not by itself provide the
uniform torsion-point or Frobenius-orbit avoidance theorem
needed to settle A3 for every curve.


\begin{claim}{A.37}{Proved; limitation of a trigonal branch-point argument}
Let \(E/L\) be an ordinary elliptic curve in characteristic two,
and let \(f:E\to\mathbf P^1\) be a separable degree-three map
whose ramification consists of three distinct simple wild points
\(R_1,R_2,R_3\). No difference \(R_i-R_j\), \(i\ne j\), is the
nonzero geometric two-torsion class. It is nevertheless possible
for all three differences to be two-primary.

An explicit example over \(\mathbf F_4\), with
\(\alpha^2+\alpha+1=0\) and \(\beta=\alpha^2\), is
\[
 E:\ y^2+xy=x^3+1,\qquad f=(y+\alpha)/x.
\]
The ramification points are
\[
 T=(0,1),\quad R=(\beta,0),\quad R'=(\beta,\beta),
\]
and the pair-difference orders are \(8,8,4\).
The unramified residual points in their degree-three fibres are
\[
 S_\infty=\infty,\qquad S_\beta=(1,1),\qquad S_\alpha=(1,0),
\]
whose pair-difference orders are \(4,4,2\).
The quadratic resolvent is the genus-two curve
\[
 C:\quad y^2+y=z+\frac{\beta}{z^2+z}+\alpha.
\]
Its constant pointed cubic cover has
\(L(\chi,T)=1+3T+4T^2\) for both nontrivial characters and
relative class number \(64\).

For a pointed dihedral cubic cover \(D\to C\), twisting by the
reflection fixing the chosen lift leaves its elliptic quotient
unchanged. Thus this twist cannot reduce the two-primary order
of the elliptic group in the proposed argument.
These statements exhibit a limitation of that route, not a
collision on \(C\).
\end{claim}
\begin{proof}
Let \(\eta\ne0\) be the unique geometric two-torsion point of
\(E\). The locally exact differential line on an ordinary
elliptic curve is the nontrivial two-torsion line bundle
corresponding to \(\eta\).
If \(D_\infty=f^*(\infty)\), the ramification formula gives
\[
 \operatorname{div}(df)
        =2(R_1+R_2+R_3-D_\infty),\qquad
 [R_1+R_2+R_3-D_\infty]=\eta.
\]
For completeness the class of half this exact differential
cannot be zero: otherwise \(df=g^2\omega\) for a nonzero
regular differential \(\omega\); applying Cartier gives
\(0=g\,\mathcal C(\omega)\), contrary to ordinariness.
If \(R_1-R_2=\eta\), cancellation would give
\(D_\infty\sim2R_2+R_3\).
The actual fibre at \(f(R_2)\) is \(2R_2+S_2\).
Cancelling \(2R_2\) and using the injective Abel map of \(E\)
gives \(S_2=R_3\). That fibre would contain both ramification
points with multiplicity at least two, contradicting degree
three. This proves the order-two assertion.

For the example, direct enumeration gives
\(E(\mathbf F_4)\simeq\mathbf Z/8\).
The pole divisor of \(f\) is \(2T+\infty\), so its degree is
three, and with \(\omega=dx/x\) one has
\[
                        df=\frac{x^2+\alpha}{x}\,\omega.
\]
The divisor of this differential and the pole divisor give
different exponent two at exactly \(T,R,R'\).
Their branch values are respectively \(\infty,\beta,\alpha\).
The stated residual points follow by evaluating \(f\);
exact elliptic addition gives the six difference orders in
the claim. In particular the distinction between order two
and a higher two-power is essential.

Writing \(f=t\), so \(y=tx+\alpha\), gives the cubic
\[
                x^3+(t^2+t)x^2+\alpha x+\alpha=0.
\]
For a separable cubic \(X^3+aX^2+bX+c\) in characteristic two,
its quadratic resolvent has Artin--Schreier class
\((b^3+a^3c+c^2)/(ab+c)^2\).
Here this simplifies to
\[
              \beta(t^2+t)+\frac{\beta}{t^2+t+1}.
\]
The change \(w\mapsto w+\alpha t\), followed by \(t=z+\alpha\),
gives the displayed \(C\).
The Galois closure has group \(S_3\), since the cubic is
irreducible over \(L(t)\) by its degree-three function-field
map and the nontrivial quadratic resolvent is connected.
All three branch inertia groups have order two, so the
degree-three map from this closure to \(C\) is unramified.
The rational residual points give rational lifts of the three
branch points of \(C\), making this constant cubic cover
pointed. Alternatively its unique pointed character is
identified directly in \(J_C(\mathbf F_4)\), which has order
twelve. The three rational points of \(C\) are its branch
points and have character occupancies \((3,0,0)\).
A.7 therefore gives \(1+3T+4T^2\) and relative value
\(8^2=64\). All the finite-field, group, resolvent and
character calculations are executed in
\path{trigonal_two_primary_obstruction.py}.

Finally, if \(D\to C\) has deck group \(\langle\delta\rangle\)
of order three and \(\tau\) is its pointed reflection,
then a \(\tau\)-fixed lift over a branch point maps to the
unramified residual point \(S_i\) of
\(E=D/\langle\tau\rangle\to\mathbf P^1\), not to \(R_i\).
Indeed the total inertia has order two and is already
accounted for in \(D\to E\). The fibre relation is
\[
                     2R_i+S_i\sim D_\infty,\qquad
                     S_i-S_j=-2(R_i-R_j).
\]
Twisting \(D\) by \(\tau\) does not twist \(E\), because
\(\tau\) acts trivially on its quotient. Every other reflection
\(\delta^j\tau\) is conjugate to \(\tau\) by a rational deck
element, so it gives the same elliptic quotient up to
\(L\)-isomorphism when the deck group is constant.
A nontrivial pure deck twist destroys the rational lift over
the prescribed branch point. Thus neither operation supplies
the missing small-two-primary-group hypothesis.
\end{proof}


\begin{claim}{A.38}{Proved partial reformulation; a finite etale torsor in the scalar case}
Let \(C/\mathbf F_q\) be a genus-two curve in characteristic
two, with a rational base point and the corresponding point-rigidified
Poincare family of degree-zero line bundles, defined over the base field. Suppose arithmetic Frobenius
on \(J_C\) is the scalar \([s]\), where \(s^2=q\).
On a nonempty rational open of \(J_C\), choose a rational basis
\(e_a\) of its one-dimensional nontrivial character cohomology,
and write the \(q\)-power coefficient as
\[
                     F_{\rm abs}^{\log_2q}e_a
                              =A(a)e_{[q]a}.
\]
For the models of A.30, \(A\) is given explicitly by its
finite Cartier--Mumford product.
On the open \(U\) where this basis and \(A\) are invertible,
\[
                         \mathcal Y:\quad z^{q-1}=A
\]
is a finite etale \(\mu_{q-1}\)-torsor.
Its isomorphism class is independent of the basis on common
domains of definition.

If \(a\in U(\mathbf F_{q^m})\), the linear coherent return
scalar for the pointed character defined by \(a\) is
\[
                    \lambda(a)=
                    N_{\mathbf F_{q^m}/\mathbf F_q} A(a).
\]
This is exactly the arithmetic Frobenius value of
\(\mathcal Y\) at \(a\), with degree \(m\).
If \(a\) has odd prime order \(\ell\) and every nonzero multiple
of \(a\) lies in \(U\), its pointed cyclic cover over
\(\mathbf F_{q^m}\) satisfies (R) exactly when all these
Frobenius values differ from one.

The torsor furnished by \(C:y^2+y=x^5\) over \(\mathbf F_{16}\)
is geometrically nontrivial. This construction supplies a
finite etale object detecting the single-character return
condition, but does not prove that its nontrivial Frobenius
locus contains a whole nonzero prime cyclic subgroup.
\end{claim}
\begin{proof}
The scalar hypothesis gives the identity of Jacobian
endomorphisms
\[
                             [q]=F_J^2.
\]
Since the chosen rational functions are defined over
\(\mathbf F_q\), one has \(A(F_Ja)=A(a)^q\).
For a point \(a\) rational over \(\mathbf F_{q^m}\),
iterating the absolute Frobenius formula gives
\[
\begin{aligned}
 A_m(a)
  &=\prod_{j=0}^{m-1}
      A([q]^j a)^{q^{m-1-j}}\\
  &=A(a)^{\,q^{m-1}(1+q+\cdots+q^{m-1})}\\
  &=N_{\mathbf F_{q^m}/\mathbf F_q}A(a).
\end{aligned}
\]
The last expression lies in \(\mathbf F_q\), so the inverse
coefficient-field conjugation in A.30 does not alter it.
Also \(F_J^m(a)=a\) implies \([q^m]a=a\).
If \(a\) has order \(n\), then \(n\mid s^m-1\), so
\(\mu_n\subset\mathbf F_{q^m}\).
The point-normalized line bundle therefore defines the
constant pointed cyclic cover over this same field.

For a different rational basis \(e'_a=b(a)e_a\), the coefficient
changes to
\[
 A'(a)=A(a)\frac{b(a)^q}{b([q]a)}
       =A(a)b(a)^{q-q^2}
       =A(a)\bigl(b(a)^{-q}\bigr)^{q-1}.
\]
This is precisely a change of Kummer generator, proving
basis independence. The relevant open is defined over
\(\mathbf F_q\) and is stable under \(F_J\), whose underlying
map is Frobenius; thus the return computation does not
require removing a new infinite collection of divisors.

Over a closed point with residue field \(\mathbf F_{q^m}\),
arithmetic Frobenius sends a root \(z\) to \(z^{q^m}\), so
its ratio is
\[
             z^{q^m-1}
              =A(a)^{(q^m-1)/(q-1)}
              =N_{\mathbf F_{q^m}/\mathbf F_q}A(a).
\]
This proves the exact comparison. For a prime-order
point, all nonzero multiples have the same field of
definition: multiplication by any \(1\le i<\ell\)
has an inverse on that cyclic subgroup.
Applying A.20 or A.30 to all its nontrivial characters
gives the stated criterion.

For the example of A.34, \(q=16\), and the degree-three
closed point \(a\) has return value a nontrivial cube
root in \(\mathbf F_4^\times\).
The program now independently checks that each
\(4096\)-power scalar is the \(4096/16\) norm of its
\(16\)-power Cartier coefficient.
If the \(\mu_{15}\)-torsor were geometrically trivial,
its arithmetic Frobenius at a degree-three point would
be the cube of an element of \(\mu_{15}\), hence have
order dividing five. The observed order three is
impossible. Thus the torsor has nontrivial geometric
monodromy.

For general nonscalar Jacobian Frobenius the corresponding
periodicity condition is \(V_q^m(a)=a\), since on prime-to-two
torsion \(\sigma_{q^m}(a)=q^m a\) is equivalent to that
condition for the \(q\)-Verschiebung.
The simplification \([q]=F_J^2\), essential to the norm
and the Kummer gauge identity above, is then unavailable.
Even in the scalar case, ordinary Chebotarev applied to
\(U\) controls closed points of \(U\); it does not impose
prime order on \(a\) and simultaneous avoidance of one
on every nonzero multiple. Those are the precise additional
assertions still needed for A3.
\end{proof}


\begin{claim}{A.39}{Proved: an exact nonabelian bounded-exponent example}
Let \(L=\mathbf F_{4096}\), let \(D/L\) be the Fermat curve
\[
                         X^{39}+Y^{39}+Z^{39}=0,
\]
and choose a primitive thirteenth root \(\zeta\in L\).
The automorphisms
\[
 \sigma(X:Y:Z)=(\zeta X:\zeta^3Y:\zeta^9Z),
 \qquad \tau(X:Y:Z)=(Y:Z:X)
\]
generate a constant group \(G\simeq C_{13}\rtimes C_3\)
acting freely on \(D\). For \(C=D/G\), one has
\[
                  g(D)=703,\qquad g(C)=19.
\]
For every odd positive integer \(m\),
\[
 \begin{aligned}
 J_D(\mathbf F_{4096^m})[2^\infty]
                  &\simeq(\mathbf Z/2)^{24},\\
 J_C(\mathbf F_{4096^m})[2^\infty]
                  &\simeq(\mathbf Z/2)^4.
 \end{aligned}
\]
Thus this cover has relative class-number valuation
\[
             v_2\!\left(\frac{h(D)}{h(C)}\right)=20
\]
over every such odd extension: it fails (R).
Nevertheless every \(P\in C(\mathbf F_{4096^m})\), with
\(m\) odd, has a rational lift to this same cover after
an additional extension of degree in \(\{1,3,13\}\),
and the bounded-exponent nonabelian criterion applies there,
since \(39>(2-1)^2\).

For an even extension degree \(m\), the exact exponent of
\(J_D(\mathbf F_{4096^m})[2^\infty]\) is
\(2^{v_2(m)+2}\). In particular it is eight for \(m=2\);
the numerical inequality for this order-39 cover then fails.
The example distinguishes two sufficient tests on a specified
cover; it does not assert that \(C\) has no other (R) witness.
\end{claim}
\begin{proof}
The relation is \(\tau\sigma\tau^{-1}=\sigma^3\).
Every nonidentity element of \(\langle\sigma\rangle\)
has three distinct projective eigenvalues, and its fixed
points are the coordinate vertices, none on \(D\).
Every remaining nonidentity element of \(G\) has order
three and is conjugate to \(\tau\) or \(\tau^2\).
Their projective fixed points have coordinates
\((1:\omega:\omega^2)\), with \(\omega^3=1\),
and at such a point the Fermat expression is
\(1+1+1=1\). Thus the action is free.
The plane genus formula and the etale Riemann--Hurwitz
formula give the two stated genera.

Here are the complete spectral data, so that the claim
about the integral two-primary exponent is independently
checkable. Put \(q=4096\), choose a multiplicative
character \(\chi:L^\times\to\mu_{39}\), and set
\[
 J(a,b)=\sum_{x\in L\setminus\{0,1\}}
                   \chi(x)^a\chi(1-x)^b,
 \qquad \alpha_{a,b}=-J(a,b).
\]
The indices are \(1\le a,b\le38\), \(a+b\not\equiv0\pmod{39}\).
For the diagonal group \(\mu_{39}^2\), each such character
occurs exactly once in \(H^1\), and its Frobenius scalar is
\(\alpha_{a,b}\). This assertion, including its sign and its
validity for both prime-to-two etale and crystalline cohomology,
is exactly \cite[Corollary~2.4, p.~181]{AC_Katz1981}.
The integral crystalline version is explained on p.~183;
Proposition~4.2 on pp.~185--186 also transfers the character
polynomial identities to the Jacobian endomorphism ring.

The accompanying program evaluates all \(1406\) sums in
\(\mathbf Z[\zeta_{39}]\), using the binary field
\[
              L=\mathbf F_2[t]/(t^{12}+t^3+1)
\]
and the primitive multiplicative generator encoded by the
integer \(3\). It checks each equality
\(\alpha\bar\alpha=4096\), each invariance under
\(\zeta_{39}\mapsto\zeta_{39}^2\), and the following table.
The conductor is \(39/\gcd(a,b,39)\); the second column
is \(\alpha+\bar\alpha\), including twice the scalar in
the rational cases.
\[
\begin{array}{r|r|r}
 \text{conductor}&\alpha+\bar\alpha&\text{character pairs}\\ \hline
 3&128&2\\
 13&-128&132\\
 39&-128&360\\
 39&-80&216\\
 39&-28&72\\
 39&28&432\\
 39&80&144\\
 39&115&48
\end{array}
\]
In particular the possible trace \(-115\) does \emph{not}
occur. As an independent sign check, direct enumeration
of projective points gives \(\#D(L)=30537\), and the
sum of the displayed eigenvalues is \(-26440=q+1-30537\).
Consequently the zeta numerator is
\[
\begin{aligned}
 P_D(T)={}&(1-64T)^2(1+64T)^{492}\\
 &\cdot(1+80T+4096T^2)^{108}
        (1+28T+4096T^2)^{36}\\
 &\cdot(1-28T+4096T^2)^{216}
        (1-80T+4096T^2)^{72}\\
 &\cdot(1-115T+4096T^2)^{24}.
\end{aligned}
\]

It is essential to prove an exponent bound on the actual
Tate lattice, rather than infer one from an unspecified
isogeny decomposition. Let \(F\) denote the \(L\)-Frobenius
endomorphism of \(J_D\). It is central in the semisimple
algebra \(\operatorname{End}_L^0(J_D)\); hence its minimal
polynomial over \(\mathbf Q\) is square-free. By the
faithful prime-to-two Tate realization, its roots occur
among the reciprocal roots of the displayed zeta numerator. The action
of \(F\) on \(V_2^{\rm et}(J_D)\) remains semisimple,
and all its eigenvalues are two-adic units, since \(F\)
is an automorphism of its integral etale Tate lattice.

There is only one unit root among the corresponding
Frobenius polynomials: the root \(\beta\in\mathbf Z_2^\times\) of
\[
                       X^2-115X+4096=0.
\]
The other root \(\gamma\) has valuation twelve.
For the traces \(\pm80\) and \(\pm28\), both roots
have positive valuation, as do the rational roots
\(\pm64\). Thus semisimplicity gives
\[
               F=\beta\,\operatorname{id}
               \quad\text{on }V_2^{\rm et}(J_D),
\]
and therefore also on the actual integral Tate lattice.
This argument excludes any hidden lattice extension.
Since
\[
 (1-\beta)(1-\gamma)=4097-115=3982=2\cdot1991,
\]
one has \(v_2(\beta-1)=1\). The rational two-primary
subgroup is therefore elementary abelian. Its order is
\(2^{24}\), directly from \(P_D(1)\), proving its exact
structure without needing to infer it merely from
Newton slopes.

For completeness, the quotient calculation is also made
at the level of individual characters. Write
\(c=-a-b\pmod{39}\). The \(\sigma\)-invariant characters
are exactly those with
\[
                    a+3b+9c\equiv5a+7b\equiv0\pmod{13}.
\]
There are \(110\) of them. Under the cycle
\((a,b,c)\mapsto(b,c,a)\), two are fixed,
namely \((13,13,13)\) and \((26,26,26)\), and the
other \(108\) form \(36\) orbits of length three.
Each length-three orbit contributes one invariant line,
with its common Jacobi eigenvalue. The two fixed lines
factor through the Fermat cubic. The coordinate cycle
acts freely on that genus-one curve and hence by a
translation, so it acts trivially on these two lines.
The invariant eigenvalues are \(64\) twice, \(-64\)
twenty-eight times, and the two roots of
\(X^2-115X+4096\) four times each. Thus
\[
        P_C(T)=(1-64T)^2(1+64T)^{28}
                           (1-115T+4096T^2)^4.
\]
The same scalar argument, or the pullback embedding
on two-adic Tate modules for the odd-degree cover,
gives the elementary abelian group of order \(2^4\).
In particular \(f_2(D)=24\), \(f_2(C)=4\), and
\(v_2(P_D(1)/P_C(1))=24-4=20\).

The field claims now follow with no change of cover.
The root \(\beta\) is congruent to \(115\pmod{4096}\),
since \(\beta=115-\gamma\). Therefore
\[
                   v_2(\beta-1)=1,\qquad
                   v_2(\beta+1)=2.
\]
For odd \(m\), the factor
\(1+\beta+\cdots+\beta^{m-1}\) is a unit, so
\(v_2(\beta^m-1)=1\). For even \(m\), the elementary
two-adic lifting formula gives
\[
                  v_2(\beta^m-1)=v_2(m)+2.
\]
As \(F\) is scalar on the lattice, these are exact
exponents of the finite rational groups, not only bounds.

For any rational point of \(C\) over an odd-degree
extension of \(L\), its geometric fibre is a \(G\)-torsor.
Frobenius acts there as multiplication by some element
of \(G\), up to conjugacy; its order is one, three or
thirteen. Thus an additional extension of that odd
degree makes the fibre split. The exponent stays two.
The nonabelian annihilator argument of A.53, with
\(a=2\cdot1\in\mathbf Z[G]\), now excludes a collision
because its degree condition is \(39>(2-1)^2\).
For even extension degree two the exponent is eight
and \(39\le(8-1)^2=49\), which explains the stated
limitation of this particular numerical test.
\end{proof}

\paragraph{Scope of the example.}
There is a separate geometric reason that this particular
quotient does not provide a new unresolved instance of A1.
The two power maps from the Fermat curve of degree \(39\)
to the Fermat curves of degrees three and thirteen are
jointly injective on geometric points: in any common
nonzero coordinate chart, equality of both images forces
ratios whose third and thirteenth powers are both one,
hence forces equal ratios. Both target Jacobians have
two-rank zero, as is already shown by the conductor-three
and conductor-thirteen entries of the Jacobi table.
A two-primary point difference on \(D\) therefore maps
to zero in both Jacobians; the Abel--Jacobi embeddings
and joint injectivity force the two points to coincide.
The collision-lifting statement A.4 transfers this
exclusion to \(C\), over every finite extension. The
point of A.39 is the exact nonabelian exponent test,
including a cover for which (R) fails; it is not a claim
that all-extension injectivity of this example required
that numerical test.

\begin{claim}{A.40}{Proved: a free order-117 group and its precise limitation}
On the Fermat curve \(D/\mathbf F_{4096}\) of A.39,
choose a primitive thirty-ninth root \(\xi\) and put
\[
 s(X:Y:Z)=(X:\xi Y:\xi^{17}Z),
 \qquad \tau(X:Y:Z)=(Y:Z:X).
\]
The constant nonabelian group
\[
                 G'=\langle s,\tau\rangle
                    \simeq C_{39}\rtimes_{16}C_3
\]
has order \(117\) and acts freely. Its quotient \(C_7=D/G'\)
has genus seven and zeta numerator
\[
 P_{C_7}(T)=(1-64T)^2(1+64T)^4
                         (1-115T+4096T^2)^4.
\]
Thus \(J_{C_7}(\mathbf F_{4096})[2^\infty]
\simeq(\mathbf Z/2)^4\), and \(D/C_7\) again has relative
class-number valuation twenty.
The bounded-exponent criterion for this cover applies
after splitting the pointed fibre over every extension
of \(\mathbf F_{4096}\) of odd or twice-odd degree:
the relevant exponents are respectively two and eight,
and \(117>(8-1)^2\).

There is also a more direct (R) witness for this particular
quotient. The subgroup \(\langle s^3,\tau\rangle\) is
normal of index three and is the order-39 group of A.39,
after choosing its generator of order thirteen.
Writing \(C_{19}=D/\langle s^3,\tau\rangle\), the cover
\(C_{19}\to C_7\) is etale cyclic of degree three, with
relative zeta numerator and relative class number
\[
                  (1+64T)^{24},\qquad 65^{24}.
\]
Its new two-rank is zero, so (R) persists over every
extension, and every prescribed base point becomes
split after an additional extension of degree at most
three. Accordingly, the smaller genus of this quotient
does not enlarge the set of curves reached by (R).
\end{claim}
\begin{proof}
Projectively,
\[
       \tau s\tau^{-1}=s^{16},\qquad
       16^3\equiv1\pmod{39},\qquad
       1+16+16^2\equiv0\pmod{39}.
\]
The diagonal generator has order \(39\), and no nonidentity
power of \(\tau\) is diagonal, giving the stated order
and semidirect product. The pairwise differences of the
three diagonal exponents \(0,1,17\) are \(1,16,17\),
all units modulo \(39\). Hence a nonidentity diagonal
element has distinct eigenvalues and fixes only the
coordinate vertices.

For an element \(s^a\tau^b\), with \(b=1\) or \(2\),
use its displayed monomial matrix. Its cube is
\(\xi^{18a}I\). If \(\lambda\) is an eigenvalue,
\[
              \lambda^3=\xi^{18a},\qquad
              \lambda^{39}=\xi^{234a}=1.
\]
All eigenvector coordinates are nonzero; their ratios
are products of powers of \(\lambda\) and \(\xi\),
so their thirty-ninth powers are equal. The Fermat
expression at an eigenpoint is therefore a nonzero
scalar times \(1+1+1=1\). There are no fixed points.
Riemann--Hurwitz gives
\(g(C_7)=1+(703-1)/117=7\).

For the Jacobi character indexed by \((a,b,c)\),
\(c=-a-b\), the diagonal invariance condition is
\[
                   b+17c\equiv22a+23b\equiv0\pmod{39}.
\]
The augmented exact program finds \(38\) characters:
the same two cycle-fixed characters as in A.39 and
twelve cycles of length three. The quotient eigenvalues
are \(64\) twice, \(-64\) four times, and the two
roots of \(X^2-115X+4096\) four times each.
This proves the displayed numerator. The actual
integral exponent follows from the scalar unit-root
argument of A.39.

Every element of \(G'\) has order in
\(\{1,3,13,39\}\). A rational fibre therefore splits
after an odd extension of one of those degrees.
Such an extension does not change the two-adic valuation
of the original constant-field degree. Hence the exact
exponent remains two for odd degrees and eight for
twice-odd degrees, as stated.

Finally, the map \(G'\to C_3\) sending \(s\) to a
generator and \(\tau\) to zero is well-defined because
\(16\equiv1\pmod3\). Its kernel is
\(\langle s^3,\tau\rangle\); the order-thirteen
generator \(s^6\), after multiplying coordinates by
a common scalar, is the diagonal generator in A.39.
Dividing the two explicitly proved zeta numerators gives
\((1+64T)^{24}\). Its roots all have positive two-adic
valuation; alternatively its value over every extension
is an odd power of \(1-(-64)^m\). Thus (R) persists
under the extension used to split any given fibre.
\end{proof}

\begin{claim}{A.41}{Proved: when scalar Frobenius forces the arithmetic test to be geometric}
Let \(C/k\) be a genus-two curve in characteristic two,
let \(P\in C(k)\), and suppose Jacobian Frobenius is the
scalar \([s]\), where \(s=\varepsilon2^r\),
\(\varepsilon\in\{1,-1\}\), and \(\#k=s^2\).
Let \(\ell\) be an odd prime and put
\[
 f=\operatorname{ord}_\ell(2),\qquad
 m=\operatorname{ord}_\ell(s),\qquad
 d=\frac{f}{\gcd(f,m)}.
\]
Every geometric cyclic degree-\(\ell\) cover has a
\(P\)-pointed descent to \(k\), possibly with nonconstant
deck group, and its first constant-deck field has degree
\(m\). Suppose
\[
                  d=1,
        \quad\text{equivalently}\quad
          \langle2\rangle\subseteq\langle s\rangle
                         \text{ in }\mathbf F_\ell^\times.
\]
Then all its geometric new two-torsion points are rational over
that first constant-deck field. In particular this
cover satisfies (R) over a field with constant deck
group if and only if its geometric new two-rank is zero.
This equivalence holds individually for every cyclic
cover, and persists over all further constant fields.
No assumption that two is primitive modulo \(\ell\)
is required.

The condition has the following explicit forms:
\[
\begin{array}{c|c}
 s&\text{condition for }d=1\\ \hline
 +2^r&\gcd(r,f)=1\\
 -2^r,\ f\text{ even}&\gcd(r+f/2,f)=1\\
 -2^r,\ f\text{ odd}&\gcd(r,f)=1.
\end{array}
\]
When \(d>1\), the same argument yields a norm in
\(\mathbf F_{2^d}^\times\), which need not be one.
It therefore gives no converse obstruction in that case.
\end{claim}
\begin{proof}
The scalar action preserves every geometric degree-\(\ell\)
cyclic quotient of the prime-to-two geometric fundamental
group. The rational point supplies its pointed descent,
with arithmetic Frobenius acting on the deck quotient
by the scalar \(s\), up to the harmless simultaneous
inverse convention. Thus the deck group first becomes
constant after degree \(m\).

Consider the new part of geometric etale cohomology
\[
 W=H^1_{\rm et}(\overline D,\mathbf F_2)_{\rm new},
 \qquad R=\mathbf F_2[X]/\Phi_\ell(X).
\]
The averaging idempotent is defined because \(\ell\)
is odd, so \(W\) is an \(R\)-module.
The Artin--Schreier comparison identifies its scalar
extension with the Frobenius-stable part of the
corresponding coherent cohomology. Every nontrivial
character of the geometric cyclic cover occurs in
\(H^1(\overline D,\mathcal O_{\overline D})\) with
dimension \(g(C)-1=1\). Consequently \(W\) has rank
zero or one over each simple field factor of \(R\).
This is the only step using genus two; compare the
functorial stable-part comparison in A.20.

Arithmetic Frobenius \(F\) is an automorphism of \(W\)
and is semilinear for
\(\sigma_s:R\to R\), \(X\mapsto X^s\).
If \(W\ne0\), its support is a union of \(\sigma_s\)-orbits
of simple factors. Thus \(W\) is free of rank one over a
\(\sigma_s\)-stable direct factor \(R'\) of \(R\).
Choose a basis on that factor; then
\[
                 F=c\sigma_s\quad(c\in(R')^\times),
 \qquad
                 F^m=\prod_{j=0}^{m-1}\sigma_s^j(c).
\]
The last product belongs to \(((R')^{\sigma_s})^\times\).

Here is the exact fixed-ring calculation. The primitive
\(\ell\)-th roots are indexed by
\(\mathbf F_\ell^\times\); their two-Frobenius orbits
have length \(f\), while multiplication by \(s\) has
order \(m\). As both acting subgroups lie in the same
cyclic group, their intersection has order \(\gcd(f,m)\).
It follows that
\[
 R^{\sigma_s}\simeq
 \left(\mathbf F_{2^{\,f/\gcd(f,m)}}\right)^{
                    (\ell-1)/\operatorname{lcm}(f,m)}.
\]
For \(d=1\) its unit group, and that of every direct
factor such as \((R')^{\sigma_s}\), is trivial.
Therefore \(F^m=1\) on \(W\).
The same holds on the dual geometric new two-torsion
of the Jacobian. If \(W\ne0\), criterion (R) fails at
the first constant field and at every later one. If
\(W=0\), the geometric new two-primary group is already
zero, so (R) holds over every extension. This proves the
equivalence, including its field quantifier.

For \(s=2^r\), one has \(m=f/\gcd(r,f)\), giving
\(d=\gcd(r,f)\). If \(s=-2^r\) and \(f\) is even,
then \(-1=2^{f/2}\pmod\ell\), so replace \(r\) by
\(r+f/2\). If \(f\) is odd, inversion is outside
\(\langle2\rangle\), and
\(m=2f/\gcd(r,f)\); because \(f\) is odd this gives
\(d=\gcd(r,f)\) again.
\end{proof}

\paragraph{Relation to the ordinary-new counterexample.}
For \(C:y^2+y=x^5\) over \(\mathbf F_{16}\), one has
\(s=-4\) and \(r=2\). At \(\ell=7\), \(f=3\) and
\(d=1\); at \(\ell=11\), \(f=10\) and
\(d=\gcd(2+5,10)=1\). For those primes the arithmetic
criterion cannot improve on the geometric new-rank-zero
criterion. At \(\ell=13\), however, \(f=12\) and
\(d=\gcd(2+6,12)=4\). The semilinear norm lies in
\(\mathbf F_{16}^\times\), and A.34 exhibits the
nontrivial values of order three which give an ordinary-new
rational (R) witness. Thus A.41 explains exactly why
the unrestricted scalar-Frobenius assertion is false,
while recovering a substantial range in which it is true.

\begin{claim}{A.42}{Proved: exact finite-level meaning of hereditary relative tests}
Let \(G\) be a finite solvable group of order prime to \(p\),
and let \(V\) be a finite \(p\)-primary abelian group with a
\(G\)-action. For a subgroup \(H\), write
\[
                 e_H=\frac1{|H|}\sum_{h\in H}h
                      \quad\text{on }V.
\]
The following conditions are equivalent:
\begin{enumerate}
\item For every pair of normal subgroups \(H<K\) of \(G\),
      one has \(V^H\supsetneq V^K\).
\item There is \(v\in V\) such that
      \(\operatorname{Stab}_G(e_Hv)=H\) for every normal
      subgroup \(H\) of \(G\).
\end{enumerate}

In particular, let \(D/k\to C/k\) be a constant solvable
prime-to-\(p\) Galois etale cover, with a rational chosen
point over \(P\), and put \(V=J_D(k)[p^\infty]\).
All relative class-number tests for
\[
                    D/H\longrightarrow D/K,
                    \qquad H<K\text{ normal in }G,
\]
fail if and only if there is a rational \(p\)-primary
class on \(D\) whose norms have free deck orbits on
every normal quotient of this particular cover.
This is a finite-level statement about classes, with
no assertion that they are differences of points.
\end{claim}
\begin{proof}
Necessity is immediate: if \(V^H=V^K\), then \(e_Hv\)
is \(K\)-fixed, so cannot have stabilizer exactly \(H\).
For sufficiency, induct on \(|G|\). The trivial group
causes no difficulty. Choose a nontrivial abelian normal
subgroup \(N\) of the solvable group \(G\), and split
\[
                  V=U\oplus W,\qquad
                  U=e_NV=V^N,\quad W=(1-e_N)V.
\]
Both summands are \(G\)-stable. The strict-invariant
condition descends to the \(G/N\)-module \(U\), so induction
gives \(u\in U\) with the required property for all normal
quotients of \(G/N\). For any normal \(H\) of \(G\),
\(e_Hu=e_{HN}u\), whence
\[
                     \operatorname{Stab}_G(e_Hu)=HN.
\]

We choose \(w\in W\) simultaneously for all the normal
subgroups \(H\). The algebra \(\mathbf Z_p[N]\) is finite
etale and commutative; it is a product of unramified
discrete valuation rings, and \(N\) acts through a
prime-to-\(p\) character on each factor. The projectors
\(e_H\) are central in \(\mathbf Z_p[G]\). They commute
with this algebra and with each other, so decompose
\(W\) further into their finitely many joint image/kernel
summands. Choose one nonzero element in every nonzero
joint summand and add these elements to obtain \(w\).

If a prime-to-\(p\) root of unity is not one, its
difference from one is a unit in the relevant
unramified valuation ring. Thus every nonzero element
of a character summand has precisely that character's
kernel as its stabilizer in \(N\), even though the
module need not be free over that ring. Our choice
therefore ensures, simultaneously for all normal \(H\),
\[
 \operatorname{Stab}_N(e_Hw)
       =\ker\bigl(N\longrightarrow
                         \operatorname{Aut}(W^H)\bigr).
\]
Denote this kernel by \(K_H\). It is normal in \(G\),
because \(N\) and \(H\) are normal and \(W^H\) is
\(G\)-stable. It contains \(H\cap N\). If this
containment were strict, then \(HK_H>H\); both would
be normal in \(G\), while \(K_H\) acts trivially on
both \(U^H\) and \(W^H\). This would give
\(V^H=V^{HK_H}\), contrary to the hypothesis.
Consequently \(K_H=H\cap N\).

Set \(v=u+w\). Since \(U\) and \(W\) are \(G\)-stable,
an element stabilizing \(e_Hv\) must stabilize each
component separately. From the first component it
belongs to \(HN\). Modulo \(H\) it is therefore
represented by some \(n\in N\); from the second
component this \(n\) lies in \(H\cap N\).
The stabilizer is exactly \(H\), as required.

For the geometric application, the prime-to-\(p\)
pullback and norm identities identify
\(J_{D/H}(k)[p^\infty]\) with \(V^H\). Therefore
\[
 v_p\!\left(\frac{h(D/H)}{h(D/K)}\right)>0
                  \quad\Longleftrightarrow\quad
                  V^H\supsetneq V^K.
\]
The pullback of the actual norm of a class \(v\)
is \(|H|e_Hv\). Multiplication by \(|H|\) is
invertible on \(V\), so it does not change its
stabilizer. The abstract equivalence gives exactly
the asserted finite-level application.
\end{proof}

\paragraph{Why tests only over the original base are weaker.}
Let \(p=2\), \(G=C_7\rtimes C_3\), and let \(V\) be
the additive group of \(\mathbf F_4\), with the action
inflated from \(G^{\rm ab}=C_3\) by multiplication by
a primitive cube root. The proper normal subgroups
are \(1\) and \(C_7\), and
\[
                     V^G=0,\qquad
                     V^1=V^{C_7}=V\ne0.
\]
Thus every nontrivial normal quotient over the original
base has a nonzero new module, and all norm maps have
their usual prime-to-two splittings. Nevertheless every
vector is fixed by \(C_7\), so no free orbit exists.
The omitted relative test is exposed by
\(V^1=V^{C_7}\). This is an algebraic countermodel to
the inference from those module conditions alone;
it is not asserted to be the full Jacobian module of
a specified curve or a counterexample to A1 or A10.

\paragraph{Remaining inverse-system issue.}
A.42 treats a finite cover with constant deck group.
At a fixed finite field, such covers need not be
cofinal in the geometric prime-to-\(p\) tower.
For nonconstant deck groups, the rational subgroup
\(J_D(k)[p^\infty]\) need not be \(G\)-stable, so the
commuting-projector argument cannot simply be repeated
on that subgroup. Making the deck group constant
produces regular classes over an extension, but does
not make them fixed by the original Frobenius.
Moreover regularity on a cofinal selection of levels
does not imply regularity after norms to every
geometric normal quotient. A semilinear descent and
compatibility statement resolving these two issues
is still required before A.42 could imply a
Frobenius-fixed element of the full ghost module.

\begin{claim}{A.43}{Proved: semilinear simultaneous regularity from all relative tests}
Let \(G\) be a finite solvable group of order prime to \(p\),
let \(V\) be a finite-dimensional \(\mathbf F_p\)-vector space
with a \(G\)-action, and let \(F\in\operatorname{GL}(V)\) and
\(\phi\in\operatorname{Aut}(G)\) satisfy
\[
                         F(gv)=\phi(g)F(v).
\]
Assume that for every positive integer \(r\) and every
\(\phi^r\)-stable pair of subgroups \(A\triangleleft B\)
with \(B\triangleleft G\) and \(B/A\) cyclic of prime order, one has
\[
                  (V^A)^{F^r}\supsetneq(V^B)^{F^r}.
                  \tag{A.43.1}
\]
Then there is \(v\in V^F\) such that
\[
                    \operatorname{Stab}_G(e_Hv)=H
                    \quad\text{for every normal }H\triangleleft G.
                    \tag{A.43.2}
\]
The subgroups \(H\) in this conclusion need not be
\(\phi\)-stable.

For a pointed geometrically Galois odd prime-to-\(p\)
etale cover \(D/k\to C/k\), take
\(V=J_D[p](\overline k)\), with its actual arithmetic
Frobenius. Failure of every relative cyclic-prime
class-number test over intermediate pointed bases which are
geometrically normal over \(C\), and over every finite
constant extension, implies (A.43.1).
Consequently it supplies the class in (A.43.2).
The rational class need not be a point difference.
\end{claim}
\begin{proof}
We induct on \(|G|\). Choose a minimal nontrivial subgroup
\(N\subseteq G\) which is normal in \(G\) and stable under
\(\phi\). It is elementary abelian of order a power of
some prime \(\ell\ne p\). Indeed its derived subgroup,
its Sylow subgroups, and its subgroup of \(\ell\)-th
powers would otherwise give smaller characteristic
subgroups; solvability rules out the nonabelian
characteristically simple case.

The averaging projector gives a \(G,F\)-stable splitting
\[
                         V=U\oplus W,\qquad
                         U=V^N,\quad W=(1-e_N)V.
\]
The hypothesis descends to the \(G/N\)-module \(U\),
so induction gives \(u\in U^F\) satisfying (A.43.2)
for that quotient. For every normal \(H\) in \(G\),
\[
                         \operatorname{Stab}_G(e_Hu)=HN.
                         \tag{A.43.3}
\]

Decompose \(W\) by the commuting projectors \(e_H\),
for all normal subgroups \(H\) of \(G\), and by the
simple factors of the commutative semisimple algebra
\(\mathbf F_p[N]\). The former joint blocks are
indexed by kernels \(L\triangleleft G\) of irreducible
\(G\)-constituents: on a block with kernel \(L\),
\[
                         e_H=1\quad\Longleftrightarrow\quad H\subseteq L.
\]
This follows from Maschke's theorem, since for normal
\(H\) the \(H\)-invariants of an irreducible constituent
are either the whole constituent or zero. We group
together all constituents with the same kernel; no
choice of a splitting field is needed for that grouping.
The further \(N\)-character factors correspond to
\(p\)-power orbits of nontrivial characters of \(N\).
All characters in such a factor have the same kernel.

The operator \(F\) permutes these refined blocks.
For each orbit of blocks whose direct sum has a
nonzero \(F\)-fixed space, choose a nonzero fixed
vector on that orbit, and add these choices to obtain
\(w\in W^F\). Each chosen vector has a nonzero component
in every block of its orbit. On an \(N\)-character
factor, a nontrivial character value minus one is
invertible; thus every nonzero component has exactly
the corresponding character kernel as its stabilizer
in \(N\). Consequently, simultaneously for all normal
\(H\), our \(w\) realizes the pointwise kernel:
\[
 \operatorname{Stab}_N(e_Hw)=
 \ker\bigl(N\text{ acting on the set }e_H(W^F)\bigr).
                         \tag{A.43.4}
\]
Here \(e_H(W^F)\) need not be a \(G\)-submodule or
even equal \((W^H)^{F^r}\) for a field of definition
of \(H\). We do not make either of these replacements.

We prove that the kernel in (A.43.4) is \(H\cap N\).
Suppose that it contains \(n\in N\setminus H\).
If \(n\) acted trivially on \(V^H\), so would all its
\(G\)-conjugates, because \(V^H\) is \(G\)-stable. Put
\[
 B=H\langle gng^{-1}:g\in G\rangle.
\]
Then \(B\triangleleft G\), \(B/H\) is nontrivial
elementary abelian of exponent \(\ell\), and
\(V^H=V^B\). Choose a hyperplane \(A/H\) in \(B/H\).
It gives \(A\triangleleft B\), \(B/A\simeq C_\ell\),
and \(V^A=V^B\), contradicting (A.43.1) after a
common power making \(F\) and \(\phi\) the identity.
Thus \(n\) does not act trivially on \(V^H\).
Since \(N\) acts trivially on \(U\),
there is an \(N\)-character \(\chi\) in \(W^H\)
with \(\chi(n)\ne1\).

Among all normal constituent kernels \(L\supseteq H\)
supporting a nonzero character in the line
\[
                 \{\chi^a: a\in\mathbf F_\ell^\times\},
\]
choose one maximal under inclusion. Put
\(K=\ker\chi\subset N\). A surviving constituent with
kernel \(L\) shows \(L\cap N\subseteq K\), so
\[
                         LK\triangleleft LN,\qquad
                         LN/LK\simeq C_\ell.
\]
Modulo \(L\), the first subgroup is
\(K/(L\cap N)\) inside the abelian group
\(N/(L\cap N)\), which proves the asserted normality.
Consider the relative module
\[
                         V^{LK}/V^{LN}.
                         \tag{A.43.5}
\]
Its constituents have a normal \(G\)-kernel containing
\(L\), and a nontrivial \(N\)-character in the displayed
line. Maximality therefore forces their \(G\)-kernel
to be exactly \(L\). Thus (A.43.5) consists precisely
of the \(L\)-kernel blocks with nontrivial \(N\)-character
kernel \(K\).

No full \(F\)-orbit meeting this relative module can
have an \(F\)-fixed vector. If it did, our choice of
\(w\) would put a nonzero component in its selected
block. Since \(H\subseteq L\), that component survives
\(e_H\); every nonzero multiple of \(\chi\) detects
\(n\). This contradicts \(n\in\operatorname{Stab}_N(e_Hw)\).

Let \(r\) be the least positive integer for which
\(\phi^r(L)=L\) and \(\phi^r(K)=K\). The pair \((L,K)\)
is determined by each refined block in (A.43.5).
Hence \(r\) divides the length \(t\) of its full
\(F\)-orbit of blocks. The part of that orbit appearing
in (A.43.5) is exactly one cycle under \(F^r\).
A fixed vector of \(F^r\) on this cycle is equivalent
to a fixed vector of \(F^t\) on one block, and hence
to a fixed vector of \(F\) on the full orbit.
There is none by the preceding paragraph.

The averaging splitting of (A.43.5) is \(F^r\)-stable,
so we conclude
\[
                         (V^{LK})^{F^r}=(V^{LN})^{F^r}.
\]
This contradicts (A.43.1) for the cyclic-prime pair
\(LK\triangleleft LN\). Therefore
\(\operatorname{Stab}_N(e_Hw)=H\cap N\) for every normal
\(H\). Combining this with (A.43.3) in the
\(G\)-stable direct sum \(U\oplus W\) proves that
\(v=u+w\) satisfies (A.43.2), exactly as in A.42.

For the geometric application, odd-order groups are solvable
by the verified Feit--Thompson theorem used in A.6.
A \(\phi^r\)-stable
subgroup defines its pointed intermediate curve over
\(k_r\). Prime-to-\(p\) norm and pullback identify
the relevant \(p\)-torsion invariants with
\(V^A,V^B\); their strictness is exactly the presence
of rational new \(p\)-torsion, hence failure of (R)
for \(D/A\to D/B\) over \(k_r\). The deck group of
this relative cyclic-prime cover may be nonconstant.
If (R) succeeded, A.23 would make that cyclic group
constant after its exact Frobenius orbit length,
raising the relative class number to that power and
preserving its \(p\)-unit property. Thus failure of
all constant-deck relative tests over every extension
indeed implies the stated hypothesis.
\end{proof}

\paragraph{The normal-level distinction remains necessary.}
The conclusion concerns all normal quotients of \(D/C\).
The relative cyclic cover isolated in the proof can
have source \(D/(LK)\) which is not normal over \(C\),
although its base \(D/(LN)\) is normal over \(C\).
Accordingly the compactness consequence for the
normal-level ghost module is a one-way implication
from hereditary failure of all relative tests.
The reverse assertion needs freeness at the nonnormal
source as well; it does not follow merely from
regularity at every normal quotient.


\subsection{The hyperelliptic refinement of the scalar obstruction}

\begin{claim}{A.44}{Proved; an all-extension obstruction using the pointed involution}
Let $C/k$ be a genus-two curve over a finite field of
characteristic two, and let $P\in C(k)$ be a Weierstrass
point. Fix a nontrivial geometric connected cyclic etale
cover of odd degree $n$, with a chosen point above $P$.
Let $k_d/k$ be the first constant extension over which
its geometric isomorphism class is stable, and use its
unique pointed descent to $k_d$. Suppose Frobenius on
its geometric cyclic deck group is multiplication by
$a\in(\mathbf Z/n)^\times$, of order $m$. Put
\[
 U=(\mathbf Z/n)^\times,\qquad
 T=\langle2\rangle\subset U,\qquad
 H=\langle a,-1\rangle\subset U.
\]
If $T\subset H$, every geometric two-torsion point of
the primitive factor is rational over the first
constant-deck field $k_{dm}$. Consequently, for this
pointed cover, the primitive class number is odd over
some constant-deck field if and only if its geometric
primitive two-rank is zero. When $n=\ell$ is prime,
this is the whole relative criterion \textup{(R)}.

More generally, the possible primitive Frobenius
multipliers over $k_{dm}$ belong to the units of a
direct factor of
\[
 \left(\mathbf F_{2^{\delta}}\right)^c,
 \qquad
 \delta=\frac{|T|}{|T\cap H|},\qquad
 c=\frac{|U|}{|TH|}.                              \tag{A.44.1}
\]
Thus the preceding sufficient condition is exactly
$\delta=1$ for this fixed-ring argument.

In particular, suppose the Jacobian Frobenius of $C/k$
is the scalar $[\varepsilon2^r]$, with
$\varepsilon\in\{1,-1\}$ and $\#k=2^{2r}$.
For an odd prime $\ell$, put $f=\operatorname{ord}_\ell(2)$.
The sufficient condition becomes, independently of the
sign $\varepsilon$,
\[
 \begin{cases}
  \gcd(r,f)=1,&f\text{ odd},\\
  \gcd(r,f/2)=1,&f\text{ even}.
 \end{cases}                                      \tag{A.44.2}
\]
This sharpens A.41, and also applies to nonscalar
Jacobian Frobenius whenever the displayed deck-group
condition holds over the first field of definition of
the chosen cyclic cover.
\end{claim}
\begin{proof}
The hyperelliptic involution $\iota$ fixes $P$ and acts
as minus one on $J_C$. It therefore preserves the
isomorphism class of every cyclic cover. It has a
unique lift $\tau$ fixing the chosen point above $P$:
two such lifts differ by a deck transformation fixing
that point, hence are equal. The lift inverts the
cyclic deck group, and its square is the identity.
Uniqueness also shows that the lift is defined over
$k_d$ and commutes with its arithmetic Frobenius $F$.
This uses the pointed descent; an arbitrary unpointed
twist need not preserve this commutation relation.

Let
\[
 R=\mathbf F_2[X]/\Phi_n(X),\qquad
 W=H^1_{\mathrm{et}}(D_{\overline k},\mathbf F_2)_{\mathrm{prim}}.
\]
Since $n$ is odd, $R$ is finite etale. Every nontrivial
coherent character space of a genus-two cyclic cover
has dimension one. The geometric etale part therefore
has rank zero or one over every field factor of $R$;
equivalently, $W$ is free of rank one over a direct
factor $R'$ of $R$. This is the same multiplicity
argument used in A.41. The support $R'$ is stable
under $X\mapsto X^a$ and $X\mapsto X^{-1}$, because
$F$ and $\tau$ act invertibly on $W$.

The action of $F$ is semilinear for the deck
automorphism $X\mapsto X^a$ (or its inverse under the
dual cohomology convention). Since $a^m=1$, $F^m$ is
$R'$-linear, and hence is multiplication by a unit
$b\in R'$. Commuting $F^m$ with $F$ shows that
$b$ is fixed by $a$. Commuting $F^m$ with $\tau$
shows that $b$ is fixed by inversion. Thus
\[
                         b\in(R'^H)^\times.           \tag{A.44.3}
\]
This argument does not assume that the primitive
factor is ordinary; rank zero factors have simply
been omitted from $R'$.

Over $\overline{\mathbf F}_2$, the primitive character
set is $U$. The field factors of $R$ are the orbits
of multiplication by two, that is, the cosets of $T$.
Taking $H$-invariants groups these into cosets of $TH$.
The stabilizer $T\cap H$ acts on each corresponding
field of degree $|T|$, leaving a field of degree
$|T|/|T\cap H|$. This proves (A.44.1). A stable
direct factor $R'$ merely removes some of these
field factors.

If $T\subset H$, the fixed ring is a product of copies
of $\mathbf F_2$, whose unit group is trivial. Hence
$b=1$ and $F^m$ is the identity on $W$. The same holds
on the dual geometric two-torsion points of the
primitive abelian factor. A nonzero such point then
remains rational over every further constant-deck
extension. Conversely, geometric two-rank zero
excludes rational two-torsion over every field, so
its class number is always odd. Any field over which
this pointed geometric cover has constant decks
contains $k_{dm}$: first the isomorphism class must
be stable, requiring a multiple of $d$, and then
the deck scalar must become the identity, requiring
a further multiple of $m$.

For scalar Jacobian Frobenius $[s]$ with
$s=\varepsilon2^r$ and $s^2=\#k$, every cyclic
isomorphism class is stable over $k$ and the pointed
deck scalar is $s$. Equivalently, on its Kummer line
the scalar is $s$, and Cartier duality gives the
deck scalar $(\#k)/s=s$. Therefore
$H=\langle2^r,-1\rangle$. For prime $\ell$, if
$f$ is odd then $-1\notin T$, and
$H\cap T=\langle2^r\rangle$; the fixed-field degree
is $\gcd(r,f)$. If $f$ is even, $-1=2^{f/2}$ in
$T$, so $H=\langle2^{\gcd(r,f/2)}\rangle$ and
the fixed-field degree is $\gcd(r,f/2)$. This proves
(A.44.2).

For composite $n$ the same inclusion and fixed-ring
formulas remain valid. The simple prime formula
must then be replaced by the actual subgroup
calculation in $(\mathbf Z/n)^\times$; in particular,
an even order of two does not by itself imply
$-1\in\langle2\rangle$.
\end{proof}

\paragraph{Scope for the counterexample pencil.}
For a member of the supersingular pencil over
$\mathbf F_{2^{38}}$, the scalar is $-2^{19}$.
If it has no nontrivial cyclic primitive-zero cover
as in A.81--A.82, (A.44.2) excludes every rational
prime-cyclic \textup{(R)} witness whose prime $\ell$
satisfies $19\nmid\operatorname{ord}_\ell(2)$,
over every finite extension. In particular, every
such witness has $\ell\ge191$: for $\ell<191$,
divisibility by19 would force $\ell=19j+1$ with
$j\in\{2,4,6,8\}$, and the four values
$39,77,115,153$ are composite. This is a necessary
condition on witnesses, not a claim that degree191
succeeds or that the arithmetic criterion A3 fails.



\begin{claim}{A.45}{Proved obstruction to a finite-type-only return bound}
Fixed algebraic coefficient functions evaluated at prime-to-$p$ torsion
points, even when their Frobenius periods are prime to $p$, do not by
themselves give a bounded valuation for the return unit minus one.
More precisely, put $p=2$, let $\omega$ be a primitive cube root of
unity, and set $K=\mathbf Q_2(\omega)$. For $n\ge1$ choose a primitive
$3^n$-th root $\zeta_n$ satisfying
\[
 m_n=3^{n-1},\qquad \zeta_n^{m_n}=\omega.
\]
These roots are prime-to-two Teichmuller points. The fourth-power
Frobenius orbit of $\zeta_n$ has the odd length $m_n$. For the single
fixed degree-one coefficient function $A(X)=1-2X$, its orbit return is
\[
 U_n=\prod_{j=0}^{m_n-1}A(\zeta_n^{4^j})
     =N_{K(\zeta_n)/K}(1-2\zeta_n)
     =1-2^{m_n}\omega.                              \tag{A.45.1}
\]
In particular
\[
 v_2(A(\zeta_n)-1)=1,\qquad v_2(U_n)=0,
 \qquad v_2(U_n-1)=m_n=3^{n-1}.                    \tag{A.45.2}
\]
For the full square-Frobenius orbit over $\mathbf Q_2$, the analogous
return is
\[
 N_{\mathbf Q_2(\zeta_n)/\mathbf Q_2}(1-2\zeta_n)
 =\Phi_{3^n}(2)=1+2^{m_n}+2^{2m_n},                 \tag{A.45.3}
\]
again a unit whose difference from one has valuation $m_n$.
Consequently the annihilator sizes $2^{v_2(U_n-1)}$ grow faster than
$3^{gn}$ for every fixed positive integer $g$.

There is an integral semilinear module realizing these annihilator
sizes. Let $L_n=K(\zeta_n)$, put $M_n=\mathcal O_{L_n}$, and let
$\sigma_n$ be the fourth-power Frobenius of $L_n/K$. Give $M_n$ the
operator
\[
 \mathcal F_n(x)=(1-2\zeta_n)\sigma_n(x).
\]
It normalizes the odd cyclic deck action $\delta_n(x)=\zeta_nx$ by
$\mathcal F_n\delta_n\mathcal F_n^{-1}=\delta_n^4$, and
\[
 \mathcal F_n^{m_n}=U_n,\qquad
 M_n/(1-\mathcal F_n)M_n
       \simeq\mathcal O_K/(2^{m_n}),\qquad
 \exp\bigl((L_n/M_n)^{\mathcal F_n}\bigr)=2^{m_n}.
                                                        \tag{A.45.4}
\]
This is an algebraic Frobenius-cocycle countermodel to a proposed
intermediate bound. It is not asserted to arise from the Jacobians
in A.107, or from pure cohomology of Weil weight one. It therefore
does not refute the conditional geometric route in A.107.
Indeed $|U_n|^2=4^{m_n}+2^{m_n}+1>4^{m_n}$ under a complex
embedding, so the required weight-one equality already fails.
\end{claim}
\begin{proof}
The identity $v_3(4-1)=1$ and the elementary binomial valuation formula
show that $4$ has order $3^{n-1}$ modulo $3^n$. Its powers exhaust
the subgroup of units congruent to one modulo three. Thus the roots
\[
 \{\zeta_n^{4^j}:0\le j<m_n\}
\]
are exactly the roots of $X^{m_n}-\omega$. All prime-to-two roots of
unity belong to unramified extensions of $\mathbf Q_2$, and the
Frobenius over its unramified quadratic extension $K$ is fourth
power. The field degree $[K(\zeta_n):K]$ is therefore $m_n$.
Evaluating the displayed monic polynomial at $X=1/2$ and multiplying
by $2^{m_n}$ proves (A.45.1). Since $\omega$ and $\zeta_n$ are
two-adic units, (A.45.2) follows immediately.

The two conjugates over $\mathbf Q_2$ are
$1-2^{m_n}\omega$ and $1-2^{m_n}\omega^2$. Their product is
\[
 (1-2^{m_n}\omega)(1-2^{m_n}\omega^2)
   =1+2^{m_n}+2^{2m_n}.
\]
This is also $\Phi_{3^n}(2)$ by
$\Phi_{3^n}(X)=1+X^{m_n}+X^{2m_n}$. Its difference from one
is $2^{m_n}(1+2^{m_n})$, so the asserted valuation is exact.
Finally $3^{n-1}/n\to\infty$, proving the comparison with $3^{gn}$.

For the integral statement, $M_n$ is free of rank $m_n$ over
$\mathcal O_K$. After extension of scalars to $L_n$, its $m_n$
embeddings diagonalize multiplication, and $\sigma_n$ cyclically
permutes them. Thus $\mathcal F_n$ becomes a weighted cyclic
permutation, with product of its weights $U_n$. Its characteristic
polynomial over $K$ is $T^{m_n}-U_n$, so
\[
                 \det(1-\mathcal F_n)=1-U_n.
\]
Modulo two the operator $\mathcal F_n$ is just the fourth-power
automorphism of $\mathbf F_{4^{m_n}}$. Its fixed space has dimension
one over $\mathbf F_4$. The Smith normal form over the discrete
valuation ring $\mathcal O_K$ therefore has exactly one nonunit
invariant factor, whose valuation equals $v_2(1-U_n)=m_n$.
This proves the cokernel assertion. Applying $1-\mathcal F_n$ to
$0\to M_n\to L_n\to L_n/M_n\to0$, and using its invertibility on
$L_n$, identifies that cokernel with $(L_n/M_n)^{\mathcal F_n}$.
Since $\mathcal O_K$ is unramified quadratic, this finite abelian
group is $(\mathbf Z/2^{m_n}\mathbf Z)^2$.

The coefficient $A$ is invertible at every point under consideration
and has constant reduction one. Thus bounded degree of $A$, bounded
valuation of $A-1$, and periods prime to two all hold in this model.
The growing quantity comes from taking a norm over an increasingly
long orbit. A fixed-locus torsion-distance theorem cannot be applied
to these norm polynomials as though their degrees were bounded.
The exact identities for $1\le n\le6$ are independently checked in
$\mathbf Z[X]/(\Phi_{3^n}(X))$ by
\path{frobenius_norm_obstruction.py}.
\end{proof}


\begin{claim}{A.46}{Proved limit of a rank-one overconvergence approach}
Let $U/\mathbf F_q$ be smooth and geometrically connected, and let
$\mathcal L$ be a rank-one unit-root overconvergent $F$-isocrystal
on $U$, with coefficients in a fixed finite extension of
$\mathbf Q_p$. Write $b_x$ for its Frobenius eigenvalue at a closed
point $x$. For every fixed integer $c\ge0$, the finite valuations
\[
 \{v_p(b_x-1):x\in|U|,\ v_p(\deg x)\le c,\ b_x\ne1\}
                                                        \tag{A.46.1}
\]
are uniformly bounded. This assertion would apply to a character
family only after its return eigenvalues had been identified with
such an overconvergent object.

There is a concrete obstruction to that identification in the
scalar example $C_0:y^2+y=x^5$ over $\mathbf F_{16}$. Its ordinary
character unit roots at the order-three parameters and the order-thirteen
parameter of A.34 cannot be interpolated by one rank-one coefficient
object with finite geometric monodromy on a common geometrically
connected base over $\mathbf F_{16}$ containing those parameter
points with their indicated closed-point degrees two and three.
In particular no rank-one
unit-root overconvergent $F$-isocrystal can have all those prescribed
return eigenvalues. This conclusion concerns the proposed interpolation;
it does not assert that an intrinsic arithmetic $F$-isocrystal on
the ordinary character locus has been constructed.
\end{claim}
\begin{proof}
By \cite[Proposition~8.6, p.~42]{AC_HartlPal2025}, a tensor power
$\mathcal L^{\otimes N}$ is constant. If its constant Frobenius
eigenvalue is $c_0$, then
\[
                         b_x^N=c_0^{\deg x}.           \tag{A.46.2}
\]
All these quantities are $p$-adic units. Suppose first that $c_0$
is not a root of unity. In a finite splitting field, let $d$ be the
order of its residue. If $d\nmid m$, then
$v_p(c_0^m-1)=0$. Otherwise write $m=dp^a u$ with $p\nmid u$.
The factorization of $X^u-1$ at
$X=c_0^{dp^a}\equiv1$ gives
\[
 v_p(c_0^m-1)=v_p(c_0^{dp^a}-1).
\]
When $v_p(m)\le c$, the latter runs through a finite set of finite
values. Moreover
\[
 v_p(b_x-1)\le v_p(b_x^N-1)=v_p(c_0^{\deg x}-1),
\]
because the other factor of $b_x^N-1$ is integral. If $c_0$ is a
root of unity, (A.46.2) puts every $b_x$ in one fixed finite set
of roots of unity, which also proves (A.46.1). This establishes
the asserted conditional bound.

We now prove the obstruction using two genuine character fibers.
On $J_{C_0}$, Frobenius over $\mathbf F_{16}$ is $[-4]$.
Every nonzero point $a_3\in J_{C_0}[3]$ therefore has degree two.
For the Weierstrass point $P=\infty$, the Raynaud theta curve is
$C_0-P$ and contains no nonzero three-torsion: a function with
divisor $3Q-3P$ would have pole order three at $P$, whereas
$L(3P)=\langle1,x\rangle$. The corresponding pointed cubic
cover over $\mathbf F_{256}$ consequently has ordinary new part,
by the rank-one character criterion. The pointed hyperelliptic
involution commutes with Frobenius and interchanges the two
nontrivial cubic characters. Their common polynomial thus has
coefficients in
$\mathbf Z[\zeta_3]\cap\mathbf R=\mathbf Z$. It is
\[
                         1+s_3T+256T^2,
 \qquad s_3\in\mathbf Z,
\]
and has an ordinary unit root $\beta$ satisfying
$\beta^2+s_3\beta+256=0$. The Weil bound and ordinariness imply
$s_3^2<1024$, so $\mathbf Q(\beta)$ is imaginary quadratic.

Let $a_{13}$ be the explicit order-thirteen parameter of A.34.
Its degree over $\mathbf F_{16}$ is three. A character unit root
$\alpha$ of its pointed cover over $\mathbf F_{4096}$ satisfies
\[
 P(X)=X^4+139X^3+12473X^2+569344X+16777216=0.       \tag{A.46.3}
\]
More precisely, put $K_0=\mathbf Q(\sqrt{13})$. Its character
quadratic is
\[
 X^2-tX+4096,
 \qquad t=\frac{-139\pm13\sqrt{13}}2,
\]
with discriminant
\[
 d=t^2-16384=\frac{-22009\mp1807\sqrt{13}}2,
 \qquad
 N_{K_0/\mathbf Q}(d)=110486961=39^2\cdot17\cdot4273.
                                                        \tag{A.46.4}
\]
Both real conjugates of $d$ are negative, so $d$ is not a square
in the real field $K_0$. Also $t=\alpha+4096/\alpha$ belongs to
$\mathbf Q(\alpha)$ and is irrational. Hence
$[\mathbf Q(\alpha):\mathbf Q]=4$ and $K_0$ is a subfield.
The rational number in (A.46.4) is neither a rational square nor
thirteen times a rational square: its valuation at seventeen is
one. A rational number is a square in $\mathbf Q(\sqrt{13})$
only in those two square classes. Thus $N(d)$ is not a square
in $K_0$. In particular the nontrivial automorphism of $K_0$
cannot extend to $K_0(\sqrt d)$, since such an extension would
make $\overline d/d$, and hence $d\overline d$, a square in $K_0$.
This proves that $\mathbf Q(\alpha)$ is not Galois. It has no
quadratic subfield other than $K_0$: two different quadratic
subfields would generate a biquadratic degree-four field.

Complex conjugation in this CM field fixes $K_0$ and sends
$\alpha$ to $4096/\alpha$. Under the chosen two-adic embedding,
\[
             v_2(\alpha)=0,
 \qquad v_2(4096/\alpha)=12.
\]
No positive power of $\alpha$ can therefore belong to $K_0$.
Since any proper subfield of a degree-four field has degree one
or two, every positive power of $\alpha$ has degree four.

If one rank-one object with finite geometric monodromy had the
specified eigenvalues at these degree-two and degree-three points,
some tensor power would be constant, giving
\[
                     \beta^N=c_0^2,
 \qquad \alpha^N=c_0^3,
 \qquad \alpha^{2N}=\beta^{3N}.
\]
The right-hand side has degree at most two over $\mathbf Q$,
whereas the left-hand side has degree four. This contradiction
proves the interpolation obstruction. For an overconvergent
rank-one unit-root object, the required finite geometric
monodromy follows from the theorem already cited.

The distinction between a full geometric cohomology object and
its unit-root part is essential. The ordinary Legendre family
has overconvergent full cohomology but a unit-root subobject
which is not overconvergent; see
\cite[Theorem~5.4.3.5, Corollary~5.4.3.10 and
Remark~5.4.3.11, pp.~49--52]{AC_Etesse2010}.
In that example the Frobenius coefficient alone can be
overconvergent while its connection is not. Therefore neither
geometric origin nor a finite algebraic formula for the Cartier
coefficient supplies the missing overconvergence hypothesis.
\end{proof}

% A.47: full Frobenius-packet obstruction for the explicit A6 pencil.
\subsection{Two remaining conductors after simultaneous packet constraints}

\begin{claim}{A.47}{Proved by an exact finite certificate}
\label{clm:ACpacketconductors}
For any member, with \(t\notin\mathbb F_4\), of the genus-two pencil of
Claim~\ref{clm:AGprimitivecounterexample} over
\(\mathbb F_{2^{38}}\), a nontrivial odd conductor of a
primitive two-rank-zero cyclic cover must belong to
\[
                              \{191,229\}.
\tag{A.47.1}
\]
Consequently the incidence estimate of Claim~A.82 improves to
\[
 \#\{\text{bad parameters}\}
       \le133086336+229171540=362257876,
\]
and at least
\[
              274877906944-4-362257876=274515649064
\tag{A.47.2}
\]
parameters have no nontrivial odd primitive-zero cover.
These are counts of parameters, not of isomorphism classes.
The assertion retains the degree-one qualification of the A6
counterexample and does not exclude rational class-number witnesses.
\end{claim}

\begin{proof}
Let \(n\) be a candidate conductor, and write
\[
 T=\langle2\rangle,\qquad
 U=\langle2,-1\rangle,\qquad
 H=\langle2^{19},-1\rangle
       \subseteq(\mathbb Z/n)^\times.
\tag{A.47.3}
\]
The theta curve contains the origin, is symmetric, and is stable
under the scalar Frobenius \([-2^{19}]\). Therefore a primitive
theta point \(aQ\), for a chosen generator \(Q\) of the cyclic
subgroup, forces the entire packet \(aHQ\) to lie on theta.
Primitive two-rank zero requires a theta point in every
multiplication-by-two orbit of primitive points. By symmetry it
thus requires at least one \(H\)-packet in each coset of \(U\).
This formulation also handles \(n=191\), where \(-1\notin T\)
and symmetry joins two different doubling orbits.

For any such choice of packets, the set of residues
\[
                  S=\{0\}\cup\bigcup_{aU} a_{aU}H
\tag{A.47.4}
\]
must satisfy
\[
 \#\{(x,y)\in S^2:x-y=d\}\le2
                 \qquad(0\ne d\in\mathbb Z/n).
\tag{A.47.5}
\]
Indeed these pairs give distinct points of the intersection of
theta with its translate by \(dQ\). The translation is nontrivial,
the principal polarization has trivial translation stabilizer,
and the intersection number is two. Additional theta points cannot
repair a violation of (A.47.5), so selecting just one required
packet per \(U\)-coset gives a necessary test.

The single-packet test of Claim~A.85 leaves only
\(191,229,457,1257\). At \(n=1257=3\cdot419\), the groups
in (A.47.3) have \(|U|=418\) and \(|H|=22\), while
\(\varphi(n)=836\). Hence two different packets are required.
Each packet has eleven elements in each nonzero residue class
modulo three, since negation exchanges the classes. Thus (A.47.4)
has occupancies \((1,22,22)\) modulo three and supplies
\[
                   2\cdot22\cdot21=924
\]
ordered pairs whose difference is a nonzero multiple of three.
There are only \(418\) such differences modulo \(1257\).
Equation (A.47.5) permits at most \(836\) pairs, a contradiction.

At \(n=457\), one has \(|U|=76\), \(|H|=4\), and six
\(U\)-cosets. Choose one theta point as the generator \(Q\),
which normalizes its packet to \(H\) in the identity coset.
Each other \(U\)-coset contains exactly nineteen possible
\(H\)-packets. The accompanying standard-library program
\texttt{primitive\_packet\_sidon.py} enumerates these packets by
group closure and tests all possible extensions of the normalized
packet. It first discards any candidate already violating
(A.47.5) with \(\{0\}\cup H\); the five remaining candidate
counts are
\[
                            17,19,17,18,18.
\]
It then performs an exhaustive depth-first search. At each step
it adds every ordered difference involving a new packet and
rejects the branch as soon as a multiplicity exceeds two.
The numbers of surviving partial sets containing respectively
one through six packets are exactly
\[
                         1,17,226,879,262,0.
\tag{A.47.6}
\]
There is therefore no possible full required set for \(457\).
The JSON certificate retains the complete normalization, packet
indices, rejected first-packet witnesses, and search counts;
the source performs the complete search without a heuristic cutoff.
For \(191\) and \(229\) it also retains compatible necessary
sets, making explicit that this test gives no exclusion there.
This proves (A.47.1).

The incidence estimate is now needed only at these two conductors.
The minimum numbers of primitive theta points are respectively
\(10\) and \(12\), and the numbers of exact-order points in
the ambient product of elliptic curves are
\(J_4(191)=1330863360\) and \(J_4(229)=2750058480\).
Each ambient nonzero torsion point lies on at most one pencil
member. Dividing these two incidence bounds by \(10\) and
\(12\), respectively, gives (A.47.2). No individually chosen
parameter is asserted to have been certified by this counting
argument.
\end{proof}


\subsection{The fixed norm module and full arithmetic sections}

\begin{claim}{A.48}{Proved: section origin is weaker than point origin}
\label{clm:ACnormsectionorigin}
Let $C/k$ be smooth, proper and geometrically connected of genus
at least two over a finite field of characteristic $p$, and let
$P\in C(k)$. Use the full etale fundamental groups in the
first row below, and put
\[
 \Omega=\pi_1(C_{\bar k},\bar P),\qquad
 \Delta=\Omega^{(p')},\qquad
 N=\ker(\Omega\longrightarrow\Delta),\qquad
 V=N^{\mathrm{ab},p}.
                                                        \tag{A.48.1}
\]
Here $N^{\mathrm{ab},p}$ is the maximal abelian pro-$p$
quotient. The actual point section $s_P$ acts on these
groups; write $F$ for arithmetic Frobenius and also for
its induced automorphism of $V$. For the pointed norm
module $M_P$ of A.100 there is a natural exact sequence
of compact abelian pro-$p$ groups
\[
 0\longrightarrow V\xrightarrow{\ F-1\ }V
      \xrightarrow{\ \varepsilon\ }M_P^F
      \longrightarrow0.
                                                        \tag{A.48.2}
\]

Every $m\in M_P^F$ is the compatible $p$-adic
abelianized difference from $s_P$ of a continuous section
of the \emph{full} arithmetic etale fundamental group
of $C$. That section projects to exactly $s_P$ in the
geometrically prime-to-$p$ quotient. If $m\ne0$, A1
implies that this full section is not geometrically
conjugate to the section of any point of $C(k)$.

Projection is surjective at every $F$-stable pointed
level, in particular
\[
                    M_P^F\twoheadrightarrow J_C(k)[p^\infty].
                                                        \tag{A.48.8}
\]
Thus nonzero section-origin classes actually exist
whenever $J_C(k)[p^\infty]\ne0$.

The statement supplies no element with free deck orbit
at every level. In particular it neither proves nor
refutes A10(c) or the nontrivial-cover existence
assertion A3.
\end{claim}

\begin{proof}
\emph{1. The integral norm limit.}
Take the open normal subgroups $H\triangleleft\Delta$
that are stable under $F$. They are cofinal: every
finite geometric cover has a model over a finite
constant extension, hence its defining subgroup has
a finite Frobenius orbit, whose intersection is a
stable normal refinement. Let $U_H\subset\Omega$
be the inverse image of $H$. The point section
gives a model $D_H/k$ with compatible rational lift
$P_H$. Write
\[
                    T_H=T_pJ_{D_H}.
\]
The Abel map based at $P_H$ identifies
$U_H^{\mathrm{ab},p}$ with $T_H$. This comparison
includes the characteristic prime: see
\cite[Proposition~9.1, Lemmas~9.4--9.5 and
Corollary~9.6]{AC_MilneJac1986}.
If $H'\subset H$, the map induced by
$U_{H'}\hookrightarrow U_H$ is the norm map on
Tate modules. Indeed the pointed Abel maps commute
with the morphism of curves and its norm homomorphism
on Jacobians. Equivalently, its dual is pullback on
first etale cohomology, not the transfer multiplied
by the covering degree.

There is an $F$-equivariant isomorphism
\[
                 V\xrightarrow{\sim}\varprojlim_H T_H.
                                                        \tag{A.48.3}
\]
Here is an elementary continuity verification. A
continuous character $c:N\to\mathbf Q_p/\mathbf Z_p$
has finite image. Choose an open normal subgroup
$O\triangleleft\Omega$ with $N\cap O\subset\ker c$.
Since $[O,N]\subset O\cap N$, the formula
\[
               \widetilde c(no)=c(n),\qquad n\in N, o\in O,
\]
defines a character on the open subgroup $NO$.
Its quotient $\Omega/(NO)$ has prime-to-$p$ order,
since it is a finite quotient of $\Delta$.
Restricting to a sufficiently small $F$-stable normal
subgroup containing $N$ therefore extends $c$ to
one of the $U_H$. Conversely, a character on $U_H$
vanishing on $N$ factors through the pro-prime-to-$p$
group $H$ and is zero. Thus
\[
 \operatorname{Hom}_{\mathrm{cont}}(N,\mathbf Q_p/\mathbf Z_p)
  =\varinjlim_H
    \operatorname{Hom}_{\mathrm{cont}}(U_H,\mathbf Q_p/\mathbf Z_p).
\]
Pontryagin duality gives (A.48.3), with the transition
maps already identified as norms. No comparison with
the full prime-to-$p$ configuration fiber is involved.

\emph{2. Frobenius coinvariants and rational torsion.}
For every $H$ there is an exact sequence
\[
 0\longrightarrow T_H\xrightarrow{\ F-1\ }T_H
     \xrightarrow{\ \varepsilon_H\ }
       J_{D_H}(k)[p^\infty]\longrightarrow0.
                                                        \tag{A.48.4}
\]
To check it, put $W_H=T_H\otimes_{\mathbf Z_p}\mathbf Q_p$.
The quotient $W_H/T_H$ is the group of geometric
etale $p$-primary points of $J_{D_H}$. The fixed
subgroup of $T_H$ is zero: otherwise it would give
a nonzero compatible infinitely $p$-divisible
system in the finite group $J_{D_H}(k)$.
Since $T_H$ has finite rank, $F-1$ is therefore
invertible on $W_H$, and
\[
 \varepsilon_H(a)=(F-1)^{-1}a\pmod {T_H}
                                                        \tag{A.48.5}
\]
gives (A.48.4). Its image is precisely the
$F$-fixed subgroup of $W_H/T_H$, which is
$J_{D_H}(k)[p^\infty]$. This is the full $p$-adic
Kummer comparison over a finite field; it does not
treat multiplication by $p$ as an etale isogeny.
The maps in (A.48.4) commute with norms.

For the cofinal stable levels, being fixed by the
action on the entire pointed inverse system means
exactly that each coordinate is rational over $k$.
Consequently
\[
           M_P^F=\varprojlim_H J_{D_H}(k)[p^\infty].
                                                        \tag{A.48.6}
\]
This also equips $M_P^F$ with its compact pro-$p$
topology, regardless of the unbounded-exponent issue
for the complete module $M_P$.

We verify the inverse-limit exactness explicitly.
For a compatible family $m_H$ on the right of
(A.48.4), each fiber
$\varepsilon_H^{-1}(m_H)\subset T_H$ is nonempty
and compact. Any finite collection of norm
compatibility equations can be satisfied by choosing
a preimage at one common refinement and taking its
norms to the finitely many levels in the equations.
The remaining coordinates may be chosen arbitrarily
in their nonempty fibers. The finite-intersection
property gives a compatible lift in
$\varprojlim_H T_H$. This proves surjectivity in
(A.48.2), without a uniform denominator or an
unjustified assertion about restricted transition maps.

If a compatible family $a_H$ maps to zero, write
$a_H=(F-1)b_H$ using (A.48.4). Each $b_H$ is
unique because $F-1$ is injective on $T_H$.
Commutation with norms then makes the $b_H$
compatible. This proves the kernel equality in
(A.48.2). Its left injectivity follows at each
level as well. Using (A.48.3) completes the proof
of that exact sequence.

The asserted surjectivity of projection has an
additional elementary proof. For a transition
$D_{H'}\to D_H$ over $k$, norm composed with
pullback is multiplication by its degree $d$.
Since $d$ is prime to $p$, multiplication by $d$
is an automorphism of $J_{D_H}(k)[p^\infty]$.
Norm on these rational $p$-primary groups is
therefore split surjective. Their finite inverse
system has surjective projection to every
coordinate, proving (A.48.8). No assertion that
freeness survives these projections is made.

\emph{3. Realization by a full arithmetic section.}
Let $m\in M_P^F$. Choose $v\in V$ with
$\varepsilon(v)=m$, and then choose $n\in N$
mapping to $v$. In the full arithmetic group
\[
        1\longrightarrow\Omega\longrightarrow\pi_1(C,\bar P)
          \longrightarrow G_k\longrightarrow1,
\]
let $F=s_P(\operatorname{Frob}_k)$ and form
$nF$. Every element of a profinite group defines
a continuous procyclic power homomorphism from
$\widehat{\mathbf Z}$. The power homomorphism for
$nF$, followed by the projection to
$G_k\simeq\widehat{\mathbf Z}$, is the identity:
it sends the dense generator $1$ to $1$.
It therefore defines a continuous section $s_n$.
Its prime-to-$p$ geometric projection is exactly
$s_P$, since $n\in N$.

At the pointed level $D_H$, the element $n$ belongs
to $U_H$. With the additive convention for
$T_H$, the abelianized cocycle difference satisfies
\[
      (s_n-s_P)(\operatorname{Frob}_k)=a_H,
                                                        \tag{A.48.7}
\]
where $a_H$ is the image of $v$ in $T_H$.
Our Kummer convention is
$\delta(x)(F)=F(\widetilde x)-\widetilde x$;
thus its cohomology class is
$a_H\bmod(F-1)T_H$, corresponding to $m_H$
under (A.48.5). This fixes the sign and verifies
the asserted compatible section realization.

\emph{4. Point origin and the centralizer.}
Suppose $m\ne0$ and $s_n$ were geometrically
conjugate to the section of $Q\in C(k)$.
Projection to the geometrically prime-to-$p$
group and A1 imply $Q=P$.
If $\omega\in\Omega$ conjugates $s_n$ to $s_P$,
its image in $\Delta$ centralizes $s_P(G_k)$,
because these two sections have identical
prime-to-$p$ projections. That geometric
centralizer is trivial by
\cite[Lemma~1.7]{AR_ST2009}. Hence $\omega\in N$.
Its conjugation changes the abelianized difference
by an element of $(F-1)V$, so (A.48.2) would
give $m=0$, a contradiction.

Thus nonzero norm-module classes have actual full
section realizations that are not point sections.
The argument says nothing further about their
deck stabilizers. In particular it does not replace
the missing point-origin assertion by a freeness
assertion or manufacture a counterexample to A3.
\end{proof}

\paragraph{Relation to the cuspidal proof.}
The point-recognition argument in A.88 starts from
finite neighborhoods that are known to contain
actual rational points. A section $s_n$ constructed
above does not carry that property merely because
it is a continuous splitting. Consequently its
existence supplies no input to that recognition
step. The exact sequence distinguishes the
Frobenius-fixed Tate module, which is zero, from
the Frobenius coinvariants, which form the
rational torsion norm module. Confusing those two
groups would incorrectly remove the A10 gap.

\paragraph{A small actual-curve example.}
Over $\mathbf F_2$, take the smooth proper model of
\[
                   C:\quad y^2+y=x+\frac1x+\frac1{x+1}.
                                                        \tag{A.48.9}
\]
The three simple poles $0,1,\infty$ give genus two
and two-rank two by the Artin--Schreier genus and
Deuring--Shafarevich formulas. Each pole supplies
one rational ramification point. There are no
other $\mathbf F_2$-points. Over $\mathbf F_4$,
the two remaining possible $x$-values give
Artin--Schreier right sides of absolute trace one,
so there are still exactly three points. Hence
\[
 \#C(\mathbf F_2)=\#C(\mathbf F_4)=3,
 \qquad P_C(T)=1-T^2+4T^4,
 \qquad \#J_C(\mathbf F_2)=4.
                                                        \tag{A.48.10}
\]
The exact elementary program
\texttt{norm\_module\_small\_curve.py} verifies
both point counts and the polynomial. In fact
$J_C(\mathbf F_2)\simeq(\mathbf Z/2)^2$:
writing the three ramification points as
$W_0,W_1,W_\infty$, the divisors of $x$ and
$x+1$ are $2(W_0-W_\infty)$ and
$2(W_1-W_\infty)$. The corresponding classes
are distinct and nonzero, because a difference
of distinct points on a positive-genus curve
cannot be principal.

For any choice of rational $P$, the fixed module
$M_P^F$ therefore surjects onto a group of order
four and has nonzero classes represented by
non-point-theoretic full sections. All three
actual point sections are distinct by A1.
No all-level free-orbit property is claimed for
these nonzero module classes.


\paragraph{A fixed geometric curve defeats every finite prime list.}
Let $(C,P)$ be the explicitly specified characteristic-two
curve with its point at infinity in
Claim~\ref{clm:AGexplicitprimitivecounterexample}.
For every finite set $S$ of odd primes there is a finite
extension $L$ of its definition field such that no
$P$-split cyclic cover of degree $\ell\in S$ satisfies
(R) over $L$, or over any further finite extension.
Indeed, there are only finitely many geometric cyclic
covers of each fixed prime degree. Since no nontrivial
odd primitive cover of this curve has primitive two-rank
zero, the new abelian variety of each prime-degree cover
has a nonzero geometric two-torsion point. Choose one
such point for each cover. A common finite extension
makes all the covers and their deck groups defined,
makes a selected point above $P$ rational, and makes
all the selected new two-torsion points rational.

Over this field every geometric cover has its unique
$P$-split twist: constant cyclic twists are parametrized
by the Frobenius translation on the fibre over $P$, and
exactly the zero translation splits that fibre.
The chosen rational new two-torsion persists after every
further extension. Thus every cover in the finite list
has even relative class number throughout those
extensions. Conversely any $P$-split cover of one of
these degrees is the unique pointed form of a member
of the same finite geometric list, so none is omitted.
Consequently there is no finite list of witness primes
that works uniformly over all finite definition fields
even for this single geometric pointed curve. This
statement permits witnesses at other primes and does
not refute A3.

% Source: agent_geometric/geometric_section.tex
% Integration fragment. Root may renumber AG.* into Part A.
\subsection{Correspondences, character orbits, and integral group-ring tests}

\begin{claim}{A.49}{Proved; a precise rational correspondence statement}
Let $L$ be a finite field of characteristic $p$, let $C/L$ be smooth,
proper and geometrically connected of genus at least two, and let
$P\ne Q$ in $C(L)$ have conjugate sections in the geometrically
prime-to-$p$ arithmetic fundamental group. Let
$\pi:D\to C$ be a connected finite etale cover with prime-to-$p$ geometric
Galois closure, let $P'\in D(L)$ lie above $P$, and let
$\phi:D\to X_0$ be an $L$-morphism to a smooth proper geometrically
connected curve of genus at least one. Assume
$J_{X_0}(L)[p]=0$. Then a lift $Q'\in D(L)$ supplied by R1 satisfies
$\phi(Q')=\phi(P')$. Moreover,
\[
 [\phi_*\pi^*(Q)-\phi_*\pi^*(P)]=0\quad\hbox{in }J_{X_0}(L).
\]
If $\pi$ is a torsor under a constant finite group $G$ over $L$,
then $P'$ makes the entire $P$-fibre rational, the $Q$-fibre is also
rational, and the stronger equality of effective divisors holds:
\[
             \phi_*\pi^*(Q)=\phi_*\pi^*(P).
\]
For a general cover the same divisor equality holds if the geometric
$p$-rank of $J_{X_0}$ is zero, or if $\deg(\pi)$ is strictly less
than the geometric gonality of $X_0$.
\end{claim}
\begin{proof}
R1 gives $[Q'-P']\in J_D(L)[p^\infty]$. The homomorphism
$\phi_*:J_D\to J_{X_0}$, defined over $L$ by pushforward of
degree-zero divisors, sends this class to zero, since a finite
abelian group with no element of order $p$ has no nonzero
$p$-primary element. On a smooth proper curve of positive genus,
the equality of the degree-one divisor classes of two points
forces the points to coincide: otherwise a function with divisor
their difference gives a degree-one map to $\mathbf P^1$.
This proves the first assertion.

On the base curve, R1 already gives $[Q-P]\in J_C(L)[p^\infty]$.
Apply the $L$-homomorphism $\phi_*\pi^*:J_C\to J_{X_0}$ to
obtain the asserted equality of divisor classes. For the constant
$G$-torsor, choose the lifts $P',Q'$ with equal section classes.
Their translates $gP',gQ'$ are $L$-rational and have $p$-primary
difference for every $g\in G$. Applying the preceding argument
to each pair and summing gives equality of effective divisors.
For geometric $p$-rank zero we may instead pass to a finite field
over which the Galois closure and all lifts are rational; the
target still has no geometric $p$-primary points. Equality of
the two geometric divisors descends to $L$.
Finally, distinct effective linearly equivalent divisors of degree
$N$ give, after cancelling their common part, a nonconstant function
of degree at most $N$. This proves the asserted gonality alternative.

The distinction between divisor classes and effective divisors
is necessary in the proof. One rational lift and the condition
$J_{X_0}(L)[p]=0$ do not assert that
$J_{X_0}(L')[p]=0$ for residue extensions $L'/L$ of the other
lifts. We therefore do not use that argument for the unrestricted
divisor assertion in lead L5.

The map $P\mapsto\phi_*\pi^*P$ to $\operatorname{Sym}^{\deg\pi}X_0$
is a morphism: etale locally on $C$, the finite etale fibre splits
into sections, whose images give an ordered tuple; quotienting by
the permutation action descends it. If $\phi$ is nonconstant, this
map is nonconstant, since a constant effective divisor would contain
every point of $\phi(D)$ in its finite support. Its equal-fibre locus
is consequently a proper closed subset of $C\times C$. This alone
does not exclude a hypothetical collision lying in that subset.
\end{proof}

\paragraph{Consequences of point separation and the finite spectral test.}
By A.88--A.90, A1 holds, so the unrestricted collision-dependent
implications in L5 hold as well: there is no pair of distinct
points satisfying their premise. In particular the unqualified
effective-divisor conclusion is no longer an open assertion.
The rationality qualifications above describe what the direct
correspondence argument proves without that theorem; they do
not weaken the now-established implication.

For clarity, the proposed spectral test has the following exact
meaning. Let $D/\mathbf F_q$ be smooth, proper and geometrically
connected of genus $h\ge1$, and write
\[
 P_D(T)=\det(1-TF\mid H^1_{\mathrm{et}}
       (D_{\overline{\mathbf F}_q},\mathbf Q_\nu)),
 \qquad
 H_D(X)=X^{2h}P_D(X^{-1}),\quad \nu\ne p.
\]
Then $J_D$ has a nonzero supersingular isogeny factor if and
only if one root $\alpha$ of $H_D$ has
$\alpha/\sqrt q$ a root of unity. Indeed, on a simple
isogeny factor this is exactly the supersingular
Newton-slope criterion; see
\cite[\S\S21.20--21.21(3)]{AC_Oort2008}. If one root has
this property, every conjugate in its simple factor does,
so that factor is supersingular. Conversely, a
supersingular factor supplies such roots.

This is a finite exact test: it is equivalent to
\[
 \deg\gcd\bigl(H_D(X),X^{2n}-q^n\bigr)>0
 \quad\hbox{for some }1\le n\le8h^2.
\]
To prove the bound, suppose $\beta=\alpha^2/q$ has exact
root-of-unity order $n$. Then
$\varphi(n)=[\mathbf Q(\beta):\mathbf Q]\le2h$.
For $n=\prod_\ell\ell^{a_\ell}$,
\[
 \frac{\varphi(n)^2}{n}
       =\prod_{\ell\mid n}\ell^{a_\ell-2}(\ell-1)^2
       \ge\frac12.
\]
Every odd-prime factor is at least one, and the factor
at two is at least $1/2$. Thus $n\le8h^2$.
Conversely, any nonconstant displayed greatest common
divisor gives $\alpha^{2n}=q^n$, which is precisely the
required root-of-unity condition.

For a constant abelian cover $D\to C$, the factorization
$P_D=\prod_\chi L(C,\chi,T)$ makes this the character
test in L5, with the trivial character included when
testing the whole Jacobian. Omitting it tests only the
new factor. A constant nonabelian cover requires all
irreducible Artin factors with their multiplicities;
a non-Galois cover requires its actual permutation
representation, or simply the polynomial $P_D$ itself.
The fixed-cover test supplies no existence theorem for
an etale cover having such a factor and no assertion
that the resulting maps jointly separate points.
Universal prime-to-$p$ supersingular correspondences
therefore remain the separate existence problem in A5.

\begin{claim}{A.50}{Proved; the closure operations in L6}
Let $\mathcal N$ be the class of smooth proper curves of genus at
least two over $\overline{\mathbf F}_p$ for which the point-section
map is injective over every finite field of definition and every
finite extension thereof. It has the following closure properties.
If $D\to C$ is a connected finite etale cover with prime-to-$p$
geometric Galois closure and $D\in\mathcal N$, then
$C\in\mathcal N$. If $C$ has a family of morphisms to curves
in $\mathcal N$ which jointly separates geometric points, then
$C\in\mathcal N$. One may also use positive-genus targets
with geometric $p$-rank zero.
\end{claim}
\begin{proof}
Suppose a collision exists on $C$ over a finite field. The finite
cover and a chosen geometric lift of $P$ are defined over a finite
extension $L$. Equality of the section classes restricts along
the injective map $G_L\hookrightarrow G_k$. In the finite
prime-to-$p$ fibre functor, equality of the two sections identifies
the Galois actions on the fibres. The fixed lift of $P$ therefore
gives a fixed lift of $Q$; choosing the corresponding paths gives
equal section classes upstairs. Equivalently this is the section
lifting step in R1, before abelianization. The geometric group
upstairs is the open subgroup with its full prime-to-$p$ topology:
cores of finite-index subgroups have prime-to-$p$ index because
the geometric Galois closure does. Thus the equality is in the
full geometrically prime-to-$p$ group of $D$, not merely a
quotient. The two lifts remain distinct because their images are
distinct. This contradicts $D\in\mathcal N$.

For the second assertion, a morphism of curves induces a
continuous map of fundamental extensions over the same Galois
group. It carries conjugate sections to conjugate sections.
Every target therefore identifies the two images of a hypothetical
collision. Joint separation contradicts their distinctness.
For a target with geometric $p$-rank zero, the Abel map and R1
give the required target injectivity directly.

These operations can be iterated in the tower. They do not
assert that every curve is obtained from the given seeds.
Indeed, the universal point-separating-cover statement of A5 is
equivalent to A1: this proof gives one implication, and, assuming
A1, the identity cover and identity morphism give the other.
The later A.88--A.90 prove A1, so the general
point-separating-cover assertion of A5 follows by this
identity-cover argument. The asserted universal correspondence
to a supersingular curve is a stronger sufficient construction,
whose existence is a separate question.
\end{proof}

\begin{claim}{A.51}{Proved; the exact Frobenius-orbit formula}
Let $C$ be a smooth proper genus-$g$ curve over an algebraically
closed field $\Omega$ of characteristic $p>0$, and let $L$ be a
line bundle of exact order $n>1$, with $(n,p)=1$. Let
$D\to C$ be the associated connected cyclic etale cover. For
$i\in\mathbf Z/n\mathbf Z$ put $V_i=H^1(C,L^i)$, with a fixed
trivialization of $L^n$. Absolute Frobenius gives $p$-semilinear
maps $F_i:V_i\to V_{pi}$. For an orbit $O$ of multiplication
by $p$ on $(\mathbf Z/n\mathbf Z)^\times$, choose $i\in O$
and let $r_O$ be the eventual rank of the iterates of
$F^{|O|}:V_i\to V_i$. Then
\[
                  f_p(P_n(D))=\sum_O |O|r_O.
\]
In particular the primitive part has $p$-rank zero if and only
if every such orbit composite is nilpotent. Here
$\dim_\Omega V_i=g-1$ for every nontrivial $i$.
\end{claim}
\begin{proof}
The finite etale algebra of the cyclic cover gives
\[
 H^1(D,\mathcal O_D)=\bigoplus_{i\in\mathbf Z/n\mathbf Z}V_i,
\]
and taking $p$th powers sends index $i$ to $pi$. Changing the
chosen trivialization only changes the displayed identifications
and does not change stable ranks. Riemann--Roch and
$H^0(C,L^i)=0$ give $\dim V_i=g-1$ for $i\ne0$.
For a finite-dimensional space with a Frobenius-semilinear
endomorphism, the stable image is the maximal subspace on which
Frobenius is bijective. Over $\Omega$, its dimension equals the
$\mathbf F_p$-dimension of the fixed space; this follows by
solving the semilinear equation $F(v)=v$ on the bijective part.
The Artin--Schreier sequence identifies that fixed space with
$H^1_{\mathrm{et}}(D,\mathbf F_p)$, whose dimension is $f_p(D)$.
On a character orbit the stable image has the same dimension
$r_O$ at each index, giving the factor $|O|$.

Since $n$ is invertible in characteristic $p$, the character
projectors split the cohomology. The endomorphism
$\prod_{\ell\mid n}(1-\sigma^{n/\ell})$ is nonzero precisely
on primitive characters. Its image therefore has exactly the
primitive stable ranks in the displayed formula. A descending
sequence of images of a semilinear endomorphism stabilizes in
finite dimension; its stable image is zero exactly when some
iterate is zero. This proves the criterion.

In genus two and characteristic two, the locally exact
differential sequence identifies $\ker F_i$ with
$H^0(C,\mathcal B\otimes L^i)$. Each $V_i$ is a line, so an
orbit composite is nilpotent precisely when one arrow is zero.
This is the theta-orbit condition in R5 and R8. For $g\ge3$,
noninvertibility of a single arrow is weaker than nilpotence
of the orbit composite.
\end{proof}

\begin{claim}{A.52}{Refuted; nonordinariness alone is not an annihilation argument}
The implication from a nonordinary character-cohomology block
to absence of surviving $p$-primary vectors is false. The failure
persists for a nonabelian group and a vector with free deck orbit
in the abstract module.
\end{claim}
\begin{proof}
Already the semilinear map
$(x,y)\mapsto(x^p,0)$ on $\overline{\mathbf F}_p^2$ has rank
one and stable rank one. It is noninvertible but is not
nilpotent. For a group example take
$G=\mathbf F_8^\times\rtimes\operatorname{Gal}(\mathbf F_8/\mathbf F_2)$
acting diagonally on $W=\mathbf F_8^2$. The vector $(1,\alpha)$,
where $\alpha^3+\alpha+1=0$, has trivial stabilizer: a group
element $(u,b)$ fixing it has $u=1$ from the first coordinate,
then $\alpha^{2^b}=\alpha$, whence $b=0$.
On $W\oplus\mathbf F_8$ take the Frobenius-semilinear map
which is the identity on the $\mathbf F_2$-structure of $W$
and zero on the last summand. It is $G$-equivariant, has
dimension nine, stable rank six, and retains the fixed vector
with free $G$-orbit. The program
\texttt{code/geometric\_checks.py} verifies the group law,
all 9261 associativity triples, the free orbit, and the ranks.
This is a refutation of the proposed linear weakening, not a
curve with a collision. The correct condition is stable-rank
zero on the relevant summand, as in A.51.
\end{proof}

\begin{claim}{A.53}{Proved; a small integral group-ring obstruction}
Let $\Omega$ be algebraically closed of characteristic $p$, and
let $\pi:D\to C$ be a connected finite etale Galois cover of
smooth proper curves, with $g(C)\ge2$ and group $G$ of order
$m$ prime to $p$. Let
\[
 a=\sum_{g\in G}a_g g\in\mathbf Z[G]\setminus\{0\},
 \qquad d=\sum_g|a_g|.
\]
Assume $(d,m)=1$ and $m>(d-1)^2$. If $G$ is nonabelian,
then for any points $P',Q'\in D(\Omega)$ lying above distinct
points of $C$ one has
\[
                         a[Q'-P']\ne0\quad\hbox{in }J_D(\Omega).
\]
\end{claim}
\begin{proof}
Suppose the class is zero. Since the two free $G$-orbits are
disjoint, the divisor $\sum a_g(gQ'-gP')$ is a nonzero
principal divisor with positive and negative degrees both
equal to $d$. Let $f\in\Omega(D)$ have that divisor, and put
\[
 M=\Omega(\gamma f:\gamma\in G),\qquad
 e=[\Omega(D):M].
\]
Let $Z$ be the smooth proper curve with function field $M$.
Every generator has degree $d$. Poonen's Lemma 2.2 and
Corollary 2.3, with Proposition 2.1
\cite{AG_Poonen2007}, give
\[
                 e\mid d,\qquad g(Z)\le(d/e-1)^2.
\]
For clarity, the induction proving the bound is as follows.
If the current field has index $d/s$ and genus at most
$(s-1)^2$, adjoining another degree-$d$ function with relative
degree $b$ gives index $d/(bs)$ and genus at most
\[
 b(s-1)^2+(b-1)(bs-1)
   =(bs-1)^2-bs(b-1)(s-1)\le(bs-1)^2.
\]
This is Castelnuovo--Severi for the two generating subfields.
The perfectness hypothesis holds because $\Omega$ is
algebraically closed. The result allows arbitrary rational
functions, so no unproved separability assertion is needed.

Since $g(D)=m(g(C)-1)+1\ge m+1>(d-1)^2$, we have $e>1$.
The field $M$ is $G$-stable. The kernel of the action on $M$
has order dividing both $e$ and $m$, so the action on $Z$ is
faithful. It is free: if a point stabilizer contained a subgroup
of prime order $\ell$, that subgroup would act freely on the
degree-$e$ fibre divisor in $D$, partitioning its points into
orbits of size $\ell$ with constant multiplicity. This would
force $\ell\mid e$, contradicting $(e,m)=1$.

Consequently $Z\to Z/G$ is etale and
$g(Z)-1=m(g(Z/G)-1)$. The numerical bound excludes
$g(Z)\ge m+1$ and genus zero is impossible in that equality;
hence $g(Z)=1$. Choose an origin on $Z$. An automorphism is
$x\mapsto t+u(x)$ with $u$ an origin-preserving automorphism.
If $u\ne1$, the nonzero endomorphism $1-u$ of the elliptic
curve is surjective, so this automorphism has a fixed point.
Every element of a free action is therefore a translation.
Translations commute, contradicting that $G$ is nonabelian.
\end{proof}

\begin{claim}{A.54}{Proved; rational and geometric nonabelian criteria}
In A.53 suppose all cover data are defined over a finite field
$L$, $G$ is constant over $L$, and $P'\in D(L)$ lies above
$P\in C(L)$. Put $B_a=\operatorname{im}(a:J_D\to J_D)$,
an abelian subvariety over $L$. If $B_a(L)[p]=0$, then $P$
collides with no other point of $C(L)$. The same conclusion
holds over all finite fields of definition if $f_p(B_a)=0$
geometrically. Thus an odd nonabelian group can use, for
example, $a=(1-\delta)(1-\epsilon)$ whenever the four support
elements are distinct and $m>9$.
\end{claim}
\begin{proof}
A collision lifts by R1 to a rational nonzero $p$-primary
class $T=[Q'-P']$. Its image $aT$ is a rational $p$-primary
point of $B_a$, hence zero under the rational hypothesis.
This contradicts A.53. In the geometric version one first
defines the cover, its automorphisms, and the lift over a
finite extension; $f_p(B_a)=0$ ensures that the torsion
vanishing persists over that extension. For the displayed
four-term element, $d=4$ and an odd group has $(d,m)=1$.

More generally, for any collection of integral endomorphisms
$a_j$, a joint map $D\to\prod_j B_{a_j}$ separating the two
fibres gives a collision obstruction if every target has no
rational $p$-torsion over the field where the lifts are rational.
A.53 is a concrete condition that supplies this separation
for one endomorphism. It does not supply existence of a
cover with the required torsion vanishing.
\end{proof}

\begin{claim}{A.55}{Proved; the Conway--McKay--Trojan test is outside the category}
The polynomial covers $t=x^{24}+x$ in characteristic two and
$t=x^{12}+x$ in characteristic three from
Conway--McKay--Trojan \cite{AG_CMT2008} are finite etale over
the affine line but are wildly ramified on their proper
compactifications. Their root-space representations do not
give instances of the prime-to-$p$ etale criteria above.
\end{claim}
\begin{proof}
Both affine derivatives are one. At infinity the ramification
indices are respectively 24 and 12, divisible by the
characteristic. A Galois closure therefore has nontrivial
wild inertia and its group has order divisible by $p$.
The source and base of the non-Galois polynomial maps are
projective lines, outside the genus hypothesis as well.

For the characteristic-two polynomial put
\[
 L(x)=x^{2048}+t^{64}x^{512}+t^8x^{16}+t^{16}x^8+x.
\]
Reduction modulo $x^{24}+x+t$ gives $L(x)=t+t^{24}$.
Consequently the roots satisfy the additive polynomial
$L(x)^2+(t+t^{24})L(x)$. It is separable of degree $2^{12}$,
because its coefficient of $x$ is the nonzero rational
function $t+t^{24}$. Its root space has dimension 12 over
$\mathbf F_2$, and the Galois action on the subspace spanned
by the original roots is faithful. The script
\texttt{code/geometric\_checks.py} verifies the polynomial
remainders and all additive exponents exactly. None of these
linear spaces is identified with the $p$-torsion of a
Jacobian. In particular faithfulness of this representation
cannot substitute for the geometric point-difference
condition in A.53.
\end{proof}

\begin{claim}{A.56}{Known, checked; the orbifold theorem covers the required class}
The pro-$\mathscr S$ finite-subgroup theorem used in AP-A.29l
applies to $\mathscr S$ equal to all finite groups of order
prime to $p$. Thus this particular cited step in R4 has the
required scope.
\end{claim}
\begin{proof}
Boggi--Zalesskii Definition 1.1 defines a class by closure under
subgroups, quotients, and extensions. Their Lemma 3.1 assumes
a hyperbolic Fuchsian group $\Lambda$ with a torsion-free
normal subgroup $H$ and $\Lambda/H\in\mathscr S$. It gives
the malnormality of the nontrivial cyclic cone groups and
containment of every finite nontrivial subgroup of the
pro-$\mathscr S$ completion in a conjugate cone group
\cite{AG_BoggiZalesskii2019}. All three closure conditions
hold for prime-to-$p$ groups.

For a prime-to-$p$ Galois map $C\to B$ with $g(C)\ge2$,
the orbifold presentation has cone orders equal to the
ramification indices. Its surjection to the deck group is
injective on each cone group, so its kernel is the
torsion-free surface group of the topological covering.
Its orbifold Euler characteristic is
$(2-2g(C))/|G|<0$. These are precisely the lemma's
hypotheses. We have read the proof, which uses the
$\mathscr S$-solenoid and acyclicity at primes occurring
in $\mathscr S$; no prime-$p$ acyclicity is needed.
This verifies the completion comparison, not every other
step of R4 by mere citation.
\end{proof}

\begin{claim}{A.57}{Open for supersingular correspondences; nontrivial primitive-cover existence is refuted}
The universal supersingular-correspondence assertion of A5
is not established here. The cited geometric existence
theorems do not imply the universal primitive-zero cover
assertion A6; the later explicit argument A.80--A.81
refutes its degree-$n>1$ version in genus two and
characteristic two. The literal version permitting the
identity cover remains open for curves of positive $p$-rank.
\end{claim}
\begin{proof}[Exact obstruction retained]
The nilpotence requirement is A.51. Raynaud's Theorem 2
provides a solvable prime-to-$p$ cover with a nonordinary
representation \cite{AG_Raynaud2000}; noninvertibility
does not imply nilpotence, as A.59 proves on an actual cover. The theta torsion
density of Yang Proposition 2.3 and Theorem 2.4
\cite{AG_Yang2016}, and the related Pop--Saidi propositions
\cite{AG_PopSaidi2003}, do not prescribe a torsion order for
which every primitive multiplication-$p$ orbit is met.
In genus two, characteristic two, a single orbit can be
killed while other primitive orbits survive. In higher
genus even one orbit can retain positive stable rank after
some Frobenius arrows become singular.

Ozman--Pries Theorem 1.1 concerns generic curves and ordinary
Pryms; Example 9.4 excludes degree-three witnesses for
nonordinary genus-two curves in characteristic two
\cite{AG_OzmanPries2019}. Generic counterexamples over a
field of positive transcendence degree do not give a
counterexample over $\overline{\mathbf F}_p$: the latter
has only countably many moduli points, and a countable
union of proper closed loci can contain all of them.
No countability argument is used in the reverse direction.

Bogomolov--Tschinkel Theorem 1.7 assumes a hyperelliptic
source and $p\ge5$; their Proposition 1.8 assumes
characteristic different from two and three
\cite{AG_BogomolovTschinkel2002}. Poonen Theorem 1.7 has
a prime-to-$p$ Galois action with quotient genus at most
two, together with one of its additional conditions
\cite{AG_Poonen2005}. These hypotheses exclude a general
curve, and an arbitrary etale cover in those statements
must additionally be checked for prime-to-$p$ Galois
closure before applying the collision lifting argument.
The searches and subsequent attempts are recorded in the
source audit; these citations are not substitutes for the
missing existence proof.
\end{proof}

\begin{claim}{A.58}{Proved; explicit cyclic zero-rank families in characteristics two and three}
Let $p=2$ or $3$, and let $n>1$ be odd and prime to $p$, with
$(p-1)(n-1)\ge4$. Suppose that
\[
                    (p-1)n\mid p^r+1
\]
for some positive integer $r$. The smooth proper curve
\[
                         C_{p,n}:y^p-y=x^n
\]
has a connected cyclic etale cover of degree $n$ whose entire
Jacobian, and therefore its primitive conductor-$n$ factor,
has $p$-rank zero. For $p=2$ the cover is the Fermat curve
$u^n+z^n=1$; for $p=3$ it is the smooth model of
$u^n=z^{2n}-1$. In both cases the map is
\[
                         x=uz,\qquad y=z^n.
\]
Thus these are actual cases of A6 in both small characteristics,
without a separability or Galois-closure qualification left over.
\end{claim}
\begin{proof}
The equation for $C_{p,n}$ has one point at infinity. At that
point $x$ has pole order $p$ and $y$ pole order $n$:
the valuation equation gives $p\,v(y)=n\,v(x)$, and the
Artin--Schreier extension has degree at most $p$, forcing
degree and ramification index $p$. Its different exponent is
$(p-1)(n+1)$, hence its genus is $(p-1)(n-1)/2$.
At $P=(0,0)$, the variable $x$ is a parameter and $y$ has
order $n$. Consequently
\[
                          \operatorname{div}(y)=n(P-\infty).
\]
The class $[P-\infty]$ has exact order $n$. Indeed, if its
order were a proper divisor $d<n$ of $n$, a function with
only pole $d\infty$ would exist. Every function regular away
from infinity is a sum of monomials $x^a y^b$ with $a\ge0$
and $0\le b<p$. Their pole orders $pa+nb$ are pairwise
distinct when grouped with fixed $b$ and then compared
modulo $p$. A pole order less than $n$ must therefore be
divisible by $p$, whereas $d\mid n$ is prime to $p$.

Adjoining $z$ with $z^n=y$ thus gives a connected cyclic
degree-$n$ extension. It is unramified because every
valuation of $y$ is divisible by $n$ and $n$ is invertible;
locally a unit has an $n$th root in the strict henselization.
Putting $u=x/z$ gives
\[
                           u^n=z^{(p-1)n}-1.
\]
The constant deck group $\mu_n$ acts by
$(u,z)\mapsto(\zeta^{-1}u,\zeta z)$ after a finite field
extension containing $\mu_n$. This is a cyclic geometric
Galois cover of order prime to $p$, so its own Galois closure
already has the required order.

We give the zero-rank argument to avoid relying on an
unstated supersingularity criterion for Fermat curves. Put
$Q=p^r$. The degree-$Q+1$ Fermat curve is isomorphic over
the algebraic closure to
\[
                        w^Q-w=t^{Q+1}.
\]
To see this, choose $\beta\in\mathbf F_{Q^2}\setminus\mathbf F_Q$
and $\gamma\in\mathbf F_{Q^2}$ with $\gamma^{Q+1}=-1$.
The substitutions
$U=W+\beta Z$ and $V=\gamma(W+\beta^Q Z)$ give
\[
 U^{Q+1}+V^{Q+1}
       =(\beta-\beta^Q)(W^QZ-WZ^Q).
\]
Rescaling the remaining coordinate gives the displayed
Artin--Schreier form. Its projection to the $t$-line has
additive group $\mathbf F_Q$ and is etale away from infinity.
The valuation relation at infinity, with
$\gcd(Q,Q+1)=1$, proves irreducibility, degree $Q$, and
total ramification there. Crew's Corollary 1.8
\cite{AG_Crew1984}, for a finite $p$-group acting on a
smooth proper curve over an algebraically closed field,
states the Deuring--Shafarevich formula
\[
 f_p(H)-1=|G|\bigl(f_p(H/G)-1\bigr)
                       +\sum_{R\in H}(e_R-1).
\]
Here it gives $f_p(H)-1=-Q+(Q-1)=-1$.

Set $N=(p-1)n$. Since $N\mid Q+1$, the coordinate-power
map from the degree-$Q+1$ Fermat curve to the degree-$N$
Fermat curve is finite separable of degree
$((Q+1)/N)^2$, prime to $p$. The latter Fermat curve maps
to our cover by
\[
                     (w,z)\longmapsto (w^{p-1},z)
\]
in the equation $w^N=z^N-1$. Its degree is $p-1$, also
prime to $p$. Pullback followed by norm is multiplication
by the degree; on geometric $p$-torsion this is invertible.
Zero $p$-rank therefore descends through both maps.
Every abelian subvariety of this zero-rank Jacobian has
zero $p$-rank, including the primitive factor.

The restrictions describe infinite families: one may take
$n=2^r+1$ for $p=2$, $r\ge2$, and
$n=(3^{2s}+1)/2$ for $p=3$, $s\ge1$. They do not cover
arbitrary curves of those genera.
\end{proof}

\begin{claim}{A.59}{Refuted on an actual curve cover: nonordinary primitive blocks need not be nilpotent}
In characteristic two, let
\[
 C:y^2+y=x^7,\qquad
 D:u^7+z^7=1,\qquad
 (x,y)=(uz,z^7).
\]
Then $D\to C$ is connected cyclic etale of degree seven,
$g(C)=3$, $g(D)=15$, and $f_2(C)=0$, $f_2(D)=6$.
Each nontrivial deck-character space has coherent
dimension two and stable Frobenius dimension one. In
particular both primitive doubling-orbit composites are
noninvertible, but neither is nilpotent. The primitive
factor has dimension twelve and two-rank six.
\end{claim}
\begin{proof}
The construction and connectedness were proved in A.58
before its congruence hypothesis was used. Alternatively,
the action
$(u:z:w)\mapsto(\zeta^{-1}u:\zeta z:w)$ on the Fermat
curve is free: in the affine chart a nonidentity element
cannot fix both nonzero coordinates, and at infinity
both $u,z$ are nonzero and fixedness requires
$\zeta^{-1}=\zeta$, impossible for nontrivial seventh
roots. Its invariants $x=uz$, $y=z^7$ give precisely the
displayed quotient. The genus formulas follow either
from these equations or from etale Riemann--Hurwitz.

We compute the Cartier operator directly. A basis of
$H^0(D,\Omega_D^1)$ is
\[
 \omega_{a,b}=u^{a-1}z^{b-7}\,du,\qquad
      a,b\ge1,\quad a+b<7.
\]
This is the usual adjunction basis
$u^i z^j\,du/z^6$, $i,j\ge0$, $i+j\le4$;
there are fifteen forms. For the separating variable $u$,
Cartier sends $(A^2+B^2u)\,du$ to $B\,du$.
Using $z=(1+u^7)/z^6$ gives
\[
 \mathcal C(\omega_{a,b})=
 \begin{cases}
 0,&a,b\text{ both odd},\\
 \omega_{a',b'},&\text{otherwise},
 \end{cases}
 \qquad
 a'\equiv4a,\quad b'\equiv4b\pmod7,
\]
where $a',b'$ are in $\{1,\ldots,6\}$. In the second
case their sum is less than seven. For example, if
$a,b$ are both even the target is $(a/2,b/2)$;
if $a$ is odd and $b$ even it is $((a+7)/2,b/2)$;
the remaining case is $(a/2,(b+7)/2)$.

The stable basis consists exactly of
\[
             (1,2),(2,4),(4,1),\qquad
             (2,1),(4,2),(1,4),
\]
which form two Cartier cycles. All other basis vectors
eventually map to zero. The deck character of
$\omega_{a,b}$ is $b-a$ modulo seven. The invariant
space has dimension three and stable rank zero; every
nonzero character has dimension two and stable rank one.
Cartier and coherent Frobenius are dual, with opposite
character conventions, so they have the same orbitwise
stable ranks. The Cartier fixed-space description of
geometric two-torsion, equivalently the Artin--Schreier
description in A.51, now gives all stated two-ranks.

The exact graph of all fifteen basis vectors and both
multiplicity tables are generated and checked by
\path{code/fermat_cyclic_checks.py}, with its output
attached. This counterexample concerns the proposed
weakening in L8. It does not assert that $C$ has no other
cyclic cover satisfying A6.
\end{proof}

\begin{claim}{A.60}{Proved; a primitive-only zero-rank example}
In characteristic two, the cyclic etale degree-$35$ cover
\[
 u^{35}+z^{35}=1\ \longrightarrow\
       C:y^2+y=x^{35},\qquad (x,y)=(uz,z^{35}),
\]
has $g(C)=17$, $g(D)=561$ and $f_2(D)=6$, while its
primitive conductor-$35$ factor has dimension $384$ and
two-rank zero.
\end{claim}
\begin{proof}
The connected etale construction is again the one proved in
A.58. The Cartier basis and its calculation in A.59 apply
with every occurrence of seven replaced by any odd $n$.
Explicitly, for $a,b\ge1$ and $a+b<n$ the operator is zero
if both indices are odd, and otherwise replaces them by
their least positive halves modulo $n$. The deck weight
is $b-a$ modulo $n$. This also proves the finite formula
\[
 f_2(P_n)=\#\left\{(a,b):
 \begin{array}{l}
 1\le a,b<n,\quad \gcd(b-a,n)=1,\\
 {}[2^j a]_n+[2^j b]_n<n
       \text{ for every }0\le j<\operatorname{ord}_n(2)
 \end{array}\right\},
\]
where $[\,\cdot\,]_n$ is the representative in
$\{1,\ldots,n-1\}$. A Cartier basis vector belongs to the
stable image exactly when its entire finite halving orbit
stays in the holomorphic basis.

For $n=35$, the complete graph of the $561$ basis vectors
has stable set
\[
 (5,10),(10,20),(20,5),\qquad
 (10,5),(20,10),(5,20).
\]
The attached program verifies every transition and iterates
the entire image set until it stabilizes; its output includes
the full graph, so the finite assertion is reproducible.
Every surviving weight is divisible by five and is therefore
nonprimitive modulo $35$. A.51 gives primitive two-rank
zero. Its dimension is
$\varphi(35)(g(C)-1)=24\cdot16=384$.
This example also distinguishes primitive vanishing from
vanishing of the complete relative factor. It does not reach
a pair beyond A3: all surviving characters have conductor
seven, so the intermediate cyclic degree-five cover has
zero full new two-rank. After defining this cover and a
chosen lift over a finite field, it is itself an L1 witness.
\end{proof}

\begin{claim}{A.61}{Proved; the elliptic rational lever in characteristics two and three}
Let $E/L$ be an elliptic curve over a finite field of
characteristic two. Then $E(L)[2]=0$ if and only if $E$ is
supersingular. Thus an ordinary elliptic target does not
weaken the geometric zero-rank requirement in L5 when $p=2$.
In characteristic three, if $E$ is ordinary and $E(L)[3]=0$,
this vanishing persists over every odd-degree extension of
$L$ and fails over every even-degree extension.
\end{claim}
\begin{proof}
For an ordinary elliptic curve in characteristic $p$, the
geometric group of $p$-torsion points is cyclic of order $p$.
For $p=2$ its unique nonzero point is fixed by every
automorphism of the algebraic closure over $L$, so it is
rational. A supersingular elliptic curve has no nonzero
geometric $p$-torsion points. For $p=3$, Frobenius on the
one-dimensional $\mathbf F_3$-space is either $1$ or $-1$.
The stated rational vanishing selects $-1$, whose $d$th
power is $1$ exactly for even $d$.
\end{proof}

\begin{claim}{A.62}{Proved; a genus-lowering closure step in characteristic two}
Let $C/\overline{\mathbf F}_2$ be a hyperelliptic curve of
genus $g\ge3$. There is a connected cyclic etale cover
$D\to C$ of degree three and a smooth proper curve $B$
of genus $g-1$ with two morphisms $D\to B$ which jointly
separate geometric points. Consequently, if every genus-$(g-1)$
curve has no collision over its finite fields of definition,
then $C$ has no collision. In particular universal A1 in
genus two and characteristic two would imply A1 for every
hyperelliptic genus-three curve in that characteristic.
\end{claim}
\begin{proof}
Choose a nontrivial order-three character of the geometric
abelianized fundamental group of $C$ and let $D\to C$ be
its connected cyclic cover, with generator $\sigma$.
Choose a hyperelliptic ramification point $W$ and a lift
$W'$ to $D$. The hyperelliptic involution $\iota$ acts by
$-1$ on $J_C$, so it preserves the character kernel and
acts by inversion on the cyclic quotient. It lifts to an
automorphism $\tau$ fixing $W'$. Its square is a deck
transformation fixing $W'$, hence is the identity.
Thus
\[
            \tau^2=1,\qquad \tau\sigma\tau=\sigma^{-1}.
\]
Let $B=D/\langle\tau\rangle$. On the three-point fibre
above each hyperelliptic ramification point, $\tau$ is an
affine inversion and has exactly one fixed point. It has
no fixed points above unramified points of $C\to\mathbf P^1$.
Because $D\to C$ is etale, the completed local rings at
these fixed points, together with the involution, identify
with the completed local rings downstairs. Hence the
different for $D\to B$ has the same total degree as the
different for $C\to\mathbf P^1$, namely $2g+2$.
Etale Riemann--Hurwitz gives $g(D)=3g-2$, and the
separable degree-two quotient formula then gives
\[
 6g-6=2(2g(B)-2)+(2g+2),\qquad g(B)=g-1.
\]
This computation retains the wild different in characteristic
two; tame ramification counting would be incorrect here.

Let $\phi:D\to B$ be the quotient map. The maps $\phi$
and $\phi\circ\sigma$ jointly separate points. Indeed, if
$R\ne S$ have the same image under both, then
$S=\tau R=\sigma^{-1}\tau\sigma R$. The product of
these two distinct involutions is a nonidentity power of
$\sigma$, which would fix $R$, contradicting the etaleness
of the cyclic cover. Every object and map descends after
a finite extension of a finite field. Apply A.50 to the
two maps and then to $D\to C$.

For $g=2$ the same construction has an elliptic target.
It only supplies the desired vanishing when that elliptic
target has two-rank zero, as A.61 explains; the reduction
does not by itself settle the genus-two base case.
\end{proof}

\begin{claim}{A.63}{Proved; exact small-prime theta tests in a supersingular genus-two family}
Over $\overline{\mathbf F}_2$ let
\[
                      C_a:y^2+y=x^5+a x^3.
\]
Every genus-two curve of two-rank zero is geometrically
isomorphic to some $C_a$.
Its Raynaud theta divisor, based at its point $\infty$,
is $\Theta_a=\{[R-\infty]:R\in C_a\}$.
It has no nonzero point of order three. It has a point
of exact order five if and only if $a=0$, and a point
of exact order seven if and only if $a^5=1$.
Consequently:
\begin{enumerate}
\item for $a=0$ a degree-five cyclic etale cover has
      zero primitive two-rank;
\item for $a^5=1$ a degree-seven cyclic etale cover has
      zero primitive two-rank, while every degree-three
      or degree-five cyclic cover is new-ordinary;
\item if $a\ne0$ and $a^5\ne1$, every geometric cyclic
      cover of degree three, five or seven is new-ordinary.
\end{enumerate}
These are geometric assertions. The last assertion alone
does not exclude rational L1 witnesses of those degrees
over a particular field.
\end{claim}
\begin{proof}
For the normal-form assertion, the genus-two hyperelliptic
map is a separable quadratic map to $\mathbf P^1$, hence
has an Artin--Schreier equation. Put its rational function
in reduced form, with odd pole orders. The
Deuring--Shafarevich formula makes the two-rank one less
than the number of poles. Thus there is one pole, which
Riemann--Hurwitz forces to have order five. Move it to
infinity, scale $x$, and remove the constant and even
powers by changing $y$; the equation becomes
$y^2+y=x^5+a x^3+b x$. Translation $x=u+r$ followed
by changes of $y$ by a polynomial of degree two removes
the coefficient of $u$ whenever
\[
                   r^{16}+a^4r^8+a^2r^2+r=b^4.
\]
This monic polynomial has a root over the algebraic
closure. The coefficient of $u^3$ remains $a$,
proving the assertion.

The Artin--Schreier map has a single pole, of order five,
so the genus is two and the two-rank is zero. The different
at infinity has order six; hence $\operatorname{div}(dx)=2\infty$.
Since $dx$ is locally exact, its half-divisor identifies
the Raynaud bundle with $\mathcal O(\infty)$, giving
the asserted theta origin.

Every function regular away from infinity is
$A(x)+B(x)y$, and $x,y$ have pole orders two and five.
If $[R-\infty]$ has exact odd order $\ell$, it is nonzero,
$R$ is affine, and a function $F$ has
$\operatorname{div}(F)=\ell R-\ell\infty$.
There is no function with unique pole order three, proving
the first assertion.

For $\ell=5$, normalize $F=A+y$ with $\deg A\le2$.
Its hyperelliptic norm is
\[
             A^2+A+x^5+a x^3=(x+r)^5,\qquad r=x(R).
\]
The coefficient of $x^3$ forces $a=0$. Conversely,
when $a=0$, the divisor of $y$ is
$5(0,0)-5\infty$, so that point gives exact order five.

For $\ell=7$, write
\[
 F=A+(x+d)y,\qquad
 A=c_3x^3+c_2x^2+c_1x+c_0.
\]
The leading coefficient of $x y$ is nonzero because the
pole order is seven, so this normalization loses no
possible function. The norm condition is
\[
 A^2+(x+d)A+(x+d)^2(x^5+a x^3)=(x+r)^7.
\]
Write $a=t^2$ and $r=s^2$, uniquely in the algebraic
closure. Comparing coefficients from degree six down to
degree three gives
\[
 c_3=s,\quad d=s^2+t,\quad
 c_2=s^8+t^2s^4+s^3+st+t^4,
\]
and
\[
             E_1=s^{16}+t^4s^8+t^2s^2+s+t^8=0.
\]
The coefficients of degrees two and one give
$c_1^2+c_1=r^5+c_2d$ and $c_0=r^6+c_1d$.
Substituting these in the constant coefficient gives
\[
             E_2=s^{24}+t^2s^{10}+s^9+t^6s^2+t^4s+t^7=0.
\]
The complete elimination is the polynomial identity
\[
                     E_2+(s^8+t^4)E_1=t^7(t^5+1).
\]
It is verified symbolically in
\path{code/supersingular_small_primes.py}; the output also
checks the resultant
$\operatorname{Res}_s(E_1,E_2)=t^{112}(t^5+1)^{16}$.

If $t\ne0$, these equations have a solution precisely
when $t^5=1$: necessity follows from the displayed
identity, and sufficiency follows by taking any root of
the monic degree-sixteen polynomial $E_1$ and then solving
the Artin--Schreier equation for $c_1$. Moreover
$(x+d)|_{x=r}=t\ne0$. The norm has its only finite zero
at $r$, but $F$ cannot vanish at both of the two points
over $r$, because its two values differ by $r+d=t$.
Its divisor is therefore exactly $7R-7\infty$ for one
of those points. This proves existence of a point of
order seven when $a^5=1$.

If $t=0$, the coefficient equations instead give $d=r$.
The norm condition then forces $A(r)=0$, so $F$ vanishes
at both hyperelliptic conjugate points above $r$. They
are distinct, since the affine Artin--Schreier map is
etale. Such an $F$ cannot have divisor $7R-7\infty$.
This excludes the extraneous zero of the resultant at $a=0$.

Finally theta is symmetric: the hyperelliptic involution
sends $[R-\infty]$ to its negative. For degree five,
two generates $(\mathbf Z/5)^\times$, so one theta point
kills the sole primitive doubling orbit. For degree
seven its two doubling orbits are exchanged by negation.
A theta point and its negative therefore meet both
orbits. A.51 proves zero primitive rank in the asserted
cases. Absence of theta at one of these prime orders
makes every relevant Frobenius arrow invertible, proving
the stated new-ordinariness assertions.
\end{proof}

\begin{claim}{A.64}{Proved; the degree-eleven theta locus in the same family}
For $C_a$ in A.63, the Raynaud theta divisor contains
a point of exact order eleven if and only if
\[
                           a^{35}+a^{10}+1=0.
\]
At each such parameter there is a connected cyclic
etale degree-eleven cover with zero full new two-rank.
At every other parameter all degree-eleven cyclic
covers are new-ordinary.
\end{claim}
\begin{proof}
At an affine point, translate its $x$-coordinate to
$u=0$ and its $y$-coordinate to zero. Removing the
even powers by adding a quadratic polynomial to $y$
gives the local and global affine equation
\[
                        Y^2+Y=u^5+a u^3+c u.
\]
The pole orders of $u,Y$ at infinity are still two
and five. Every $c$ can occur for fixed $a$: if the
translated point has original coordinate $r$, the
coefficient calculation is
$c^4=r^{16}+a^4r^8+a^2r^2+r$, a monic degree-sixteen
polynomial in $r$. Its value can be any prescribed
element over the algebraic closure.

On the branch with $Y(0)=0$, the formal expansion to
the necessary order is
\[
 Y=cu+c^2u^2+a u^3+c^4u^4+u^5
                       +a^2u^6+c^8u^8+u^{10}+O(u^{11}).
\]
A function with sole pole order eleven has the form
$A(u)+B(u)Y$, with $\deg A\le5$ and
$B=u^3+b_2u^2+b_1u+b_0$. Choose $A$ to cancel terms
of degree at most five. Writing $v=c^4$, the necessary
and sufficient conditions for vanishing to order at
least eleven are the five coefficients
\[
 \begin{aligned}
 b_0a^2+b_1+b_2v+a&=0,&
 b_1a^2+b_2+v&=0,\\
 b_0v^2+b_2a^2+1&=0,&
 b_1v^2+a^2&=0,&
 b_0+b_2v^2&=0.
 \end{aligned}
\]
If $v=0$, the fourth equation gives $a=0$ and the
third becomes $1=0$. Thus $v\ne0$. Eliminate
$b_1=a^2/v^2$, $b_2=a^4/v^2+v$ and
$b_0=a^4+v^3$. The remaining equations are
\[
 \begin{aligned}
 E_1&=v^7+a^4v^4+a^2v^3+v^2+a^6=0,\\
 E_2&=a^2v^5+v^4+(a^6+a)v^2+a^4v+a^2=0.
 \end{aligned}
\]
The exact resultant is
\[
              \operatorname{Res}_v(E_1,E_2)
                        =a^4(a^{35}+a^{10}+1).
\]
The program \path{code/supersingular_eleven.py}
verifies all five substitutions and computes this
resultant over $\mathbf F_2[a]$, with output attached.
When $a=0$, $E_2=v^4$ excludes every nonzero $v$.
When $a\ne0$, both displayed leading degrees are
unchanged, so resultant zero is equivalent to a common
root. The constant coefficient $a^6$ of $E_1$
ensures that any such root is nonzero.

For each solution, the constructed function has a
pole of exactly order eleven, and a zero of order
at least eleven at the selected point. Thus its
divisor is precisely eleven times the difference
between that point and infinity. The point is
distinct from infinity, so its theta class has
exact prime order eleven. Conversely any such
theta point supplies the coefficient equations
above. This proves the equivalence, including
the extraneous resultant factor at $a=0$.
Since two has order ten modulo eleven, A.51 gives
the final two-rank conclusions.
\end{proof}

\begin{claim}{A.65}{Proved; even deck groups in odd characteristic}
Let $\Omega$ be algebraically closed of odd characteristic,
and let $D\to C$ be a connected etale Galois cover of
smooth proper curves with $g(C)\ge2$ and geometric group
$G$ of order $m$ prime to the characteristic. Suppose
$m\nmid4$. For points $P',Q'\in D(\Omega)$ in distinct
fibres, the class $[Q'-P']$ has trivial stabilizer in $G$.
In particular this holds for every nonabelian $G$,
including groups of even order.

Consequently, over a finite field of odd characteristic,
the full relative class-number criterion L1 holds for a
pointed geometrically Galois etale cover with such a group.
The deck group need not be constant: the geometric average
is the rational pullback--norm projector. This supplies
an even-group version of the nonabelian criterion, without
asserting that its existence reaches a pair outside A3.
\end{claim}
\begin{proof}
If $1\ne\delta\in G$ fixes $[Q'-P']$, then
\[
           \delta Q'+P'-\delta P'-Q'
\]
is principal. Its four points are distinct, since the
deck action is free and the two original points lie in
distinct fibres. Thus its function has degree two.
Since $g(D)\ge2$, this is the hyperelliptic map. It is
separable because the characteristic is odd.

The hyperelliptic pencil is unique, so every automorphism
of $D$ preserves its ramification-point set $W$.
There are exactly $2g(D)+2$ such geometric points in odd
characteristic. The action of $G$ on $W$ is free, because
$D\to C$ is etale. Hence
\[
 m\mid\#W=2g(D)+2=2m\bigl(g(C)-1\bigr)+4,
\]
contradicting $m\nmid4$. This proves the stabilizer assertion.

For the rational criterion, a collision lifts to a rational
$p$-primary point-difference class $T$ on $D$ by L0,
or by its pointed nonconstant-deck form A.21. The
prime-to-$p$ isogeny projector
$m^{-1}\pi^*\pi_*$ splits its invariant and new parts.
If $p$ does not divide $h_L(D)/h_L(C)$, the rational new
part is zero. Thus $T$ is fixed by every geometric deck
element, contrary to the stabilizer assertion. The
pullback--norm operator is defined over $L$, so the same
argument covers nonconstant geometric deck groups.

The divisibility restriction in the geometric argument is
real: a hyperelliptic curve can carry a free group of
order two or four. No claim about arbitrary point-difference
stabilizers in those degrees is being inferred. For odd
groups the Feit--Thompson reduction A.23 applies; an even
group need not have an odd-order abelian quotient, so that
specific reduction does not extend automatically.
\end{proof}

\begin{claim}{A.66}{Proved; primitive-only vanishing already in genus two}
For the characteristic-two family $C_a$ in A.63,
the Raynaud theta divisor contains a point of exact
order nine if and only if $a^5=1$.
At each such parameter the associated connected
cyclic etale degree-nine cover has genus ten and
two-rank two. Its primitive conductor-nine factor
has dimension six and two-rank zero, whereas its
degree-three intermediate cover is new-ordinary.
Thus primitive-only vanishing can occur with genus-two
base, not just in the genus-seventeen example A.60.
\end{claim}
\begin{proof}
Use the local equation
$Y^2+Y=u^5+a u^3+c u$ at a selected affine point as
in A.64; every value of $c$ can occur for fixed $a$.
A function with sole pole order nine has the form
\[
 F=A(u)+(u^2+b_1u+b_0)Y,\qquad \deg A\le4.
\]
Choose $A$ to cancel the coefficients through degree
four. With $v=c^4$, vanishing through degree eight
is precisely the following complete system:
\[
 \begin{aligned}
 b_0+b_1v+a&=0,& b_0a^2+b_1+v&=0,\\
 b_1a^2+1&=0,& b_0v^2+a^2&=0.
 \end{aligned}
\]
The third equation excludes $a=0$ and gives
$b_1=a^{-2}$. The first then gives
$b_0=v/a^2+a$. The remaining two equations become
\[
                  a^5+1=0,
       \qquad v^3+a^3v^2+a^4=0.
\]
These exact substitutions are independently checked in
\path{code/supersingular_nine.py}. If $a^5=1$, the
second polynomial is monic and has a root over the
algebraic closure. Construct $c$, the selected point,
and $F$ from that root. The function has pole exactly
nine at infinity and vanishes at least nine times at
the selected point. Its divisor is therefore exactly
$9R-9\infty$. The class has order nine: it is nonzero,
and order three is impossible by A.63. This also
proves the necessity of the condition from any
order-nine theta point.

Multiplication by two is transitive on the six primitive
characters modulo nine. Their coherent eigenspaces are
one-dimensional; the theta point makes an arrow zero,
so A.51 gives zero primitive rank. The sole remaining
nontrivial conductor is three. Its character subgroup
does not meet theta by A.63, so the degree-three
intermediate new factor is ordinary of dimension two.
The base has two-rank zero. Hence the covering curve
has two-rank exactly two, and Riemann--Hurwitz gives
$g(D)=9(g(C_a)-1)+1=10$.

These parameters also admit the degree-seven whole-new
zero-rank covers of A.63. Thus, as with A.60, this
example separates conditions on the displayed cover;
it does not place the base curve outside the A3 class.
\end{proof}

\begin{claim}{A.67}{Proved; the full-relative criterion for every deck order}
Let $D\to C$ be a connected etale geometrically Galois
cover of smooth proper curves over a finite field $L$
of characteristic $p$, with $g(C)\ge2$, nontrivial
geometric group $G$ of order prime to $p$, and a rational
point $P'$ above $P\in C(L)$. If
\[
                       p\nmid h_L(D)/h_L(C),
\]
then $P$ collides with no other point of $C(L)$.
There is no odd-order restriction and no constant-deck
restriction. For $|G|>2$ the proof is elementary once
the rational projector is established. The degree-two
case uses the verified odd-characteristic hyperelliptic
case of R4.
\end{claim}
\begin{proof}
We first prove the exact geometric fact needed here.
If a finite group acts freely on a smooth proper curve
$D$ of genus at least two, and $P',Q'$ lie in distinct
orbits, then
\[
             \#\operatorname{Stab}_G[Q'-P']\le2.       \tag{A.67.1}
\]
Indeed, a nonidentity stabilizer element $\delta$ gives
the principal degree-two divisor of A.65. The resulting
map to $\mathbf P^1$ is separable: the only possible
inseparable case is a purely inseparable extension of a
rational function field of degree two in characteristic
two, and that extension is again rational, contrary to
$g(D)\ge2$. Thus $D$ has its unique hyperelliptic
involution $\iota$. The unequal but linearly equivalent
degree-two divisors
$\delta Q'+P'$ and $\delta P'+Q'$ belong to its pencil.
Consequently
\[
                   \delta Q'=\iota P',\qquad
                   \delta P'=\iota Q'.
\]
There is at most one such nonidentity $\delta$, because
the action is free and its value on $Q'$ has now been
specified. This proves (A.67.1) in every characteristic.

If $P,Q$ collide, the pointing lifts them to rational
$P',Q'$ with a nonzero $p$-primary difference $T$ and
equal full prime-to-$p$ section classes upstairs. As in
A.21, the rational projector
$|G|^{-1}\pi^*\pi_*$ and the hypothesis kill the new
part of $T$. Thus the whole geometric group fixes $T$.
For $|G|>2$ this contradicts (A.67.1).

If $|G|=2$, the characteristic is odd. The same
degree-two divisor argument makes $D$ hyperelliptic.
The lifted points still collide on $D$ by the full
subgroup-topology statement in A.50. But R4, specifically
AP-A.17--A.19 with the induction of AP-A.15, proves
injectivity for every hyperelliptic curve in odd
characteristic over every finite extension. This is the
remaining contradiction.

This extends the sufficient criterion; it does not turn
an even-order witness into an odd-order witness. The
existence assertion A3 and its odd-group reductions
therefore retain their own quantifiers.
\end{proof}

\begin{claim}{A.68}{Proved; an ordinary elliptic-pullback barrier in characteristic two}
Let $C$ be an ordinary smooth proper curve of genus at
least two over $\overline{\mathbf F}_2$, with a separable
degree-two map $C\to E$ to an ordinary elliptic curve.
For every connected odd-degree etale isogeny $E'\to E$,
the pullback $D=C\times_E E'$ is ordinary. In particular
no such $D$ admits a nonconstant morphism to a
positive-genus curve of two-rank zero, or a nonzero
homomorphism from its Jacobian onto an abelian variety
of two-rank zero.

This applies to both elliptic quotients of the ordinary
genus-two curves in A.28. Thus those isogeny-pullback
families supply neither the rational new-part vanishing
of A.27 nor the geometric zero-rank targets in L5.
It is a restriction on these specified families, not on
all odd-degree covers of $C$.
\end{claim}
\begin{proof}
Write $n=\deg(E'\to E)$ and let $r$ be the number
of geometric ramification points of $C\to E$.
The quadratic extension is ramified, because
Riemann--Hurwitz and $g(C)\ge2$ exclude an unramified
cover of an elliptic curve. It is linearly disjoint from
the odd-degree isogeny, so $D$ is connected.
Since $D\to C$ is etale of degree $n$,
\[
                      g(D)-1=n\bigl(g(C)-1\bigr).
\]
The elliptic curve $E'$ is ordinary because it is
isogenous to $E$. The quadratic map $D\to E'$ has
exactly $nr$ geometric ramification points; etale base
change preserves their local quadratic extensions.
The Deuring--Shafarevich formula
\citep[Corollary~1.8]{AG_Crew1984} gives
\[
 f_2(C)-1=2(f_2(E)-1)+r=r,
 \qquad f_2(D)-1=2(f_2(E')-1)+nr=nr.
\]
Ordinariness of $C$ makes $r=g(C)-1$, so these
equalities prove $f_2(D)=g(D)$.

Every abelian quotient of an ordinary abelian variety
is ordinary: Poincare reducibility expresses the source
up to isogeny as the quotient times a complementary
abelian variety, and both dimension and two-rank are
additive and isogeny invariant. A nonconstant morphism
of positive-genus curves gives a surjective Jacobian
pushforward. Its target therefore cannot have positive
dimension and two-rank zero. This proves the assertion.
\end{proof}

\begin{claim}{A.69}{Proved; the complete degree-thirteen theta locus}
For the characteristic-two family $C_a$ of A.63, put
$U=a^5$. Its Raynaud theta divisor contains a point
of exact order thirteen if and only if
\[
 (U+1)(U^3+U+1)(U^4+U+1)(U^6+U^5+1)=0.              \tag{A.69.1}
\]
Each such parameter admits a connected cyclic etale
degree-thirteen cover with zero full new two-rank.
At every other parameter all degree-thirteen covers
are new-ordinary. In particular every degree-thirteen
cover of $C_0:y^2+y=x^5$ has ordinary new part.
\end{claim}
\begin{proof}
Use the translated local equation of A.64,
$Y^2+Y=u^5+a u^3+c u$, and write $v=c^4$.
A function with pole exactly thirteen at infinity has
the form
\[
 F=A(u)+(u^4+b_3u^3+b_2u^2+b_1u+b_0)Y,
                  \qquad \deg A\le6.
\]
The polynomial $A$ cancels terms through degree six.
The remaining coefficients through degree twelve give
exactly the following six equations:
\[
 \begin{aligned}
 b_1a^2+b_2+b_3v+a&=0,&
 b_0v^2+b_2a^2+b_3+v&=0,\\
 b_1v^2+b_3a^2+1&=0,&
 b_0+b_2v^2+a^2&=0,\\
 b_1+b_3v^2&=0,&
 b_0a^4+b_2+v^2&=0.
 \end{aligned}
\]
The fifth and fourth equations give
$b_1=b_3v^2$ and $b_0=b_2v^2+a^2$.
Put $w=v^4+a^2$. The third equation gives
$b_3w=1$, so $w\ne0$. The first gives
$b_2=a+(a^2v^2+v)/w$. Substitution in the second
now gives $aw^2=1$, in particular $a\ne0$.

Write $a=t^8$ with $t\ne0$, uniquely in the algebraic
closure. The preceding equations force
\[
 \begin{gathered}
 v=t^4+t^{-1},\qquad b_3=t^4,\qquad b_1=t^{12}+t^2,\\
 b_2=t^{28}+t^{18}+t^3,
 \qquad b_0=t^{36}+t^{11}+t.
 \end{gathered}
\]
The last coefficient equation, multiplied by $t^2$,
is $P(t^5)=0$, where
\[
 P(U)=U^{14}+U^9+U^7+U^6+U^4+U^2+U+1.
\]
Since $a^5=(t^5)^8$ and $P$ has coefficients in
$\mathbf F_2$, this is equivalent to $P(a^5)=0$.
Direct multiplication gives the factorization in
(A.69.1). The exact program
\path{code/supersingular_thirteen.py} checks all six
substitutions and the factorization.

Conversely, a root of (A.69.1) has $a\ne0$; take
its unique eighth root $t$ and the displayed coefficients.
All six equations hold. The local parameter $c$ and
point on $C_a$ exist by A.64. The resulting function
has pole exactly thirteen and zero at least thirteen
at that point, so its divisor is $13R-13\infty$.
Its nonzero class has exact prime order thirteen.
Finally two generates $(\mathbf Z/13)^\times$,
so the one-dimensional Frobenius-orbit criterion A.51
gives the last assertions.
\end{proof}

\begin{claim}{A.70}{Refuted; scalar base Frobenius does not identify the rational and geometric tests}
Let $C_0/\mathbf F_{16}$ be $y^2+y=x^5$, with
$P=\infty$. The Frobenius of $J_{C_0}$ is the scalar
$[-4]$. Over $L=\mathbf F_{4096}$ there exists a
constant cyclic degree-thirteen $P$-split cover whose
relative class number is odd, although its entire new
factor is ordinary of dimension twelve.
Moreover $L$ is the first extension of $\mathbf F_{16}$
over which any constant cyclic degree-thirteen cover
can occur. Thus the distinction already occurs at the
minimal constant-deck field under scalar base Frobenius.
\end{claim}
\begin{proof}
The exact point counts are
\[
            \#C_0(\mathbf F_{16})=33,
             \qquad \#C_0(\mathbf F_{4096})=4353.
\]
Both are independently enumerated in
\path{code/scalar_frobenius_count.py} by the absolute
trace criterion for $y^2+y=x^5$; the complete output
is attached. The first achieves the upper Weil bound,
so all four Frobenius eigenvalues are $-4$.
Frobenius is semisimple on a prime-to-two Tate module,
and that realization is injective on endomorphisms
\citep[Sections~1--2, pp.~134--139]{AG_Tate1966}.
Consequently the endomorphism itself is $[-4]$.
Its cube is $[-64]$, which also yields the second count
from the trace formula, and
\[
 J_{C_0}(L)=J_{C_0}[65](\overline{\mathbf F}_2)
                      \simeq(\mathbf Z/65)^4.
\]
Set $H=13J_{C_0}(L)=J_{C_0}[5]$, so $|H|=625$,
and put $A=J_{C_0}(L)/H\simeq(\mathbf Z/13)^4$.
Write $N_b$ for the number of $R\in C_0(L)$ with
$[R-P]$ in the coset $b\in A$.
The zero coset has $N_0=33$: $[5][R-P]=0$ is
equivalent to the class being fixed by $[-4]$, hence,
by the Abel embedding based at $P$, to
$R\in C_0(\mathbf F_{16})$.

Suppose every nontrivial quotient $A\to\mathbf Z/13$
failed the pointed rational criterion. Since two is
primitive modulo thirteen, failure at one nontrivial
character forces failure at all its nontrivial powers.
The genus-two formula in A.8 therefore gives
\[
        \sum_{b\in A}N_b\psi(b)=1\quad\hbox{in }
        \overline{\mathbf F}_2
        \qquad(1\ne\psi\in\widehat A).
\]
The trivial character has the same value because
$4353$ is odd. Fourier inversion, valid since $|A|$
is odd, gives $N_b\equiv0\pmod2$ for every $b\ne0$.
Every one of the $4353-33=4320$ points in nonzero
cosets must therefore share its coset with another point.

For any nonzero $h\in J_{C_0}(L)$ there are at most
two ordered pairs of distinct curve points $(R,S)$
with $[R-S]=h$. Indeed an equality of two such
differences is an equality of effective degree-two
divisor classes; the unique hyperelliptic pencil is
the only possibility for unequal effective divisors,
giving the two exchanged ordered pairs. Therefore
the total number of ordered distinct-point pairs in
the same $H$-coset is at most
$2(|H|-1)=1248$. This contradicts $4320>1248$.
Some prime cyclic quotient is a pointed rational
witness. This is the specified application of the
simultaneous-quotient bound in A.35; the criteria
agent's exact certificate also supplies a concrete
Mumford divisor and its full $L$-polynomials.

A.69 shows that theta has no thirteen-torsion when
$a=0$, so every such cover has ordinary new factor,
of dimension twelve. Finally $-4$ has order three
modulo thirteen. The scalar Frobenius on the geometric
abelianization has a nonzero constant cyclic quotient
only after an extension of degree divisible by three.
This proves the minimal-field assertion.
\end{proof}

\subsection{Exact plane-quartic character computation}

\begin{claim}{A.71}{Proved; a finite divisor algorithm for pointed plane quartics}
Let $C/\mathbf F_q$ be a smooth plane quartic, let
$P\in C(\mathbf F_q)$, and let $K$ be its canonical
divisor class, represented by a line section. From the closed
points of degrees one, two and three, their local expansions,
and linear algebra over $\mathbf F_q$, one can compute the
entire group $J_C(\mathbf F_q)$ and every pointed cyclic
prime-order character. The algorithm below gives the exact
twisted $L$-polynomial for every such character; it does not
replace the Jacobian group law by the numerical zeta polynomial.
\end{claim}
\begin{proof}
Put $B=3P$. Every degree-three divisor class has an effective
representative, by Riemann--Roch. Moreover an effective
degree-three divisor $D$ satisfies
\[
             h^0(D)=1+h^0(K-D).
\]
The degree-one class $K-D$ is effective precisely when it is
the class of a point $R\in C(\mathbf F_q)$. In this case
$D\sim K-R$ and $h^0(D)=2$; otherwise $h^0(D)=1$.
Distinct points give distinct degree-one classes, since a
function with divisor $R-S$ would have degree one on a
genus-three curve. Thus all effective degree-three divisors
represent different classes, except that, for each rational
point $R$, the $q+1$ members of $|K-R|$ represent one class.
This is an exact and exhaustive equivalence test.

The test is implemented using lines. A nonzero line section
vanishing along $D$ exists precisely in the special case;
its residual intersection is the point $R$. No two independent
line sections can vanish along a degree-three divisor on a
smooth plane quartic: their common zero in the plane is a
single point, and at most one line is tangent there. Thus
the computed residual point is unambiguous. A unique
representative for every resulting class identifies
$\operatorname{Pic}^3(C)(\mathbf F_q)$ with
$J_C(\mathbf F_q)$ by $D\mapsto[D-B]$.

Here is the group law, including nonreduced divisors and
nonrational closed points. Given effective representatives
$D,E$ of degree three, choose a nonzero cubic section
vanishing along $D+E+B$. Its existence follows from
\[
 \deg(3K-D-E-B)=3,\qquad h^0(3K-D-E-B)\ge1.
\]
All sections of $3K$ are restrictions of homogeneous cubics:
both spaces have dimension ten, and a nonzero cubic cannot
vanish identically on a smooth quartic. The residual divisor
$R$ has degree three and
\[
                         R\sim3K-D-E-B.
\]
Choose a second cubic vanishing along $R+2B$, with residual
divisor $S$ of degree three. Then
\[
 S\sim3K-R-2B\sim D+E-B,
 \qquad [S-B]=[D-B]+[E-B].
\]
In particular the result is independent of both cubic choices.
All the residual closed points have degree at most three,
so the initial finite list contains every point that can occur.

For completeness the local multiplicity conditions are exact
linear conditions, not evaluations at the support alone. At
each closed point of degree $d$, choose an affine chart and
a coordinate with nonzero partial derivative. Take the other
coordinate as a uniformizer and solve the curve equation
recursively as a formal series through order $12/d$.
The nonzero derivative makes every coefficient uniquely
determined. A multiplicity $m$ imposes the vanishing of
the first $m$ coefficients, each in $\mathbf F_{q^d}$,
or exactly $md$ linear conditions over $\mathbf F_q$.
Intersection orders are read from the same series. B\'ezout's
degree bound is four for a line and twelve for a cubic;
the program checks that all residual degrees are exactly
one or three as appropriate.

For every prime $\ell\ne\operatorname{char}(\mathbf F_q)$
dividing the computed group order, form the image
$\ell J_C(\mathbf F_q)$, its cosets, and a basis of the
elementary abelian quotient. The nonzero linear forms on
that quotient, modulo multiplication by
$\mathbf F_\ell^\times$, give exactly every prime cyclic
quotient, once each. The pointed Lang construction of A.5
gives a geometrically connected constant cyclic cover with
a rational lift above $P$ for each quotient.

Let $\chi$ denote a selected character and put
$\kappa=[K-4P]$. Its polynomial is
\[
 \begin{split}
 L_P(\chi,T)={}&1+c_1(\chi)T+c_2(\chi)T^2
       +q\chi(\kappa)c_1(\chi^{-1})T^3
       +q^2\chi(\kappa)T^4,\\
 c_i(\chi)={}&\sum_{D\ge0,\,\deg D=i}\chi([D-iP])
                                      \quad(i=1,2).
 \end{split}                                                   \tag{A.71.1}
\]
The first two sums are directly enumerable from the closed
points already listed. The degree-three coefficient follows
also without a functional-equation convention: every
degree-three class has one effective representative, except
$K-R$, which has $q+1$. The uniform contribution cancels
for a nontrivial character; the extra contribution is
$q\sum_R\chi([K-R-3P])$. Similarly every degree-four
class has $q+1$ effective representatives except $K$, which
has $q^2+q+1$. The extra contribution is
$q^2\chi(\kappa)$. For larger degrees Riemann--Roch makes
the number of effective representatives constant in every
class, so all subsequent coefficients vanish. This proves
(A.71.1) and its exact normalization at $P$.

Finally evaluate at $T=1$ and take the norm from
$\mathbf Q(\zeta_\ell)$. This is the relative class number.
Its reduction modulo the characteristic is computed exactly
as a polynomial resultant with
$1+X+\cdots+X^{\ell-1}$; nonvanishing is also checked by
the independent Euclidean gcd test. Every conjugate
$L_P(\chi^i,T)$ is given by the substitution
$\zeta_\ell\mapsto\zeta_\ell^i$ in the retained coefficient
vectors. These steps compute characters of the actual
finite Jacobian, including the chosen base-point twist.
\end{proof}

\begin{claim}{A.72}{Proved by exact computation; every pointed plane quartic over $\mathbf F_3$}
There are exactly $3001$ isomorphism classes of pointed
smooth plane quartics $(C,P)$ over $\mathbf F_3$.
For every one of them the program
\path{code/plane_quartic_F3.cpp} computes all constant
odd-prime cyclic pointed tests over $\mathbf F_3$.
It retains the exact $L$-polynomials for all $2692$ tested
quotients, including failures. There is a successful
base-field witness for $2368$ pointed classes.
For $637$ classes a degree-five witness succeeds; five is
therefore their least admissible odd prime over all finite
extensions. For the other $2364$ classes this calculation
does not determine the all-extension minimum.

The $633$ classes without a base-field witness are not
asserted to lack witnesses after extension. Likewise,
the smallest successful prime in the base-field table
is not called the all-extension minimum when it exceeds
five. The $\mathbf F_5$ quartic calculation is given in
A.73--A.74, with the same distinction between base-field
and all-extension minima.
\end{claim}
\begin{proof}
Move $P$ to $(0:0:1)$ and its tangent to $Y=0$.
After scaling the defining equation, every pointed quartic
has the normalized form
\[
 F=YZ^3+(aX^2+bXY+cY^2)Z^2
       +(dX^3+eX^2Y+fXY^2+gY^3)Z+Q_4(X,Y),
\]
with twelve arbitrary coefficients in $\mathbf F_3$.
There are $3^{12}=531441$ such forms. The complete
stabilizer of the point and tangent in
$\operatorname{PGL}_3(\mathbf F_3)$ consists of
\[
 (X,Y,Z)\longmapsto(AX+BY,CY,DX+EY+Z),
 \quad A,C\in\mathbf F_3^\times,
 \quad B,D,E\in\mathbf F_3,
\]
and has order $108$. The program applies every one of
these transformations, rescales the $YZ^3$ coefficient
to one, and marks every orbit. Since a smooth plane
quartic is canonically embedded, every isomorphism of
pointed curves is induced by a projective linear
transformation. Thus these are exactly the pointed
isomorphism classes, including classes with automorphisms.
Pointless quartics contribute no pointed pair; each rational
point on every other quartic is covered up to pointed
isomorphism by this table.

Smoothness is decided over the algebraic closure. The three
partial derivatives of $F$ are homogeneous cubics. As
$4$ is invertible in $\mathbf F_3$, their common
projective zero locus is exactly the singular locus.
They have no projective common zero if and only if
\[
  S_4^{\oplus3}\longrightarrow S_7,
  \qquad(A,B,C)\longmapsto AF_X+BF_Y+CF_Z
\]
is onto, where $S=\mathbf F_3[X,Y,Z]$.
Indeed, in the absence of a projective zero the three
cubics form a regular sequence, whose quotient has
Hilbert series $(1+T+T^2)^3$ and hence vanishing
degree-seven part. Conversely surjectivity excludes a
common projective zero by evaluation. This is a rank
$36$ test on a $36$ by $45$ matrix. The complete
enumeration gives $315360$ smooth normalized forms,
partitioned into the $3001$ displayed orbits.

For each representative, A.71 is implemented with
$\mathbf F_9=\mathbf F_3[t]/(t^2+1)$ and
$\mathbf F_{27}=\mathbf F_3[t]/(t^3+2t+1)$.
The field tables check an inverse for every nonzero
element. Closed points are full Frobenius orbits,
with projective coordinates normalized by their first
nonzero coordinate. The coefficient order for a quartic
is $X^iY^jZ^{4-i-j}$, first $i=0,\ldots,4$, then
$j=0,\ldots,4-i$; all files use that order.

The program independently counts the curve over
$\mathbf F_3,\mathbf F_9,\mathbf F_{27}$, obtains
the degree-six zeta numerator by Newton identities,
and checks that its value at one equals the number
of computed degree-three divisor classes. It also
checks that the number of effective degree-three
divisors is $h(C)+3\#C(\mathbf F_3)$, that every
special class has exactly four representatives,
that the two-cubic group law has the correct identity,
and that the group order annihilates every class.
Associativity is checked on a deterministic set of
triples for each curve. Its validity in general follows
from the divisor-class calculation in A.71.

For every character the degree-three coefficient is
additionally enumerated directly. As a vector indexed
by $\mathbf Z/\ell$, its difference from the reciprocal
formula in (A.71.1) is exactly
$(h(C)/\ell)(1+X+\cdots+X^{\ell-1})$; the program
checks every coefficient of this identity. Resultant
and gcd tests agree for all $2692$ quotients.
The separate program \path{code/review_quartic_norms.py}
also computes the full positive integer resultant for
every quotient. All $2692$ exact relative class numbers
are retained in
\path{code/A4_quartic_F3_exact_norms.csv}, and each has
the same residue modulo three as the finite-field
algorithm. This independently checks the resultant sign,
including every failed test.
The complete audited run gives the counts in the
claim. The JSON lines retain all closed points,
all Jacobian representatives, quotient bases,
coordinates of every class in each quotient,
character vectors, exact cyclotomic coefficient
vectors and relative class numbers modulo three.
The companion CSV records every test and its result.
Since three is the characteristic and the criterion
here retains odd degree, five is the smallest allowed
prime; hence the $637$ degree-five successes have
the stated global minimum.
\end{proof}

\begin{claim}{A.73}{Proved; exhaustive small slices for the $\mathbf F_5$ quartic computation}
Every pointed smooth plane quartic over $\mathbf F_5$
has a normalized representative of precisely one of the
following three types. Write $u,v$ for the coefficients
of $X^4,X^3Y$, respectively, and retain the notation
of A.72 for the other coefficients:
\[
 \begin{array}{c|c|c|c}
 \text{Contact}&\text{Fixed coefficients}
       &\text{Free}&\text{Weight}\\ \hline
 2&a=1,\ b=c=d=0&8&2\\
 3&a=0,\ d=1,\ b=c=e=0&7&3\\
 4&a=d=0,\ u=1,\ b=c=v=0&6&4.
 \end{array}
\]
In a row of weight $w$, two normalized forms are pointed
isomorphic exactly when related by
$(X,Y,Z)\mapsto(AX,A^wY,Z)$ with
$A\in\mathbf F_5^\times$, followed by division by $A^w$.
Thus $5^8+5^7+5^6=484375$ forms, each under a group of
order four, give an exhaustive census. The exact geometric
smoothness test of A.72 gives $94283$ smooth pointed
isomorphism classes, corresponding to $187512000$ of
the original $5^{12}$ normalized coefficient forms.
\end{claim}
\begin{proof}
The tangent section is
\[
                  F(X,0,Z)=aX^2Z^2+dX^3Z+uX^4.
\]
If $a=d=u=0$, the line $Y=0$ is a component, impossible
for a smooth plane quartic. The three displayed cases
therefore exhaust all possibilities and their contact
orders distinguish them intrinsically.

Apply the general point-and-tangent stabilizer from A.72,
now over $\mathbf F_5$, and divide the transformed equation
by $C$. Its first relevant coefficients are
\[
 \begin{aligned}
 a'&=aA^2/C,&
 b'&=3D+2aAB/C+bA,\\
 c'&=3E+aB^2/C+bB+cC,&
 d'&=2aA^2D/C+dA^3/C.
 \end{aligned}                                                \tag{A.73.1}
\]
If $a\ne0$, choose, for any nonzero $A$,
\[
 C=aA^2,\qquad D=-dA/(2a),\qquad
 B=-A(3D+bA)/2,
\]
and then the unique $E$ making $c'=0$. This gives
the first row.

If $a=0$ and $d\ne0$, take
\[
 C=dA^3,\qquad D=-bA/3.
\]
The transformed coefficient of $X^2YZ$ is then
\[
                   e'=3D^2+2bAD+3B/A+eA^2.
\]
Choose the unique $B$ making this zero, followed by
the unique $E$ making $c'=0$. This gives the second row.

Finally let $a=d=0$ and $u\ne0$. Take
$C=uA^4$ and $D=-bA/3$. The transformed coefficient
of $X^3Y$ is
\[
                 v'=D^3+bAD^2+eA^2D+4B/A+vA^3.
\]
There is a unique $B$ making it zero, and then a unique
$E$ making $c'=0$. This gives the third row. Every
division used here is legitimate in characteristic five.

Conversely, a stabilizer element carrying one form of a
fixed row to another must have $C=A^w$ and successively
$D=B=E=0$, by exactly the displayed coefficients.
The residual group is therefore precisely the stated
diagonal group. Its action on a free coefficient of
$X^iY^jZ^{4-i-j}$ is multiplication by
$A^{i+wj-w}$, which is what the program enumerates.

The full stabilizer has order
$(5-1)^2\,5^3=2000$. Each normalized slice orbit
of size $s$ consequently corresponds to $500s$
forms in the original coefficient space. Enumerating
the three slices, marking every diagonal orbit, and
performing the algebraic-closure smoothness rank test
gives the stated $94283$ classes and weighted count.
These are finite exact calculations retained in
\path{code/plane_quartic_F5_census.csv}; the subsequent
Jacobian tests use the same representatives.
\end{proof}


\begin{claim}{A.74}{Proved by exact computation; every pointed plane quartic over $\mathbf F_5$}
For all $94283$ pointed smooth plane-quartic isomorphism
classes over $\mathbf F_5$, the program
\path{code/plane_quartic_F5.cpp} computes every constant
odd-prime cyclic pointed test over the base field.
There are $134422$ such quotient tests, and their exact
cyclotomic $L$-polynomials are retained, including failures.
A successful base-field witness exists for $89906$ classes;
$4377$ classes have no such base-field witness.
For $39917$ classes a degree-three witness succeeds.
Consequently three is their least admissible odd prime
over all finite extensions. The all-extension least prime
for each of the other $54366$ classes remains undetermined
by this calculation, even when its base-field table has
a larger successful prime.
\end{claim}
\begin{proof}
The exhaustive representatives and the algebraic-closure
smoothness test are those of A.73. Apply the exact group
and character algorithm A.71 to each representative,
with the explicit fields
\[
 \mathbf F_{25}=\mathbf F_5[t]/(t^2+2),\qquad
 \mathbf F_{125}=\mathbf F_5[t]/(t^3+t+1).
\]
Both displayed polynomials are irreducible over
$\mathbf F_5$: the quadratic has no root, and direct
substitution of $0,1,2,3,4$ gives no root for the cubic.
The finite-field tables additionally check every nonzero
inverse. The initial projective closed-point lists have
$31$, $310$ and $5240$ points of degrees one, two and three,
respectively, as follows either from their explicit Frobenius
orbits or from the counts of $\mathbf P^2$ over the three
fields. Every local multiplicity, cubic complement,
Jacobian operation and character value is computed over
these finite fields, with no floating-point arithmetic.

The implementation verifies, for every curve, that the
computed number of Jacobian classes agrees with the zeta
polynomial from the three point counts. Writing
$s_i=5^i+1-\#C(\mathbf F_{5^i})$, its first coefficients
are
\[
 a_1=-s_1,\qquad
 a_2=(s_1^2-s_2)/2,\qquad
 a_3=(-s_3-a_1s_2-a_2s_1)/3,
\]
and
\[
                   h(C)=126+26a_1+6a_2+a_3.
\]
The divisions by two and three here are divisions in the
integers; the program checks their exactness. It also checks
that there are $h(C)+5\#C(\mathbf F_5)$ effective
degree-three divisors, exactly six representatives for each
special class, the correct group identity, annihilation of
every class by $h(C)$, and the same deterministic
associativity checks as in A.72. The proof of the group
law itself is A.71.

Every odd prime divisor of $h(C)$ except the characteristic
five is used. For each prime the actual quotient
$J_C(\mathbf F_5)/\ell J_C(\mathbf F_5)$ is computed,
and every linear form up to scalar is tested. The direct
degree-three coefficient enumeration agrees coefficient by
coefficient with the reciprocal formula (A.71.1), with
exact uniform difference $(h(C)/\ell)\Phi_\ell$.
The resultant and Euclidean-gcd tests agree on all
$134422$ characters. The independent integer-arithmetic
program \path{code/review_quartic_norms.py} additionally
computes the full positive relative class number for every
one of these tests. All $134422$ integer values agree
modulo five with the finite-field algorithm, including every
failed test; they are retained in
\path{code/A4_quartic_F5_exact_norms.csv}. Thus the
computed norms are checked as integers, as well as by
two independent finite-field algorithms.
All coefficient vectors, Jacobian
representatives, quotient coordinates, base-point shifts,
and both outcomes are retained in
\path{code/plane_quartic_F5_tests.jsonl}; the companion
CSV lists every quotient and test outcome.

The run is partitioned into four disjoint index classes
modulo four. The merging program
\path{code/launch_quartic_F5.py} checks that their census
rows are exactly the original $94283$ representatives,
with no omissions or repetitions, before accepting the
aggregate. The complete exact output is
\path{code/plane_quartic_F5_summary.json}. It gives the
counts in the claim, together with the distribution of
least successful base-field primes. Since three is the
smallest allowed odd prime in characteristic five, each
degree-three success establishes the all-extension minimum
without any need to exclude smaller extensions or primes.
The other records retain their explicitly limited scope.
\end{proof}


\begin{claim}{A.75}{Proved; universal degree obstructions for geometric primitive vanishing}
Let $C/\overline{\mathbf F}_2$ have genus two and let
$n>1$ be odd. Put
\[
 U_n=\langle2,-1\rangle\subset(\mathbf Z/n)^\times,
 \qquad r_n=\frac{2\varphi(n)}{|U_n|}.
\]
If a connected cyclic etale degree-$n$ cover of $C$
has primitive conductor-$n$ factor of two-rank zero, then
\[
                   r_n(r_n-1)\le2(n-1).                 \tag{A.75.1}
\]
If $C$ is nonordinary, the stronger necessary condition is
\[
                   r_n(r_n+1)\le2(n-1).                 \tag{A.75.2}
\]
Consequently degree $127$ and degree $257$ cannot have
primitive new two-rank zero for any genus-two base in
characteristic two. More generally, no degree
$n=2^d-1$ with $d\ge48$ can have this property.
These restrictions concern geometric two-rank, and assert
no corresponding exclusion of the rational class-number
criterion. They also do not exclude the possibility that
every curve admits some other degree as in A6.
\end{claim}
\begin{proof}
In characteristic two the Raynaud bundle is a degree-one
theta characteristic $B$ on the Frobenius twist of $C$.
Its theta divisor is the embedded curve
\[
 \Theta=\{L\in J_C:h^0(B\otimes L)>0\}=C-B.
\]
Since $B^{\otimes2}\simeq\omega_C$, Serre duality shows
$\Theta=-\Theta$. The harmless Frobenius twist does not
change any geometric torsion or intersection assertion
used here. The point zero belongs to $\Theta$ if and
only if $C$ is nonordinary: its multiplicity space
$H^0(B)$ is the kernel of Frobenius on $H^1(\mathcal O_C)$,
which vanishes exactly in the ordinary case.

Let $\Gamma\subset J_C$ be the cyclic order-$n$ group
of character line bundles for the cover. Every nontrivial
coherent character space has dimension one. By A.51,
primitive rank zero requires each orbit of multiplication
by two on the primitive elements of $\Gamma$ to meet
$\Theta$. Write $f=\operatorname{ord}_n(2)$. If
$-1\in\langle2\rangle$, each such orbit must contain
at least two theta points, since a nonzero point of odd
order is distinct from its negative. There are
$\varphi(n)/f$ orbits, giving at least
$2\varphi(n)/f=r_n$ primitive theta points. If
$-1\notin\langle2\rangle$, the doubling orbits are
paired by negation. Each pair contributes at least two
points, or $\varphi(n)/f=r_n$ points in total.

For any nonzero $h\in J_C$, there are at most two
ordered pairs of distinct points of $\Theta$ with
difference $h$. This is the genus-two difference bound
proved in A.70: unequal effective degree-two divisors
in one class belong to the unique hyperelliptic pencil,
which supplies at most the exchanged pair. Equivalently,
distinct translates of the theta curve have intersection
number two. Thus any $m$ distinct theta points belonging
to $\Gamma$ satisfy
\[
                         m(m-1)\le2(n-1).
\]
Use $m\ge r_n$ for (A.75.1). In the nonordinary case,
the additional point zero is distinct from all primitive
points, so $m\ge r_n+1$ gives (A.75.2).

For $127$, two has order seven and its subgroup does
not contain $-1$. Hence $|U_{127}|=14$ and $r_{127}=18$,
whereas $18\cdot17=306>252$. For $257$, two has
order sixteen and contains $-1$; therefore $r_{257}=32$
and $32\cdot31=992>512$. All group computations,
the necessary inequalities, and all odd degrees through
$2001$ are checked in
\path{code/primitive_degree_obstructions.py}, with its
complete exact output attached.

For the infinite family, an elementary bound valid for
odd $n$ is
\[
                       \varphi(n)\ge\frac23 n^{2/3}.
\]
Indeed, for every prime $t\ge5$ one has
$t-1\ge t^{2/3}$, and this gives
$\varphi(t^e)\ge(t^e)^{2/3}$. The possible factor
$3^e$ costs at most the constant $2/3$, since
$2\cdot3^{e-1}\ge(2/3)3^{2e/3}$ for $e\ge1$.
Multiply these prime-power bounds.
If $n=2^d-1$, then $|U_n|\le2d$, and consequently
\[
 r_n\ge\frac{\varphi(n)}d
      \ge\frac{2n^{2/3}}{3d}.
\]
For $d\ge48$ we have $n^{1/6}\ge3d$: it suffices
that $2^{(d-1)/6}\ge3d$, whose seed at $d=48$
and subsequent monotonicity follow from the exact integer
inequalities $2^{47}\ge144^6$ and
$2\cdot48^6>49^6$. The ratio
$2^{(d-1)/6}/d$ increases thereafter.
It follows that $r_n\ge2\sqrt n$, contradicting
(A.75.1). This proves the stated infinite family of
geometrically excluded degrees.
\end{proof}


\begin{claim}{A.76}{Proved; a linear rational-exponent criterion at a fixed field}
Let $L=\mathbf F_q$ have characteristic $p$, let $C/L$
be smooth proper and geometrically connected with $g(C)\ge2$,
and let $P\in C(L)$. Let $\pi:D\to C$ be a connected
finite etale Galois cover with constant group of order $m$
prime to $p$, and choose $P'\in D(L)$ over $P$.
Write $p^e$ for the exponent of $J_D(L)[p^\infty]$,
with $p^e=1$ when this group is zero. If
\[
                       (q-1)p^e<2(m-1),                \tag{A.76.1}
\]
then $P$ collides with no other point of $C(L)$.
The group need not be nonabelian, and its order need not
be odd. More precisely, any collision forces
\[
 \operatorname{ord}[Q'-P']\ge\frac{2(m-1)}{q-1}
\]
for the rational lifts given by the collision.

Consequently, a tower of constant pointed prime-to-$p$
Galois covers over one fixed $L$, of degrees $m_i\to\infty$
and rational $p$-primary exponents $b_i=o(m_i)$, would
exclude every collision with $P$. Existence of such a tower
for every pointed curve is not established here.
\end{claim}
\begin{proof}
Suppose $P$ and a distinct $Q\in C(L)$ collide.
A.1 gives a rational lift $Q'\in D(L)$ and a
nonzero $p$-primary class $T=[Q'-P']$ in $J_D(L)$.
Since the deck group is constant and the action is free,
every translate of $P'$ and of $Q'$ is rational.
The two fibres are disjoint, so $D(L)$ contains at least
$2m$ points.

Put $b=\operatorname{ord}T$, which divides $p^e$.
Over the algebraic closure there is a rational function
$f$ with
\[
                         \operatorname{div}(f)=bQ'-bP'.
\]
It can be chosen over $L$. Indeed the divisor is rational,
so every Galois conjugate of $f$ differs from it by a
constant. These constants form a multiplicative cocycle
in a finite extension over which $f$ is defined; Hilbert 90
rescales $f$ to an $L$-rational function with the same
divisor. The corresponding map $D\to\mathbf P^1_L$
has degree $b$, a single geometric point above zero,
and a single geometric point above infinity. Each of the
other $q-1$ rational values has at most $b$ rational
preimages, whether or not the map is separable. Therefore
\[
                        2m\le\#D(L)\le2+(q-1)b.
\]
This gives the sharper order bound and then
$2(m-1)\le(q-1)p^e$, contrary to (A.76.1).

In the tower assertion $q$ is fixed, so $b_i=o(m_i)$
eventually implies (A.76.1). This is a conditional
consequence, not a proof of the required growth hypothesis.
Enlarging the constant field changes both $q$ and the
rational torsion exponent; neither is held fixed silently.
\end{proof}


\begin{claim}{A.77}{Proved; cyclotomic exponent bounds and further quartic certificates}
Let $D\to C$ be a constant pointed cyclic etale cover
of prime degree $\ell\ne p$ over $L=\mathbf F_q$,
with $g(C)\ge2$. Put
\[
 h=\#J_C(L),\quad h_{\rm rel}=\#J_D(L)/h,
 \quad f=\operatorname{ord}_\ell(p),\quad v=v_p(h_{\rm rel}).
\]
Choose a faithful character $\chi$ and write
$a=L_P(\chi,1)\in\mathbf Z[\zeta_\ell]$ as a polynomial
in $\zeta_\ell$. Let $b$ be the degree of its gcd with
$\Phi_\ell$ modulo $p$. Define
\[
 t=\begin{cases}
 0,&v=0,\\
 1+(v-b)/f,&v>0.
 \end{cases}
 \qquad B=p^{\max\{v_p(h),t\}}.                       \tag{A.77.1}
\]
Then $t$ is a nonnegative integer and
\[
             \exp J_D(L)[p^\infty]\mid B.
\]
Consequently $(q-1)B<2(\ell-1)$ is a sufficient
no-collision criterion at $P$ by A.76.
The weaker bound with $t=v/f$ is valid as well.

Applying this bound to the complete, already certified
quartic polynomials and integer norms of A.72 and A.74
gives the following exact base-field conclusions:
\[
 \begin{array}{c|r|r|r|r}
 q&\text{Pointed classes}&\text{L1 successes}
   &\text{Additional successes}&\text{Total certified}\\ \hline
 3&3001&2368&33&2401\\
 5&94283&89906&123&90029.
 \end{array}
\]
Every additional class has failed every base-field
constant odd-prime L1 test, yet has a successful exponent
certificate. There are $52$ individual covers over
$\mathbf F_3$ and $272$ over $\mathbf F_5$ which fail
L1 and satisfy this exponent criterion. This is a strict
enlargement of the verified sufficient tests over the
fixed base field. It makes no assertion that those
pointed curves lack an L1 witness after extension.
\end{claim}
\begin{proof}
Let $P_D=\ker(J_D\to J_C)^0$ be the connected Prym.
The pullback--norm isogeny
$J_C\times P_D\to J_D$ has kernel killed by a power
of $\ell$, so it induces an isomorphism on $p$-divisible
groups over $L$. In particular its rational $p$-primary
group splits into the base and Prym parts. The base part
has exponent dividing $p^{v_p(h)}$.

Let $\sigma$ generate the deck group. The relation
\[
 \Phi_\ell(\sigma)=1+\sigma+\cdots+\sigma^{\ell-1}
                       =\pi^*\pi_*
\]
vanishes on $P_D$. Therefore the full integer ring
$\mathcal O_K=\mathbf Z[\zeta_\ell]$, with
$K=\mathbf Q(\zeta_\ell)$, acts integrally on $P_D$,
not merely up to isogeny. Frobenius $F$ commutes with
this action because the deck group is constant over $L$.
Its $K$-linear characteristic polynomial $f_F(Z)$
has degree $2g(C)-2$ and coefficients in $\mathcal O_K$;
its conjugates are the reciprocal character polynomials
$Z^{2g(C)-2}L_P(\chi^i,Z^{-1})$. One may choose the
embedding of $K$ so the identity character corresponds
to the displayed $\chi$; reversing the deck generator
only permutes these embeddings. Integrality follows
also because its roots are Frobenius algebraic integers.

Cayley--Hamilton on an auxiliary prime-to-$p\ell$
Tate module, followed by faithfulness of that realization,
gives the endomorphism identity $f_F(F)=0$ on $P_D$.
On $P_D(L)[p^\infty]$ the endomorphism $F$ is one,
so $f_F(1)=a$ annihilates this finite group. Moreover
\[
                      N_{K/\mathbf Q}(a)=h_{\rm rel}.
\]
These statements use the actual Frobenius endomorphism
and the integral cyclotomic action; there is no inference
of an integral group structure from the numerical zeta
polynomial alone.

Because $p\ne\ell$, the algebra
\[
 \mathcal O_K\otimes\mathbf Z_p=\prod_{i=1}^{(\ell-1)/f}O_i
\]
is a product of unramified discrete valuation rings,
each of residue degree $f$. Write $t_i=v_p(a_i)$ for
the image of $a$ in $O_i$. The corresponding factor
of the finite Prym group is annihilated by $p^{t_i}$:
$a_i$ is $p^{t_i}$ times a unit of $O_i$.
Thus its exponent divides $p^{\max_i t_i}$, and
\[
                       v=f\sum_i t_i.
\]
Exactly $b/f$ of the integers $t_i$ are positive.
Indeed $\Phi_\ell$ modulo $p$ is square-free, with
irreducible factors all of degree $f$, and its common
factors with $a$ are precisely the residue factors
where $a_i$ is zero. If $v=0$, every $t_i$ is zero.
Otherwise
\[
 \max_i t_i\le\sum_i t_i-\bigl(b/f-1\bigr)
                         =1+(v-b)/f=t.
\]
This proves (A.77.1), the divisibility assertion,
and the simpler bound $\max_i t_i\le v/f$.

The exact script \path{code/quartic_rational_exponent.py}
reads the retained Jacobian orders, mod-$p$ gcd degrees,
and independently computed integer relative class numbers.
It verifies $f\mid v$, $f\mid b$, $v\ge b$, and
$v=0\Longleftrightarrow b=0$ for every quotient,
then applies (A.77.1) and the strict inequality of A.76.
The complete successful certificate tables are
\path{code/plane_quartic_F3_rational_exponent_certificates.csv}
and
\path{code/plane_quartic_F5_rational_exponent_certificates.csv},
with companion JSON summaries. Each certificate records
its pointed model, prime quotient character, both exact
class numbers, $f,v,b$, the exponent upper bound,
and the two sides of the strict inequality. Comparing
with every L1 test for the same pointed model gives
exactly the additional counts in the table. The weaker
norm-only scan and the initially tried total-class-number
bound are also retained, with their outputs, so that no
failed or superseded attempt is suppressed.
\end{proof}


\begin{claim}{A.78}{Proved; the geometric limitation of two ordinary elliptic targets}
Let $\Omega$ be algebraically closed, let $D/\Omega$
be a smooth proper curve of genus at least two, let $A$
be an abelian surface, and let $u:A\to A'$ be an isogeny
with a nonzero point in its geometric kernel.
For every morphism $\phi:D\to A$, the composite
$u\phi$ fails to be injective on geometric points.
In particular, if $A=E_1\times E_2$ is a product of
ordinary elliptic curves in characteristic two, then
no joint map
\[
 D\longrightarrow E_1\times E_2
       \xrightarrow{[2^{e_1}]\times[2^{e_2}]}
                    E_1\times E_2,
 \qquad e_1,e_2\ge1,
\]
is geometrically point-separating.

Thus replacing each of two ordinary elliptic targets by
multiplication by its rational two-primary exponent cannot
supply a geometrically point-separating family for L6.
This does not exclude separation on rational points over
a specified finite field: the geometric intersection points
in the proof may only appear after extension, when the
rational exponents themselves can increase.
\end{claim}
\begin{proof}
If $\phi$ already identifies two geometric points, so does
$u\phi$. Suppose therefore that $\phi$ is injective on
geometric points. It is nonconstant and finite onto its
irreducible image curve $Z\subset A$. Its function-field
extension over the normalization $\widetilde Z$ is purely
inseparable: a nontrivial separable degree would give
multiple geometric points in a general fibre. A finite
purely inseparable map between smooth curves over a perfect
field is a power of relative Frobenius followed by an
isomorphism. Hence
\[
                g(\widetilde Z)=g(D)\ge2.
\]
The image may be singular, but its arithmetic genus is
at least this normalization genus. Adjunction on the smooth
abelian surface, whose canonical bundle is trivial, gives
\[
                   Z^2=2p_a(Z)-2>0.
\]
In fact $Z$ is ample. Every irreducible curve $Y$ on an
abelian surface has $Y^2\ge0$, by intersecting it with a
distinct translate. If $Z\cdot Y=0$, the Hodge index
theorem and $Z^2>0$ would force $Y^2<0$, since $Y$
is not numerically zero (an ample divisor has positive
intersection with $Y$). This is impossible. Thus $Z$
meets every curve positively, and the surface
Nakai--Moishezon criterion proves ampleness. Only the
already established inequality $Z^2>0$ is needed below.

Choose $0\ne t\in\ker u(\Omega)$. If $Z=Z-t$,
any point $z\in Z$ gives two distinct points $z,z+t$
on $Z$. Otherwise the two curves $Z$ and $Z-t$ have
intersection number $Z^2>0$, because translations preserve
numerical equivalence. They therefore meet; a point in
their intersection again gives distinct $z,z+t\in Z$.
Choose preimages of these two points in $D$. Their images
under $u\phi$ agree, contradicting injectivity.

For ordinary elliptic curves in characteristic two,
$E_i[2](\Omega)$ is nontrivial, so the displayed product
isogeny has a nonzero geometric kernel. Moreover the unique
nonzero geometric two-torsion point of each ordinary
elliptic curve over a finite field is rational; hence its
rational two-primary exponent is always at least two,
as noted in A.61. The general result applies.

Finally, merely knowing that two unmultiplied maps
$\phi_i$ separate points imposes no new torsion constraint
on a hypothetical collision. Once
$T=[Q'-P']$ is $p$-primary, every induced Jacobian
homomorphism automatically sends $T$ to a $p$-primary
point of its target. If target exponents $b_i$ are used,
it is the maps $[b_i]\phi_i$, rather than the original
$\phi_i$, whose point separation would contradict the
collision. The theorem explains exactly why this geometric
replacement fails for two ordinary elliptic targets.
\end{proof}


\begin{claim}{A.79}{Proved; a limit of normal-chain tests with nonconstant deck Frobenius}
There is a finite odd solvable group $G$, a finite
$\mathbf F_2[G]$-module $V$, and an automorphism $F$ of $V$
semilinear for an automorphism of $G$, such that
\[
 (V^H)^F\supsetneq(V^K)^F
 \quad\text{for every proper inclusion }H<K
       \text{ of normal subgroups of }G,
\]
but no $F$-fixed vector has trivial $G$-stabilizer.
Thus the constant-deck simultaneous-vector assertion cannot
be extended to nonconstant Frobenius using only these
normal-chain tests, even if they are imposed over every
finite constant-field extension.
This is an abstract module counterexample, without a claim
that the module is realized by a curve or its Jacobian.
\end{claim}
\begin{proof}
Let $N=(\mathbf F_{25},+)$ and let $c$ of order three act
on $N$ by multiplication by $\omega$, where
$\omega^2+\omega+1=0$ over $\mathbf F_5$.
The polynomial is irreducible over $\mathbf F_5$,
so $N$ is an irreducible two-dimensional $\mathbf F_5$
module for $\langle c\rangle$. Put
\[
             G=N\rtimes\langle c\rangle,
             \qquad |G|=75.
\]
Its normal subgroups are exactly $1,N,G$.
Indeed, for a normal subgroup $L$, the intersection
$L\cap N$ is either $1$ or $N$ by irreducibility.
In the first case $[L,N]\subset L\cap N=1$,
so $L\subset C_G(N)=N$ and $L=1$.
In the second case the quotient $G/N$ has prime order,
so $L=N$ or $G$.

Take the left regular module $V=\mathbf F_2[G]$,
choose $0\ne t\in N$, and set $F=\rho(t)$,
where $\rho$ denotes the left regular representation.
Then
\[
 F\rho(g)F^{-1}=\rho(tgt^{-1}),
\]
so $F$ is semilinear for the inner deck automorphism
$g\mapsto tgt^{-1}$. For every subgroup $H$,
\[
 (V^H)^F=V^{\langle H,t\rangle},\qquad
 \dim_{\mathbf F_2}(V^H)^F
       =\frac{75}{|\langle H,t\rangle|}.
\]
The dimensions for $H=1,N,G$ are respectively
$15,3,1$. Every required normal-chain inclusion
is therefore strict. This remains true for every power $F^a$:
the dimensions are $15,3,1$ when $5\nmid a$, and $75,3,1$
when $5\mid a$. Nevertheless every vector in
$V^F$ is fixed by the nonidentity element $t$.
The indicator vector of the subgroup $\langle t\rangle$
has stabilizer exactly that subgroup, so the minimum
stabilizer order on $V^F$ is exactly five.

The missing intermediate test is explicit:
$1\triangleleft\langle t\rangle$ is stable under the
deck Frobenius, and
\[
            (V^1)^F=(V^{\langle t\rangle})^F.
\]
Here $\langle t\rangle$ is not normal in $G$.
There is also a missing test with a $G$-normal upper group:
if $A_0\subset N$ is an $\mathbf F_5$-line not containing $t$,
then $A_0\triangleleft N$, both are $F$-stable, and
$\dim(V^{A_0})^F=\dim(V^N)^F=3$.
Consequently a hypothesis imposed on all relevant
intermediate relative covers can detect this obstruction,
even though the normal-subgroup-only hypothesis cannot.
The standard-library certificate
\path{code/nonconstant_normal_chain_counterexample.py}
enumerates the complete normal-subgroup lattice,
checks the semilinear identity and all fixed dimensions,
and computes the displayed vector's exact stabilizer.
Its full output is retained in the companion JSON file.
\end{proof}


% A.80--A.81: full direct proof; independently approved by root and criteria.
\begin{claim}{A.80}{Proved; an effective conductor bound in the scalar supersingular case}
Let $C/\mathbf F_{2^{2r}}$ be a genus-two curve with two-rank zero,
with rational Weierstrass point $P$, and suppose Jacobian Frobenius
is $[\varepsilon2^r]$, where $\varepsilon\in\{1,-1\}$.
If a nontrivial connected cyclic \'etale cover of odd degree $n$
has primitive part of two-rank zero, put
\[
 f=\operatorname{ord}_n(2),\qquad d=\gcd(f,r),\qquad
 M=\frac{\varphi(n)}d.
\]
Then
\[
             M(M+1)\le 2(n-1),\qquad n<10r^3.       \tag{A.80.1}
\]
Consequently the geometric primitive-cover question for this
fixed scalar-Frobenius curve has only finitely many possible
conductors, with the displayed explicit bound.
\end{claim}
\begin{proof}
The Raynaud theta divisor is the symmetric Abel curve
$\Theta=C-P$, and contains zero. In characteristic two its
nontrivial character cohomology has dimension $g-1=1$.
The orbit criterion A.51 therefore says that primitive
rank zero requires at least one theta point in every
multiplication-by-two orbit on the primitive characters
of the cyclic subgroup of order $n$.
There are $\varphi(n)/f$ such orbits, each of size $f$.

Frobenius preserves $\Theta$, and symmetry removes the sign
$\varepsilon$. Thus multiplication by $2^r$ preserves its
geometric points. A theta point in one of the required
orbits brings with it its entire multiplication-by-$2^r$
orbit, of size $f/d$. The primitive theta intersection
therefore has at least $M=\varphi(n)/d$ points.
Together with zero there are at least $M+1$ points of
$\Theta$ in the cyclic subgroup. For any nonzero difference
there are at most two ordered pairs of distinct points of
a genus-two Abel curve, since two distinct translates of
its principal theta divisor have intersection number two.
The exact difference count used in A.75 now gives
$M(M+1)\le2(n-1)$.

For completeness, the elementary uniform estimate
\[
                  \varphi(n)\ge\frac23 n^{5/6}
                       \quad(n\text{ odd})             \tag{A.80.2}
\]
follows prime-power by prime-power. For primes $\ell\ge5$,
$(\ell-1)^6\ge\ell^5$: it holds at five, and
$(x-1)^6/x^5$ increases for $x>1$. Hence
$\varphi(\ell^a)^6/\ell^{5a}\ge1$ for every $a\ge1$.
The only remaining prime, three, contributes a factor at
least $64/243$, which is larger than $(2/3)^6$.
Multiplying proves (A.80.2). Since $d\le r$, the first
inequality of (A.80.1) implies
\[
 \frac{4}{9r^2}n^{5/3}\le M^2<2n,
 \qquad
 n<\left(\frac92\right)^{3/2}r^3<10r^3.
\]
The strict inequality is preserved because $n>1$.
\end{proof}

\begin{claim}{A.81}{Refuted: universal nontrivial primitive-cover existence}
\label{clm:AGprimitivecounterexample}
There are supersingular genus-two curves over
$\mathbf F_{2^{128}}$ for which every nontrivial connected
cyclic \'etale cover of odd degree has primitive part of
positive two-rank. Thus the version of the existence statement
A6 requiring degree $n>1$ is false, already for genus two
in characteristic two.

A completely specified family containing these examples is
\[
 C_t:\qquad y^2+y=
       \left(\frac{u^2+t^2u+t}{t^3+1}\right)^3,
 \qquad
 t\in\mathbf F_{2^{128}}\setminus\mathbf F_4,       \tag{A.81.1}
\]
where $C_t$ means the smooth proper model. The proof gives
an explicit finite incidence set of excluded parameters
and a positive exact lower bound for its complement.
The program certifies this nonemptiness and the displayed
algebraic identities; it does not claim to have enumerated
one particular parameter in the complement.

All these curves themselves satisfy A1: their Jacobians
have geometric two-rank zero. The refutation concerns
existence of nontrivial primitive-zero covers, and gives no
counterexample to A1 or to the arithmetic criterion A3.
The prompt does not explicitly require $n>1$. With its
convention $P_1=J_C$, the identity cover is a witness for
every curve constructed here. Therefore these examples
do not refute the literal degree-one-permitted version
of A6; that version remains open for curves of positive
$p$-rank.
\end{claim}
\begin{proof}
We first establish the family and its torsion incidence
bound directly. It is the characteristic-two pencil
studied by Moret--Bailly in \emph{Polarisations de degr\'e 4
sur les surfaces ab\'eliennes}, the theorem, parts~1(c) and~3
\citep{AG_MoretBailly1979}; the explicit calculations below
supply the required properties independently.

\emph{1. Explicit genus-two curves and elliptic maps.}
Let $k$ be a finite field containing $\mathbf F_4$, put
$E:y^2+y=x^3$, $A=E\times E$, and take
$t\in k\setminus\mathbf F_4$. Write $d=t^3+1$, so
$t\ne0$ and $d\ne0$. If $X=u^2+t^2u+t$, expansion gives
\[
 X^3=u^6+t^2u^5+(t+t^4)u^4+t^6u^3
                +(t^5+t^2)u^2+t^4u+t^3.
\]
Over the perfect field $k$, successive Artin--Schreier
changes remove the even powers $u^6,u^4,u^2$. The highest
remaining odd power is $u^5$, with nonzero coefficient
$t^2/d^3$. Thus (A.81.1) has genus two and two-rank zero,
and has exactly one point $P$ above infinity; that point
is $k$-rational and is its unique Weierstrass point.
In particular $\omega_C\simeq\mathcal O_C(2P)$.

Define two morphisms to $E$ by
\[
\begin{aligned}
 x_1&=\frac{u^2+t^2u+t}{d}, & y_1&=y,\\
 x_2&=\frac{tu^2+u+t^2}{d}, & y_2&=y+\frac{u^3+1}{d}.
\end{aligned}                                                   \tag{A.81.2}
\]
The first elliptic equation is the defining equation of
$C_t$. For the second, put $Z=tu^2+u+t^2$ and $Y=u^3+1$.
The polynomial identity
\[
                 X^3+Z^3=dY^2+d^2Y
\]
proves $y_2^2+y_2=x_2^3$. Properness extends both maps
across $P$, and both take $P$ to the elliptic origin.
The joint map $i_t:C_t\to A$ is birational onto its image,
because
\[
                         u=x_2+t x_1,
                         \qquad y=y_1.                \tag{A.81.3}
\]

\emph{2. The pencil and its only geometric base point.}
On $A$ consider
\[
 f=\frac{y_2+y_1+1+x_1x_2}{x_2+x_1^2},\qquad
 D_\infty=E\times\{0\}+\Gamma_{F_E},                 \tag{A.81.4}
\]
where $F_E:E\to E$ is the degree-two Frobenius morphism.
The polar divisor of $f$ is exactly $D_\infty$.
Here is the local check. In the affine part, the denominator
vanishes on the two graphs $\Gamma_{F_E}$ and
$\Gamma_{-F_E}$. On the first graph its numerator is one,
so the pole is simple. On the second the numerator
vanishes and the factor cancels. Along $E\times\{0\}$,
the orders of $y_2$ and $x_2$ are minus three and minus two,
so again the pole is simple with nonzero leading coefficient.
Along $\{0\}\times E$, the numerator has order minus three
and the denominator minus four, so $f$ vanishes rather
than having a pole. These account for every possible
polar divisor.

The two components of $D_\infty$ meet geometrically only
at the origin: the kernel of $F_E$ has only its
identity as a geometric point. At every nonzero point of
the components of $D_\infty$, the preceding nonzero pole
coefficient shows that the section $f$ of
$\mathcal O_A(D_\infty)$ does not vanish. Hence the pencil
generated by $1$ and $f$ has no nonzero geometric base point.
It has the origin as a base point, because (A.81.2) sends
$P$ to zero for every admissible $t$.

Substitution in (A.81.4) gives $f\circ i_t=t$.
The divisor $D_t=(f-t)+D_\infty$ is precisely the image
$i_t(C_t)$, with multiplicity one. To check that no
component is being suppressed, write $x=x_1,z=x_2$.
Eliminating $y_2$ from the numerator equation for $f=t$
and the two elliptic equations gives
\[
 (z+x^2)\bigl(z^2+t^2z+t^2x^2+x+t\bigr)=0.
\]
The first factor is the cancelled graph. The second is
exactly (A.81.1) after $u=z+tx$, as (A.81.2)--(A.81.3)
show. It is irreducible because its smooth proper model
has the Artin--Schreier equation with nonzero reduced
$u^5$ coefficient. It has multiplicity one: away from
the denominator, differentiating on $A$ gives
\[
                   df=dx+
                  \frac{z^2+x+f}{z+x^2}\,dz,
\]
which is nonzero along the generic point of this component.
The boundary component $\{0\}\times E$ is a zero fibre
only for $t=0$, excluded here, and the other boundary
component is polar. Therefore $D_t\sim D_\infty$.

\emph{3. The induced isogeny and the theta origin.}
The elliptic components in $D_\infty$ have self-intersection
zero and intersect with length $\deg F_E=2$. Consequently
\[
                         D_t^2=4.                     \tag{A.81.5}
\]
The map $i_t$, based at $P$, induces a homomorphism
$v_t:J_{C_t}\to A$. Its image is all of $A$: otherwise
the genus-two image curve would lie in a translate of an
elliptic subvariety and would itself have normalization
genus one. Thus $v_t$ is an isogeny.
For the principal Abel theta divisor $\Theta=C_t-P$,
birationality gives $(v_t)_*\Theta=D_t$.
For an isogeny of abelian surfaces, pullback and pushforward
on the rational N\'eron--Severi group satisfy
$v_t^*(v_t)_*=(\deg v_t)\operatorname{id}$, and pullback
multiplies intersection numbers by $\deg v_t$.
Using $\Theta^2=2$ therefore gives
\[
                       D_t^2=2\deg v_t.
\]
Equation (A.81.5) proves $\deg v_t=2$.
In particular it induces an isomorphism on geometric
$n$-torsion for every odd $n$.

There is no unspecified translation of Raynaud theta in
this argument. Since $C_t$ has two-rank zero,
$J_{C_t}[2](\overline k)=0$. Its theta characteristic is
therefore unique; as $\omega_C=\mathcal O_C(2P)$, the
Raynaud theta characteristic is $\mathcal O_C(P)$.
The corresponding theta divisor is exactly $C_t-P$.
Thus a nonzero theta point killed by an odd $n$ maps to
a nonzero point of $A[n]$ on $D_t$.

\emph{4. Uniform torsion incidence and scalar Frobenius.}
Every nonzero geometric point of $A$ belongs to at most
one member of the pencil, by its single base point.
Since $\#A[n](\overline k)=n^4$ for odd $n$, there are
at most $n^4-1$ parameters over $\overline{\mathbf F}_2$
for which $\Theta_{C_t}$ contains a nonzero $n$-torsion point.
This bounds all geometric torsion points, regardless of
which extension makes them rational.

For $k=\mathbf F_{4^e}$, Frobenius on $E$ is $[(-2)^e]$:
indeed $\#E(\mathbf F_2)=3$, so its Frobenius
polynomial is $X^2+2$, and Cayley--Hamilton for the
endomorphism gives $F_E^2=[-2]$.
The isogeny $v_t$ is defined over $k$, so the same exact
scalar formula holds on $J_{C_t}$. The conductor bound
A.80 applies with $r=e$. If a primitive-zero cyclic cover
exists, its odd conductor therefore satisfies $3\le n<10e^3$,
and its theta criterion supplies a nonzero theta point
of exact order $n$.

\emph{5. A finite, nonempty counterexample set.}
Take $e=64$, $q=4^{64}=2^{128}$, and $B=10\cdot64^3=2621440$.
Define a finite subset of $\mathbf F_q$ by
\[
 \mathcal B=\mathbf F_4\ \cup\
 \left(\mathbf F_q\cap
 \bigcup_{\substack{3\le n<B\\ n\text{ odd}}}
 \{f(a):a\in A[n](\overline{\mathbf F}_2)\setminus\{0\},
                         \ f(a)\text{ finite and defined}\}\right).
                                                               \tag{A.81.6}
\]
Here $f(a)$ means the value of the pencil morphism
$A\setminus\{0\}\to\mathbf P^1$, including extension across
removable zeros of the displayed numerator and denominator.
This is an exact algebraic definition using the displayed
rational function on the fixed elliptic surface $E\times E$.
It specifies a finite set, without an assertion that the
large set has been enumerated. The incidence bound gives
\[
 \#\mathcal B\le4+
       \sum_{\substack{3\le n<B\\ n\text{ odd}}}(n^4-1)<4+B^5<q.
\]
Every $t\in\mathbf F_q\setminus\mathcal B$ defines a
smooth curve (A.81.1), and no odd conductor allowed by
A.80 can supply a primitive-zero cover. Larger conductors
are impossible by that same bound. This proves the
asserted counterexample existence over the stated finite
field, with no restriction on the field of definition of
the hypothetical cover.

The exact counting certificate actually gives
\[
 \#(\mathbf F_q\setminus\mathcal B)\ \ge\
       340282354541538070615576657943687331836>0.
\]
The standard-library program
\path{code/supersingular_primitive_counterexample.py}
verifies the polynomial identities of (A.81.2)--(A.81.4),
the Artin--Schreier coefficient expansion, the elementary
integer bounds, and the exact odd fourth-power sum.
Its complete output is retained in the companion JSON.
The program is a finite algebra and nonemptiness certificate,
not an enumeration of an individual parameter or a test
of infinitely many covers.

Finally, $J_{C_t}$ is isogenous to $E^2$ and hence has
geometric two-rank zero. R1 therefore already excludes a
collision on each $C_t$. The two supersingular elliptic
maps (A.81.2) jointly separate geometric points as well:
$v_t$ has degree two and a connected kernel, since
$J_{C_t}[2](\overline k)=0$. It is therefore radicial on
geometric points, while the Abel map is an
embedding. This is consistent with A.78, whose hypothesis
requires a nonzero geometric kernel point.
\end{proof}

\begin{claim}{A.82}{Proved; nontrivial primitive-cover counterexamples over $\mathbf F_{2^{38}}$}
For the family $C_t$ of A.81 over $\mathbf F_{2^{38}}$,
at least $59927733332$ parameters $t$ define curves having
no nontrivial connected cyclic \'etale cover of odd degree
with primitive two-rank zero. This counts parameters,
without asserting that the corresponding curves are
pairwise nonisomorphic or selecting one individual parameter.
The degree-one qualification at the end of A.81 applies.
\end{claim}
\begin{proof}
There is a useful refinement of the torsion incidence bound.
For an integer $e\ge1$ let
\[
\begin{split}
 \mathcal S_e={}&\{n:\ 3\le n<10e^3,\ n\text{ odd},\
       M_e(n)(M_e(n)+1)\le2(n-1)\},\\
 M_e(n)={}&\frac{\varphi(n)}{\gcd(\operatorname{ord}_n(2),e)},\qquad
 J_4(n)=n^4\prod_{\ell\mid n}(1-\ell^{-4}).
\end{split}                                                   \tag{A.82.1}
\]
The integer $J_4(n)$ is the number of points of exact order
$n$ in $A[n]\simeq(\mathbf Z/n\mathbf Z)^4$.
For $n\in\mathcal S_e$ define $\mathcal B_{e,n}$ to be
the parameters $t\in\mathbf F_{4^e}\setminus\mathbf F_4$
for which the pencil member $D_t$ contains at least
$M_e(n)$ distinct points of exact order $n$ in $A$.
Every nonzero point belongs to at most one pencil member.
Counting incidences therefore gives
\[
               \#\mathcal B_{e,n}
                  \le\left\lfloor\frac{J_4(n)}{M_e(n)}\right\rfloor.
                                                               \tag{A.82.2}
\]
If $C_t$ admits a nontrivial primitive-zero cover of
conductor $n$, A.80 puts $n$ in $\mathcal S_e$ and gives
at least $M_e(n)$ exact-order-$n$ points of its theta
divisor. The odd-torsion isomorphism induced by $v_t$
sends them injectively to points counted in
$\mathcal B_{e,n}$. Thus the explicitly defined set
\[
 \mathbf F_{4^e}\setminus
       \left(\mathbf F_4\cup
          \bigcup_{n\in\mathcal S_e}\mathcal B_{e,n}\right)
                                                               \tag{A.82.3}
\]
consists entirely of parameters with the asserted property,
and its size is at least
\[
       4^e-4-
       \sum_{n\in\mathcal S_e}
          \left\lfloor\frac{J_4(n)}{M_e(n)}\right\rfloor.
                                                               \tag{A.82.4}
\]
The definition of $\mathcal B_{e,n}$ uses the threshold
$M_e(n)$; the stronger bound (A.82.2) is not being asserted
for parameters with only one theta point of order $n$.

For $e=19$, exact enumeration below the proved global
cutoff $68590$ gives precisely
\[
 \mathcal S_{19}=\{191,229,361,419,457,571,573,647,687,
                         955,1083,1257,1371\}.
\]
The certificate records all factorizations, totients,
multiplicative orders, values of $M_{19}(n)$, exact-order
point counts, and integer quotients. Its sum is
\[
 \sum_{n\in\mathcal S_{19}}
      \left\lfloor\frac{J_4(n)}{M_{19}(n)}\right\rfloor
       =214950173608.
\]
As $4^{19}=274877906944$, formula (A.82.4) becomes
\[
             274877906944-4-214950173608
                     =59927733332>0.
\]
The root program \path{code/A6_scalar_conductor_bounds.py}
and its complete JSON output give the global enumeration.
The independent review
\path{agent_geometric/code/primitive_counterexample_e19_review.py}
factors every odd integer below the cutoff by trial division
and recomputes every relevant multiplicative order without
using the root program's sieve or cache. Its complete JSON
agrees on every $e=19$ record. These are exhaustive finite
arithmetic certificates for the displayed existence theorem,
not an enumeration of a parameter in (A.82.3).
\end{proof}

\begin{claim}{A.83}{Proved: geometric origin alone does not bound unit-root returns}
Consider the fixed smooth elliptic family over the multiplicative group
\[
 \mathcal E\longrightarrow\mathbf G_{m,\mathbf F_4},
 \qquad E_a:\ y^2+xy=x^3+a.                         \tag{A.83.1}
\]
For every positive odd integer $m$, there is a closed parameter $a$
of exact degree $m$ over $\mathbf F_4$ such that
\[
 \#E_a(\mathbf F_{4^m})=4^m,
 \qquad E_a(\mathbf F_{4^m})\simeq\mathbf Z/2^{2m}\mathbf Z.
                                                        \tag{A.83.2}
\]
If $\beta_a\in\mathbf Z_2^\times$ is its unit-root Frobenius
eigenvalue, then
\[
 \beta_a^2-\beta_a+4^m=0,
 \qquad v_2(\beta_a-1)=2m.                         \tag{A.83.3}
\]
The parameter is a point of odd order in
$\mathbf G_m(\overline{\mathbf F}_2)$, and its Teichmuller lift
is prime-to-two torsion. Thus a single fixed geometric family,
ordinary fibers, integral cohomology of Weil weight one, and
prime-to-two Frobenius periods do not give a uniform bound for
the return unit minus one, even at prime-to-two torsion parameters.
The assertion does not put these parameters in the powers of one
fixed odd prime, and does not identify this family with the Abel
tower of A.107. Those restrictions remain essential unresolved
inputs to that route.
\end{claim}
\begin{proof}
The discriminant and $j$-invariant of (A.83.1) are $a$ and $1/a$,
respectively. Hence the displayed projective Weierstrass model is
a smooth elliptic scheme over $\mathbf G_m$. Each fiber is ordinary:
the geometric point $(0,\sqrt a)$ is a nonzero point of order two,
whereas an elliptic curve has at most one independent geometric
two-torsion point in characteristic two.

Fix odd $m$ and put $q=4^m$. Waterhouse's classification
\citep[Theorem~4.1, case~(1), p.~536]{AG_Waterhouse1969}
asserts that every integer $t$ with $|t|\le2\sqrt q$ and
$(t,2)=1$ is the Frobenius trace of an ordinary elliptic curve
over $\mathbf F_q$. Apply it to $t=1$. The resulting elliptic
curve $E/\mathbf F_q$ has $q$ rational points.
Since $4\mid q$, Proposition~5.1 of
\citep[p.~311]{AG_AuerTop2002} identifies $E$ over $\mathbf F_q$
with a member $E_a$ of (A.83.1).

For completeness, the normal-form step can be seen directly.
An ordinary elliptic curve over the perfect field $\mathbf F_q$
has an equation
\[
               y^2+xy=x^3+b x^2+a,\qquad a\ne0.
\]
Its unique nonzero geometric two-torsion point is
$T=(0,\sqrt a)$. As $E(\mathbf F_q)$ is a two-group and
the geometric two-primary torsion has rank one, it is cyclic;
since its order is $q\ge4$, the point $T$ has a rational half.
The duplication formula for a point $R$ with nonzero
$x$-coordinate is
\[
                 x(2R)=x(R)^2+\frac a{x(R)^2}.
\]
Thus a half of $T$ must have $x$-coordinate $r=a^{1/4}$.
Writing its $y$-coordinate as $rz$, the equation becomes
\[
                       z^2+z=b+r+r^2.
\]
It is solvable over $\mathbf F_q$ precisely when
$\operatorname{Tr}_{\mathbf F_q/\mathbf F_2}(b)=0$.
Consequently $b=s^2+s$ for some $s\in\mathbf F_q$, and
the substitution $y\mapsto y+sx$ removes the $x^2$ term.
This proves the required isomorphism to $E_a$ and also confirms
the rational four-torsion step in the cited proposition.

The characteristic polynomial of $E_a$ is $X^2-X+q$.
Modulo two it has distinct roots zero and one. Its unit root
$\beta_a$ therefore belongs to $\mathbf Z_2^\times$, and the
other root is $q/\beta_a$, of valuation $2m$. Evaluating at one
gives
\[
              (1-\beta_a)(1-q/\beta_a)=q.
\]
The second factor is a two-adic unit, proving (A.83.3).
The geometric two-primary torsion of an ordinary elliptic curve
is $\mathbf Q_2/\mathbf Z_2$, so its finite rational two-primary
subgroup is cyclic. Its full rational order here is $q$, which
proves the group assertion in (A.83.2).

It remains to prove that the parameter has exact degree $m$;
the existence of a curve over $\mathbf F_{4^m}$ alone would not
establish the assertion about periods. Let
$d=[\mathbf F_4(a):\mathbf F_4]$, so $d\mid m$. The displayed
model $E_a$ is defined over $\mathbf F_{4^d}$. Write
$\gamma\in\mathbf Z_2^\times$ for its unit root over that field.
The unit root over $\mathbf F_{4^m}$ is
$\gamma^{m/d}$. Because $m/d$ is odd,
\[
 v_2(1-\gamma^{m/d})=v_2(1-\gamma);
\]
the quotient $(1-\gamma^{m/d})/(1-\gamma)$ is a sum of an
odd number of odd two-adic integers. The nonunit Frobenius root
does not contribute to the two-primary part of the rational
point count, so
\[
             v_2\bigl(\#E_a(\mathbf F_{4^d})\bigr)=2m.
\]
If $d<m$, Hasse's bound contradicts this divisibility:
\[
 4^m\le\#E_a(\mathbf F_{4^d})
       \le(2^d+1)^2<4^m.
\]
Hence $d=m$. Finally $a\in\mathbf F_{4^m}^\times$ has order
dividing the odd number $4^m-1$, and its fourth-power Frobenius
orbit has length exactly $m$. Teichmuller lifting preserves that
order and that orbit length.

The small exact certificate
\path{agent_geometric/code/geometric_unitroot_growth.py}
constructs examples for $m=1,3,5$, records the finite fields,
the parameter, its order and degree, and the point count, and
checks the unit-root congruence. This finite check illustrates
the theorem; its assertion for every odd $m$ follows from the
classification and the degree argument above.
\end{proof}

% A.84: root and criteria independently audited.
\subsection{A geometric unit-root obstruction along a fixed prime-power sequence}

\begin{claim}{A.84}{Proved}\label{clm:AGfixedthreepower}
For each \(n\geq1\), choose a primitive \(3^n\)-th root
\(a_n\in\overline{\mathbb F}_2\), and put
\[
 d_n=2\cdot3^{n-1},\qquad k_n=\mathbb F_{2^{d_n}},\qquad
 E_n:\ y^2+xy=x^3+a_n.
 \tag{A.84.1}
\]
Let \(\beta_n\in\mathbb Z_2^\times\) be the unit root of
\(X^2-t_nX+2^{d_n}\), where \(t_n\) is the Frobenius trace
of \(E_n/k_n\). Then
\[
 \lim_{n\to\infty}v_2(\beta_n-1)=+\infty,
 \qquad
 \exp E_n(k_n)[2^\infty]=2^{v_2(\beta_n-1)}.
 \tag{A.84.2}
\]
Every \(E_n\) is a fibre of the single smooth ordinary elliptic
family \(y^2+xy=x^3+a\) over \(\mathbb G_{m,\mathbb F_2}\).
The parameter \(a_n\) has exact degree \(3^{n-1}\) over
\(\mathbb F_4\), an odd degree, and has exact multiplicative
order \(3^n\). Thus a fixed geometric family, weight-one Weil
purity, integrality, ordinariness, a fixed prime-power torsion
sequence, and bounded \(2\)-part of the parameter degrees do
not imply a uniform bound for the return valuations.

The conclusion asserts divergence, without a quantitative
comparison between \(2^{v_2(\beta_n-1)}\) and a prescribed
power of \(3^n\). It neither identifies this family with the
relative factors of an Abel tower nor disproves the weaker
polynomial growth bound required by that route.
\end{claim}

\begin{proof}
We first prove a finite-cover statement. If \(Y\to
\mathbb A^1_{\mathbb F_2}\) is a finite etale torsor under a
finite abelian \(2\)-group and its fibre at \(0\) is split,
then its fibre at every primitive \(3^n\)-th root is split over
\(k_n\) for all sufficiently large \(n\).
It is enough to prove this for a cyclic quotient of order
\(2^r\), and then take the maximum of finitely many bounds.
The Artin--Schreier--Witt exact sequence and affine vanishing
of \(H^1(\mathbb A^1,W_r\mathcal O)\) give
\[
 H^1_{\mathrm{\acute et}}(\mathbb A^1_{\mathbb F_2},
                         \mathbb Z/2^r)
 =W_r(\mathbb F_2[a])/(F-1)W_r(\mathbb F_2[a]).
 \tag{A.84.3}
\]
Here \(F\) is Witt Frobenius. These standard exact statements
are recorded explicitly in \citet[Theorem~1.2 and the following
cohomology sequence, p.~2]{AG_Sekiguchi2001}; compare the
field-theoretic treatment and Frobenius trace formula in
\citet[Theorem~3.2, Section~5.4, p.~29]{AG_KostersWan2016}.
Affine vanishing also follows directly by induction from the
filtration of \(W_r\mathcal O\) with additive-group quotients.

Choose a representative \(w(a)\) in (A.84.3).
On \(W_r(\mathbb F_2)\), Frobenius is the identity. Since
splitting at \(0\) says that the specialized cohomology class
vanishes, it follows that \(w(0)=0\), as an actual Witt vector.
There is a finite additive expansion
\[
                 w(a)=\sum_{i=0}^{r-1}\sum_{j\in J_i}
                                  V^i[a^j],
                 \qquad J_i\subset\mathbb Z_{>0}\text{ finite}.
 \tag{A.84.4}
\]
Repeated terms may be retained in this expression. To obtain
it, subtract the sum of the Teichmuller monomials of the first
Witt coordinate. The difference has first coordinate zero,
so is uniquely \(V\) of a shorter Witt vector with polynomial
coordinates. This shorter vector still specializes to zero
at \(a=0\). Induction proves (A.84.4), with no constant
monomial and only finitely many exponents. Let \(D\) be the
largest exponent occurring, taking \(D=0\) if the class is
zero.

The order of \(2\) modulo \(3^n\) is \(d_n\): the order
is even and \(v_3(4^m-1)=1+v_3(m)\). Hence the absolute
Frobenius conjugates of \(a_n\) are precisely all primitive
\(3^n\)-th roots. In the full Witt ring of \(k_n\), the
Teichmuller lift \(z=[a_n]\) has exact order \(3^n\), and
\[
 \sum_{s=0}^{d_n-1}z^{j2^s}
 =\sum_{\zeta\text{ primitive of order }3^n}\zeta^j=0
 \qquad(0<j<3^{n-1}).
 \tag{A.84.5}
\]
Indeed the last sum is the sum over all \(3^n\)-th roots
minus the sum over all \(3^{n-1}\)-th roots; both geometric
sums vanish for the stated exponents. Since \(F\) commutes
with \(V\), (A.84.4)--(A.84.5) show that
\[
 \operatorname{Tr}_{W_r(k_n)/W_r(\mathbb F_2)}w(a_n)=0
                       \qquad\text{if }3^{n-1}>D.
 \tag{A.84.6}
\]
For \(F(b)-b=w(a_n)\), the action of the degree-\(d_n\)
Frobenius on \(b\) is translation by this trace, by telescoping
\(F^{d_n}(b)-b\). Thus (A.84.6) is exactly the splitting
condition. This proves the finite-cover statement.

We now apply it to the elliptic family. Its discriminant is
\(a\); each fibre over \(\mathbb G_m\) is ordinary and
has geometric etale \(2\)-rank one. For each \(r\geq1\),
let
\[
 Q_r=\ker\bigl(V^r:E^{(2^r)}\longrightarrow E\bigr).
 \tag{A.84.7}
\]
This is a finite etale group scheme over \(\mathbb G_m\),
with geometric fibres cyclic of order \(2^r\). Its frame
torsor has group \((\mathbb Z/2^r)^\times\), a finite
abelian \(2\)-group. We show directly that this torsor is
trivial over the local field \(K=\mathbb F_2((a))\).
This statement concerns the etale group (A.84.7), and does
not presume rationality of the non-etale group scheme
\(E[2^r]\) over the imperfect field \(K\).

In the purely inseparable extension \(K^{1/2}\), the point
\(R_1=(0,a^{1/2})\) has exact order two. Suppose a point
\(R=(x,y)\) of order \(2^s\) is defined over
\(K^{1/2^s}\), and both coordinates have positive valuation
(the zero coordinate is allowed). The convergent series
\[
                  \lambda=\sum_{i\geq0}x^{2^i}
 \tag{A.84.8}
\]
belongs to the same complete local field, has positive
valuation, and satisfies \(\lambda^2+\lambda=x\).
Set
\[
 u^2=y+(\lambda+1)x,\qquad v=u(\lambda+u).
 \tag{A.84.9}
\]
These coordinates belong to \(K^{1/2^{s+1}}\) and have
positive valuation. The curve equation for \(R\) gives
\(u^4+xu^2=a\); in particular \(u\ne0\).
It follows directly that \(v^2+uv=u^3+a\), and the
elliptic doubling formulas give \(2(u,v)=R\).
Thus \((u,v)\) has exact order \(2^{s+1}\).
Induction constructs a generator of the geometric
\(2^r\)-torsion in \(E(K^{1/2^r})\).
Applying relative Frobenius \(F^r\) gives a \(K\)-rational
generator of \(Q_r\). Consequently \(Q_r/K\) is constant.

The frame torsor is therefore unramified at \(a=0\),
with trivial residue Frobenius there. Its normalization over
\(\mathbb A^1_{\mathbb F_2}\) is finite etale and has split
fibre at \(0\). Apply the finite-cover statement to this
extension. For each fixed \(r\), and all sufficiently large
\(n\), the group \(Q_r\) at \(a_n\) is constant over
\(k_n\). Relative Frobenius identifies its geometric points
Galois-equivariantly with those of \(E_n[2^r]\), so the
unit root satisfies \(\beta_n\equiv1\pmod{2^r}\).
This proves the limit in (A.84.2).

Finally, the rational \(2\)-primary subgroup of an ordinary
elliptic curve over a finite field is cyclic. The other
Frobenius root \(2^{d_n}/\beta_n\) is divisible by two,
and therefore
\[
 v_2\#E_n(k_n)
 =v_2\bigl((1-\beta_n)
              (1-2^{d_n}/\beta_n)\bigr)
 =v_2(1-\beta_n).
\]
This proves the exponent formula. The exact parameter degree
over \(\mathbb F_4\) is \(\operatorname{ord}_{3^n}(4)
=3^{n-1}\), as asserted.
\end{proof}

\paragraph{Finite certificates and quantitative scope.}
The independent complete point-count program
\texttt{agent\_geometric/code/kloosterman\_three\_power\_test.cpp}
and its JSON output use the irreducible polynomial
\(\Phi_{3^n}(X)=X^{2\cdot3^{n-1}}+X^{3^{n-1}}+1\)
for \(n=1,2,3\). They verify, respectively,
\[
\begin{array}{c|r|r|r|r}
 n&\#k_n&\#E_n(k_n)&t_n&v_2(\beta_n-1)\\\hline
 1&4&4&1&2\\
 2&64&64&1&6\\
 3&262144&262912&-767&8
\end{array}
\]
The trace-one pattern of the first two examples already
fails at the third. The proof of divergence does not use
these finite data. Its threshold for a fixed torsion level
is the finite degree bound \(3^{n-1}>D\) in (A.84.6),
where \(D\) comes from representatives of that level's
frame-torsor characters. No bound for the growth of these
representative degrees with the torsion level is asserted.

% Supplement to A.82: an exact reduction for any individual-parameter test.
\subsection{Four remaining conductors in the explicit pencil}

\begin{claim}{A.85}{Proved}\label{clm:AGsidonfour}
For every member \(C_t\) of the explicit family in
Claim~A.81 over
\(\mathbb F_{2^{38}}\), any nontrivial odd conductor of a
primitive two-rank-zero cyclic cover belongs to
\[
                         \{191,229,457,1257\}.
                                                        \tag{A.85.1}
\]
More precisely, its theta curve has no point of exact order
\(361,419,571,573,647,687,955,1083\), or \(1371\).
This assertion is a necessary-condition reduction; it does not select
an individual parameter from the nonempty complement in Claim~A.82.
\end{claim}

\begin{proof}
The scalar Frobenius is \([-2^{19}]\), and the theta curve
\(\Theta=C_t-P\) is symmetric. If \(Q\in\Theta\) has odd
exact order \(n\), put
\[
 H_n=\langle-2^{19},-1\rangle\subseteq(\mathbb Z/n)^\times,
 \qquad S_n=\{0\}\cup H_n.
\]
Every point \(aQ\), \(a\in S_n\), lies on \(\Theta\).
For each nonzero \(d\in\mathbb Z/n\), the ordered pairs
\((a,b)\in S_n^2\) with \(a-b=d\) give distinct geometric
points in \(\Theta\cap(\Theta+dQ)\). The two theta curves
are distinct: translation stabilizing the divisor of a principal
polarization belongs to the trivial kernel of that polarization.
Their intersection number is \(\Theta^2=2\), so there are at
most two such pairs, even if some intersections are nonreduced.

Apply this test to the thirteen conductors already proved exhaustive
in Claim~A.82. The complete finite calculation gives
\[
\begin{array}{r|r|r}
 n&|H_n|&\max_{d\ne0}\#\{(a,b)\in S_n^2:a-b=d\}\\\hline
 191&10&2\\229&4&2\\361&18&4\\419&22&6\\457&4&2\\
 571&6&4\\573&20&10\\647&34&6\\687&8&4\\955&40&30\\
 1083&18&4\\1257&22&2\\1371&8&4
\end{array}
\]
The standard-library program
\texttt{primitive\_counterexample\_sidon\_review.py}
constructs each subgroup by closure under its two specified generators
and enumerates every ordered pair. Its retained JSON output records
all subgroups and, for each failing difference, every representing
ordered pair. The nine failures exclude even a single theta point of
the indicated exact order. A primitive two-rank-zero conductor must
have such a theta point, by Claim~A.80. This proves (A.85.1).
\end{proof}


\subsection{An individually specified finite counterexample}

\begin{claim}{A.86}{Refuted, with an explicit finite counterexample: nontrivial primitive-cover existence}
\label{clm:AGexplicitprimitivecounterexample}
Let
\[
 \begin{split}
 M(Z)&=Z^{38}+Z^{14}+Z^{10}+Z^9+Z^8+Z^5+Z^2+Z+1,\\
 k&=\mathbb F_2[Z]/(M(Z)),\qquad t=Z\bmod M(Z),
 \qquad d=t^3+1.
 \end{split}                                                    \tag{A.86.1}
\]
The polynomial \(M\) is irreducible. Let \(C/k\) be the smooth
projective model of the affine equation
\[
             y^2+y=\left(\frac{u^2+t^2u+t}{d}\right)^3,
                                                               \tag{A.86.2}
\]
and let \(P\) be its unique point at infinity. This is a
supersingular genus-two curve, with \(P\in C(k)\).
For every connected cyclic etale cover
\(D\to C_{\bar k}\) of odd degree \(n>1\), and every generator
\(\sigma\) of its deck group, the abelian subvariety
\[
 P_n(D)=\operatorname{im}\left(
       \prod_{\ell\mid n}(1-\sigma^{n/\ell}):J_D\longrightarrow J_D
                              \right)
\]
has positive two-rank. Thus the nontrivial-cover version of A6 is
false for this completely specified curve, independently of the
counting existence argument.

The degree-one qualification is essential: \(J_C\) has two-rank
zero, so the identity cover satisfies the literal variant that
permits \(n=1\). The example also satisfies A1 and is not a
collision or a counterexample to A1.
\end{claim}

\begin{proof}
First, the standard-library bit-polynomial program
\path{primitive_parameter_field_check.py} verifies
\[
 Z^{2^{38}}\equiv Z\pmod M,
 \qquad
 \gcd(M,Z^{2^2}-Z)=\gcd(M,Z^{2^{19}}-Z)=1.
                                                               \tag{A.86.3}
\]
Since the prime divisors of \(38\) are \(2\) and \(19\),
the finite-field irreducibility criterion proves that \(M\) is
irreducible. Its retained output gives every successive squaring
residue, using bitwise addition, multiplication, and division in
\(\mathbb F_2[Z]\). In particular \(t\notin\mathbb F_4\),
and \(d\ne0\).

Claim~A.81 applies to (A.86.2). It supplies the genus,
supersingularity, unique point at infinity \(P\), which is rational,
scalar Frobenius
\([-2^{19}]\), and the pointed degree-two isogeny
\(\alpha:J_C\to E^2\), where \(E:y^2+y=x^3\). On the
embedded theta curve \(C-P\), its two coordinates have
\[
 A(u)=\frac{u^2+t^2u+t}{d},\qquad
 B(u)=\frac{tu^2+u+t^2}{d}                              \tag{A.86.4}
\]
as their elliptic \(x\)-coordinates. Both elliptic images of a
finite point are finite, and \(P\) maps to \((0_E,0_E)\).
The theta divisor relevant to the primitive two-rank criterion
is precisely \(C-P\).

By Claims~A.80, A.85, and~\ref{clm:ACpacketconductors}, the
only remaining possible conductors of a primitive-zero cover are
\(191\) and \(229\). We prove the stronger assertion that
\(C-P\) has no nonzero point killed by either integer.

Here is a direct characteristic-two verification of the elliptic
torsion polynomials used in that test. On \(E\), the tangent
and chord formulas give the rational-function identities
\[
 x([2]R)=x(R)^4,
 \qquad
 x(R+S)+x(R-S)=\frac{1}{(x(R)+x(S))^2}.
                                                               \tag{A.86.5}
\]
For the second identity, replacing \(S=(x_S,y_S)\) by
\(-S=(x_S,y_S+1)\) changes the chord slope by
\(1/(x_R+x_S)\); the other terms in the two addition formulas
cancel. The tangent slope is \(x(R)^2\), giving the first
identity. Hence the rational functions \(r_n(X)=x([n]R)\),
where \(X=x(R)\), satisfy
\[
 r_1=X,\qquad r_{2m}=r_m^4,\qquad
 r_{2m+1}=X+\frac{1}{(r_m+r_{m+1})^2}.             \tag{A.86.6}
\]
All equalities are rational-function identities, so values where
the displayed addition formula has a denominator zero are
handled by the reduced rational function. Its finite poles are
exactly the \(x\)-coordinates of nonzero points killed by \(n\).

Independently, define polynomials in \(\mathbb F_2[X]\) by
\[
 \begin{gathered}
 \psi_1=\psi_2=\psi_4=1,\qquad \psi_3=X^4+X,\\
 \psi_{2m}=\psi_m^4,\qquad
 \psi_{2m+1}=\psi_{m+2}\psi_m^3+
                       \psi_{m-1}\psi_{m+1}^3\quad(m\ge2).
 \end{gathered}                                                 \tag{A.86.7}
\]
The program \path{primitive_parameter_rational_check.py}
constructs (A.86.6) using reduced numerator--denominator pairs,
without using (A.86.7) for that construction. For both required
integers it verifies that the reduced denominator is exactly
\(\psi_n^2\). Thus no general division-polynomial formula in
characteristic different from two is being imported. The checks are
\[
\begin{array}{c|r|r|r}
 n&\deg\psi_n&\deg\operatorname{num}(r_n)&
                         \deg\operatorname{den}(r_n)\\\hline
 191&18240&36481&36480\\
 229&26220&52441&52440
\end{array}
\]
and the reduced numerators are coprime to their denominators.

Finally, exact polynomial arithmetic over the explicitly defined
field \(k\) gives polynomials \(U_n,V_n\in k[u]\) satisfying
\[
       U_n(u)\psi_n(A(u))+V_n(u)\psi_n(B(u))=1
                \qquad(n=191,229).                            \tag{A.86.8}
\]
The producer \path{primitive_parameter_gcd.py} evaluates
(A.86.7) directly at the quadratic polynomials (A.86.4), computes
the extended gcd, and checks the whole identity (A.86.8) by
polynomial multiplication. The degrees of the two evaluated
polynomials are \(36480\) and \(52440\), respectively; in
each case both Bezout coefficients have degree four less.
All four polynomials for each integer are retained as binary
coefficient arrays. Their encoding, degrees, and SHA-256 hashes
are specified in \path{primitive_parameter_gcd.json}:
coefficients are in ascending degree, each represented by an
unsigned 64-bit little-endian integer whose bit \(i\) is the
coefficient of \(t^i\). The recorded runtime dependency is
\texttt{python-flint==0.9.0}. This is a complete exact identity
certificate, not a numerical approximation or a random-point test.
An independently written loader, \texttt{A6\_individual\_curve\_verify.py},
also verifies every coefficient-array hash, degree, and encoding,
reconstructs both compositions independently, and checks both full
Bezout identities; its retained run exits successfully.

If a finite \(Q\in C(\bar k)\) had \([Q-P]\ne0\) killed
by \(n=191\) or \(229\), its image under \(\alpha\) would
be killed by \(n\) in both elliptic factors. By (A.86.5)--
(A.86.7), the finite value \(u(Q)\) would be a common zero of
\(\psi_n(A(u))\) and \(\psi_n(B(u))\), contradicting
(A.86.8). The remaining point \(P\) represents the origin
and cannot supply a primitive character. Therefore neither
remaining conductor occurs, over any finite extension or over
\(\bar k\), and the claimed assertion follows.

Finally \(J_C[2^\infty](\bar k)=0\). A collision would already
give a two-primary class \([Q-P']\) on \(C\), for the two
colliding points \(P',Q\), by R1. This class is zero, and the
Abel map of a positive-genus curve is injective on geometric
points. Hence \(Q=P'\), proving the stated A1 assertion.
\end{proof}

% Independently audited by root and the F/topos agent; integrated as A.87.
\subsection{A failure of full prime-to-\texorpdfstring{\(p\)}{p} configuration exactness}

\begin{claim}{A.87}{Proved: the full prime-to-\(p\) configuration fiber does not inject}
\label{clm:AGconfigurationodd}
Let \(C\) be a smooth proper connected complex curve of genus
\(g\geq2\), let \(P\in C\), and put
\[
 U=\{(x,z)\in C\times C:x\ne z\},\qquad F=C-\{P\}.
\]
Use the second projection \(U\to C\) to identify \(F\) with a
fiber. For every prime \(p\), write \(\Gamma^{(p')}\) for the
inverse limit of the finite quotients of a group \(\Gamma\)
whose orders are prime to \(p\). Then
\[
 \pi_1(F)^{(p')}\longrightarrow
 \pi_1(U)^{(p')}                                     \tag{A.87.1}
\]
is not injective. The same assertion holds for the geometric
prime-to-\(p\) etale fundamental groups of every smooth proper
curve of genus at least two in characteristic \(p\).
In particular it holds for the smooth projective model of
\[
                    C_0/\mathbb F_2:\quad y^2+y=x^5,
             \qquad P=\infty.                       \tag{A.87.2}
\]

This refutes the full prime-to-\(p\) configuration exactness
condition proposed as a sufficient route in
\cite[Remark 3.13]{AR_ST2009}. It does not supply two coincident
point sections and does not refute A1. The proof specifies a finite
curve and finite branched-cover construction; it does not enumerate
a particular nonidentity profinite kernel element or an individual
finite monodromy matrix witnessing the failure.
\end{claim}

\begin{proof}[Independent proof for \(p=2\)]
For this proof write \((-)^{\mathrm{odd}}=(-)^{(2')}\).
The extension to every prime is proved below by a separate
congruence argument.
\emph{An odd-degree branched family.}
Choose a connected cyclic etale cover
\(a:C'\to C\) of degree three. Put \(m=3^{2g}\), and let
\(b:B\to C\) be the connected etale cover corresponding to
\[
 \ker\bigl(\pi_1(C)\longrightarrow H_1(C,\mathbb F_3)\bigr).
                                                               \tag{A.87.3}
\]
It is Galois of degree \(m\), factors through \(C'\), and
\(b^*:H^1(C,\mathbb F_3)\to H^1(B,\mathbb F_3)\) is zero:
each homomorphism \(\pi_1(C)\to\mathbb F_3\) restricts to
zero on the subgroup in (A.87.3).

On \(Y=C'\times B\) consider the effective divisor
\[
          D=(a\times b)^{-1}(\Delta_C).
                                                               \tag{A.87.4}
\]
It consists of three disjoint graphs \(B\to C'\). Its first
Chern class is divisible by three in \(H^2(Y,\mathbb Z)\).
Indeed, in the Kunneth decomposition modulo three, the two
degree-two summands of the diagonal class pull back by maps of
degrees \(3\) and \(m\), and its cross summand vanishes because
\(b^*H^1(C,\mathbb F_3)=0\). The integral cohomology of a
product of compact oriented surfaces is torsion-free, so vanishing
modulo three is equivalent to integral divisibility.
The class \(c_1(D)/3\) is integral of type \((1,1)\); the
Lefschetz \((1,1)\) theorem gives a line bundle with that Chern
class. Its cube differs from \(\mathcal O_Y(D)\) by an element
of \(\operatorname{Pic}^0(Y)\). Surjectivity of multiplication
by three on this complex abelian variety therefore gives
\[
                         L^{\otimes3}\simeq\mathcal O_Y(D).
                                                               \tag{A.87.5}
\]

Take the cyclic cover \(S\to Y\) obtained by adjoining a cube
root of the canonical section of \(L^{\otimes3}\). It is a
connected smooth projective surface. The composite
\(f:S\to B\) is smooth and proper with connected fibers:
locally along any graph the cover has the form \(w^3=t\), where
\((t,z)\) are coordinates with \(z\) on \(B\); these become
the coordinates \((w,z)\) on \(S\). Each fiber is a cyclic
triple cover of \(C'\), branched at the three distinct points
\(a^{-1}(b(z))\), each with ramification index three. Consequently
\[
 g(C')=3g-2,\quad g(B)=m(g-1)+1,\quad
                  h:=g(S_z)=9g-5.                   \tag{A.87.6}
\]
This is the usual cyclic construction of a Kodaira fibration;
compare \cite[Section 1.3]{AG_Rollenske2010}. All the particular
divisibility and smoothness assertions needed here were proved
above.

The family is nonisotrivial. To see this without any genericity
assumption, suppose every fiber were isomorphic to a fixed curve
\(T\). There are only finitely many degree-three morphisms
\(T\to C'\): their Hom scheme is of finite type, and the tangent
space at a morphism \(q\) is
\(H^0(T,q^*T_{C'})=0\), since this line bundle has negative
degree. Thus there would be only finitely many possible branch
sets in \(C'\). The sets \(a^{-1}(b(z))\), however, vary with
the surjective map \(b:B\to C\). This is a contradiction.

\emph{The required comparison of completions.}
Let \(S^\circ\) be the complement of the ramification divisor
in \(S\). There is a tower of finite etale covers
\[
 S^\circ\longrightarrow Y-D
 \longrightarrow (C\times B)-\operatorname{graph}(b)
 \longrightarrow U.                                 \tag{A.87.7}
\]
The first two arrows are cyclic of degree three, and the last is
the Galois base change of (A.87.3). The Galois closure of the
composite has a three-group as its group: the Galois closure of a
tower of covers with three-group Galois groups embeds in the
corresponding iterated wreath product. In a fiber of \(U\to C\),
the cover \(S_z^\circ\to F\) has degree nine and likewise has
three-group Galois closure.

Here is the elementary topology-of-completion fact being used.
If \(H\leq E\) contains a normal subgroup of finite odd index,
then \(H^{\mathrm{odd}}\) is the open subgroup given by the
closure of \(H\) in \(E^{\mathrm{odd}}\). For any normal
odd-index subgroup \(N\) of \(H\), first intersect \(N\)
with that fixed normal subgroup of \(E\), and then take its
\(E\)-core. The resulting quotient is an extension of an odd
group by a subgroup of a finite product of odd groups. It has
odd order and its restriction to \(H\) refines the given
quotient. This proves fullness of the induced pro-odd topology.

Suppose now that (A.87.1) were injective. Apply the preceding
fact to both the total cover and its degree-nine fiber cover.
The resulting homomorphism
\[
 \pi_1(S_z^\circ)^{\mathrm{odd}}\longrightarrow
                 \pi_1(S^\circ)^{\mathrm{odd}}        \tag{A.87.8}
\]
is injective. The discrete fiber group is normal in the discrete
total group and the latter surjects onto \(\pi_1(B)\):
\(S^\circ\to B\) is a locally trivial punctured-surface
fibration with connected fiber. Thus its compact pro-odd closure
is normal in the completed total group. These assertions also
follow from the Fadell--Neuwirth sequence for \(U\to C\) and
the pullback covers; see \cite[Section 2.1]{AG_Bellingeri2004}.

For every odd prime \(\ell\), compactification of the punctured
fiber gives a quotient
\[
 \pi_1(S_z^\circ)^{\mathrm{odd}}\longrightarrow
                          H_1(S_z,\mathbb F_\ell).
                                                               \tag{A.87.9}
\]
Its kernel is stable under the discrete total group, because the
family preserves the peripheral loops, and therefore under its
pro-odd closure. The conjugation action on this finite quotient
is continuous. For example this follows by testing the images
of a finite generating set of the pro-odd fiber group. Hence its
image has odd order. Inner conjugation by the fiber group acts
trivially on (A.87.9). Since the total discrete group surjects
onto \(\pi_1(B)\), the image just obtained is precisely the
usual monodromy image on \(H_1(S_z,\mathbb F_\ell)\).
In particular, if
\[
 \Gamma=\operatorname{im}\bigl(\pi_1(B)\to
                    \operatorname{Sp}_{2h}(\mathbb Z)\bigr),
\]
then the reduction of \(\Gamma\) modulo every odd prime has
odd order. No assertion that completion is left exact for an
arbitrary surface fibration has been used.

\emph{The monodromy contradiction, in characteristic zero.}
The Feit--Thompson theorem makes each of those finite reductions
solvable \cite{AC_FeitThompson1963}. The uniform solvable-length
theorem for linear groups supplies an integer \(d(2h)\),
depending only on the dimension and not on the field, that
bounds their derived lengths; see
\cite[Section 1, printed p. 59]{AG_Newman1972}.
Every element of \(\Gamma^{(d(2h))}\) therefore becomes the
identity modulo every odd prime. An integral matrix with this
property is the identity. Thus \(\Gamma\) is solvable.

On the other hand, Deligne's theorem states that the identity
component of the Zariski closure of the monodromy of a smooth
proper complex algebraic family is semisimple
\cite[Corollary 4.2.9(a)]{AG_Deligne1971}.
The Zariski closure of a solvable group is solvable, so that
identity component must be trivial. Hence \(\Gamma\) is finite.
After a finite etale cover of \(B\) the integral local system
is constant. Its period map is then a globally defined
holomorphic map from a compact curve into Siegel upper half
space. In its bounded realization all coordinate functions are
constant. The principally polarized Jacobian is constant, and
Torelli makes the fibers of \(f\) isomorphic. This contradicts
the nonisotriviality already proved and proves (A.87.1).

\emph{Passage to characteristic two.}
This last Hodge argument is used only over \(\mathbb C\).
Let \((C,P)\) be a pointed smooth proper curve over an
algebraically closed field of characteristic two. Choose a
finitely generated subfield over \(\mathbb F_2\) over which
\((C,P)\) is defined, and replace it by its algebraic closure
\(k_0\) inside the given field. Prime-to-two fundamental groups
of these smooth proper pairs are invariant under algebraically
closed extension, so it suffices to work over \(k_0\).
A pointed smooth proper lifting over \(W(k_0)\) exists: the
obstruction group for deforming a pointed smooth curve is
\(H^2(C,T_C(-P))=0\), and a relatively ample line bundle
algebraizes the formal lifting. The countable field \(k_0\)
has a Witt fraction field of cardinality at most that of
\(\mathbb C\); an algebraic closure of this characteristic-zero
field therefore embeds abstractly into \(\mathbb C\). This
gives the complex geometric generic fiber used below.

For this lift, both
\((\mathcal C,\mathcal P)\) and
\((\mathcal C\times\mathcal C,\Delta_{\mathcal C})\)
are smooth proper pairs with relative normal-crossing boundary.
SGA~1 XIII, Corollary~2.9 and the specialization construction in
Section~2.10 identify their maximal prime-to-two fundamental
groups on the geometric generic and special fibers
\cite{AR_SGA1}. In more detail, Corollary~2.9 makes the nonabelian
\(R^1\) for every finite group of odd order locally constant;
its base-change statement and Section~2.10 give these
specialization isomorphisms. Functoriality for
\(x\mapsto(x,\mathcal P)\) identifies the fiber-inclusion maps
as well. Injectivity in characteristic two would therefore imply
injectivity for the complex curve, which was disproved above.

For the completely specified example (A.87.2), a lift requires
no deformation argument. The equation \(y^2+y=x^5\) over
\(\mathbb Z_2\), with the chart at infinity
\[
                         v^2+u^3v=u,
\]
gives the required smooth genus-two model and the section
\(u=v=0\). On the special fiber the finite chart is smooth
because its derivative with respect to \(y\) is one; the
infinity chart is smooth at that section because its derivative
with respect to \(u\) is a unit. Its generic fiber embeds into
\(\mathbb C\), so the same functorial specialization diagram
applies.
\end{proof}

\paragraph{Finite numerical checks and scope.}
For genus two, (A.87.3)--(A.87.7) have
\[
 m=81,\qquad g(C')=4,\qquad g(B)=82,
 \qquad g(S_z)=13,
\]
and the punctured fiber has three punctures and first Betti
number \(28\); its compactification has first Betti number
\(26\). The retained exact-rational program
\path{configuration_family_invariants.py} also verifies
\[
 D^2=-486,\quad K_YD=972,\quad K_Y^2=1944,
 \quad e(S)=3888,\quad K_S^2=9072,\quad
 \operatorname{sign}(S)=432.
                                                               \tag{A.87.10}
\]
The signature is an independent numerical consistency check,
not a substitute for the monodromy argument. The new obstruction
concerns completing the \emph{whole} configuration group to
pro-odd. It does not contradict the exact sequence with full
tame base and only the fiber replaced by its prime-to-two
quotient, nor the exact single-prime pro-\(\ell\) sequences.

\paragraph{A cohomological obstruction to naive cuspidal devissage.}
There is a concrete loss of information when one tries to
iterate maximal cuspidally pro-\(\ell\) reconstruction through
ramified covers. Work over an algebraically closed field of
characteristic two, or over \(\mathbb C\), and take the
pro-odd groups throughout. Let \(C\) have genus \(g\geq2\),
let \(P\in C\), and write
\[
 F=\pi_1(C-P)^{\mathrm{odd}},\quad
 G=\pi_1(C)^{\mathrm{odd}},\quad K=\ker(F\to G),
 \quad N=\ker(K\to K^{(3)}).
\]
Here \(K^{(3)}\) denotes the maximal pro-three quotient,
so \(F/N\) is the maximal cuspidally pro-three quotient over
\(G\). Choose a cyclic etale cubic \(a:C'\to C\), and a
cyclic cubic \(D\to C'\) branched exactly at the three
points of \(a^{-1}(P)\), with ramification index three at
each. Such a cover exists because the degree-three branch
divisor has a cube root in \(\operatorname{Pic}(C')\):
choose a line bundle of degree one and correct its cube
by divisibility of \(\operatorname{Pic}^0(C')\) over the
algebraically closed field.
It has genus \(h=9g-5\).

Put \(D^\circ=D-\pi^{-1}(P)\). Its degree-nine etale cover
of \(C-P\) is visible in \(F/N\), since its Galois closure
has three-group Galois group. Let \(H\leq F/N\) be its
open subgroup, and let \(\overline H\) be the quotient of
\(H\) by the closed normal subgroup generated by its three
puncture inertia groups. Then, for every odd prime
\(\ell\ne3\),
\[
 \dim_{\mathbb F_\ell}H^1(\overline H,\mathbb F_\ell)
       =2g(C')=6g-4
       <18g-10
       =\dim_{\mathbb F_\ell}H^1(D,\mathbb F_\ell).
                                                               \tag{A.87.11}
\]
Thus \(\overline H\) is a proper quotient of the full
group \(\pi_1(D)^{\mathrm{odd}}\), already detected on
finite abelian quotients. For \(g=2\), the two dimensions
are \(8\) and \(26\).

Here is a direct proof of the comparison, including its
surjectivity assertion. Let \(F'\leq F\) correspond to
\(C'-a^{-1}(P)\), and let
\(\chi:F'\twoheadrightarrow\mathbb Z/3\) define
\(D^\circ\). The group \(K\) is contained in \(F'\),
and \(\chi|_K\) is surjective, because a branch inertia
generator lies in \(K\) and maps nontrivially. If
\(I=\ker\chi\), this gives \(F'=IK\). Also \(N\subset I\),
and hence
\[
 1\longrightarrow (I\cap K)/N\longrightarrow
 H=I/N\longrightarrow \pi_1(C')^{\mathrm{odd}}
 \longrightarrow1.                                  \tag{A.87.12}
\]
The left-hand term is a closed subgroup of the pro-three
group \(K/N\). Every homomorphism from \(H\) to
\(\mathbb F_\ell\) therefore factors uniquely through
the last group when \(\ell\ne3\). The puncture inertia
groups in \(H\) lie in the left-hand term, so killing them
does not change these homomorphisms. This proves the first
equality in (A.87.11); the second is the usual surface
cohomology dimension and the computed genus of \(D\).
The natural map from \(\pi_1(D)^{\mathrm{odd}}\) to
\(\overline H\) is surjective by the construction.

This is an obstruction to the specific procedure of taking
an open subgroup of a cuspidally pro-three quotient and then
killing its puncture inertia. It is not a theorem that all
cuspidal quotients together determine insufficient information
by every possible method. In a solvable induction, the next
cuspidally pro-\(\ell\) theorem needs the \emph{full} proper
prime-to-two group of the intermediate ramified compactification.
Equation (A.87.11) shows exactly why the preceding formal
operations do not provide that input. That ramified
compactification is also absent from the original proper
unramified tower.

% Extension of A.87; intended for incorporation into that statement/proof,
% not a second independently numbered configuration claim.

\paragraph{The congruence lemma needed for every excluded prime.}
Let \(\Gamma\subset\operatorname{GL}_N(\mathbb Z[1/M])\) be
finitely generated. Suppose that the identity component \(G\) of its
Zariski closure over \(\mathbb Q\) is a nontrivial semisimple group.
For every prime \(p\) and every integer \(a\geq1\), there are
infinitely many rational primes \(\ell\ne p\), with \(\ell\nmid M\),
such that
\[
       p^a\mid
       \#\operatorname{im}\bigl(\Gamma\longrightarrow
                    \operatorname{GL}_N(\mathbb F_\ell)\bigr).
                                                        \tag{A.87.13}
\]

\begin{proof}
The subgroup \(\Gamma_0=\Gamma\cap G(\mathbb Q)\) has finite index,
is finitely generated, and is Zariski dense in \(G\). Let
\[
     1\longrightarrow Z\longrightarrow \widetilde G
          \xrightarrow{\pi}G\longrightarrow1
\]
be its simply connected central covering, and put \(c=\#Z\).
The connecting map
\[
               \delta:G(\mathbb Q)\longrightarrow H^1(\mathbb Q,Z)
\]
is a group homomorphism. Indeed, for lifts of two rational elements,
the two Galois cocycles multiply because their values lie in the
central subgroup \(Z\). The abelian group \(H^1(\mathbb Q,Z)\)
is killed by \(c\). Therefore the image of the finitely generated
group \(\Gamma_0\) under \(\delta\) is finite. Its kernel
\(\Gamma_1\) has finite index in \(\Gamma_0\), and every element
of \(\Gamma_1\) has a rational lift through \(\pi\).

Choose a finite generating set of \(\Gamma_1\), choose rational
lifts of its generators, and let \(\widetilde\Gamma\) be the
group generated by those lifts. Its image is exactly \(\Gamma_1\).
Its Zariski closure is all of \(\widetilde G\): the image of that
closure under the finite morphism \(\pi\) is closed and contains
the dense subgroup \(\Gamma_1\); equality of dimensions, and
connectedness of \(\widetilde G\), finish the assertion.

Choose a faithful rational representation of \(\widetilde G\).
After enlarging a finite set of primes \(S\), all chosen lifts
are integral over \(\mathbb Z[S^{-1}]\), and \(\pi\), the
original faithful representation of \(G\), and their central
kernel spread out over this ring. We may require that these
models are smooth semisimple group schemes, that the morphism
is a central isogeny of degree \(c\), and that the embedding
of \(G\) into the original matrix group is a closed immersion.
Thus, outside \(S\), reducing \(\pi(\widetilde\gamma)\) gives
the same matrix as reducing the original element of \(\Gamma_1\).
The kernel of this comparison on finite-field points has at
most \(c\) elements.

Nori's strong-approximation theorem, in precisely its
connected and simply connected form
\cite[Theorem 5.4, printed p. 274]{AC_Nori1987}, says that the
congruence closure of \(\widetilde\Gamma\) is open in
\(\widetilde G(\prod_{\ell\notin S}\mathbb Z_\ell)\).
Consequently, after excluding finitely many more primes,
\[
       \widetilde\Gamma\bmod\ell
                      =\widetilde G(\mathbb F_\ell).
\]
Here smoothness gives surjectivity of reduction from
\(\mathbb Z_\ell\)-points, and an open subgroup of a product
contains the full product outside finitely many factors.

Choose a finite Galois number field \(K\) which splits
\(\widetilde G\) and contains the roots of unity of order
\(p^{a+v_p(c)}\). Exclude also the finitely many rational
primes needed to spread the chosen splitting isomorphism
over the ring of integers of \(K\). Infinitely many rational primes split
completely in \(K\). One may use the usual splitting-prime
theorem, or the following elementary sufficient form. If
\(K=\mathbb Q(\alpha)\) with \(\alpha\) integral and minimal
polynomial \(f\), infinitely many primes divide nonzero values
\(f(n)\). Indeed, if all such prime divisors belonged to a
finite set with product \(R\), choose \(b\) with
\(d=f(b)\ne0\). The integer polynomial
\(f(b+dRt)/d\) is congruent to one modulo \(R\) and has
nonconstant absolute value, a contradiction for sufficiently
large integers \(t\). Away from the discriminant and the
index of \(\mathbb Z[\alpha]\), a root of \(f\) modulo a
prime gives a residue-degree-one prime of \(K\). Since
\(K/\mathbb Q\) is Galois, that rational prime splits completely.

Discard the finite exceptional set, and choose any such
\(\ell\ne p\). Then \(\widetilde G_{\mathbb F_\ell}\)
is a nontrivial split semisimple group and
\(p^{a+v_p(c)}\mid\ell-1\). Its split maximal torus has
positive rank, so
\[
       v_p\bigl(\#\widetilde G(\mathbb F_\ell)\bigr)
                          \geq a+v_p(c).
\]
Its image in the original \(\operatorname{GL}_N(\mathbb F_\ell)\)
loses at most \(v_p(c)\) from this valuation, and that image is
the reduction of \(\Gamma_1\). It is a subgroup of the
reduction of \(\Gamma\). This proves (A.87.13).
\end{proof}

\paragraph{Replacement of the pro-odd restriction in A.87.}
Fix an arbitrary prime \(p\), and write \((-)^{(p')}\) for
completion with respect to all finite groups of order prime
to \(p\). For every smooth proper complex curve \(C\) of
genus \(g\geq2\), every \(P\in C\), and
\(U=(C\times C)-\Delta_C\), the map
\[
       \pi_1(C-\{P\})^{(p')}
                 \longrightarrow \pi_1(U)^{(p')}
                                                        \tag{A.87.14}
\]
is not injective. The corresponding geometric prime-to-\(p\)
map is likewise noninjective for every smooth proper curve of
genus at least two over an algebraically closed field of
characteristic \(p\).

\begin{proof}
In the explicit family construction of A.87 choose a prime
\(r\ne p\): one may take \(r=2\) for odd \(p\), and
\(r=3\) for \(p=2\). Replace its cyclic triple cover by a
cyclic degree-\(r\) etale cover \(a:C'\to C\), and take
\(b:B\to C\) to be the full mod-\(r\) homology cover,
of degree \(r^{2g}\). This kills the cross term of the
diagonal class modulo \(r\), and both remaining Kunneth
terms vanish modulo \(r\). The same integral-cohomology and
Picard-divisibility argument therefore gives
\[
       L^{\otimes r}\simeq
       \mathcal O_{C'\times B}\bigl((a\times b)^{-1}\Delta_C\bigr).
\]
The resulting cyclic cover is a smooth proper family
\(S\to B\). Each fiber is a degree-\(r\) cover of \(C'\)
branched at the \(r\) distinct points \(a^{-1}(b(z))\), and
\[
   g(C')=1+r(g-1),\qquad g(B)=1+r^{2g}(g-1),\qquad
   g(S_z)=1+r^2(g-1)+\frac{r(r-1)}2.
\]
It is nonisotrivial by the same finite-Hom argument: a fixed
source curve has only finitely many degree-\(r\) maps to the
fixed hyperbolic curve \(C'\), while their required branch
sets vary through the fibers of the surjection \(a\).

Every cover in the punctured total-space tower has an
\(r\)-group Galois group, and its composite has an
\(r\)-group Galois closure by the wreath-product argument.
The punctured fiber has the same property. An open subgroup
containing a normal subgroup of finite prime-to-\(p\) index
has its full pro-prime-to-\(p\) topology induced from the
ambient group: taking normal cores produces an extension
of a prime-to-\(p\) group by a subgroup of a product of such
groups. This uses no solvability assumption.

Consequently, if (A.87.14) were injective, the completed
punctured-fiber group would be a closed normal subgroup of
the completed total group of \(S^\circ\to B\). For every
\(\ell\ne p\), its compactified homology
\(H_1(S_z,\mathbb F_\ell)\) is a finite quotient with
monodromy-invariant kernel. Conjugation is continuous, the
fiber acts trivially on this abelian quotient, and the
discrete total group surjects onto \(\pi_1(B)\). Therefore
every reduction modulo \(\ell\ne p\) of the integral
monodromy group
\[
       \Gamma\subset\operatorname{Sp}_{2g(S_z)}(\mathbb Z)
\]
would have order prime to \(p\). These are reductions of
the actual integral homology local system, rather than
abstract finite quotients of an unrelated group.

Deligne's semisimplicity theorem used in A.87 makes the
connected Zariski closure of \(\Gamma\) semisimple. It is
nontrivial: a trivial identity component would make
\(\Gamma\) finite, so after a finite etale base change the
period map would be a holomorphic map from a compact curve
to the bounded realization of Siegel space. It would be
constant; Torelli would make the family isotrivial. The
congruence lemma, with \(a=1\), now contradicts the asserted
prime-to-\(p\) orders of all the reductions. This proves
(A.87.14) over \(\mathbb C\).

For characteristic \(p\), lift the pointed smooth curve
as in A.87 and apply the prime-to-\(p\) specialization
isomorphisms for the two smooth proper pairs
\((\mathcal C,\mathcal P)\) and
\((\mathcal C\times\mathcal C,\Delta_{\mathcal C})\).
SGA~1 XIII, Corollary~2.9 and Section~2.10 apply because
all the finite groups in question have order invertible
on the base. Functoriality identifies the two fiber maps.
Thus injectivity of the special-fiber map would imply
injectivity of its complex counterpart, which has just
been contradicted. The Hodge and period-map argument has
been used only in characteristic zero.
\end{proof}

The extension eliminates the full configuration exactness
shortcut for every excluded prime. It does not produce a
collision, and by itself does not decide A1.
The later A.88--A.90 settle A1; universal relative-witness
existence remains a separate question. The characteristic-two proof in A.87
via Feit--Thompson and bounded solvable length remains an
independent alternative to the strong-approximation step.

\paragraph{What the same-prime quotient does retain.}
The preceding cross-prime rank loss must not be confused
with a loss of the entire zeta function. There is a general
elementary complement. Let \(F\) be a profinite group,
let \(K\triangleleft F\) be closed, fix a prime \(r\),
and put \(N=\ker(K\to K^{(r)})\). For every open
subgroup \(I\leq F\) containing \(N\), the natural map
induces an isomorphism
\[
                         I^{(r)}\simeq (I/N)^{(r)}.
                                                        \tag{A.87.15}
\]
Here the superscript denotes the maximal pro-\(r\)
quotient. In particular, the assertion is not restricted
to open subgroups having an \(r\)-group Galois core.

To prove it, take an open normal subgroup \(J\triangleleft I\)
such that \(I/J\) is an \(r\)-group, and set
\(I_0=\operatorname{core}_F(I)\).
The group \(J\cap I_0\) is normal in \(I_0\), and
\(I_0/(J\cap I_0)\) embeds in \(I/J\).
Let \(M=\operatorname{core}_F(J\cap I_0)\).
Then \(I_0/M\) embeds in a finite product of
\(r\)-groups, so it is an \(r\)-group. Since
\(N\triangleleft F\) and \(N\subset I\), we have
\(N\subset I_0\); consequently the image of \(K\)
in \(F/I_0\) is an \(r\)-group. Its image in
\(F/M\) is an extension of an \(r\)-group by a
subgroup of \(I_0/M\), and is therefore an
\(r\)-group as well. The defining property of \(N\)
now gives \(N\subset M\subset J\). Thus every
finite \(r\)-quotient of \(I\) kills \(N\),
which proves (A.87.15).

Apply this to the geometric punctured-curve group and
its cuspidal kernel. For every finite cover visible in
the maximal cuspidally pro-\(r\) quotient, quotienting
by the puncture inertia and then taking the maximal
pro-\(r\) quotient recovers the \emph{full} proper
pro-\(r\) group of its smooth compactification.
Indeed, formation of a maximal pro-\(r\) quotient
commutes with quotienting by the closed normal subgroup
generated by the specified inertia groups. This applies
in particular to the ramified cover \(D\) in
(A.87.11), even though its other-prime cohomology has
been lost.

If the curves and cover are defined over a finite field
of characteristic different from \(r\), this comparison
is compatible with the arithmetic Frobenius action.
It therefore retains \(H^1(D_{\bar k},\mathbb Z_r)\)
with Frobenius, and hence the exact integral Weil
polynomial and zeta function of \(D\). Equality of the
\(r\)-adic characteristic polynomials is equality of
their integer polynomials, since
\(\mathbb Z[T]\hookrightarrow\mathbb Z_r[T]\).
Thus the obstruction in (A.87.11) concerns the full
proper prime-to-\(p\) group needed by a subsequent
different-prime reconstruction theorem; it does not
prevent transfer of these zeta functions. No inference
from that transfer to point separation is made here.
In particular, the Weil polynomial does not by itself
specify the integral \(p\)-primary rational torsion
exponent, or supply R1-compatible point lifts on these
ramified covers.


\subsection{A conditional cuspidal criterion for point separation}

\begin{claim}{A.88}{Proved: a conditional cuspidal separation criterion}
\label{clm:AGcuspidalcriterion}
Let \(k\) be a finite field of characteristic \(p\), let
\(C/k\) be smooth, proper and geometrically connected of
genus at least two, and let \(P,Q\in C(k)\).
Write \(\Pi_C\) for the arithmetic fundamental group with
maximal geometrically prime-to-\(p\) quotient. Fix a prime
\(r\ne p\). Let \(G_P\) and \(G_Q\) be the maximal
cuspidally pro-\(r\) quotients of the corresponding
arithmetic prime-to-\(p\) groups of \(C-P\) and \(C-Q\),
relative to their maps to \(\Pi_C\).

Suppose there is a continuous isomorphism
\[
                  \beta:G_P\xrightarrow{\sim}G_Q
\]
over \(\operatorname{id}_{\Pi_C}\) which carries the
cuspidal decomposition group at \(P\) to a conjugate of
the cuspidal decomposition group at \(Q\). Then \(P=Q\).
The isomorphism is in particular over the identity of
\(G_k\), so it preserves arithmetic Frobenius.
\end{claim}

This is an implication from the existence of the specified
isomorphism. Its proof does not use a configuration-space
fiber exactness theorem. The required construction from
conjugate proper point sections is supplied separately by
Claims~\ref{clm:AGfixedrconfiguration} and
\ref{clm:AGcuspidaltransport}.

\begin{proof}
Assume \(P\ne Q\). We give the argument in five steps.

\emph{1. The point difference is \(p\)-primary.}
Projection of the cuspidal decomposition groups to the
proper quotient shows that the sections of \(P\) and
\(Q\) in \(\Pi_C\) are conjugate. The conjugator
may be made geometric: multiply it by an element of
the target point decomposition group with the same
arithmetic projection, which does not change the
conjugated subgroup. Thus these sections are
geometrically conjugate. Passing
to the geometric abelian quotient and the Jacobian
\(J_C\), their difference has zero Kummer class for
every prime \(\ell\ne p\). The Kummer map on the
\(\ell\)-primary part of \(J_C(k)\) is injective:
the usual Kummer exact sequences, followed by inverse
limit, give the injection of the \(\ell\)-adic
completion of \(J_C(k)\) into
\(H^1(G_k,T_\ell J_C)\). Since \(J_C(k)\) is finite,
this completion is precisely its \(\ell\)-primary part.
Consequently \([P-Q]\) has \(p\)-power order. There are
an integer \(e\geq0\) and \(f\in k(C)^\times\) with
\[
                  \operatorname{div}(f)=p^e(P-Q).
                                                        \tag{A.88.1}
\]
Indeed this divisor is principal over \(\bar k\), and
its rational generator descends by Hilbert~90, since
the divisor is \(k\)-rational.

\emph{2. Every visible finite cover has its zeta function
and its number of missing points preserved.}
For the geometric group of \(C-P\), let \(K_P\) be
the kernel over the full proper geometric prime-to-\(p\)
group and let \(N_P=\ker(K_P\to K_P^{(r)})\).
Equation (A.87.15) applies to every open subgroup
containing \(N_P\). Thus, for every finite cover
visible in \(G_P\), killing its cusp inertia and then
taking its maximal pro-\(r\) quotient gives the full
proper pro-\(r\) group of its smooth compactification.
This construction is natural for the arithmetic action.
The same holds for \(G_Q\).

The isomorphism \(\beta\) identifies the relevant
inertia groups, because it identifies the cuspidal
decomposition groups and their geometric kernels.
It therefore identifies the compactified
\(H^1(-,\mathbb Z_r)\) representations, including
Frobenius. The integral Weil polynomials, hence the
point counts of the compactifications over every
finite constant extension, agree. The missing points
also correspond: their decomposition groups and
residue degrees are given by the cuspidal
double-cosets and their arithmetic projections in
the respective open subgroups. It follows that
corresponding \emph{open} covers have the same
number of rational points over their constant
fields and over every further finite extension.

\emph{3. Noncuspidal point sections are transported.}
Extend constants temporarily so that a chosen
noncuspidal point \(R\) is rational. Let
\(D_R^P\subset G_P\) be its decomposition subgroup.
It maps isomorphically to the absolute Galois group
of the new constant field. Every open subgroup
\(H'\subset G_Q\) containing \(\beta(D_R^P)\)
has full arithmetic projection. The source cover
for \(\beta^{-1}(H')\) has a rational noncuspidal
point over \(R\), so the target open cover has a
rational noncuspidal point by Step~2.

The finite sets of rational points of these target
covers form a directed inverse system with the
finite intersection property: any finite collection
of conditions can be checked in the cover for the
intersection of the corresponding open subgroups.
Compactness gives a compatible point of the
pro-cover corresponding to the closed subgroup
\(\beta(D_R^P)\). Choose a geometric lift in the
universal pro-cover defined by \(G_Q\). Its
stabilizer is the image in \(G_Q\) of an actual
target decomposition subgroup, and is contained in
\[
        \bigcap_{H'\supset\beta(D_R^P)}H'
                                      =\beta(D_R^P).
\]
Both groups map isomorphically onto the constant-field
Galois group, so this containment is equality.
Equivalently, \(\beta\) carries \(D_R^P\), up to
conjugacy, to the decomposition subgroup of an actual
noncuspidal target point.

For the full proper prime-to-\(p\) group, the
exceptional geometric point set is finite
\cite[Proposition~1.8(vi)]{AR_ST2009}.
Outside that set, a proper decomposition subgroup
has a unique point label. Since \(\beta\) is over
the identity of \(\Pi_C\), the target point in
Step~3 is therefore \(R\) itself whenever \(R\)
is outside this finite set and \(\{P,Q\}\).
This statement applies to every closed point after
its finite constant extension.

\emph{4. A difference cocycle turns point transport
into constant Kummer evaluations.}
Let \(M=\mathbb Z_r(1)\). Denote by
\[
  1\longrightarrow M\longrightarrow E_P
     \longrightarrow\Pi_C\longrightarrow1,
  \qquad
  1\longrightarrow M\longrightarrow E_Q
     \longrightarrow\Pi_C\longrightarrow1
                                                        \tag{A.88.2}
\]
the maximal geometrically cuspidally central
quotients of \(G_P,G_Q\). Their geometric
extension classes are the degree-one cycle classes
of \(P\) and \(Q\) in
\(H^2(\Delta_C,M)\simeq\mathbb Z_r\).
Here is a construction avoiding any assertion
about a completed configuration or torsor fiber.
Use the Kummer cycle class of
\(\mathcal O_C(P)\) in
\(H^2(\Pi_C,M)\) to define the extension.
Its restriction to \(C-P\) is canonically
split by the canonical section of that line
bundle, giving a homomorphism from the
punctured-curve group to the extension.
The local residue maps the puncture inertia
onto the canonical generator of \(M\).
Conversely, the maximal geometrically central
cuspidal kernel is generated topologically by
this one inertia image: the full cuspidal
kernel is normally generated by the puncture
inertia, and geometric conjugation has been
killed. Hence the displayed homomorphism
identifies this central kernel with \(M\).
The same construction applies to \(Q\).
The geometric degree can also be checked directly
from the punctured surface presentation
\cite[XIII, Corollary~2.12]{AR_SGA1}: the puncture
word is the product of the \(g\) handle commutators.
After geometric cuspidal centralization it is one central
generator. Its image has infinite \(r\)-adic order:
send the first handle generators to \(1+e_{12}\)
and \(1+e_{23}\) in the upper-unitriangular
Heisenberg group over \(\mathbb Z_r\), and the
remaining handle generators to one. The puncture word
maps to \(1+e_{13}\). The resulting central relator
is the standard degree-one orientation extension.
This verifies the geometric degree-one class directly,
in addition to the cycle-class construction.
These are one-dimensional Kummer and localization
constructions, with no configuration-space fiber
injectivity input.

The isomorphism \(\beta\) induces an extension
isomorphism \(\bar\beta:E_P\to E_Q\).
Its multiplier on the kernel \(M\) is one:
it is over the identity on \(\Delta_C\), and
both degree-one extension classes are the same
generator of \(H^2(\Delta_C,M)\).
On the fiber product \(E_P\times_{\Pi_C}E_Q\)
define
\[
                 \chi(u,v)=\bar\beta(u)v^{-1}\in M.
                                                        \tag{A.88.3}
\]
This is a continuous crossed homomorphism for the
cyclotomic action. Pull it back along the natural
map from the arithmetic prime-to-\(p\) group of
\(C-\{P,Q\}\) to that fiber product. Its residues
at \(P,Q\) are respectively \(+1,-1\), with
the canonical Kummer orientations. Therefore
\[
                         \kappa_r(f)-p^e\chi
                                                        \tag{A.88.4}
\]
has zero residues, kills both inertia groups, and
descends to a class in \(H^1(\Pi_C,M)\).

This last cohomology group consists of constant
Kummer classes. For completeness, the
Hochschild--Serre exact sequence gives
\[
 H^1(G_k,M)\xrightarrow{\sim}H^1(\Pi_C,M),
                                                        \tag{A.88.5}
\]
because \(H^1(\Delta_C,M)^{G_k}=0\): its
Frobenius eigenvalues have nonzero Weil weight,
so none is one. Finite-field
Kummer theory represents the class of (A.88.4) by
some \(a\in k^\times\).

We check that the evaluation of \(\chi\) at
each point \(R\) from Step~3 is zero as a
cohomology class. Work over its finite residue
field so that \(R\) is rational. The target
proper point section satisfies
\[
                       N_{\Pi_C}(D_R)=D_R.
                                                        \tag{A.88.6}
\]
Indeed a normalizer element can be divided by
the unique element of \(D_R\) with the same
arithmetic projection; the quotient is geometric
and centralizes \(D_R\), hence is trivial by
the standard Frobenius-weight centralizer
argument (equivalently
\cite[Lemma~1.7]{AR_ST2009}).

Choose a conjugator between
\(\beta(D_R^P)\) and the target subgroup
\(D_R^Q\). Its projection to \(\Pi_C\)
lies in the normalizer in (A.88.6). Remove that
projection by multiplying by the corresponding
element of \(D_R^Q\). The resulting conjugator
lies in the cuspidal kernel. Its image in
\(E_Q\) lies in \(M\); conjugation by it
changes a section only by a coboundary.
Consequently (A.88.3) restricts to a coboundary
on the actual point section of \(R\).
Evaluating (A.88.4) thus gives
\[
       \kappa_{r,k(R)}\bigl(f(R)/a\bigr)=0
                                                        \tag{A.88.7}
\]
for every closed \(R\) outside the fixed finite
exceptional set and \(\{P,Q\}\).

\emph{5. Constant Kummer evaluations force a
constant function.}
If \(f\) were nonconstant, its map from
\(C_{\bar k}\) to \(\mathbb P^1_{\bar k}\)
would be surjective. Choose a nontrivial
\(r\)-power root of unity \(\zeta\) such that
\(a\zeta\) avoids the finitely many images
of the excluded points. There are infinitely
many choices because \(r\ne p\).
Take a closed point \(R\) with
\(f(R)=a\zeta\). Its finite residue field
contains \(\zeta\), whose nontrivial
\(r\)-primary Kummer class in that field
contradicts (A.88.7). Thus \(f\) is constant,
contradicting (A.88.1) for \(P\ne Q\).
This proves the conditional criterion.
\end{proof}

\paragraph{The separate existence construction and source scope.}
Claims~\ref{clm:AGfixedrconfiguration} and
\ref{clm:AGcuspidaltransport} below supply the required
cuspidal isomorphism from conjugate proper point sections,
using a direct one-prime relative configuration argument.
They do not use the full prime-to-\(p\) configuration
fiber assertion excluded by A.87.
For clarity about the earlier literature route,
Saidi--Tamagawa's Corollary~2.3 constructs the compatible
cuspidally abelian comparison only when the chosen cusp
coordinates lie outside its exceptional sets; its
Theorem~2.7 refers to that comparison. The full
configuration-fiber assertion in the cited Mochizuki
author versions requires a quotient-convention or
source-scope reconciliation with A.87. Neither an
unverified exceptional-coordinate extension of those
statements nor that full configuration assertion is an
input to A.88--A.90.


% Integrated after independent root, criteria, and geometric audits, 2026-11-16.
% Original author draft retained in agent_F/F_fixed_r_configuration_draft.tex.
\subsection{A fixed-prime repair of configuration fiber exactness}

\begin{claim}{A.89}{Proved: fixed-prime cuspidal configuration exactness}
\label{clm:AGfixedrconfiguration}
Let \(p\) and \(r\) be distinct primes. Let \(C\) be a
smooth proper connected complex curve of genus at least two,
let \(x\in C\), and set
\[
 U=(C\times C)-\Delta_C,\qquad V=C-\{x\},\qquad
 \Delta=\pi_1(C)^{(p')},\quad F=\pi_1(V)^{(p')},\quad
 E=\pi_1(U)^{(p')}.
                                                        \tag{A.89.1}
\]
The fiber embedding uses the second projection \(U\to C\).
Define
\[
\begin{aligned}
 K_F&=\ker(F\to\Delta),
 &K_E&=\ker(E\to\Delta\times\Delta),\\
 F_{\mathrm{c}r}&=F/\ker(K_F\to K_F^{(r)}),
 &E_{\mathrm{c}r}&=E/\ker(K_E\to K_E^{(r)}).
\end{aligned}
                                                        \tag{A.89.2}
\]
Then the fiber map induces an exact sequence
\[
 1\longrightarrow F_{\mathrm{c}r}
 \longrightarrow E_{\mathrm{c}r}
 \xrightarrow{\mathrm{pr}_2}\Delta\longrightarrow1.
                                                        \tag{A.89.3}
\]
The same statement holds for geometric etale prime-to-\(p\)
groups over every algebraically closed field of characteristic
\(p\), with the corresponding characteristic relative quotients.

More precisely, every finite quotient of \(F_{\mathrm{c}r}\)
is dominated, on the fiber, by a finite quotient of
\(E_{\mathrm{c}r}\). The finite quotient of the total group
can be chosen to have a normal \(r\)-subgroup whose quotient
is a finite prime-to-\(p\) quotient of
\(\Delta\times\Delta\).
\end{claim}

The distinction from full prime-to-\(p\) exactness is essential.
Here the entire cuspidal kernel, on both sides, is replaced by
its maximal pro-\(r\) quotient for one fixed prime \(r\).
No injection of \(F\) into \(E\) is assumed or concluded.

\begin{proof}
\emph{1. A pro-\(r\) monodromy lemma without a section.}
Let \(1\to L\to H\to B\to1\) be an exact sequence of
discrete groups, where \(L\) is a free group of finite rank
at least two. Suppose the image of conjugation on
\(L/[L,L]L^r\) is an \(r\)-group. Then
\[
                         L^{(r)}\longrightarrow H^{(r)}
                    \quad\hbox{is injective}.             \tag{A.89.4}
\]
Indeed, conjugation by \(H\) acts canonically on the
characteristic completion \(L^{(r)}\). The kernel of
\[
 \operatorname{Aut}(L^{(r)})\longrightarrow
 \operatorname{Aut}(L^{(r)}/\Phi(L^{(r)}))
                                                        \tag{A.89.5}
\]
is pro-\(r\). Here is the finite argument, including the
topological passage. For a finite \(r\)-group \(P\), fix
an ordered basis of \(P/\Phi(P)\), of size \(d\).
There are exactly \(|\Phi(P)|^d\) tuples lifting that basis,
and every such tuple generates \(P\), by the elementary
Frattini criterion. The kernel of the automorphism action
on \(P/\Phi(P)\) acts freely on these tuples: an
automorphism fixing a generating tuple is the identity.
Its order therefore divides \(|\Phi(P)|^d\), so it is an
\(r\)-group. The iterated Frattini subgroups of the
finitely generated pro-\(r\) group \(L^{(r)}\) are
characteristic, open, and cofinal: their images in every
finite \(r\)-quotient eventually vanish. Applying the
finite argument to these quotients proves (A.89.5).

It follows that the closure of the conjugation image of
\(H\) in \(\operatorname{Aut}(L^{(r)})\) is pro-\(r\).
The homomorphism from \(H\) to that closure extends to
\(H^{(r)}\). On \(L^{(r)}\) its composite is the inner
automorphism map. A nonabelian free pro-\(r\) group has
trivial center, so this composite is injective; see, for
example, \cite[Theorem~5.1, proof Case~I, and
Corollary~5.3]{FF_KochloukovaZalesskii2011}.
This proves (A.89.4). In particular no splitting of
\(H\to B\), topological section, or lift of an outer
action is needed.

The compact embedding in (A.89.4) induces the full
pro-\(r\) topology on the fiber. Explicitly, if
\(W\triangleleft L^{(r)}\) is open, there is an open
normal subgroup \(N\triangleleft H^{(r)}\) with
\(N\cap L^{(r)}\subset W\). Thus the finite quotient
\(H^{(r)}/N\) dominates the prescribed quotient of the
fiber; equality with that quotient is unnecessary.

\emph{2. Finite proper covers make the monodromy an
\(r\)-group modulo \(r\).}
Let \(a:C'\to C\) be a connected finite Galois etale
cover with group \(A\) of order prime to \(p\). Form
\[
 Y=(C'\times C')-(a\times a)^{-1}(\Delta_C).
                                                        \tag{A.89.6}
\]
The map \(Y\to U\) is Galois with group \(A\times A\).
The second projection \(Y\to C'\) has fiber
\(V'=C'-a^{-1}(x)\), after choosing a point over \(x\)
in the second factor. Its removed divisor is the disjoint
union of the graphs of all deck transformations of \(a\).
Thus every puncture is labeled throughout this family.
The topological fibration is locally trivial; since the
closed surface \(C'\) is aspherical, its homotopy sequence
gives
\[
 1\longrightarrow\pi_1(V')\longrightarrow\pi_1(Y)
 \longrightarrow\pi_1(C')\longrightarrow1.
                                                        \tag{A.89.7}
\]
The fiber group is free of rank
\(2g(C')+|A|-1\), in particular of rank at least two.

In homology with coefficients in \(\mathbb F_r\), the
boundary loops span a subspace \(W\) and give an exact
sequence
\[
 0\longrightarrow W\longrightarrow H_1(V',\mathbb F_r)
 \longrightarrow H_1(C',\mathbb F_r)\longrightarrow0.
                                                        \tag{A.89.8}
\]
The boundary labels and complex orientations make monodromy
the identity on \(W\). The compactified family is the
product \(C'\times C'\to C'\), so monodromy is also
the identity on the last term. Therefore its image lies
in the additive shear group
\(\operatorname{Hom}_{\mathbb F_r}(H_1(C',\mathbb F_r),W)\),
a finite \(r\)-group. Fiber conjugations already act
trivially on homology. Step~1 consequently proves
\[
             \pi_1(V')^{(r)}\hookrightarrow\pi_1(Y)^{(r)}.
                                                        \tag{A.89.9}
\]

\emph{3. Every finite relative fiber quotient is detected.}
Take a continuous surjection
\(q:F_{\mathrm{c}r}\to Q\), with \(Q\) finite, and
put \(R=q(K_F^{(r)})\). This is a normal \(r\)-subgroup
of \(Q\); the quotient \(A=Q/R\) is a finite quotient
of \(\Delta\). Use its Galois proper cover \(C'\to C\)
in Step~2. On discrete fundamental groups write
\(\Gamma=\pi_1(U)\), \(\Gamma_0=\pi_1(Y)\), and
\(\Lambda=\pi_1(V)\), \(\Lambda_0=\pi_1(V')\).
Then \(\Gamma_0\triangleleft\Gamma\),
\(\Gamma/\Gamma_0=A\times A\), and
\(\Lambda\cap\Gamma_0=\Lambda_0\).
The restriction of \(q\) to \(\Lambda_0\) has image
in \(R\), so factors through a finite \(r\)-quotient
of \(\Lambda_0^{(r)}\).

By (A.89.9) and the final paragraph of Step~1 there is
a finite \(r\)-quotient of \(\Gamma_0\), with kernel
\(J\triangleleft\Gamma_0\), such that
\[
                         J\cap\Lambda_0\subset\ker(q).
                                                        \tag{A.89.10}
\]
Take its normal core in \(\Gamma\):
\[
                  J^*=\bigcap_{g\in\Gamma/\Gamma_0}gJg^{-1}.
                                                        \tag{A.89.11}
\]
This finite intersection is legitimate because \(J\) is
normal in \(\Gamma_0\). The quotient
\(\Gamma_0/J^*\) embeds in a finite product of
\(r\)-groups, hence is an \(r\)-group. Consequently
\(\Gamma/J^*\) is a finite prime-to-\(p\) group with
normal \(r\)-subgroup \(\Gamma_0/J^*\) and quotient
\(A\times A\). Moreover
\[
                 J^*\cap\Lambda\subset\ker(q),
                                                        \tag{A.89.12}
\]
because \(J^*\subset\Gamma_0\) and (A.89.10) applies.
Thus this total-space quotient dominates \(q\) on the
fiber. The inclusion of finite kernels continues to hold
after prime-to-\(p\) completion: it can be checked in
the joint finite image of the two quotient maps, in which
the discrete group has dense, hence surjective, image.

The quotient \(\Gamma\to\Gamma/J^*\) factors through
\(E_{\mathrm{c}r}\), not merely through \(E\).
Indeed its image of \(K_E\) lies in
\(\Gamma_0/J^*\), since \(\Gamma/\Gamma_0\) comes
from the proper product cover. That image is therefore an
\(r\)-group, and the kernel of \(K_E\to K_E^{(r)}\)
is killed. This proves the claimed cofinal finite-quotient
statement.

\emph{4. Injectivity and exactness.}
The fiber map factors through \(F_{\mathrm{c}r}\),
because it takes \(K_F\) to the pro-\(r\) kernel in
\(E_{\mathrm{c}r}\). Every finite quotient of that fiber
group is dominated by one from Step~3, so its kernel is
trivial. This proves the first arrow of (A.89.3).

The ordinary topological fibration
\(C-\{x\}\to U\to C\) is right exact on fundamental
groups. Right exactness survives maximal prime-to-\(p\)
completion: quotienting \(E\) by the closed image of
the fiber has exactly the finite prime-to-\(p\) quotients
of \(\pi_1(C)\), hence is \(\Delta\). The subgroup
killed to form \(E_{\mathrm{c}r}\) is contained in
this fiber kernel. The image of the compact group
\(F_{\mathrm{c}r}\) is closed and is therefore exactly
the kernel of \(\mathrm{pr}_2\). This proves (A.89.3).

\emph{5. Characteristic-\(p\) specialization.}
Spread out the pointed curve over a finitely generated
subfield, pass to a countable algebraic closure, and lift
the smooth pointed proper curve to a mixed-characteristic
complete discrete valuation ring. Curve deformations are
unobstructed, and proper formal lifts algebraize. An
algebraic closure of the fraction field embeds abstractly
in \(\mathbb C\). Both
\((\mathcal C,\mathcal x)\) and
\((\mathcal C\times\mathcal C,\Delta_{\mathcal C})\)
are smooth proper pairs with relative normal-crossing
boundary. SGA~1 XIII, Corollary~2.9 and Section~2.10 give
functorial isomorphisms of their geometric prime-to-\(p\)
groups \cite{AR_SGA1}; this is the same comparison used
in Claim~\ref{clm:AGconfigurationodd}.
These isomorphisms identify \(F,E,\Delta\), their proper
quotients and the fiber maps. They consequently identify
the characteristic relative quotients (A.89.2).
The complex result therefore proves the asserted
characteristic-\(p\) result.
\end{proof}

\paragraph{Arithmetic consequence, with the section specified.}
Let \(C/k\) instead be a smooth proper geometrically
connected curve of genus at least two over a finite field
of characteristic \(p\). Write \(\Pi_C\) for its
arithmetic geometrically prime-to-\(p\) group and let
\(\mathcal E\) be the corresponding group of
\(U=(C\times C)-\Delta_C\). Form its cuspidally
pro-\(r\) quotient relative to
\[
                \Pi_C\times_{G_k}\Pi_C.
                                                        \tag{A.89.13}
\]
The kernel relative to this product is geometric, so this
quotient has geometric subgroup \(E_{\mathrm{c}r}\)
from (A.89.2). If \(P\in C(k)\) and
\(s_P:G_k\to\Pi_C\) is its actual point section,
the natural map from the arithmetic cuspidally pro-\(r\)
group of \(C-P\) identifies it with
\[
       \mathcal E_{\mathrm{c}r}
                  \times_{\Pi_C,\mathrm{pr}_2,s_P}G_k.
                                                        \tag{A.89.14}
\]
Indeed on geometric kernels this is (A.89.3), and on
arithmetic quotients it is the identity of \(G_k\).
The elementary extension comparison therefore makes it
an isomorphism. Since a section is injective, the right
side can equivalently be viewed as the inverse image of
\(s_P(G_k)\) under the second projection. Choices of
basepoint paths give the usual geometric inner ambiguity;
no section of the configuration fibration was postulated.

\paragraph{The boundary comparison.}
Claim~\ref{clm:AGcuspidaltransport} identifies the actual
puncture decomposition group with the intersection with
the global diagonal boundary decomposition group. This
is a separate tame local calculation; it is not silently
included in the geometric fiber theorem.

% Integrated after independent root, criteria, and geometric audits, 2026-11-16.
% Original author draft retained in agent_F/F_fixed_r_boundary_draft.tex.
\subsection{The diagonal boundary and transport of cuspidal groups}

Retain the arithmetic notation of
\ref{clm:AGfixedrconfiguration}: \(C/k\) is smooth proper
of genus at least two over a finite field of characteristic
\(p\), \(r\ne p\), and
\(\mathcal E_{\mathrm{c}r}\) is the cuspidally pro-\(r\)
quotient of the geometrically prime-to-\(p\) group of
\(U=(C\times C)-\Delta_C\), relative to
\(\Pi_C\times_{G_k}\Pi_C\).

\begin{claim}{A.90}{Proved: diagonal boundary and cuspidal transport}
\label{clm:AGcuspidaltransport}
There is a diagonal boundary decomposition subgroup
\(\mathcal D\subset\mathcal E_{\mathrm{c}r}\) with
an exact sequence
\[
 1\longrightarrow\mathbb Z_r(1)\longrightarrow\mathcal D
 \longrightarrow\Pi_C\longrightarrow1,
                                                        \tag{A.90.1}
\]
whose map to the proper product is the diagonal map on
the last term. For every rational point \(P\), choose
compatible geometric paths through this boundary and put
\(D_P=s_P(G_k)\subset\Pi_C\) and
\(\mathcal E_P=\mathrm{pr}_2^{-1}(D_P)\).
The fixed-\(r\) configuration theorem identifies
\(\mathcal E_P\) with the arithmetic cuspidally
pro-\(r\) group of \(C-P\), and under this
identification its puncture decomposition group is
\[
                           \mathcal D\cap\mathcal E_P.
                                                        \tag{A.90.2}
\]

If two rational points \(P,Q\) have geometrically conjugate
proper sections in \(\Pi_C\), there is an isomorphism
between the actual cuspidally pro-\(r\) groups of
\(C-P\) and \(C-Q\), over the identity of
\(\Pi_C\), which carries the first puncture decomposition
group to a conjugate of the second. Consequently
\(P=Q\): rational point sections of the full proper
geometrically prime-to-\(p\) fundamental group are
injective up to geometric conjugacy.
\end{claim}

\begin{proof}
\emph{Construction of the boundary group.}
We give the local tame-cover construction explicitly to
isolate it from any full configuration-fiber assertion.
Let \(X=C\times C\), \(D=\Delta_C\), and
\(N=\mathcal O_X(D)|_D\), the normal line bundle.
For each positive integer \(n\) prime to \(p\), form
the root stack \(X_n=\sqrt[n]{(X,D)}\). Locally where
\(D\) has parameter \(t\), this is the quotient of
\(\operatorname{Spec}(A[z]/(z^n-t))\) by \(\mu_n\).
The reduced inverse image \(D_n\) of \(D\) is the
\(\mu_n\)-gerbe of \(n\)-th roots of \(N\).

A finite etale cover of \(U\) with prime-to-\(p\)
geometric monodromy has tame inertia along \(D\).
Take \(n\) divisible by its finitely many inertia
orders. The absolute Abhyankar lemma, SGA~1 XIII,
Proposition~5.2, says that after the local substitution
\(z^n=t\), its normalization extends uniquely to a
finite etale cover of the whole Kummer chart
\cite{AR_SGA1}. The cited statement already includes
extension over the entire chart; its proof applies
purity after the codimension-one tame calculation.
The unique extensions and their deck
descent data glue to a finite etale cover of \(X_n\).
Restriction to \(D_n\) gives a finite etale cover of
that gerbe. All constructions commute with increasing
\(n\), with field extension, and with pullback
transverse to \(D\). They therefore define the boundary
homomorphism to the tame fundamental group of \(U\),
and hence to its geometrically prime-to-\(p\) quotient.

Here are the exact group assertions needed from this
construction. The fundamental group of the gerbe \(D_n\)
maps onto \(\pi_1(D)\), and its kernel is the image
of \(\mu_n\). This is the elementary gerbe exact
sequence: etale locally on \(D\) the gerbe is
\(B\mu_n\); a cover on which the stabilizer acts
trivially descends to \(D\), and covers of \(D\)
pull back fully faithfully. The inertia is normal,
with the cyclotomic descent action. No injection of
\(\mu_n\) is required here. Passing to compatible
root orders gives the tame boundary group \(\mathcal T\)
and a right exact sequence
\[
 \widehat{\mathbb Z}^{(p')}(1)\longrightarrow\mathcal T
              \longrightarrow\pi_1(D)\longrightarrow1.
                                                        \tag{A.90.3}
\]
The exactness follows either from the same descent
criterion or by inverse limits of the compact gerbe
sequences; the transition maps are surjective, as is
seen in the local maps \(\mu_m\to\mu_n\) for
\(n\mid m\).

Now take the maximal geometrically prime-to-\(p\)
quotients. On geometric groups, right exactness of
completion shows that quotienting the completed
boundary group by the image of inertia gives
\(\Delta_C\). On arithmetic groups the quotients
over \(G_k\) are unchanged. Thus
\[
 \widehat{\mathbb Z}^{(p')}(1)\longrightarrow
 \mathcal T^{\mathrm{geom},p'}\longrightarrow
 \Pi_C\longrightarrow1                             \tag{A.90.4}
\]
is right exact. This step explicitly rules out an extra
kernel arising from the discarded geometric base kernel.

Let \(\mathcal D\) be the image of
\(\mathcal T^{\mathrm{geom},p'}\) in
\(\mathcal E_{\mathrm{c}r}\). Its projection to the
proper product factors through the injective diagonal
\(\Pi_C\to\Pi_C\times_{G_k}\Pi_C\).
Consequently the kernel of \(\mathcal D\to\Pi_C\)
is exactly the image of the inertia in (A.90.4).
Since it is contained in the cuspidal pro-\(r\) kernel,
this inertia image is a quotient of \(\mathbb Z_r(1)\).

This quotient is the full \(\mathbb Z_r(1)\).
To verify this without using a global boundary theorem,
take a geometric fiber \(C-\{P\}\). Its cuspidal
pro-\(r\) group maps onto its maximal pro-\(r\)
fundamental group, which is free on the usual
\(2g\) surface generators. Its puncture inertia
generator is, up to inversion, the product
\(\prod_{i=1}^g[a_i,b_i]\). Map this free group to
the pro-\(r\) upper-unitriangular Heisenberg group
over \(\mathbb Z_r\) by
\[
 a_1\longmapsto 1+e_{12},\qquad
 b_1\longmapsto 1+e_{23},\qquad
 a_i,b_i\longmapsto1\quad(i>1).
                                                        \tag{A.90.5}
\]
The puncture generator maps to \(1+e_{13}\), of
infinite \(r\)-adic order. Thus its \(r\)-inertia
is \(\mathbb Z_r(1)\). The fixed-\(r\) fiber
injection carries this inertia injectively into
\(\mathcal E_{\mathrm{c}r}\). In the local Kummer
charts the fiber meets the diagonal transversely
with multiplicity one, so this is exactly the
diagonal inertia in (A.90.4), with its ordinary
Kummer generator. This proves (A.90.1).

Finally, the same transverse local chart sends the
actual puncture decomposition group of \(C-P\)
into \(\mathcal D\). It lies over \(D_P\), and
it maps isomorphically on inertia and on \(D_P\).
The inverse image of \(D_P\) in (A.90.1) is an
extension of those same two groups. The resulting
map is therefore an isomorphism onto that inverse
image, proving (A.90.2). Compatible paths can be chosen
simultaneously using the common boundary gerbes;
other choices give the permitted geometric inner
conjugacies.

\emph{Transport from conjugate proper point sections.}
Suppose \(P,Q\in C(k)\) have geometrically conjugate
proper sections. Choose
\(\gamma\in\Delta_C\) with
\(\gamma D_P\gamma^{-1}=D_Q\).
By (A.90.1), choose a geometric
\(d\in\mathcal D\) mapping to \(\gamma\).
Conjugation by \(d\) identifies
\(\mathcal E_P\) with \(\mathcal E_Q\), and by
(A.90.2) carries their puncture decomposition groups
onto each other. On the first proper coordinate
it induces \(\operatorname{Inn}(\gamma)\).

The first projection
\(\mathcal E_Q\to\Pi_C\) is surjective, and is
surjective on geometric kernels: it is the proper
quotient of the punctured curve group, by the
fixed-\(r\) theorem. Choose a geometric
\(e\in\mathcal E_Q\) mapping to \(\gamma^{-1}\)
under that first projection. Then
\[
            \beta=\operatorname{Inn}(e)\circ
                         \operatorname{Inn}(d):
                 \mathcal E_P\xrightarrow{\sim}\mathcal E_Q
                                                        \tag{A.90.6}
\]
is over the identity of \(\Pi_C\), is over the
identity of \(G_k\), and carries the first puncture
decomposition group to a conjugate of the second
inside \(\mathcal E_Q\). These are exactly the
hypotheses of Claim~\ref{clm:AGcuspidalcriterion}.
The construction uses the one-prime relative fiber
comparison of Claim~\ref{clm:AGfixedrconfiguration}
and the tame boundary comparison just proved; it does
not invoke full prime-to-\(p\) configuration fiber
exactness. Claim~\ref{clm:AGcuspidalcriterion} now
implies \(P=Q\), as asserted. The argument applies
unchanged after every finite extension of the constant
field.
\end{proof}


% Source: sections/A_computations.tex
\subsection{Exact finite-field computations and their scope}

\begin{claim}{A.91}{Proved by exhaustive exact computation; the genus-two base-field census}
For $q=4,8,16$, the supplied programs enumerate every
$\mathbf F_q$-isomorphism class of smooth proper genus-two curves,
every rational point of each, and every prime-order constant-deck
pointed cyclic cover over $L=\mathbf F_q$, up to automorphism of its
cyclic deck group.  The exact twisted $L$-polynomial and its relative
class number modulo two are recorded for every test.  The totals are
\[
\begin{array}{c|r|r|r|r}
q&\text{curve classes}&\text{pointed pairs}&\text{tests}&
       \text{pairs with a base-field witness}\\\hline
4&146&730&968&613\\
8&1094&9846&14854&8978\\
16&8470&143990&263765&138769
\end{array}
\]
These are all tests over the displayed base field.  For pairs without
a proved least-prime result, the output makes no assertion of a least
prime over arbitrary extensions.  This is the ``list of all tests
tried'' alternative in A4 for these three genus-two censuses.
\end{claim}
\begin{proof}
We give the completeness argument as well as the reproducible
calculation.  In characteristic two every genus-two curve has a
separable hyperelliptic quotient.  A reduced Artin--Schreier equation
has pole orders $(5)$, $(3,1)$, or $(1,1,1)$, by the different formula
$2g-2=-4+\sum(m_i+1)$.  The first two pole sets are rational; the
last set has split, quadratic, or cubic Frobenius type.  We use
respectively
\[
 ax^5+bx^3+cx+d,\quad ax^3+bx+c/x+d,\quad
 ax+b/x+c/(x+1)+d,
\]
and $ax+(bx+c)/(x^2+x+\tau)+d$ or $n(x)/h(x)+d$ in the last
two cases.  Here $\operatorname{Tr}(\tau)=1$, $h$ is a fixed
irreducible monic cubic, $d\in\{0,\tau\}$, and all required
leading residues are nonzero.  The remaining base transformations
are the full affine group for the first type, scalings for the
second, the permutation group of three branch points in the split
case, the order-two stabilizer of the quadratic set, and the
order-three stabilizer of the cubic set.  A generator change
$y\mapsto y+v(x)$ removes even poles; after reduction it leaves
only the constant Artin--Schreier equivalence.  The program takes
exactly these orbits and does not divide by coefficient Frobenius.

This is the classification in \cite[Propositions~8--16,
Theorem~17]{AR_CNP2005}; as an independent check, our type counts are
\[
\begin{array}{c|rrrrr}
q&f_2=0&f_2=1&\text{split}&\text{quadratic}&\text{cubic}\\\hline
4&14&24&14&48&46\\
8&30&112&140&462&350\\
16&70&480&1240&3930&2750
\end{array}
\]
Their sums agree with that theorem's formula
$2q^3+q^2+q-2+8[4\mid\log_2q]$.  The pole calculations also verify
smoothness of each smooth completion and its genus; no test merely
for rational singularities is used.

For an odd-degree model $Y^2+hY=f$ with unique infinity, we enumerate
all reduced Mumford pairs $(U,V)$, with $U$ monic of degree at most
two, $\deg V<\deg U$, and
$U\mid V^2+hV+f$.  For a fixed $U$ the last condition is an affine
linear system over $\mathbf F_2$ in the coefficients of $V$;
\texttt{finite\_jacobian.py} solves it exactly.  Cantor composition
and reduction implement addition.  Each group element has a unique
such representative, so enumeration is exhaustive.

For an ordinary cubic-pole model with a rational point, we move that
point's image to infinity and translate the Artin--Schreier
coordinate.  The resulting model has $h$ monic of degree three,
$\deg f\le5$, and two rational infinities $O_+,O_-$.
The canonical class is $K=O_++O_-$.  Each nonzero Jacobian class
has a unique representative $E-K$ with $E$ effective of degree two;
the zero class is represented by $K$.  We enumerate
$(U,V,a,b)$ with $\deg U+a+b=2$, allowing infinity multiplicities.
For addition the usual finite-divisor composition first cancels
hyperelliptic pairs.  If both remaining infinity multiplicities
are positive, subtract $K$.  Otherwise the reduction is a
Riemann--Roch calculation in
$\langle1,x,x^2,x^3,Y\rangle=L(3K)$: the code solves the vanishing
conditions on the conjugate finite divisor and the required pole
bounds at $O_+,O_-$.  Dividing the norm of the resulting function
by $U$ gives the complementary divisor.  Its Laurent orders give
the two infinity multiplicities.  This derives the implemented
group law directly from principal divisors.  The script
\texttt{check\_even\_jacobian.py} verifies closure, inverses,
commutativity, and sampled associativity on four independent models.
Pointless models have no pointed tests and are explicitly retained.

For every odd prime $\ell\mid\#J_C(\mathbf F_q)$, the program
forms $J/\ell J$, determines a basis, and enumerates every nonzero
linear functional modulo scalar.  These are exactly all the
surjective prime-order characters, with duplicate generators
removed.  If $v$ is its coordinate vector, we evaluate it on
every $[R-P]$.  A.7 gives the exact polynomial
\[
 1+\Bigl(\sum_R\zeta_\ell^{\chi(R-P)}\Bigr)T
       +q\zeta_\ell^{\chi(K-2P)}T^2.
\]
Its relative class number is odd precisely when
$1+\sum_R X^{\chi(R-P)}$ is coprime modulo two to
$1+X+\cdots+X^{\ell-1}$, by A.8.  The program computes that gcd
by binary-polynomial Euclidean division.  Thus every recorded
mod-two value is exact, including primes where two is not a
primitive root.  The JSON records include the quotient basis;
the CSV records give the character coordinates and cyclotomic
polynomial coefficients.  All arithmetic, not only point counting,
is reproduced by the supplied source files and outputs.
\end{proof}

\begin{claim}{A.92}{Proved by exact computation; degree seven for classes 203 and 302}
Let $\alpha^2+\alpha+1=0$, and let $C_{203}/\mathbf F_4$ have equation
\[
 Y^2+(z^2+z+\alpha)Y
  =(z^2+z+\alpha)(z^3+z^2+z+\alpha^2).
\]
Let $C_{302}$ be its coefficientwise square.  Every original
$\mathbf F_4$-rational point on either curve has an L1 witness of
degree seven over $\mathbf F_4$ or $\mathbf F_{64}$.
Together with the no-degree-three and no-degree-five proof in A.20,
the least odd prime for each of these fourteen pointed pairs is seven.
The witnesses are constant-deck Lang covers with the indicated
base-point twist, not merely geometric covers.
\end{claim}
\begin{proof}
Both curves have seven rational points and Jacobian order $28$ over
$\mathbf F_4$.  The tests are reproduced by
\texttt{A4\_203\_extensions.py}; all character bases and exact
polynomials are retained in its six JSON outputs.  For $C_{203}$,
the seven points are infinity and the two points above each of
$z=1,\alpha,\alpha^2$.  At infinity the unique prime-seven quotient
has relative class number $9409=97^2$.  At the two points above
$\alpha$ the value is $4243$.  Both values are odd.  At the four
points above $1$ and $\alpha^2$ it is $25768$, so that particular
base-field quotient fails.  Over $\mathbf F_{16}$ the corresponding
tests still fail for these four points.

Over $\mathbf F_{64}$ the Jacobian order is
$6076=4\cdot7^2\cdot31$.  In the deterministically computed basis
of $J/7J$, each of the four remaining points has a successful
character $(1,3)$; the exact norm of its displayed twisted
polynomial is $120974743723$, an odd integer.  The program also
lists the other successful directions and every unsuccessful one.
The infinity direction $(1,0)$ has norm $129436611529$.
Coefficientwise squaring identifies the entire calculation for
$C_{302}$, and the script independently repeats it.

The finite extension comparison here is explicit: each quotient is
of $J_C(\mathbf F_{64})$, and its Lang pullback is based at the
image of the \emph{original} $P$.  Thus it is a geometrically
connected constant cyclic cover in which $P$ splits, and A.5
applies over that exact field.  We never infer this rational
condition from its holding over a different field.  Finally A.20
excludes the only smaller odd primes, three and five, over all
finite extensions.  This establishes the asserted minimum, not
only a bound obtained from the search.
\end{proof}

\begin{claim}{A.93}{Proved by exact computation and Cartier certificates; all F4 pointed pairs}
Every one of the $730$ pointed pairs in the complete $\mathbf F_4$
genus-two census has a rational L1 witness over a specified finite
extension. Starting with the $117$ pairs without a base-field witness,
the successive tests over $\mathbf F_{16},\mathbf F_{64},\mathbf F_{256}$
leave $52,26,24$ pairs. For each of these last $24$, the least prime is
$11$, with a witness over $\mathbf F_{4^{60}}$. An independent
prime-$113$ witness exists over $\mathbf F_{4^{16}}$.
The full characteristic-zero twisted $L$-polynomials are supplied for
the first $706$ pairs. For the last $24$, we give exact mod-two
Cartier certificates and the relative zeta numerator modulo two;
the full characteristic-zero polynomials have not been computed.
This remaining polynomial requirement of A4 is therefore explicit.
\end{claim}
\begin{proof}
The extension scripts use the same exhaustive Jacobian and character
algorithms as A.91. They embed $\mathbf F_4$ by the first root of
$X^2+X+1$ in the specified larger binary field, map every original
point by that embedding, and base every character at that point.
The irreducible binary moduli are $19,67,283$ for sizes $16,64,256$.
Only pairs already supplied with a witness are omitted at the next
stage. The two final supersingular successes are independently
certified in A.24 by a degree-five Fermat cover.
The remaining model identifiers are
$42,43,46,47,54,55,56,57$ in the complete F4 JSON census.
In Artin--Schreier form the first four are
\[
 y^2+y=x+c/[x(x+1)]+d,
 \qquad c\in\{\alpha,\alpha^2\},\quad d\in\{0,\alpha\};
\]
the last four are the quadratic-pole forms
$y^2+y=x+c/(x^2+x+\alpha)+d$ with the same parameters.
The first four have respectively $7,3,7,3$ remaining points;
only infinity remains on each of the last four.

All eight forms become one of the two coefficient-conjugate split
forms over $\mathbf F_{4^4}$. Indeed choose $s^2+s=\alpha$, replace
$x$ by $z+s$ in the quadratic-pole forms, and remove the resulting
constant by $y\mapsto y+e$ with $e^2+e=s+d$. These choices exist
in $\mathbf F_{4^4}$. The split twists need only the latter
Artin--Schreier constant change. Both witness fields displayed in
the claim contain this field. The original Jacobian orders
$32,12,48,8$ are prime to $11$ and $113$. Consequently the difference
between any two original rational points is divisible by those
primes already in the original Jacobian, so each corresponding
Lang character has trivial base-point twist on those points.
The coefficient-conjugate form is handled by squaring coefficients.

The exact Cartier calculation, its normalization at infinity,
and the two independent certificates are proved in A.30--A.33.
For degree $11$ the relative zeta numerator reduces to
$1+T^4+T^{10}$, whose value at $1$ is one in $\mathbf F_2$.
The degree-$113$ relative numerator is likewise explicitly retained
and has value one. These are proofs of the required rational
relative-class-number parity, rather than point-count estimates.
The no-$3$, no-$5$, and no-$7$ arguments in those statements prove
the asserted minimum $11$ over all finite extensions.

All original finite-field tests, including failures, are retained
in \texttt{A4\_F4\_extensions.csv} and
\texttt{A4\_F4\_F256.csv}. The two large-field certificates retain
an irreducible field modulus, the torsion Mumford divisor, every
nonzero multiple, the tangent coefficients and the Frobenius test.
The independent pure-Python verifier checks these without the
optional compiled arithmetic helper. The combined witness table
identifies these certificates separately from full $L$-polynomial
records; it does not invent the missing characteristic-zero data.
\end{proof}

\begin{claim}{A.94}{Proved; complete comparison with the geometric R6 classes}
The $146$ arithmetic forms in the $\mathbf F_4$ census map to
exactly $64$ geometric classes: $4$ of two-rank zero, $12$ of
two-rank one, and $48$ ordinary classes. The file
\texttt{A4\_R6\_class\_comparison.csv} gives every geometric class,
every arithmetic model above it, every original rational point's
witness coverage, and the primitive annihilator reported in R6.
The distinguished ordinary classes have the following arithmetic
identifiers:
\[
\begin{array}{c|l}
203&66,67\\
302&64,65\\
122&46,47,56,57\\
133&42,43,54,55.
\end{array}
\]
This verifies the comparison and all original-point rational
certificates. The earlier primitive certificates themselves are
not asserted to have been rerun: their auxiliary coefficient files
were not among the supplied background files.
\end{claim}
\begin{proof}
We compute the geometric invariants, rather than identify arithmetic
models by a suggestive label. For the ordinary normalized equation
$y^2+y=ax+b/x+c/(x+1)$, the unordered triple $\{a,b,c\}$ is
the complete geometric datum. The key $s_1s_2s_3$ records its
elementary symmetric functions. For a quadratic-pole model
$ax+(bx+c)/(x^2+x+\alpha)+d$, choose the two roots $r,r+1$.
Translation gives the triple $(a,br+c,b(r+1)+c)$.
For a cubic-pole model $n(x)/h(x)$, with distinct roots
$r_0,r_1,r_\infty$, put
\[
 K=\frac{r_1+r_\infty}{r_1+r_0},\qquad
 \rho_i=\frac{n(r_i)}{h'(r_i)}.
\]
The substitution $x=(r_\infty z+Kr_0)/(z+K)$ gives
\[
 (a,b,c)=\left(
 \frac{\rho_\infty}{K(r_0+r_\infty)},\quad
 \frac{K\rho_0}{r_\infty+r_0},\quad
 \frac{(1+K)\rho_1}{r_\infty+r_1}\right).
\]
These expressions follow by comparing the simple pole residues at
$z=\infty,0,1$; the remaining term is constant and is removed
geometrically by an Artin--Schreier change.

For two-rank one, $ax^3+bx+c/x$ has normalized pair
$(ac^3,bc)$. For two-rank zero, $ax^5+bx^3+cx$ has invariant
$b^5/a^3$. The completeness of these invariants, including
geometric rather than arithmetic isomorphism, is precisely
\cite[Definition~19, Theorems~20 and~22, Corollary~23]{AR_CNP2005}.
The first two normalizations also follow directly by scaling and
permuting the branch coordinates as above.

The comparison program evaluates all formulas exactly in a common
field $\mathbf F_{4096}$, checks that the symmetric invariants lie
in its specified $\mathbf F_4$ subfield, and obtains the displayed
partition with $64$ keys. It then joins the arithmetic model
identifiers to the independently computed witness table; every
one of the $730$ original points has a successful record.
The R6 column is a transcription of AP-A.C2.22, AP-A.C2.25 and
AP-A.C2.26, labelled as such. In particular its annihilator $15$
for classes $122,133$ is a primitive-factor statement, whereas
the degree-$11$ certificates here are whole-relative rational
statements. Neither column is silently substituted for the other.
\end{proof}

\begin{claim}{A.95}{Proved; exact cubic classification in all three censuses}
In characteristic two, the locally exact differential theta divisor
of a genus-two curve contains a nonzero point of order three
if and only if the curve is ordinary and its normalized geometric
coefficients satisfy
\[
                         (a+b+c)^3=abc.
\]
Consequently the all-extension cubic-existence criterion A.25
is decidable from this invariant and the base-field cubic tests.
In the complete pointed censuses, the numbers are
\[
\begin{array}{c|r|r|r}
q&\text{original pairs}&\text{least prime is }3&
   \text{no cubic witness over any extension}\\\hline
4&730&264&466\\
8&9846&3240&6606\\
16&143990&40284&103706.
\end{array}
\]
For $q=4$, explicit characters and full twisted polynomials are
retained for all $264$ positive pairs, over $\mathbf F_4$,
$\mathbf F_{16}$ or $\mathbf F_{256}$. For the other fields the
table is an exact existence and minimum classification; it does
not pretend that every additional character polynomial has been
computed. Whenever a retained degree-five witness has a negative
cubic classification, five is its proved all-extension least prime.
\end{claim}
\begin{proof}
In the ordinary normalized model let $W_0,W_1,W_\infty$ be
the branch points and put $B_\Theta=W_0+W_1-W_\infty$.
The theta divisor consists of $[R-B_\Theta]$. If such a nonzero
class has order three, there is a function with divisor
$3R-3B_\Theta$. Multiplying it by $x(x+1)$ gives a function
with divisor $3R-W_0-W_1-W_\infty$. Riemann--Roch gives
\[
 L(W_0+W_1+W_\infty)=\langle1,y\rangle,
\]
since this divisor has degree three and $g=2$.
Write the function, after scaling, as $y+\beta$, and write
$t=x(R)$ and $\delta=\beta^2+\beta$. The fibre at $R$ is
unramified: a branch point would give a two-primary theta
class, which cannot have order three. Its norm identity is
\[
 b+(b+c)x+ax^2+ax^3+\delta(x^2+x)=a(x+t)^3.
\]
Comparison of coefficients gives
\[
 b=at^3,\qquad c=a(t+1)^3,\qquad \delta=a(t+1).
\]
It follows that $(a+b+c)^3=abc$.

Conversely suppose that equality holds. If $b+c\ne0$, put
$t=(a+c)/(b+c)$; substitution, using the assumed equality,
gives $b=at^3$ and $c=a(t+1)^3$. If $b=c$, the equality forces
$a=b=c$, and either root of $t^2+t+1$ has the same property.
In both cases $t\ne0,1$. Choose $\beta$ with
$\beta^2+\beta=a(t+1)$ and let $R=(t,\beta)$.
The displayed norm identity proves
\[
 \operatorname{div}\left(\frac{y+\beta}{x(x+1)}\right)
                      =3(R-B_\Theta).
\]
The class is nonzero, since $2B_\Theta$ is canonical, whereas
$R$ is not a branch point of the hyperelliptic map. It therefore
has exact order three.

In two-rank zero or one, the locally exact differential line is
$\mathcal O(W)$ for a branch point $W$: in the single-pole model
use $\operatorname{div}(dx)=2W_\infty$, and in the $(3,1)$
model use $\operatorname{div}(dx)=2W_0$.
Thus theta is $C-W$. A nonzero point of order three on theta
would give a function with sole pole of exact order three at
the Weierstrass point $W$. The gap sequence $1,3$, or directly
the basis $L(3W)=\langle1,x\rangle$ after making $W$ infinity,
excludes this. This proves the geometric criterion in every stratum.

A.25 now says that a cubic witness exists precisely when this
geometric condition holds or a base-field cubic character succeeds.
The latter are all enumerated in A.91. The program
\texttt{A4\_cubic\_classification.py} computes the invariant
for every arithmetic form and joins that decision to every
original point. Its exact field operations and its full table
give the three counts in the claim. For the ordinary quadratic
model the test is
\[
 (a+b)^3=a(\tau b^2+bc+c^2),\qquad h=x^2+x+\tau;
\]
for cubic-pole models it uses the residue normalization of A.94.
The required splitting fields have sizes at most $q^3$.

For the $59$ positive $\mathbf F_4$ pairs lacking a base-field
cubic witness, the extra exact character tests supply $49$
over $\mathbf F_{16}$ and the remaining $10$ over
$\mathbf F_{256}$. Both successful and unsuccessful tests are
retained. In particular the latter ten are not assigned the
incorrect smaller field merely from the geometric theta condition.
Finally there is no smaller admissible prime than three; after
three is excluded, a successful degree-five record proves its
minimum as stated.
\end{proof}

\begin{claim}{A.96}{Open; the remaining computational scope of A4}
The base-field pointed odd-prime tests for smooth plane quartics over
$\mathbf F_3$ and $\mathbf F_5$ are complete in A.72 and A.74,
including independent exact integer norm audits. The complete
consolidated table is \path{code/A4_quartic_witnesses.csv}.
However, least primes over arbitrary finite extensions have not been
determined for every genus-two or quartic pointed pair outside the
explicitly proved cases. The full characteristic-zero character
polynomials for the large-field degree-$11$ and degree-$113$
genus-two certificates have not been computed; their full mod-two
certificates and the precise fields are retained in A.31.
The exhaustive tables and successful finite tests do not replace
either missing minimum or missing characteristic-zero polynomial.
\end{claim}

% Source: sections/A_structure.tex
\subsection{Exceptional decomposition groups and Frobenius classes}

\begin{claim}{A.97}{Known, checked; the exact Sa\"idi--Tamagawa statement}
Let $C/k$ be proper hyperbolic over a finite field, and put
$\Sigma=\{\ell:\ell\ne p\}$.  On the universal geometrically
pro-$\Sigma$ covering, let $\widetilde E_C$ consist of points whose
decomposition subgroup is commensurate with that of another point;
let $E_C$ be its image in the set of closed points of $C$.
The following are the statements supplied by the two papers.
\begin{enumerate}
\item In \cite[Definition~1.5, Lemma~1.7, Proposition~1.8(iii)--(vii),
Definition~1.10 and Corollary~1.11]{AR_ST2009}, the commensurability
map is injective above a fixed geometric point, is injective above a
rational point, has finite fibres, and has finite exceptional image
$E_C$ when $\Sigma$ is cofinite.  Decomposition subgroups of closed
points are recovered injectively outside $E_C$.
\item For a $\Sigma$-covering $f:C'\to C$,
\cite[Proposition~2.11]{AR_ST2017} proves
$E_{C'}=f^{-1}(E_C)$ if $f$ is \'etale.  If $f$ is tame Galois,
it proves $E_{C'}\subset f^{-1}(E_C)\cup R_f$, where $R_f$ is the
non-\'etale locus.
\item With $J_C$-largeness as defined in that paper,
\cite[Proposition~2.12]{AR_ST2017} places the exceptional points over
a proper subfield $F\subset\bar k$, finite of degree at most two over
$F_0=\bar k^{\ker\rho_{J_C,\Sigma'}}$.  The degree is at most one
when $C$ is nonhyperelliptic.  Proposition~2.13 gives a prime $\ell$
and finiteness of the exceptional points rational over each finite
extension of the unique $\mathbf Z_\ell$-extension of $k$; the
precise hypothesis is almost $\Sigma$-separatedness, with the stated
$J_C$-large refinements on the choice of $\ell$.
\end{enumerate}
These results do not assert $E_C=\varnothing$.
\end{claim}


\begin{proof}
We use their commensurability equivalence relation, rather than replace
it by equality of arbitrary chosen subgroup representatives.  The
stated hypotheses apply because the complement of our $\Sigma$ is
the single prime $p$.  For the first item the crucial centralizer
statement, their Lemma~1.7, applies to a closed subgroup $Z$ whose
projection to $G_k$ is injective and open: projection embeds the entire
commensurator of $Z$ in $G_k$.  Its intersection with the geometric
kernel is consequently trivial.  This proves injectivity above a
fixed geometric point.  The finite-fibre statement above a closed
point follows by its finitely many geometric conjugates.  The proof
of global finiteness first passes to a nonhyperelliptic cover, then
uses finiteness of the appropriate torsion intersection with the
difference variety in its Jacobian.  It bounds the exceptional set;
it does not show that this intersection is empty.

For the second item the universal covering is the same after passing
to an open subgroup.  Commensurability is unchanged by intersecting
with it, giving equality of the exceptional inverse images in the
\'etale case.  The tame argument allows the stated ramification locus,
which is why equality is not claimed there.

For the third, the proof of Proposition~2.12 uses
$\delta:C\times C\setminus\Delta\to J_C$,
$(x,y)\mapsto[x-y]$.  Its geometric fibres have cardinality at most
$e$, with $e=2$ in the hyperelliptic case and $e=1$ otherwise.
An exceptional pair has difference in the $\Sigma'$-primary torsion,
which is rational over $F_0$.  Thus both coordinates are rational
over a subextension of degree at most $e$.  This is the actual
descent comparison producing $F$.  Almost separation means exactly
that the exceptional residue fields are contained in such a proper
$F$.  Choosing a nonzero $\mathbf Z_\ell$ projection of its Galois
subgroup gives Proposition~2.13's finite-intersection property and
the asserted finiteness.  None of these comparisons removes the
possible exceptional points.
\end{proof}

\begin{claim}{A.98}{Proved; A1 in the exceptional-set language}
For a fixed proper hyperbolic $C/k$, absence of collisions after
\emph{every} finite constant extension is equivalent to $E_C$ being
empty.  Universal A1 is therefore equivalent to emptiness of all these
exceptional sets.  Absence of collisions only among $C(k)$ is the
weaker assertion that no two distinct rational-point decomposition
sections become geometrically conjugate at that field.
By A.3 all these exceptional sets are empty for
$\Sigma=\{\ell:\ell\ne p\}$.  This last conclusion is a consequence
of the proof in the present return, not an assertion attributed to
either of the cited papers.
\end{claim}
\begin{proof}
After choosing compatible lifts, conjugate rational-point sections give
equal decomposition subgroups at two different geometric points.
They are commensurate, hence exceptional.  Conversely, take an
exceptional pair upstairs.  Their geometric images are distinct by
A.97.  Their decomposition subgroups have open intersection, because
they are commensurate.  Enlarge the finite constant field so that both
geometric points are rational and the arithmetic quotient of this
intersection contains the new $G_L$.  Each decomposition group
projects injectively for a proper curve.  Restricting the two inverse
sections to $G_L$ therefore gives the same subgroup and the same
section.  These are sections of distinct $L$-points.  This comparison
also explains why the single-field and all-extension statements must
be kept separate.

The ``Frobenius'' terminology of L7 is exact with this qualification:
$G_k=\widehat{\mathbf Z}$ has a distinguished arithmetic generator,
and a continuous section is determined by its value on that generator.
Geometric conjugacy of these elements is geometric conjugacy of the
sections.  Recovering the curve from an outer isomorphism of its
fundamental group supplies no extra argument for excluding the
exceptional pair just described.  The full abelian class-field map
does separate points, but its prime-to-$p$ quotient leaves precisely
the possible $p$-primary difference in A.1.
Applying A.3 after the finite constant extension used above excludes
every exceptional pair and proves the asserted emptiness.
\end{proof}

\subsection{The field-faithfulness margin of scan 36}

\begin{claim}{A.99}{Proved; the geometrically prime-to-$p$ field functor is faithful}
Let $F,K$ be fields finitely generated over $\mathbf F_p$.  Let
$E_F^{(p')}$ and $E_K^{(p')}$ denote their absolute Galois groupoids
with the kernel over $\overline{\mathbf F}_p$ replaced by its maximal
prime-to-$p$ quotient, retaining the maps to $BG_{\mathbf F_p}$.
If two $\mathbf F_p$-embeddings $u,v:F\hookrightarrow K$ induce
isomorphic morphisms of these groupoids \emph{over}
$BG_{\mathbf F_p}$, including the specified relative commutation
isomorphisms, then $u=v$.
Thus the margin in cote~149, scan~36, concerning Th\'eor\`eme~2
(scan~22), has a correct proof using the genus-zero target
$\mathbf G_m$ and its products.  This assertion does not settle A1
for proper curves.
\end{claim}
\begin{proof}
For $n$ prime to $p$, the torsor $X^n=a$ of $a\in K^\times$ has
geometric monodromy contained in $\mu_n$, hence its Kummer class
factors through $E_K^{(p')}$.  The coefficient group $\mu_n$ is
pulled back from $BG_{\mathbf F_p}$.  We claim that
\[
 K^\times\longrightarrow
       \prod_{(n,p)=1}H^1(E_K^{(p')},\mu_n)
                                                        \tag{A.99.1}
\]
is injective.  If $a$ is killed by every Kummer class, it is an
$n$th power in $K$ for every such $n$.  On a normal proper model
over the finite algebraic constant field $k_c$, every divisorial
valuation of $a$ is divisible by all positive integers prime to
$p$, hence is zero.  Therefore $a\in k_c^\times$.  An $n$th root
in $K$ of a nonzero constant is algebraic over $k_c$, and so lies
in $k_c$.  Taking $n=\#k_c^\times$, which is prime to $p$, forces
$a=1$.  If $K$ itself is finite, the last argument alone suffices.

For every $a\in F^\times$, functorial pullback of its Kummer
torsors under the two relative groupoid morphisms gives equal
classes in (A.99.1).  The relative commutation isomorphism is what
identifies their coefficient systems without an uncontrolled
cyclotomic automorphism.  Thus $u(a)/v(a)=1$ by injectivity.
Both embeddings also preserve zero, proving equality.

Geometrically this test is evaluation of the coordinate of
$\mathbf G_m=\mathbf P^1\setminus\{0,\infty\}$.  To test a finite
list of nonzero generators one takes its product.  Each multiplicative
coordinate separates by the same Kummer argument; no proper-curve
section assertion has entered the proof.

Absolute Frobenius requires care.  It is isomorphic to the identity
on the absolute \'etale topos, but the isomorphism induces arithmetic
Frobenius on the base fibre functor.  It is generally not an
isomorphism between the two \emph{specified relative} morphisms.
Erasing that commutation datum, or taking the maximal prime-to-$p$
quotient of the entire arithmetic group, would change the assertion
we have proved.  Scan~22's relative-groupoid convention supports
precisely the convention retained here.
\end{proof}

\subsection{The norm module and its necessary point-origin condition}

\begin{claim}{A.100}{Proved; construction and additivity of the norm module}
Fix $C/k$ and $P\in C(k)$.  Over $\bar k$, take all finite connected
Galois prime-to-$p$ covers $D_N\to C_{\bar k}$ indexed by open normal
subgroups $N\triangleleft\Pi_C$, with compatible lifts $P_N$.
There is a well-defined module
\[
 M_P=\varprojlim_N J_{D_N}(\bar k)[p^\infty]
                \quad\text{under norm},
 \qquad \mathbf Z_p[[\Pi_C]]\times M_P\longrightarrow M_P.
\]
Arithmetic Frobenius acts semilinearly: if $\alpha$ is its action on
$\Pi_C$, then $F(am)=\alpha(a)F(m)$.  It need not commute with the
deck action.  The module itself need not be a torsion group of bounded
exponent, although each coordinate is $p$-primary torsion.

If $P,Q,R$ have equal point-section classes, their compatible lifted
differences satisfy $T_{PR}=T_{PQ}+T_{QR}$ in this same based module.
The formula is canonical after choosing the compatible based tower;
independently based towers require their path identification first.
\end{claim}
\begin{proof}
The finite maps $D_{N'}\to D_N$, for $N'\subset N$, induce norm
homomorphisms on Jacobians.  They compose transitively and preserve
$p$-primary torsion.  Deck actions commute with the norms through the
quotient maps $\Pi_C/N'\to\Pi_C/N$.  At a fixed coordinate $m_N$
of finite order $p^a$, a completed group-ring element acts through
$(\mathbf Z/p^a)[\Pi_C/N]$.  This defines the stated action and its
compatibility, independently of the representative in the completed
ring.

The section $s_P$ acts on the entire pointed tower; it sends the
index $N$ to $\alpha(N)$ and the chosen point to the chosen point.
Galois functoriality of norm therefore defines $F$ on the inverse
limit and gives the semilinearity formula.  Each $N$ has a finite
orbit under $F$.  A power $F^m$ stabilizing $N$ makes this pointed
cover a cover over $k_m$; at that coordinate, $F^m$ is its usual
arithmetic Frobenius on the Jacobian.

A.1 and uniqueness give compatible lifts $Q_N,R_N$ whenever the
sections collide.  In the divisor class group of $D_N$ the identity
\[
 [R_N-P_N]=[Q_N-P_N]+[R_N-Q_N]
\]
is exact.  Norm sends each term to the same term at the lower level,
so the identities give the claimed one in $M_P$.  Equality of
point-section classes supplies a fixed path.  Its geometric
centralizer ambiguity is trivial by the Lemma~1.7 comparison in
A.97.  In particular no unrelated choice of origins is silently
added to the formula.
\end{proof}

\begin{claim}{A.101}{Proved; what a collision gives in the module}
A collision of $P\ne Q$ gives a nonzero $F$-fixed vector
$T_{PQ}\in M_P$ whose coordinate at every level is a nonzero
point difference.  Every odd-order geometric deck group acts freely
on that coordinate.  In characteristic two this covers every level.
In odd characteristic the same conclusion for all prime-to-$p$
levels follows using the independently established hyperelliptic
case of R4.
The converse from an arbitrary free-orbit vector in $M_P$ is not
proved: the coordinate must come from actual points, compatibly.
\end{claim}
\begin{proof}
The torsion and compatibility assertions follow from A.1.  The Abel
map of a positive-genus curve is injective on geometric points, so
none of these differences is zero.  Galois uniqueness makes the
family $F$-fixed, even at levels that Frobenius permutes.  For odd
deck groups, A.4 proves freeness.

For completeness, suppose $p$ is odd and a nonidentity deck
transformation fixes the point difference at another level.  Its
two effective degree-two divisors in A.4 are distinct and disjoint,
so $D_N$ is hyperelliptic.  The lifted points still have equal
sections in its full geometrically prime-to-$p$ extension, after a
finite field extension defining all the data.  The odd-characteristic
hyperelliptic case of R4, with its prime-to-$p$ correspondence checked
in A.56 and its source audit, excludes such a collision.  Hence the
stabilizer is trivial at these levels too.  This last step uses
that established geometric theorem, rather than incorrectly apply
the odd-order argument to an involution.

None of the module construction forces a general torsion vector to
be represented by a divisor $Q_N-P_N$.  R3's small-support function
argument depends on that very property.  The abstract finite-module
example in A.54 makes clear why a faithful or free action alone
cannot replace it in that proof.
\end{proof}

\begin{claim}{A.102}{Proved; the rational tests as module projections}
Let $D\to C_{k_m}$ be a geometrically Galois prime-to-$p$ cover with
a chosen rational lift of $P$.  The projection of $M_P^F$ to its
Jacobian is contained in $J_D(k_m)[p^\infty]$.  Every homomorphism
$h:J_D\to B$ defined over $k_m$ with $B(k_m)[p]=0$ kills that
projection.  If $D$ is a nontrivial odd-degree L1 witness, its deck
group fixes the projection of every element of $M_P^F$.
Consequently $M_P^F$ has no element with free deck orbit at
\emph{every} level whenever there is an L1 witness for $(C,P)$
over any finite extension.

In particular this nonexistence assertion holds for every original
$\mathbf F_4$-rational point on every genus-two curve over $\mathbf F_4$,
by A.91--A.93. The first requested classes $203$ and $302$ already
follow from the degree-seven witnesses of A.92 and A.20.
\end{claim}
\begin{proof}
Choose an identification of the given pointed geometric cover with
the level $N$.  Its descent is the one induced by the power $F^m$
of the chosen point section, so an $F$-fixed family has an
$F^m$-fixed coordinate.  Fixed geometric torsion points are exactly
rational torsion points over $k_m$.  The assertion for $h$ follows
because its image is a rational $p$-primary point.  For L1,
$e=|G|^{-1}\pi^*\pi_*$ is a $k_m$-defined projector on this
$p$-primary group.  The complementary part vanishes by the relative
class-number condition, exactly as in A.5.  Thus every projection
is $G$-invariant, including a possible zero projection.  Such a
coordinate cannot have a free orbit for nontrivial $G$.

This proof requires all the specified levels in the free-orbit
condition.  Although characteristic subgroups form a cofinal system
for constructing the inverse limit, freeness of a vector does not
in general descend under a norm to a quotient level.  We therefore
do not replace ``every level'' by ``a chosen cofinal system'' in
this assertion.  The projections used by L5 are likewise the typed
maps $h\circ\operatorname{pr}_D$ just proved to vanish; a
correspondence separating actual points adds the further geometric
information stated in A.49.
\end{proof}

\begin{claim}{A.103}{Open; remaining scope of the ghost-module proposal}
For every genus-two curve $C$ in characteristic two and every
rational point $P$ over a finite field of definition, $M_P^F$ has no
nonzero vector with free deck orbit at every normal level.
The implication from this assertion to A1 is proved.  A converse,
or a proof that arbitrary such vectors are compatible differences
of points, is not established.  Every original $\mathbf F_4$ pointed pair, including the first
requested classes $203$ and $302$, is settled by A.102.
\end{claim}

\subsection{A rational exponent obstruction on a nonabelian cover}

\begin{claim}{A.104}{Proved; small annihilators of the actual rational torsion group}
Let $L$ be a finite field of characteristic $p$, let $C/L$ be
smooth, proper and geometrically connected with $g(C)\ge2$, and
let $P\in C(L)$. Let $\pi:D\to C$ be a connected finite etale
Galois cover with constant nonabelian group $G$ of order $m$
prime to $p$, and choose $P'\in D(L)$ above $P$.
The group
\[
                 H=J_D(L)[p^\infty]
\]
is a finite $\mathbf Z[G]$-module, with the action induced by
the $L$-automorphisms of $D$. Suppose a nonzero element
$a=\sum_g a_g g\in\mathbf Z[G]$ annihilates $H$, and put
$d=\sum_g|a_g|$. If
\[
                    (d,m)=1,\qquad m>(d-1)^2,
\]
then $P$ collides with no other point of $C(L)$ in the full
geometrically prime-to-$p$ arithmetic fundamental group.

In particular, if $p^e$ is the exponent of $H$, with $p^e=1$
when $H=0$, it suffices that
\[
                         m>(p^e-1)^2.                 \tag{A.104.1}
\]
For characteristic two, a constant pointed nonabelian cover
of degree $21$ meets this condition whenever its rational
two-primary Jacobian group has exponent at most four.
\end{claim}
\begin{proof}
A hypothetical collision gives, by A.1 over the same field
$L$, a point $Q'\in D(L)$ above the distinct point $Q$ and
a class $T=[Q'-P']\in H$. The integral group-ring element
defines an $L$-endomorphism
$a:J_D\to J_D$, whose restriction to the finite group $H$
is zero by hypothesis. Thus $a[Q'-P']=0$. After scalar
extension to $\overline L$, all hypotheses of A.53 are
satisfied. That theorem proves that this particular divisor
class cannot be annihilated by $a$, a contradiction.

For the stated numerical consequence take $a=p^e\cdot1$.
Its action on $H$ is zero, its coefficient norm is $p^e$,
and $(p^e,m)=1$ because $m$ is prime to $p$. This proves
(A.104.1). For $p=2$, $m=21$ and $2^e\le4$, we have
$21>(4-1)^2=9$.

The exact field is essential. Passing to a larger field to
make a deck group constant can enlarge $H$ and its exponent,
so no such extension is made silently. The criterion only
requires an annihilator of the finite group of rational
torsion points; it does not require the image abelian variety
$a(J_D)$ to have no rational $p$-torsion. For the scalar
choice $a=p^e$, that image is the whole Jacobian. The
contradiction also uses that $T$ is a difference of points
in two distinct fibres. It does not apply to an arbitrary
vector in the norm module of A.100.
\end{proof}

\subsection{Hereditary relative tests and the full normal-cover inverse limit}

\begin{claim}{A.105}{Proved; a compactness consequence with exact directions}
Let $C/k$ be a smooth proper geometrically connected curve of genus
at least two over a finite field of characteristic two, and let
$P\in C(k)$. Use the full geometrically prime-to-two group and the
pointed norm module $M_P$ of A.100. Consider the following three
conditions.
\begin{enumerate}
\item There are a finite extension $L/k$ and compatible pointed
finite etale covers $Z\to Y\to C_L$ such that both $Z$ and $Y$
are geometrically Galois of odd degree over $C$, the relative
cover $Z\to Y$ is nontrivial and Galois, and
$2\nmid h_L(Z)/h_L(Y)$.
\item The module $M_P^F$ has no nonzero element whose projection
has a free geometric deck orbit at every normal level.
\item There are a finite extension $L/k$, a pointed geometrically
Galois odd-degree cover $Y\to C_L$, and a pointed cyclic cover
$Z\to Y$ of odd prime degree with constant deck group over $L$,
such that $Z\to C_L$ has geometric Galois closure of odd degree
and $2\nmid h_L(Z)/h_L(Y)$.
\end{enumerate}
Then condition~1 implies condition~2, and condition~2 implies
condition~3. The source $Z$ in condition~3 is not asserted to be
normal over $C$.

More precisely, if condition~3 fails, then $M_P^F$ contains a
nonzero element killed by two with free deck orbit at every
normal level. Such an element is not asserted to arise from
points of $C$ or of its covers. This is a statement about the
entire inverse limit, not just finitely many quotients.
\end{claim}
\begin{proof}
For the first implication, choose the two normal levels of the
pointed geometric tower corresponding to $Z$ and $Y$. A vector
in $M_P^F$ has a coordinate in $J_Z(L)[2^\infty]$, by the
pointed descent comparison of A.102. Norm and pullback for
$Z\to Y$ give the averaging projector for its geometric deck
group, even if that group is not constant over $L$. Its
complementary rational two-primary group is zero by the
relative class-number hypothesis; this is A.21 and A.67.
Consequently every element of $\operatorname{Gal}(Z_{\bar k}/Y_{\bar k})$
fixes the coordinate. This is a nontrivial subgroup of the
geometric deck group of $Z/C$, so the coordinate does not have
a free orbit. No assertion about a nonnormal level is used.

We prove the second implication by contraposition. Let
$\Pi=\pi_1(C_{\bar k},P)^{(2')}$, and let $\alpha$ be the
continuous automorphism defined by the arithmetic Frobenius
through the point section. Open normal $\alpha$-stable subgroups
$N$ are cofinal among all open normal subgroups of $\Pi$.
Indeed, a pointed finite geometric cover and its chosen point
descend over a finite constant extension, so its defining
subgroup has a finite $\alpha$-orbit; intersect that orbit.
Finite intersections also give a common refinement of any
finite list of such subgroups.

Write $D_N/k$ for the pointed model supplied by the point
section, $G_N=\Pi/N$ for its geometric deck group, and
\[
                  V_N=J_{D_N}[2](\bar k).
\]
This is a finite-dimensional $\mathbf F_2$-vector space.
Arithmetic Frobenius $F_N$ acts on it semilinearly for its
finite odd-order group $G_N$; no constant-deck assumption is
made. Each $G_N$ is solvable, by the verified odd-order theorem
used in A.6.

Failure of condition~3 implies the hypothesis of A.43 at this
level. To check the direction, any cyclic-prime test isolated
there has a normal intermediate base $D_N/B$ over $C$, a
source $D_N/A$ with odd geometric Galois closure, and the
compatible rational point over a finite constant extension.
If its relative class number were odd with a nonconstant
cyclic deck group, A.23 would make that relative group constant
over its exact Frobenius orbit extension while preserving the
oddness. This would give condition~3. We do not make an
additional, uncontrolled extension to constantize the deck
group of $D_N/B$ over $C$.

For each $N$ let $S_N$ be the finite set of $v\in V_N^{F_N}$
such that the norm of $v$ to $D_N/H$ has free orbit under
$G_N/H$ for every normal $H\triangleleft G_N$. The cover
$D_N/H$ and its torsion can be considered geometrically;
$H$ need not be Frobenius-stable. Pullback on two-torsion is
injective, because the degree is odd, and
\[
 (D_N\to D_N/H)^*\operatorname{Nm}_{D_N/(D_N/H)}(v)
                         =|H|e_Hv.
\]
Its equivariance and the invertibility of $|H|$ identify the
required stabilizer with $H$. Thus A.43 proves $S_N\ne\varnothing$.

If $N'\subset N$, norm is a map of finite sets
\[
        S_{N'}\longrightarrow S_N,
        \qquad v\longmapsto\operatorname{Nm}_{D_{N'}/D_N}(v).
\]
It commutes with arithmetic Frobenius and preserves
two-torsion. For every normal quotient of $D_N/C$, its
composite with the next norm is the norm to that same normal
quotient of $D_{N'}/C$. Hence it preserves every freeness
condition defining $S_N$.

The inverse limit of these nonempty finite sets is nonempty.
Here is the precise compactness argument, which requires no
surjectivity of the restricted norm maps. In the compact
product $\prod_N S_N$, each compatibility equation defines a
closed subset. Any finite list of equations is satisfied by
choosing one element at a common refinement and taking its
norms to all coordinates occurring in the list. The remaining
coordinates can be chosen arbitrarily. The finite-intersection
property therefore gives a compatible family $(v_N)_N$.

This family determines an element of $M_P[2]$, by cofinality,
and it is fixed by $F$ because every displayed coordinate is
fixed. Given an arbitrary normal geometric level, choose an
$\alpha$-stable refinement. Its norm to that quotient has
free deck orbit by the definition of $S_N$, so freeness holds
at every normal level, including ones that Frobenius permutes.
There is a nontrivial odd-degree geometric cover of $C$:
for any odd prime $\ell$, its geometric abelian fundamental
group has a quotient $\mathbf Z/\ell$. At a stable refinement
dominating that quotient, a free orbit cannot be the orbit
of zero. Thus the inverse-limit element is nonzero.

The first implication only treats normal sources. In the
contrapositive construction, A.43 isolates a cover with
normal base but possibly nonnormal source. Freeness at every
normal level does not imply freeness after norm to that
nonnormal source. We therefore claim precisely the two
implications stated, and neither an equivalence of all three
conditions nor a point-origin assertion.
\end{proof}

\subsection{A high-rank abelian tower and the missing exponent estimate}

\begin{claim}{A.106}{Proved; the elliptic translation obstruction}
In A.53, its conclusion holds whenever the prime-to-$p$
group $G$ cannot be generated by two elements as an abelian
group. More explicitly, it holds if $G$ is nonabelian, or
if $G$ is abelian and some Sylow subgroup needs at least
three generators. Under the same group hypothesis, the
rational annihilator criterion A.104 holds with no change
in its field, pointing, or numerical conditions.
For a scalar annihilator $p^e\cdot1$, it is enough that
the cover is geometrically Galois, with all other data
and the chosen lift defined over $L$; rational deck
transformations are unnecessary.
\end{claim}
\begin{proof}
We retain the notation of A.53. If $a[Q'-P']=0$, the
nonzero divisor supported on the two disjoint free orbits
gives a degree-$d$ function $f$ on $D$. The field generated
by its $G$-translates defines a smooth proper curve $Z$ and
a finite map $D\to Z$ of degree $e$, with
\[
 e\mid d,\qquad g(Z)\le(d/e-1)^2.
\]
These are exactly the Castelnuovo--Severi induction and
comparison of function fields proved in A.53. The kernel
of the $G$-action on $Z$ has order dividing $e$ and $m$,
so is trivial. A nontrivial point stabilizer would have
a prime-order subgroup partitioning a fibre divisor of
degree $e$ into equally weighted free orbits; this would
make that prime divide both $e$ and $m$. Thus the action
on $Z$ is free. The inequalities and Riemann--Hurwitz now
give $g(Z)=1$, just as in A.53.

Every element of a free finite action on an elliptic
curve is a translation: writing it as $z\mapsto t+u(z)$,
the case $u\ne1$ would give a fixed point because $1-u$
is a surjective nonzero elliptic endomorphism. Hence $G$
embeds in the group of geometric points of an elliptic
curve. Its exponent $n$ is prime to $p$, and its image
lies in the finite group
\[
                  E[n](\overline k)\simeq
                  (\mathbf Z/n\mathbf Z)^2.
\]
A subgroup of this group is abelian and its Sylow
subgroups need at most two generators. This contradicts
the hypothesis on $G$ and proves the geometric claim.
For A.104, a collision supplies the rational point
difference $T$, and the asserted annihilator of the
actual finite group $J_D(L)[p^\infty]$ kills $T$.
The geometric result just proved excludes that equality.
In particular the scalar annihilator $p^e$ works when
$m>(p^e-1)^2$. For this scalar case the deck group need
not be constant over $L$: the scalar kills the rational
lifted difference, and all subsequent translates and
function-field comparisons take place over the algebraic
closure. R1 supplies that lifted difference also for a
pointed cover which is only geometrically Galois. No action
of the nonrational deck transformations on $J_D(L)$ is
being asserted.
\end{proof}

\begin{claim}{A.107}{Proved; an exact conditional abelian-tower route}
Let $C/k$ be smooth, proper and geometrically connected,
where $k=\mathbf F_q$ has characteristic $p$, let
$g(C)=g\ge2$, and let $P\in C(k)$. Fix an odd prime
$\ell\ne p$. Write $J=J_C$ and let
$j_P:C\to J$ be the Abel morphism $x\mapsto[x-P]$.
For $n\ge1$ form the fibre product over $k$
\[
 D_n=C\mathbin{\times}_{J,[\ell^n]}J,
 \qquad \pi_n:D_n\longrightarrow C.
                                                        \tag{A.107.1}
\]
This is a geometrically connected finite etale cover
of degree $\ell^{2gn}$, pointed by $(P,0)$. Let $m_n$
be the order of arithmetic Frobenius on $J[\ell^n]$,
put $L_n=\mathbf F_{q^{m_n}}$, and let $E_n$ be the
exponent of the finite group $J_{D_n}(k)[p^\infty]$,
with exponent one for the zero group. If for one $n$
\[
                       E_n-1<\ell^{gn},                 \tag{A.107.2}
\]
then $P$ collides with no other point of $C(k)$ in
the full geometrically prime-to-$p$ arithmetic group.
In particular an estimate $E_n=o(\ell^{gn})$ would
suffice for this $(C,P,\ell)$.

There is an integer $t\ge1$ such that
\[
           m_n=m_1\ell^{\max(0,n-t)}\quad(n\ge1).
                                                        \tag{A.107.3}
\]
Consequently the exponents of $J(L_n)[p^\infty]$ are
uniformly bounded. This last bound concerns the base
Jacobian only; no bound for the relative parts of
$J_{D_n}$, and hence no universal instance of
(A.107.2), is asserted here. If $E_n^{\rm const}$
is the exponent over the constant-deck field $L_n$,
then $E_n\mid E_n^{\rm const}$, so the analogous
bound over $L_n$ would also suffice, but that field
enlargement is unnecessary for the scalar criterion.
\end{claim}
\begin{proof}
Multiplication by $\ell^n$ on $J$ is finite etale,
so its pullback has the displayed degree and is etale.
The Abel morphism induces the usual isomorphism on
geometric abelian pro-$\ell$ fundamental groups.
Its monodromy onto $J[\ell^n]$ is therefore surjective,
which proves geometric connectedness. The pair $(P,0)$
is $k$-rational. Over $L_n$ the deck group is the
constant group
\[
                  G_n\simeq(\mathbf Z/\ell^n)^{2g},
\]
and the complete fibre over $P$ is rational. This
group needs $2g>2$ generators. Without extending $k$,
R1 supplies a $k$-rational lifted difference killed by
$E_n$: its hypotheses only require a geometrically
prime-to-$p$ Galois closure and the chosen rational
lift. Apply the nonconstant-deck scalar form of A.106 with
$a=E_n\cdot1$, $d=E_n$ and $m=\ell^{2gn}$.
Their greatest common divisor is one, and (A.107.2)
is exactly the required strict inequality.

Let $A\in\operatorname{GL}_{2g}(\mathbf Z_\ell)$
be arithmetic Frobenius on $T_\ell J$. The matrix
$A^{m_1}-1$ is nonzero: otherwise every eigenvalue
of Frobenius would be a root of unity, contrary
to its complex absolute value $q^{1/2}>1$.
Let $t$ be the minimum $\ell$-adic valuation of its
entries. Then $t\ge1$. For an odd prime $\ell$, the
binomial formula gives
\[
 v_\ell\bigl((1+\ell^t B)^{\ell^r}-1\bigr)=t+r
 \quad(r\ge0)
\]
when at least one entry of $B$ is a unit; matrix
valuation here means the minimum entry valuation.
Indeed the first term at each powering has that
valuation and every later term has strictly larger
valuation. Raising to an integer prime to $\ell$
does not change it. An exponent making $A$ the
identity modulo $\ell^n$ must first be a multiple
of $m_1$, and the same calculation then proves
(A.107.3). In particular $v_p(m_n)=v_p(m_1)$.

For completeness we prove the asserted base bound.
On the geometric etale $p$-divisible group of $J$,
Frobenius is represented by a matrix
$B\in\operatorname{GL}_{f_p(J)}(\mathbf Z_p)$.
Its eigenvalues $\alpha_i$ are algebraic $p$-adic
units, and none is a root of unity. The rational
$p$-primary group over $\mathbf F_{q^b}$ is the
kernel of $B^b-1$ on $(\mathbf Q_p/\mathbf Z_p)^{f_p(J)}$.
Since $B^b-1$ is invertible over $\mathbf Q_p$,
its order has $p$-valuation
\[
                  \sum_i v_p(\alpha_i^b-1).             \tag{A.107.4}
\]
Choose a finite $p$-adic splitting field. If the
residue of $\alpha_i$ has order $d_i$, then the
summand is zero unless $d_i\mid b$. Otherwise
write $b=d_i p^c u$, with $p\nmid u$. Factoring
$X^u-1$ at $X=\alpha_i^{d_i p^c}\equiv1$ shows
\[
 v_p(\alpha_i^b-1)=
 v_p(\alpha_i^{d_i p^c}-1),
\]
because the second factor reduces to the unit $u$.
For $b=m_n$ the integer $c=v_p(m_1)$ is fixed.
Thus every summand in (A.107.4) belongs to a fixed
finite set of finite numbers. The orders, and
therefore the exponents, of the base groups are
uniformly bounded. The case $f_p(J)=0$ is immediate.

The comparison does not identify these base groups
with the entire rational torsion of $J_{D_n}$. The
new character components are the precise missing
input in the proposed route; omitting them would
invalidate the conclusion about a collision.
\end{proof}

% Source: agent_F/F_section.tex
% Part F contribution.  Local labels FF.* are to be renumbered by the editor.
% AI-produced by GPT-6 Astra (OpenAI, ChatGPT Pro mode). No priority claimed.
\section{Part F. Reconstruction from \'etale topoi}
\label{sec:FF}

Throughout this part, a relative geometric morphism includes its specified
commutation isomorphism with the structural morphisms to the base topos.
The adjective \emph{pinned} has a separate meaning: the induced map of
topological spaces sends closed points to closed points
\citep[Definition~2.19]{FF_CHW2024}. We retain both pieces of data.
Write \(\mathcal R(k)\) for the reconstruction property in the question.
All assertions below concern characteristic zero except where a
finite field or an arithmetic base is explicitly specified.

\subsection{Points, sections, and the qualification in scan 52}

\begin{claim}{F.1}{Proved}\label{clm:F.1}
Let \(k\) have characteristic zero and let \(X/k\) be a finite-type
scheme. The isomorphism classes of relative geometric morphisms
\(BG_k\to X_{\mathrm{et}}\) are the disjoint union, over the points
\(x\in X\), of the continuous sections of
\[
                 G_{k(x)}\longrightarrow G_k,
\]
modulo conjugation by its kernel, with compatible separable closures
understood. A point for which the displayed map is not surjective
contributes no section. The closed-point summands are exactly
\(X(k)\). In particular, for an integral curve the other summand is
the set of birational sections at its generic point.

Consequently, for every characteristic-zero field the assertion in
cote~149, scan~52, read with \emph{all} relative topos morphisms, is
false, even with universal homeomorphisms inverted.
\end{claim}
\begin{proof}
Pull back a geometric morphism along the canonical point
\(\mathrm{Set}\to BG_k\). This gives a geometric point of
\(X_{\mathrm{et}}\), hence a separably closed point above a unique
\(x\in X\). Descent from this point to \(BG_k\) is a continuous
action of \(G_k\) on its fibre functor. The automorphism group of
that functor is \(G_{k(x)}\). The relative commutation isomorphism
makes its composite to \(G_k\) the identity. Changing the compatible
geometric point acts by the geometric kernel, which gives exactly the
stated ambiguity. Conversely, such an action supplies the descent
datum and hence the geometric morphism. This is the point description
of the small \'etale topos, applied with its continuous automorphism
group; it involves the residue field at \(x\), rather than only
\(\pi_1(X)\).

If \(x\) is closed, \(k(x)/k\) is finite separable. Its absolute
Galois group maps to an open subgroup of \(G_k\) of index
\([k(x):k]\). A right inverse forces this index to be one. Conversely
a rational point gives a relative morphism, uniquely up to the
specified relative isomorphism. This also follows from
\citet[Proposition~3.6]{FF_CHW2024}.

For the counterexample put \(X=\mathbb A^1_k\). The henselian
discrete valuation field \(H=\operatorname{Frac}(k[t]^h_{(t)})\)
has residue field \(k\). In characteristic zero its decomposition
sequence splits: in a separable closure obtained by adjoining
separable constants and compatible roots of \(t\), act on the
constants and fix those roots. Thus
\(G_k\to G_H\to G_{k(t)}\) is a section over \(G_k\).
It defines a relative topos morphism whose image in \(\mathbb A^1_k\)
is the generic point. Every scheme morphism
\(\operatorname{Spec}k\to\mathbb A^1_k\) has rational, hence closed,
image. Inverting universal homeomorphisms does not change this
distinction; in particular the source field is already absolutely
weakly normal. This is the tangential construction of AP-F.1 and
AP-F.17, also illustrated in \citet[Example~2.30]{FF_CHW2024}.
\end{proof}


\begin{claim}{F.2}{Known, conditional as stated}\label{clm:F.2}
Let \(k/\mathbb Q\) be finitely generated and \(K/k\) a regular
one-variable function field. Under conjecture~a) in cote~149,
scans~126--127, each section \(G_k\to G_K\) factors through a
decomposition group \(D_v\) of a nontrivial valuation trivial on
\(k\), with residue field finite over \(k\). For an actual section
over \(G_k\), that residue field is \(k\) itself. Its centre lies
on the smooth proper model of \(K/k\), and need not lie on a
specified open curve with that function field.
\end{claim}
\begin{proof}
The manuscript expressly requires the valuation ring to contain
\(k\) and its residue field to be algebraic over \(k\). On the
smooth proper curve, the valuative criterion gives a centre.
The valuation is nontrivial, so its centre is a closed point;
the dominating valuation ring is its discrete valuation ring.
Its residue field is therefore a finite extension \(k(v)/k\).
The decomposition sequence is
\[
       1\longrightarrow I_v\longrightarrow D_v
         \longrightarrow G_{k(v)}\longrightarrow1.
\]
The image of \(D_v\to G_k\) is the open subgroup associated to
\(k(v)/k\). A section of \(G_K\to G_k\) contained in \(D_v\)
forces that image to be all of \(G_k\), so \(k(v)=k\).
The factorization itself is the conjectural assertion, not a theorem
proved here. On scan~126 the uncertain reading is
\(I_V\) versus \(N_V\); the subsequent discussion of its normalizer
and scan~127's explicit factorization support the inertia reading.
Our conditional statement uses the explicit factorization on
scan~127 and does not silently assume that every normalizer equals
the decomposition group.
\end{proof}

\subsection{Finite fields and an arithmetic base}

\begin{claim}{F.3}{Refuted}\label{clm:F.3}
Closed-point-preserving relative \'etale-topos morphisms over
\(\mathbb F_q\) need not be induced by morphisms after universal
homeomorphisms have been inverted. This already fails for
\(X=Y=\mathbb G_{m,\mathbb F_q}\). The corresponding assertion
for all finite-type schemes over \(\operatorname{Spec}\mathbb Z\)
is also false.
\end{claim}
\begin{proof}
Identify \(G_{\mathbb F_q}\) with \(\widehat{\mathbb Z}\), using
the arithmetic Frobenius. Each \(z\) in this central group gives a
2-automorphism of the structural functor
\(X_{\mathrm{et}}\to BG_{\mathbb F_q}\). Twist the relative
structure of the identity of \(X_{\mathrm{et}}\) by that
2-automorphism, choosing the sign so that its geometric-point
function is
\[
          a\longmapsto\operatorname{Frob}_q^z(a)
          \qquad(a\in\overline{\mathbb F}_q^{\times}).
\]
The underlying map of spaces is the identity, so this relative
automorphism is pinned. The point function follows either directly
by transporting the relative base identification, or from the
finite-extension point description in Proposition~3.6 of CHW.

Choose \(z\in\widehat{\mathbb Z}=\prod_\ell\mathbb Z_\ell\)
whose 2-adic component is zero and whose other components are one.
This is not an integer: its 2-adic component would force that integer
to be zero, contradicting its 3-adic component. Suppose our relative
morphism came from
\(f:X^{\mathrm{perf}}\to X\). Every unit in
\(\mathbb F_q[T^{\pm1/p^\infty}]\) is
\(cT^{m/p^r}\), with \(c\in\mathbb F_q^\times\).
The induced geometric fundamental-group homomorphism must be an
isomorphism, because our topos morphism is one. On
the maximal prime-to-\(p\) quotient
\(\pi_1(X_{\overline{\mathbb F}_q})^{(p')}=\prod_{\ell\ne p}\mathbb Z_\ell\)
it is multiplication by \(m/p^r\). Hence
\(m/p^r=\pm p^s\) for an integer \(s\). Evaluation at 1 gives
\(c=1\).

The negative sign is impossible: composing with an integral power
of Frobenius would make inversion on the multiplicative group,
extended by zero, a field automorphism. Additivity at \(u+1\)
would then give \(u^2+u+1=0\) for every
\(u\ne0,-1\), which is false in the infinite algebraic closure.
For the positive sign, write \(q=p^a\). Equality of the field
automorphisms gives \(az=s\) in \(\widehat{\mathbb Z}\), hence
\(z=s/a\in\mathbb Q\cap\widehat{\mathbb Z}=\mathbb Z\), another
contradiction. This proves nonalgebraicity even after perfection.

For the last assertion take \(q=p\) and regard \(X\) as a vertical
\(\mathbb Z\)-scheme. The closed immersion
\(\operatorname{Spec}\mathbb F_p\hookrightarrow
\operatorname{Spec}\mathbb Z\) induces a geometric embedding of
small \'etale topoi. Postcomposition with a geometric embedding is
fully faithful on relative factorization data. Thus a relative
identification over \((\operatorname{Spec}\mathbb Z)_{\mathrm{et}}\)
would give the same identification over \(BG_{\mathbb F_p}\), which
we have excluded. A morphism of these vertical schemes over
\(\mathbb Z\) is automatically a morphism over \(\mathbb F_p\).
\end{proof}

The distinction in the last proof addresses the positive-characteristic
caution attached to scan~52. Absolute Frobenius and the identity have
isomorphic \emph{absolute} \'etale-topos morphisms. Their relative
commutation isomorphisms retain different data. Inverting universal
homeomorphisms makes Frobenius invertible; it does not identify it with
the identity in the category of schemes.

\begin{claim}{F.4}{Known}\label{clm:F.4}
The finite-field counterexample does not invalidate the following
restricted positive statements. CHW Theorem~0.1 gives reconstruction
over infinite finitely generated fields, with absolute weak
normalization. Carlson--Stix Theorem~D, numbered~6.12 in the version
dated 20 September 2026, gives relative reconstruction when the base
is noetherian, Nagata and geometrically unibranch, each irreducible
component has a sub-\(p\)-adic function field, the schemes are
componentwise universally open and of finite type over the base,
and the topos morphisms send constructible subsets to constructible
subsets. Universal homeomorphisms are inverted.
\end{claim}
\begin{proof}
These are the hypotheses and conclusions of the opened statements
\citep[Theorem~0.1]{FF_CHW2024}
and \citep[Theorem~D, Theorem~6.12]{FF_CS2026}.
Carlson--Stix call their morphisms \emph{chivalrous}; their Remark~2.1
shows that these are pinned on fibres. This is a stronger and
different relative condition from merely preserving the closed
points of the total space. Their Corollary~E includes the stated
normal-source Hom conclusion. Over \(\operatorname{Spec}\mathbb Z\),
the generic field is sub-\(p\)-adic, but vertical schemes such as
those in Claim~\ref{clm:F.3} fail universal openness. Thus the
restricted theorem does not assert the statement refuted there.
\end{proof}

\subsection{The birational automorphism experiment}

\begin{claim}{F.5}{Known}\label{clm:F.5}
Let \(K_0\) be a number field and \(X/K_0\) a smooth proper
geometrically connected curve, with \(K=K_0(X)\). Every continuous
outer automorphism of \(G_K\) over \(G_{K_0}\), modulo inner
conjugation by the geometric kernel, is induced by a unique
\(K_0\)-automorphism of \(X\). Preservation of the closed-point
inertia and decomposition groups is a consequence, and need not be
assumed. In particular the experiment in cote~149, scan~52, has a
positive answer for the genus at least two curves asked about.
\end{claim}
\begin{proof}
A number field is sub-\(p\)-adic. Mochizuki's Theorem~B, equivalently
Theorem~17.1, states the birational Hom theorem for regular function
fields over a sub-\(p\)-adic base: base-linear field embeddings
correspond bijectively, contravariantly, to continuous open
homomorphisms of their \emph{full} absolute Galois groups over the
base group, modulo inner conjugation by the geometric kernel
\citep[Theorem~17.1, author pp.~87--88]{FF_Mochizuki1999}.
The title of that paper does not replace these full groups by
pro-\(p\) quotients in this theorem.

Apply the theorem to the given automorphism and to its inverse.
Its uniqueness and functoriality imply that the resulting field
embeddings are inverse, giving a \(K_0\)-automorphism of \(K\).
It extends uniquely to the smooth proper curve: the local rings at
its closed points are precisely the discrete valuation rings of
\(K/K_0\) with finite residue field, and a birational map between
regular proper curves extends across those points. The resulting
scheme automorphism transports their valuation, inertia and
decomposition groups. This proves the asserted redundancy of the
loop hypothesis. The birational argument itself does not require
the genus bound.
\end{proof}

\subsection{Two proposed field criteria}

\begin{claim}{F.6}{Proved}\label{clm:F.6}
If \(\mathcal R(k)\) holds in characteristic zero, then for every
finite extension \(E/k\),
\[
             C_{\operatorname{Aut}(\bar k)}(G_E)=1.
\]
In particular every open subgroup of \(G_k\) has trivial centre.
This proves the necessary direction of candidate F3(a).
\end{claim}
\begin{proof}
Finite-extension ascent of reconstruction follows by taking the
slice over the finite \'etale object \(\operatorname{Spec}E\):
an \(E\)-relative map is a \(k\)-relative map with that factorization,
and uniqueness over \(k\) preserves the factorization. Thus
\(\mathcal R(E)\) holds.

Let \(\sigma\in\operatorname{Aut}(\bar k)\) centralize \(G_E\).
It preserves the fixed field \(E\). Its action on coefficients,
with the variable fixed, yields a semilinear automorphism of
\(\mathbb G_{m,E}\), whose topos morphism is relative over
\(BG_E\) because the induced conjugation on \(G_E\) is the
identity. With this relative structure its point function is
\(a\mapsto\sigma(a)\). This includes the case where \(\sigma\)
fixes \(E\): then its nontriviality is in the relative structure.
Reconstruction would make the map \(cT^n\). Its being a topos
equivalence forces \(n=\pm1\), and its value at 1 forces \(c=1\).
Inversion cannot be a field automorphism by the additivity calculation
in Claim~\ref{clm:F.3}. Hence \(\sigma(a)=a\) for all nonzero
\(a\), proving the claim. This is the explicit ghost test of
AP-F.39--F.40.
\end{proof}

\begin{claim}{F.7}{Refuted}\label{clm:F.7}
Candidate F3(b) is false. Reconstruction need not imply the Isom
form of the Grothendieck conjecture for hyperbolic curves of genus
zero, even for a pair of four-punctured projective lines over the
base field itself. This remains a counterexample if the group
isomorphisms are required to preserve the peripheral conjugacy
classes.
\end{claim}
\begin{proof}
Put \(k=\overline{\mathbb Q}(t)\). Reconstruction over this field
is AP-F.45; the dependent argument is checked in
Claim~\ref{clm:FFaudit} below. Consider
\[
 U_2=\mathbb P^1_k\setminus\{0,1,\infty,2\},\qquad
 U_3=\mathbb P^1_k\setminus\{0,1,\infty,3\}.
\]
An isomorphism extends to their smooth proper completions, hence
is a fractional linear transformation carrying one four-element
set to the other. The cross-ratio orbits are respectively
\[
 \{2,\tfrac12,-1\},\qquad
 \{3,\tfrac13,-2,-\tfrac12,\tfrac32,\tfrac23\}.
\]
They are disjoint. The two curves are therefore not isomorphic
over \(k\), nor over an algebraically closed characteristic-zero
extension. The exact rational computation is reproduced in
\texttt{F\_cross\_ratio.py}.

On the other hand these are constant curves over
\(\overline{\mathbb Q}\). The characteristic-zero K\"unneth
isomorphism of SGA~1, Expos\'e~XIII, Proposition~4.6, applies to
this smooth curve and to \(\operatorname{Spec}k\). Resolution in
dimension one suffices for its hypothesis. Together with Expos\'e
XIII, Corollary~2.12, it gives
\[
 \pi_1(U_2)\simeq\widehat F_3\times G_k
          \simeq\pi_1(U_3)
\]
as groups over \(G_k\) \citep{FF_SGA1}. The projection is the
arithmetic structural map. A homeomorphism of the complex
four-punctured spheres gives an isomorphism of the free profinite
factors preserving their peripheral classes; the product with
the identity of \(G_k\) supplies the required relative isomorphism.
Thus the proposed equivalence already fails in its necessary
direction. No inference from a full topos to its fundamental
group is being reversed here.
\end{proof}

\begin{claim}{F.8}{Proved}\label{clm:F.8}
For descent along a transcendental field extension, a commuting
square of small \'etale topoi with a finitely generated model is
not by itself a scalar-descent datum for a topos morphism.
Moreover, the usual model systems of a finite-type scheme over
an arbitrary field are not generally pro-objects in the category
whose transition morphisms are required to be pinned.
\end{claim}
\begin{proof}
Take \(k=\mathbb C\), a finitely generated \(F\subset\mathbb C\),
and a nonidentity \(\sigma\in\operatorname{Aut}(\mathbb C/F)\).
Act on the coefficients of \(\mathbb G_{m,\mathbb C}\) and fix
the variable. The induced topos automorphism is pinned and, since
\(BG_{\mathbb C}=\mathrm{Set}\), relative over \(\mathbb C\).
It commutes with the projection to the model over \(F\). Yet its
point function \(a\mapsto\sigma(a)\) cannot be given by a
\(\mathbb C\)-morphism of \(\mathbb G_m\), by the preceding
unit and additivity argument. Thus this commuting square does not
mean that the automorphism is the scalar extension of the identity
on the model.

For the transition issue use \(F\subset F(u)\) with \(u\)
transcendental. The closed point \(T=u\) of
\(\mathbb A^1_{F(u)}\) maps to the generic point of
\(\mathbb A^1_F\). Hence the transition is not pinned.
Likewise, the small \'etale topos of the scalar extension is
not in general the 2-fibre product of the two small topoi over
\(BG_F\): over a separably closed extension that fibre product
only makes the algebraic constant extension of the original
geometric topos, whereas transcendental scalar extension supplies
additional closed points. Candidate F3(c) must specify a genuine
descent category; the displayed commuting-square interpretation
does not suffice.
\end{proof}

\subsection{The reconstruction chain and a normal-extension enlargement}

\begin{claim}{F.9}{Proved}\label{clm:FFcore}
Let \(D/\mathbb Q\) be finitely generated, and let \(k/D\) be
an algebraic Galois extension containing all roots of unity.
If \(k\) is Hilbertian, then every finitely generated extension
\(L/k\) has \(\mathcal R(L)\).
\end{claim}
\begin{proof}
We give the common argument underlying AP-F.43 and AP-F.49,
including the enlargement of their hypotheses. For a field \(E\)
put \(D(E^\times)=\bigcap_{n\ge1}(E^\times)^n\).

\emph{The scalar kernel.}
For every finitely generated \(M/\mathbb Q\),
\(D(M^\times)=1\). Divisors on a normal proper model put a
divisible element and each of its roots in the number field of
constants; finite valuations and then the finitely generated unit
group of that number field prove the assertion.
If \(E/k\) is finite and \(a\in D(E^\times)\), choose a finite
Galois \(M/D\) with \(a\in M\) and \(E\subset Mk\).
For \(\sigma\in\operatorname{Gal}(Mk/k)\),
\(\tau\in\operatorname{Gal}(Mk/M)\), and \(z_n\in E\) with
\(z_n^n=a\), the automorphisms \(\sigma,\tau\) commute: check
this on \(M\) and \(k\), both of which they preserve.
Moreover \(\tau z_n/z_n\in\mu_n\subset k\) is fixed by
\(\sigma\). Therefore \(\sigma z_n/z_n\in M\), and
\(\sigma a/a\in D(M^\times)=1\). Thus \(a\in k\).
For a finitely generated \(E/k\), the proper-model divisor
argument first puts \(a\) and all its roots in the finite relative
algebraic closure of \(k\) in \(E\). We have proved
\[
             D(E^\times)\subset k^\times
             \quad(E/k\text{ finitely generated}).       \tag{F.9.1}
\]
Equality and radical closedness of \(k\) are not required.

\emph{Detection of arithmetic finite quotients.}
Every finitely generated extension of a Hilbertian field is
Hilbertian. For the transcendental case we use
\citet[Theorem~1.2, Corollary~3.2]{FF_BarySorokerParan2013};
finite extensions preserve it by
\citet[Corollary~3.2.2]{FF_Serre2007}. Let \(Y\to U\) be a
connected finite \'etale \(G\)-torsor over a smooth geometrically
connected curve over \(L\). A rational function of arbitrarily
large available degree \(d\) on its proper completion makes
\(Y\to\mathbb P^1_L\) finite \'etale after deleting finitely
many values. Hilbert specialization of the entire connected cover
gives a residue field of degree \(d|G|\). Its quotient gives
\(x\in U\) with \([L(x):L]=d\) and
\[
                         G_{L(x)}\twoheadrightarrow G.   \tag{F.9.2}
\]
This retains the arithmetic constant quotient; geometric
connectedness of \(Y\) is not assumed. The exact specialization
statement is \citet[Proposition~3.3.1]{FF_Serre2007}.
For higher-dimensional smooth quasiprojective \(U\), successive
hyperplane sections preserve every geometric component of the
cover and its arithmetic permutation. Poonen--Slavov Theorem~1.5
applies component by component because the covering map has
zero-dimensional nonempty fibres \citep{FF_PoonenSlavov2020}.
Combining with the smooth Bertini open gives a smooth curve,
retaining any prescribed nonempty open restriction. Thus (F.9.2)
also gives the required finite-quotient detection in that case.

\emph{One punctured-line calculation.}
Consider
\(U=\operatorname{Spec}L[T,1/(P_1\cdots P_r)]\), with distinct
monic irreducible \(P_j\), and choose \(x_0\in U(L)\).
A compatible Kummer class with actual values at all finite-extension
points, normalized by its actual value \(b_0\) at \(x_0\), is
\[
 \kappa(b_0)+\sum_j\lambda_j
             \bigl([P_j(T)]-[P_j(x_0)]\bigr),
                  \qquad\lambda_j\in\widehat{\mathbb Z}. \tag{F.9.3}
\]
This uses the vanishing of the Picard group of this localization
of a principal ideal domain.

Choose \(F/\mathbb Q\) finitely generated, containing \(D\)
and the finite displayed data, with \(L=Fk\); then \(L/F\)
is algebraic Galois. Hilbertianity gives a nontrivial finite
extension of \(L\). Descend it, and take a finite Galois extension
\(M/\mathbb Q(t_1,\ldots,t_s)\) containing \(F\), these
extension data and the roots of all \(P_j\). Choose
\(1\ne\sigma\in\operatorname{Gal}(Mk/L)\). It preserves \(M\)
and acts nontrivially on it. We require discrete valuations
\(v_i\) of \(M\), with mutually disjoint pairs
\(\{v_i,\sigma^{-1}v_i\}\), such that the relevant nonzero coefficients and
distinct root differences are units.

Here are the two exhaustive cases. If \(\sigma\) moves the
number field \(C_M=M\cap\overline{\mathbb Q}\), choose good
rational primes splitting completely in \(C_M\). Above each,
a Gauss valuation on the rational function field has an unramified
prolongation to a suitable normal model of \(M\), outside a finite
exceptional set. The restrictions of \(v_i\) and
\(\sigma^{-1}v_i\) to \(C_M\) differ. When \(s=0\), take
the residue prime sufficiently large to avoid finitely many roots.
When \(s>0\), the residue field is infinite. If \(\sigma\)
fixes \(C_M\), it is a nonidentity automorphism of the
positive-dimensional function field over \(C_M\). Choose a proper
normal \(\sigma\)-equivariant model, for example by normalizing
a proper model of the fixed field. Infinitely many prime divisors
are moved. Indeed, if the moved divisors formed a finite set \(S\),
choose \(f\) with \(\sigma f\ne f\), and a fractional linear
transform \(h\) of \(f\) for which both \(h\) and \(1+h\)
are units at every member of \(S\). Such a transform exists
because \(C_M\) is infinite and only finitely many residue values
are excluded. Their divisors are fixed, so properness gives
\(\sigma h=ah\) and \(\sigma(1+h)=b(1+h)\), with
\(a,b\in C_M^\times\). Comparing these identities forces
\(a=b=1\), contradicting \(\sigma f\ne f\). Avoiding the
finite bad divisors and the previously selected finite orbits
now gives the required disjoint moved pairs. These are the arithmetic and geometric divisor
arguments of AP-F.43 and AP-F.45 respectively.

Weak approximation at this finite set of distinct discrete
valuations chooses \(\alpha\in M\) equal to a root of \(P_i\)
plus a uniformizer modulo the square of the \(v_i\)-ideal, and
avoiding every root modulo the \(\sigma^{-1}v_i\)-ideal.
The unit derivatives and root differences give
\[
 v_i(P_j(\alpha))=\delta_{ij},\qquad
 v_i(P_j(\sigma\alpha))=0.                               \tag{F.9.4}
\]
Set \(R_j=P_j(\sigma\alpha)/P_j(\alpha)\). The actual value
\(b_\alpha\in Mk\) lies in one field \(F'=MS_0\), where
\(S_0/D\) is finite inside \(k\). This field is finite over
\(M\), \(\sigma\)-stable, and fixed independently of \(n\).
If \(z^n=a\in F'\) with \(z\in Mk\), the commutation
calculation above gives \(\sigma z/z\in F'\). Apply it to
the evaluated relation (F.9.3). For integer representatives
\(a_{j,n}\) of \(\lambda_j\bmod n\) it gives
\[
 \frac{\sigma b_\alpha}{b_\alpha}
       \prod_jR_j^{-a_{j,n}}\in(F'^\times)^n.
\]
Apply the integer-valued map \(v_i\circ N_{F'/M}\), for every
\(n\). We obtain
\[
 [F':M]\lambda_i
       =-(v_i\circ N_{F'/M})(\sigma b_\alpha/b_\alpha)
                     \in\mathbb Z\subset\widehat{\mathbb Z}.
                                                               \tag{F.9.5}
\]
An element of \(\widehat{\mathbb Z}\) with a nonzero integer
multiple in \(\mathbb Z\) is itself in \(\mathbb Z\): the
corresponding rational number is integral at every prime. Thus
all puncture exponents are integral. No valuation on the infinite
field \(Mk\) was asserted to have discrete value group.

\emph{The Albanese step.}
Let \(X/L\) be regular, affine and connected, with
\(x_0\in X(L)\), and let \(\xi\) be the Kummer class pulled
back by a pinned relative map to \(\mathbb G_m\). It has actual
values at every finite-extension point by the point-source
description. The generalized Albanese \(A\), based at \(x_0\),
identifies normalized units with \(\operatorname{Hom}(A,\mathbb G_m)\).
The comparison of \citet[Section~3, Proposition~4.1,
Corollary~4.2, Proposition~4.3]{FF_SpiessSzamuely2003} identifies
geometric \(H^1\) with the dual of \(T_\ell A\): the finite
kernel in the abelian fundamental-group comparison contributes
no map to the torsion-free group \(\mathbb Z_\ell(1)\).
The required smooth projective compactification exists in
characteristic zero.

Descend the finite geometric data to \(F\) as above and put
\(H=G_L\triangleleft G_F\), \(W=(V_\ell A)_H\). Because
\(\mu_\infty\subset L\), the target Tate twist is trivial
on \(H\). Projected point-Kummer cocycles are homomorphisms
\(c_x:G_E\to W\), for finite \(E/L\). Their images span \(W\).
Indeed, a nonzero annihilating functional, scaled to be integral,
defines a nonzero finite quotient of \(\pi_1(X)\) which is zero
on the base-point section and every closed-point decomposition
group. This contradicts (F.9.2).

Choose finitely many spanning points. Their coordinates, their
actual values, and the normalization descend to a single
finitely generated field \(M_1/\mathbb Q\) containing \(F\).
Restriction to \(G_{M_1k}\) preserves their rational spans,
because it is restriction to a subgroup of finite index.
The normalized geometric functional
\(\lambda:T_\ell A\to\mathbb Z_\ell(1)\) satisfies
\[
 \lambda(c_x(ghg^{-1}))
       =\chi_\ell(g)\lambda(c_x(h))
 \quad(g\in G_{M_1},\ h\in G_{M_1k}).
\]
Here the point and its value are \(M_1\)-rational,
\(G_{M_1k}\) is normal in \(G_{M_1}\), and the cocycle
coboundary disappears in \(H\)-coinvariants. Spanning proves
that \(\lambda\) is \(G_{M_1}\)-equivariant.
The integral Tate-Hom theorem for 1-motives over finitely generated
fields now puts it in
\(\operatorname{Hom}_{M_1}(A,\mathbb G_m)\otimes\mathbb Z_\ell\).
We use the exact integral statement recorded in
\citet[Section~1.4]{FF_Jossen2014}; its application is over
\(M_1\), not over the infinite field \(L\).

After a finite extension splitting the boundary, every normalized
geometric unit is defined over the same finite model: its conjugates
have equal divisors and value 1 at \(x_0\), hence are equal.
For every \(\ell\) we therefore obtain the completion of one
fixed integral unit lattice. On a curve, norm to a punctured
rational line isolating one boundary point at a time. Norm
multiplies residues by residue-field degree, not ramification
index. Equation (F.9.5) makes each boundary coefficient integral.
In higher dimensions use a smooth curve transverse to the chosen
boundary divisor; the local finite-level expression \(uz^a\)
has the same exponent on that curve. The injective integer divisor
matrix of a unit basis, and an invertible minor, then put all unit
exponents in \(\mathbb Q\cap\widehat{\mathbb Z}=\mathbb Z\).
We obtain a unit \(h\), normalized by \(h(x_0)=b_0\), with
\(\xi=\kappa(h)\). Equality of the remaining constant classes
follows from Hochschild--Serre and evaluation at \(x_0\).

\emph{Actual point functions.}
Equation (F.9.1) leaves \(\varphi(x)=c_xh(x)\), \(c_x\in k^\times\).
For a curve and nonconstant \(h\), choose closed points of
degree greater than \(\deg h\), using (F.9.2). Their values
lie outside \(L\), so the point function is nonconstant.
On the inverse image of \(\mathbb G_m\setminus\{1\}\), apply
the same class argument to \(T-1\), obtaining
\(\varphi(x)-1=d_xh_1(x)\), \(d_x\in k^\times\).
Choose a closed point whose degree exceeds the common pole-degree
bound of \(c_xh-d_xh_1-1\). Its vanishing forces this function
to be identically zero. Hence \(h_1=Ah+B\), with
\(A,B\in k^\times\). At every point with \(h(x)\notin k\),
linear independence of \(1,h(x)\) over \(k\) gives
\(c_x=-A/B=c\). Pulling back \(T-b\), for each
\(b\in k^\times\), gives in the same way a unit proportional
to \(ch-b\). At a point with \(h(x)\in k\), take
\(b=ch(x)\); invertibility excludes \(\varphi(x)\ne b\).
Thus \(\varphi=ch\) everywhere, and normalization gives \(c=1\).

If \(h\) is constant, pass to a nontrivial finite extension
\(L'/L\) and choose \(\beta\in L'\setminus L\).
Every value of \(\varphi\) lies in \(hk\subset L\), so
\(T-\beta\) pulls back to a unit on the whole curve. Its
associated unit has all point values in \(L'\), and the same
unbounded-degree argument makes it constant. Normalizing at
\(x_0\) gives both
\(\varphi(x)/h\in k\) and
\((\varphi(x)-\beta)/(h-\beta)\in k\). Since \(\beta/h\notin k\),
these force \(\varphi(x)=h\). For higher-dimensional \(X\),
after a finite extension join \(x_0\) to the geometric point
in question by a smooth curve. Its normalized unit has the same
Kummer class; the ratio is a divisible unit, hence constant, and
equals 1 at \(x_0\). This proves the same value equality.

We have proved precisely the point-function condition for regular
connected affine sources with a rational point in
\citet[Theorem~4.18(5)]{FF_CHW2024}. That theorem gives
\(\mathcal R(L)\) for all finite-type schemes. All arguments
apply after every finite extension of \(L\), as required by the
intermediate reductions.
\end{proof}

\begin{claim}{F.10}{Proved}\label{clm:FFnormal}
Let \(D/\mathbb Q\) be finitely generated and let
\(D\subset k\subsetneq\bar D\) be an algebraic Galois extension.
Then \(\mathcal R(L)\) holds for every finitely generated
extension \(L/k\). In particular the field
\(\mathbb Q^{\mathrm{tr}}\) of totally real algebraic numbers
reconstructs. The assertion does not claim closure under arbitrary
subfields of \(L\).
\end{claim}
\begin{proof}
First suppose \(\mu_\infty\subset k\). Choose a finite Galois
extension \(M/D\) not contained in \(k\). Then \(kM/k\) is
proper finite and \(kM/D\) remains Galois. The field \(D\) is
Hilbertian. Weissauer's theorem, in the opened stronger form
\citet[Corollary~2.3]{FF_JardenRazon2020}, applied with its
\(R=K=D\), \(M=N=k\) and \(M'=kM\), makes \(kM\)
Hilbertian. Claim~\ref{clm:FFcore} gives reconstruction over
\(LkM\), which is finitely generated over \(kM\). The extension
\(LkM/L\) is algebraic; algebraic descent in
Claim~\ref{clm:FFdescent} below gives \(\mathcal R(L)\).

For general \(k\), put \(D_{\mathrm{cyc}}=D(\mu_\infty)\)
and \(k'=kD_{\mathrm{cyc}}\). The cyclotomic extension is
abelian and proper in \(\bar D\): Hilbert specialization gives
a nonabelian symmetric extension of \(D\), which an abelian
extension cannot contain. Concretely \(X^3-X-T\) has regular
group \(S_3\) over \(D(T)\) by the Morse-polynomial theorem
\citep[Theorem~4.4.5]{FF_Serre2007}, and specialization over
the Hilbertian field \(D\) retains that group.
Both \(k\) and \(D_{\mathrm{cyc}}\)
are proper Galois extensions of \(D\). If one contains the
other, their compositum is plainly proper. Otherwise
\citet[Theorem~2.2]{FF_JardenRazon2020}, with \(R=D\), says
that their compositum is Hilbertian, so is again proper.
Thus \(k'/D\) is a proper Galois extension containing
\(\mu_\infty\). The preceding case applies to \(Lk'\);
algebraic descent gives the assertion for \(L\).

Finally \(\mathbb Q^{\mathrm{tr}}/\mathbb Q\) is Galois:
every conjugate of a totally real algebraic number is again
totally real. It is proper, since it is contained in \(\mathbb R\)
and omits \(i\). The theorem therefore applies directly.
This argument also proves that the Hilbertianity or smallness
assumption of AP-F.49 is a sufficient route, not a necessary
hypothesis for its normal-extension conclusion.
\end{proof}

\begin{claim}{F.11}{Proved}\label{clm:FFdescent}
Let \(K/k\) be algebraic in characteristic zero. If
\(\mathcal R(K)\) holds, then \(\mathcal R(k)\) holds.
\end{claim}
\begin{proof}
It is enough to verify the regular-affine point-function condition
of CHW Theorem~4.18(5). Given the pinned relative map over \(k\),
slice at each finite \'etale constant object corresponding to a
finite subextension of \(K/k\). These are the actual small \'etale
topoi of the finite scalar extensions. Their compatible inverse-image
functors pass to \(K\): every quasicompact separated \'etale chart,
every map between finitely many charts and every finite covering
family descends to a finite stage. CHW Theorem~A.8 and Corollary~A.9
ensure that the inverse-image functor preserves these representable
charts, while Stacks Tags~0EYL and~01ZM give their finite-presentation
descent \citep{FF_StacksLimits}. Thus we obtain a relative map on
the small topos of \(X_K\), not an assumed arbitrary transcendental
base-change square.

Closed points of \(X_K\) map to closed points of \(X\), because
\(K/k\) is algebraic. If the latter point has residue field
finite over \(k\), its scalar-extension fibre is zero-dimensional
over \(K\); every resulting point is closed. The lifted map
is therefore pinned. Reconstruction over \(K\) gives a unit
\(h_K\) on the regular source, with its prescribed point function.
Use one common algebraic closure \(\bar k\) of \(K\) and \(k\).
That point function is the original \(G_k\)-equivariant function
on \(X(\bar k)\). For every \(\sigma\in G_k\), the two
regular functions \(\sigma h_K\) and \(h_K\) have equal values
at every geometric point. Reducedness makes them equal.
Their finitely many coefficients therefore descend to \(k\).
CHW's criterion supplies reconstruction for all schemes over \(k\).
This verifies AP-F.32 in the direction used above.
\end{proof}

\begin{claim}{F.12}{Proved}\label{clm:FFaudit}
The dependent core of AP-F.24--F.49 used here is valid with its
stated field and morphism restrictions. Its positive conclusions
AP-F.37, F.43--F.46 and F.49 remain sufficient statements; its
negative conclusions AP-F.39--F.42 and F.47--F.48 remain obstructions.
None constitutes a necessary-and-sufficient classification.
\end{claim}
\begin{proof}
We record the checks at the places where a stronger, incorrect
replacement would break the argument.

In AP-F.24 the arithmetic wreath embedding is
\((G\times_P\cdots\times_PG)\rtimes S_n\), with one common
constant quotient \(P\). It is not the unrestricted wreath
product. The jet conditions on a high-degree pencil give simple
branch transpositions and isolated boundary branch labels.
A spanning-tree gauge puts a literal \(S_n\) in the geometric
monodromy. Remaining base labels have support at most two and
generate the geometric group in their coordinate projections.
For \(n\ge5\), \(A_n\) is contained in the stable derived
term. If \(u\) is one of those base elements and
\(\tau^{-1}(i)\) is outside its support, then
\([u,\tau]_i=u_i\). Thus the stable derived term projects onto
the whole geometric group in each coordinate. After Hilbert
specialization a prosolvable base change retains that term, and
the surviving constant quotient is still \(P\). Multiplying
by a literal even permutation returns the selected block without
changing its constant image. This proves arithmetic, rather than
merely geometric, detection. The same argument with the kernel
of the largest allowed formation quotient proves AP-F.44 and
AP-F.46. Properness of the hereditary formation excludes all
sufficiently large alternating groups, which is exactly what
this replacement needs.

AP-F.26--F.31 use the five steps exhibited in
Claim~\ref{clm:FFcore}: spanning of the coinvariant Albanese by
point cocycles, arithmetic equivariance after actual values descend,
integral Tate-Hom over a finitely generated field, integral boundary
residues, and equality of actual point functions. The norm residue
is weighted by residue degree. The principal-divisor lattice need
not be primitive; the integral Tate-Hom assertion, or equivalently
the torsion-free Picard Tate quotient, controls that denominator.
Pointwise Kummer equality leaves a divisible scalar and is not the
end of the proof. Pullbacks of \(T-b\) and closed points of
unbounded degree remove it. The final step is CHW Theorem~4.18(5),
not a claim that the fundamental group equals the full topos.

For AP-F.43--F.46, the two descent bases retain distinct roles:
the Hilbertian geometric base supplies detection, while a finitely
generated arithmetic field with normal \(G_L\) supplies
Tate-Hom equivariance. The norm field containing each actual value
is fixed before its Kummer level varies. AP-F.45 uses moved
geometric divisors when all algebraic constants are present;
AP-F.43 uses moved arithmetic primes when a number-field extension
is missing. Claim~\ref{clm:FFcore} verifies both alternatives
without imposing an unnecessary condition on which occurs.
AP-F.49 uses only the inclusion (F.9.1). Its small-group argument
uses open subgroups and quotients, not arbitrary closed subgroups.
Small extensions preserve Hilbertianity by
\citet[Theorem~4.4]{FF_BarySoroker2008}. The bounded-degree
compositum proof retains the arithmetic constant field and the
degree of the detecting point.

The algebraic descent step is Claim~\ref{clm:FFdescent}.
The subfield assertions in AP-F.37 and F.43--F.46 additionally
take relative algebraic constants in a finite-type descent field
and their normal closures; normality and regular base change
give precisely the claimed pro-formation quotient. These arguments
must retain their stated positive-transcendence restriction when
\(\overline{\mathbb Q}\) is among the constants.

For the negative statements, AP-F.39--F.40 are the centralizer
test of Claim~\ref{clm:F.6}. In AP-F.41--F.42 the Puiseux
closure uses series whose coefficients lie in a single finite
constant extension. The finite trace-dual expansion proves the
algebraic-core identity; arbitrary infinite coefficient fields
are not substituted. In AP-F.47--F.48 the compatible roots
\((1+t_1)^{1/n}\) are in the henselization by Hensel uniqueness.
The substitution \(t_1\mapsto t_1+t_1^2\) centralizes inertia
and the coefficient action on an explicitly generated algebraic
closure. A coefficient field and rationally independent values
reduce the Abhyankar equality case to an algebraic extension of
that monomial henselization. Transcendence defect destroys the
last algebraicity implication, so that proof does not cover it.

The external inputs just used were opened with their numbered
statements: SGA~1 XII~5.1 and XIII~2.12,~4.6; Serre~3.3.1,
3.4.1; Bary-Soroker--Paran~1.2,~3.2; Poonen--Slavov~1.5;
Spie\ss--Szamuely~4.1--4.3; Jossen~1.4; CHW~3.6,~4.18,
A.8--A.9; and the valuation statements listed in the bibliography.
The original Jannsen ICM article could not be reopened through
the oversized proceedings file; the exact integral theorem is
instead verified in Jossen's opened primary article. We do not
mark that unavailable original as opened.
\end{proof}

\subsection{The pro-system, with its morphisms specified}

We use the following precise relative interpretation of the
pro-form in cote~149, scans~6--8. Let \(I\) be the filtered poset
of finitely generated subfields \(F\subset k\) containing a fixed
field of definition \(F_0\) for \(X,Y\). Put
\(X_F=X_0\times_{F_0}F\), and similarly for \(Y_F\).
The transition functors are arbitrary geometric morphisms; they
are not required to be pinned. A \emph{coherent pinned relative
model morphism} is a germ, under cofinal restriction, of relative
pinned morphisms
\[
       \varphi_F:(X_F)_{\mathrm{et}}\longrightarrow
                         (Y_F)_{\mathrm{et}}
                  \quad\text{over }BG_F,
\]
together with commutation isomorphisms for every \(F\subset F'\),
over \(BG_F\), satisfying the composition coherence. Isomorphic
germs are identified. This definition includes the relative base
identifications; a family of absolute morphisms is insufficient.

\begin{claim}{F.13}{Proved}\label{clm:FFpro}
For any characteristic-zero field \(k\) and absolutely weakly
normal finite-type \(k\)-schemes \(X,Y\), their \(k\)-morphisms
are in natural bijection with coherent pinned relative model
morphisms as just defined. The assertion is independent of the
initial field and the chosen models, under cofinal refinement.
This proves a precise coherent relative interpretation of F1(iii).
It does not identify this category with an unspecified ordinary
pro-category equipped only with a condition on closed points.
\end{claim}
\begin{proof}
Finite-presentation descent gives models and descends each
\(k\)-morphism to some finitely generated field, by Stacks
Tag~01ZM \citep{FF_StacksLimits}. These models can be taken
absolutely weakly normal. In fact in characteristic zero this is
seminormality, and it descends under a faithfully flat field
extension: in an affine reduced noetherian ring, the criterion
\(b^2,b^3\in A\Rightarrow b\in A\) in its total quotient
ring descends by flatness and contraction of ideals. Apply this
to a model whose scalar extension is the given seminormal ring.
The same argument applies to every intermediate model. Thus the
algebraic morphism gives the specified coherent system.

Conversely, CHW Theorem~0.1 reconstructs each \(\varphi_F\) as
a unique morphism \(f_F:X_F\to Y_F\). We must prove that
compatibility of topoi really implies scalar compatibility,
because a transcendental transition need not be pinned.
Fix \(F\subset F'\). For a point \(x\in X_F(\bar F)\),
put \(E=F'(x)\), finite over \(F'\), in a common algebraic
closure. The two composites to \(Y_F\) have the same relative
point morphism over \(BG_F\), by coherence. Their image is the
closed point \(y=f_F(x)\). Since \(F(y)/F\) is finite
separable, relative morphisms \(BG_E\to BG_{F(y)}\) over
\(BG_F\) are exactly the \(F\)-embeddings \(F(y)\to E\).
This assertion needs no algebraicity of \(E/F\): the finite
\(G_F\)-set of embeddings is pulled back to \(G_E\), whose
fixed elements are precisely embeddings into \(E\).
Consequently \(f_{F'}(x)=(f_F)_{F'}(x)\) with their residue
embeddings, for every such \(x\).

For completeness this point argument also handles nonseparated
targets. Cover \(Y_F\) by affine opens \(V\). Coherence on
the corresponding subterminal objects says that the inverse
images of \(V_{F'}\) under the two morphisms are the same
open, namely \((f_F^{-1}V)_{F'}\). On this open the target
is affine. The \(\bar F\)-points of the model are dense after
any scalar extension: expand the finitely many coefficients of
a function in a linearly independent set over \(\bar F\).
Vanishing at all \(\bar F\)-points then makes every coefficient
function zero. Reducedness of the source therefore gives equality
of the two morphisms on that open. They glue, proving
\(f_{F'}=(f_F)_{F'}\).

Choose one index and restrict to the cofinal upper cone above it.
The entire family is now the scalar extension of its morphism
at that index; it gives a \(k\)-morphism. Uniqueness and
independence of models are again finite-presentation descent.
All comparisons retained the structural 2-cells, so this is a
relative statement.
\end{proof}

\begin{claim}{F.14}{Proved}\label{clm:FFc}
With coherent relative scalar descent understood as in
Claim~\ref{clm:FFpro}, candidate F3(c) is correct:
\(\mathcal R(k)\) holds if and only if every pinned relative
map between finite-type \(k\)-scheme topoi is the realization
of such a system of finitely generated models, after absolute
weak normalization of its source. The single commuting-square
interpretation is excluded by Claim~\ref{clm:F.8}.
\end{claim}
\begin{proof}
A finite-type target may also be replaced by its absolute weak
normalization: its small \'etale topos is unchanged, and the
right-adjoint property gives
\(\operatorname{Hom}(X^{\mathrm{awn}},Y^{\mathrm{awn}})
 =\operatorname{Hom}(X^{\mathrm{awn}},Y)\)
\citep[Theorem~1.13]{FF_CHW2024}. Thus the source and target
hypotheses of Claim~\ref{clm:FFpro} lose no maps in this criterion.
A coherent system realizes on the small \'etale sites by the
finite-presentation colimit of charts and covering families
(Stacks Tag~0EYL). Claim~\ref{clm:FFpro} shows that it is the
system of one algebraic morphism, so its realization is precisely
that morphism's small-topos map. Conversely every algebraic
morphism has such a system. Their essential image is therefore
exactly the algebraic maps. If reconstruction holds, every pinned
relative map is in that image. If every pinned relative map has
such a coherent descent, it is algebraic by the same claim.
Faithfulness holds over every characteristic-zero field: equal
relative topos maps give equal geometric-point maps, including
residue embeddings, and equality on affine opens of a reduced
source gives equality of the scheme maps. This proves both
directions. This characterization is a descent criterion; it
does not supply an intrinsic classification in terms of familiar
field invariants.
\end{proof}

\begin{claim}{F.15}{Proved}\label{clm:FFcomplex}
For \((\mathbb C,\mathbb C)\), the algebraic pro-profinite
group of cote~149, scan~8, is the system
\[
 \left\{
   G_F=\varprojlim_{N/F\ {
m finite\ Galois},\ N\subset\mathbb C}
             \operatorname{Gal}(N/F)
 \right\}_{F\subset\mathbb C, F/\mathbb Q\ {
m finitely\ generated}},
                                                               \tag{F.15.1}
\]
with restriction maps defined using the algebraic closures
\(\bar F=\{z\in\mathbb C:z\text{ is algebraic over }F\}\).
Its inverse limit is trivial, but the system is not the trivial
pro-object in profinite groups. The relative model theorem
\ref{clm:FFpro} applies to every absolutely weakly normal
finite-type complex scheme over this system of base topoi.
\end{claim}
\begin{proof}
An abstract algebraic description dispenses even with a chosen
embedding of the individual number fields in \(\mathbb C\).
Choose a set \(B\) of cardinality \(2^{\aleph_0}\) and an
algebraic closure \(\Omega\) of \(\mathbb Q(B)\).
The classification of algebraically closed fields identifies
\(\Omega\) with \(\mathbb C\), noncanonically.
A cofinal system of index fields consists of the finite Galois
extensions of \(\mathbb Q(B_0)\) inside \(\Omega\), where
\(B_0\) runs through the finite subsets of \(B\).
Indeed finitely many elements of \(\Omega\) are algebraic over
one such rational function field, and a finite normal closure
contains all of them. This description uses only independent
indeterminates, polynomial algebra, finite extensions and
automorphism groups. Different choices give isomorphic pro-objects,
with no claimed canonical identification.

Every finitely generated subfield is specified algebraically by
a finite tuple \(a\in\mathbb C^n\):
\(F=\operatorname{Frac}(\mathbb Q[T_1,\ldots,T_n]/I_a)\),
where \(I_a\) is the ideal of polynomial relations of that tuple.
Finite Galois extensions inside \(\mathbb C\), their finite
automorphism groups, composita, and restrictions define both
levels in (F.15.1); no analytic topology enters. For
\(F\subset F'\), \(\bar F\subset\bar F'\), and restriction
is a continuous group homomorphism \(G_{F'}\to G_F\).
Its image is the open subgroup corresponding to
\(F'\cap\bar F\), finite over \(F\) by Stacks
Lemma~9.26.11 \citep{FF_StacksFields}.

A compatible family of automorphisms at every \(F\) fixes each
element of \(\mathbb C\): place that element in the index field
\(\mathbb Q(a)\), where it must be fixed. Thus the inverse
limit is \(G_{\mathbb C}=1\). At the finite quotient level,
the same fact is visible because the index field may be enlarged
to contain any chosen finite Galois extension.
Nevertheless no transition to the stage \(G_{\mathbb Q}\)
is the trivial map. Its image is
\(G_{F'\cap\overline{\mathbb Q}}\), the nontrivial absolute
Galois group of a number field. The identity of a pro-object
can equal the zero morphism only if, for every stage, some later
transition to that stage is zero. This fails here. Keeping the
two levels of the pro-system, rather than replacing them by
their inverse limit or flattening the finite quotients, is
therefore essential. The final reconstruction statement is
Claim~\ref{clm:FFpro} with \(k=\mathbb C\).
\end{proof}


\subsection{Valuations: positive and negative cases, and the remaining test}

\begin{claim}{F.16}{Proved}\label{clm:FFmixed}
Residue characteristic \(p\) does not, even together with any
specified finite transcendence degree over \(\mathbb Q\), decide
reconstruction. For each prime \(p\) and integer \(d\geq0\),
there exist henselian valued fields of transcendence degree \(d\)
over \(\mathbb Q\), with residue characteristic \(p\), which
reconstruct, and others which fail reconstruction.
\end{claim}
\begin{proof}
Choose algebraically independent elements
\(a_1,\ldots,a_d\in\mathbb Q_p\), and give
\(D=\mathbb Q(a_1,\ldots,a_d)\) the restricted \(p\)-adic
valuation. Such a tuple exists: \(\mathbb Q_p\) is uncountable,
whereas a field of countable transcendence degree over the
countable field \(\mathbb Q\), and its algebraic closure, are
countable. Henselization universality embeds \(D^h\) as a
valued subfield of the henselian field \(\mathbb Q_p\).
Henselization is algebraic and immediate
\citep[Proposition~3.3.25, Corollary~3.3.26]{FF_ADH2017}.
Thus \(\operatorname{trdeg}_{\mathbb Q}D^h=d\), its value
group is \(\mathbb Z\), and its residue field is
\(\mathbb F_p\): these assertions hold for \(D\) because
\(\mathbb Q\subset D\subset\mathbb Q_p\).
The field \(D^h\) is sub-\(p\)-adic, so Carlson--Stix
Theorem~A gives \(\mathcal R(D^h)\)
\citep[Theorem~A, Theorem~4.1]{FF_CS2026}.
For \(d=0\) this is already the ordinary henselization of
\(\mathbb Q\) at \(p\), embedded in \(\mathbb Q_p\).

For the opposite verdict take an algebraic closure
\(K=\overline{\mathbb Q(t_1,\ldots,t_d)}\).
Extend the \(p\)-adic Gauss valuation on the displayed rational
function field to \(K\). Its residue characteristic remains
\(p\), and an algebraically closed valued field is henselian.
The transcendence degree is \(d\). The group \(G_K\) is
trivial while \(\operatorname{Aut}(K)\) is nontrivial; a
nonidentity automorphism of \(\overline{\mathbb Q}\) extends
after choosing a transcendence basis. The centralizer obstruction
of Claim~\ref{clm:F.6} therefore gives \(\neg\mathcal R(K)\).
This supplies both verdicts without inferring a uniform answer
from the residue characteristic.
\end{proof}

\begin{claim}{F.17}{Proved}\label{clm:FFdefect}
There is an explicit henselian field \(K/\mathbb Q\) of
transcendence degree two, residue characteristic zero and
transcendence defect one which fails reconstruction. Namely, in
\(\overline{\mathbb Q}((t))\) put
\[
 u=\sum_{n\geq1}t^{n!},\qquad
 K=\overline{\mathbb Q}(t,u)^h
\]
for the induced \(t\)-adic valuation and its henselization
inside that complete field. Then
\[
 \kappa(K)=\overline{\mathbb Q},\qquad
 v(K^\times)=\mathbb Z,\qquad G_K\simeq\widehat{\mathbb Z}.
\tag{F.17.1}
\]
The transcendence degree two is the smallest possible finite
transcendence degree for a nontrivial residue-characteristic-zero
valuation with positive transcendence defect.
\end{claim}
\begin{proof}
The power series \(u\) is transcendental over
\(\overline{\mathbb Q}(t)\). Indeed let
\(P(t,U)\in\overline{\mathbb Q}[t,U]\) be nonzero, of
\(U\)-degree \(e\) and \(t\)-degree at most \(m\), and
write \(u_N=\sum_{n=1}^N t^{n!}\). If \(P(t,u)=0\), then
\[
 v\bigl(P(t,u_N)\bigr)\geq(N+1)!,\qquad
 \deg_tP(t,u_N)\leq m+eN!
\]
whenever the polynomial on the left is nonzero. For all
sufficiently large \(N\) these inequalities are incompatible.
Consequently infinitely many distinct \(u_N\) would be roots
of the nonzero polynomial \(P(t,U)\) over
\(\overline{\mathbb Q}(t)\), a contradiction.
Henselization is algebraic and immediate, proving the asserted
transcendence degree, residue field and value group.

For compatible roots of \(t\), let
\(E=\bigcup_n K(t^{1/n})\). Each extension has degree \(n\)
by the distinct valuation classes of
\(1,t^{1/n},\ldots,t^{(n-1)/n}\); since all roots of unity
are in \(K\), it is cyclic Galois. The union is henselian,
with residue field \(\overline{\mathbb Q}\) and divisible
value group \(\mathbb Q\). In residue characteristic zero,
these two conditions make a henselian field algebraically closed
\citep[Corollary~3.3.22]{FF_ADH2017}. Thus \(E=\bar K\)
and \(G_K=\varprojlim_n\mu_n\simeq\widehat{\mathbb Z}\).
Its nontrivial centre contradicts the necessary condition in
Claim~\ref{clm:F.6}. Notice that this obstruction may fix \(K\)
pointwise: it is a nontrivial central automorphism of its algebraic
closure, detected by the retained relative structural 2-cell.

Finally the rational rank of a nontrivial ordered value group
is at least one. The residue transcendence degree is nonnegative.
Positive transcendence defect therefore requires
\(\operatorname{trdeg}_{\mathbb Q}K\geq2\). In our example
\(2-0-1=1\), as claimed. The attached finite script checks
the displayed factorial-gap inequalities for sample bounds; the
argument above proves them for all bounds and proves transcendence.
\end{proof}

\begin{claim}{F.18}{Open}\label{clm:FFremaining}
The following sufficient direction of candidate F3(a) remains
unproved and unrefuted in this work:
\[
 \bigl[C_{\operatorname{Aut}(\bar k)}(G_E)=1
          \text{ for every finite }E/k\bigr]
                  \quad\Longrightarrow\quad\mathcal R(k).
\tag{F.18.1}
\]
Nor have we decided every finite-transcendence-degree henselian
field with residue characteristic zero and transcendence defect.
In particular our arguments do not decide
\[
 K_* = \mathbb Q\bigl(t,\sum_{n\geq1}t^{n!}\bigr)^h
           \subset\mathbb Q((t)),
\tag{F.18.2}
\]
with its induced valuation. This is a specified unresolved test,
not a claim that it is the least field under inclusion or that
its centralizers satisfy (F.18.1).
The general mixed-residue-characteristic case is also not
classified. A previously unresolved degree-zero test was
\[
                  K_p=\operatorname{Frac}(\mathbb Z_{(p)}^{sh}),
                                                        \tag{F.18.3}
\]
with a chosen strict henselization inside
\(\overline{\mathbb Q}\). Its reconstruction is now
proved in Claim~\ref{clm:FFstrictreconstruct}, and
Claim~\ref{clm:FFdiscretealgebraic} treats every algebraic
extension of \(\mathbb Q\) admitting a mixed-characteristic
discrete valuation. These conclusions do not classify
arbitrary nondiscrete valuations or all finite-transcendence-degree
mixed-characteristic henselian fields.
Claims~\ref{clm:FFtamereconstruct}--\ref{clm:FFpadicsubfields}
further settle bounded wild ramification. The next
specified degree-zero test
\(U_5=T_5(\mu_{5^\infty})\) from
Claim~\ref{clm:FFcyclotomictest} is also settled
positively by Claim~\ref{clm:FFcyclotomicreconstruct},
despite the detector failure in
Claim~\ref{clm:FFcyclotomicdetector}.
Claims~\ref{clm:FFmaxabelianreconstruct} and
\ref{clm:FFnilpotentreconstruct} further settle the
maximal abelian and each fixed bounded-nilpotency
extension. A remaining specified degree-zero
test is the maximal metabelian field
\(B_2\) of Claim~\ref{clm:FFmetabeliantest},
which has the centralizer condition over every
finite extension. Its reconstruction is not decided.
Claims~\ref{clm:FFmetabelianaffine}--\ref{clm:FFhighergenusnorm}
subsequently determine affine automorphism point functions,
all abelian generic actions on rational and elliptic
Galois slices and on the specified genus-two slice,
and the portion of the etale site
generated by genus-zero covers. This includes all
abelian covers of arbitrary genus, all Galois genus-one
covers, and all degree-three genus-one covers.
Arbitrary covers and noninvertible pinned maps
remain unresolved.
\end{claim}
\begin{proof}[State of the attempted proof and its exact boundary]
Claim~\ref{clm:F.6} proves the necessary direction. The known
substitution ghosts of AP-F.41, F.42, F.47 and F.48 fail its
hypothesis, so none refutes the converse. In particular one must
not infer sufficiency from having checked all the previous
counterexamples. Attempts to replace the reconstruction argument
by this rigidity condition have not supplied the missing step:
centralizers control invertible semilinear ghosts, whereas a
general pinned map to \(\mathbb G_m\) has Kummer cocycles and
actual point values which have not been forced to descend by
centralizer triviality alone. The finite-quotient detection and
normal arithmetic action in Claim~\ref{clm:FFcore} are genuine
extra inputs; we have not deduced them from (F.18.1).

For valuations, AP-F.48 supplies a coefficient field \(F\) and
parameters with rationally independent values. Under the
Abhyankar equality the ambient field is algebraic over the
corresponding monomial henselization, so algebraic descent
transfers its failure. Positive transcendence defect removes
that algebraicity statement. It does not remove all failures,
as Claim~\ref{clm:FFdefect} shows, but that example uses
algebraically closed residue field and a central Galois group.
For (F.18.2) the residue field is \(\mathbb Q\), and the
same radical-and-constant argument gives
\(G_{K_*}\simeq\widehat{\mathbb Z}(1)\rtimes G_{\mathbb Q}\).
Its centre does not supply the previous obstruction.

There is a concrete obstruction to copying the old substitution
proof. Put \(u=\sum t^{n!}\). Then
\(v(u-t-t^2)=6\), whereas substitution
\(t\mapsto t+t^2\), \(u\mapsto u\), gives
\[
 u-(t+t^2)-(t+t^2)^2
       =-t^2-2t^3-t^4+t^6+\cdots,
\]
of valuation two. Thus that field embedding of
\(\mathbb Q(t,u)\) is not valuation preserving and does not
extend by the henselization universal property used in AP-F.48.
Allowing \(u\) to move would require an actual algebraic
self-correspondence preserving its formal branch; none has been
constructed here. Attempts to use rigid henselian fields would
also require proving the centralizer condition for every finite
extension and independently constructing a nonalgebraic pinned
map. Neither assertion follows from field rigidity alone.

Similarly, Claim~\ref{clm:FFpro} proves the pro-form for the
explicit coherent relative model category. The transition maps
of the models need not be pinned, so simply writing an ordinary
\(\operatorname{Pro}(\operatorname{Top}^{\mathrm{pin}})\)
is invalid. We have not proved that any other proposed definition
of a closed-point-preserving ordinary pro-morphism has the same
morphisms as that coherent category. The literal unspecified
pro-category reading of F1(iii), and that part of F6, is therefore
not silently marked proved. The correct coherent-descent version
of F3(c) is Claim~\ref{clm:FFc}; the bare commuting-square
version is refuted by Claim~\ref{clm:F.8}.
\end{proof}

\subsection{Function fields over unrestricted characteristic-zero constants}

\begin{claim}{F.19}{Proved}\label{clm:FFfunction}
Let \(F\) be any characteristic-zero field. Every finitely
 generated transcendental extension \(L/F\) satisfies
\(\mathcal R(L)\). In particular this applies to
\(\mathbb C(t)\), without a finite-transcendence-degree
restriction on the constant field.
\end{claim}
\begin{proof}
First suppose that the constant field \(C\) is algebraically
closed, and that \(L/C\) is finitely generated of positive
transcendence degree. The field \(L\), and every finite extension
of it, is Hilbertian by
\citet[Theorem~1.2, Corollary~3.2]{FF_BarySorokerParan2013}.
The divisor argument gives
\[
                     D(E^\times)=C^\times
        \quad(E/L\text{ finite}),                         \tag{F.19.1}
\]
since the relative constants in \(E/C\) are \(C\) itself.
Consequently arithmetic finite-quotient detection and unbounded
closed-point degrees hold exactly as in Claim~\ref{clm:FFcore}.
We verify separately the two steps of that proof which previously
used an arithmetic normal algebraic extension.

\emph{Integral puncture exponents.}
Use the punctured line, polynomials \(P_j\), actual value at
\(x_0\), and coefficients \(\lambda_j\) of (F.9.3).
Choose a nontrivial finite Galois extension \(M/L\) containing
all roots of the \(P_j\), and
\(1\ne\sigma\in\operatorname{Gal}(M/L)\).
The finitely generated function field \(M/C\) has a proper
normal model, and the automorphism \(\sigma\) fixes its
algebraically closed constant field. The moved-prime-divisor
argument in Claim~\ref{clm:FFcore} supplies the disjoint good
valuation pairs. Weak approximation supplies \(\alpha\in M\)
with the valuation pattern (F.9.4).
This time the actual value \(b_\alpha\) belongs to \(M\)
itself, because \(M/L\) is finite. At every Kummer level,
evaluate the class and apply \(\sigma\). Since the constant
normalization is fixed by \(\sigma\), we obtain
\[
 \frac{\sigma b_\alpha}{b_\alpha}
        \prod_j\left(\frac{P_j(\sigma\alpha)}
                            {P_j(\alpha)}\right)^{-a_{j,n}}
                         \in(M^\times)^n.
\]
Apply the discrete valuation \(v_i\) of \(M/C\), for every
\(n\). It gives
\(\lambda_i=-v_i(\sigma b_\alpha/b_\alpha)\in\mathbb Z\).
No norm to a finitely generated \(\mathbb Q\)-field and no
valuation on an infinite algebraic compositum are needed.

\emph{Arithmetic equivariance through a normal image.}
Let \(X/L\) be regular, affine and connected, with rational
base point \(x_0\), and let \(A\) be its generalized Albanese.
Choose a finitely generated subfield \(C_0/\mathbb Q\) of
\(C\) and a regular finitely generated function field
\(F_0/C_0\), with compatible models for this finite list of
geometric data, such that
\[
 L=F_0C,\qquad F_0\text{ and }C\text{ are linearly disjoint
 over }C_0.                                                \tag{F.19.2}
\]
These are ordinary models of the geometrically integral
\(C\)-variety with function field \(L\); enlarge \(C_0\)
to include all coefficients. Geometric integrality is checked
after the algebraically closed extension \(C/C_0\), so the
model has the required regularity. Put
\(\bar C_0=\{c\in C:c\text{ is algebraic over }C_0\}\).
The algebraic closure of \(F_0\) inside \(L\) is
\(F_0\bar C_0\). To see this after adjoining \(\bar C_0\),
a nontrivial finite extension of the function field of a normal
geometrically integral variety cannot acquire a rational generic
section merely by extending its algebraically closed constant
field: a connected normal cover over that field remains
geometrically connected. This also proves linear disjointness
from every finite algebraic extension not already in the constant
extension. Thus restriction has image
\[
 \operatorname{im}(G_L\longrightarrow G_{F_0})
           =H:=G_{F_0\bar C_0}\triangleleft G_{F_0}.       \tag{F.19.3}
\]
This is a quotient image, not an identification of \(G_L\)
with a subgroup of \(G_{F_0}\). Surjectivity onto the indicated
closed subgroup follows equivalently from the fixed field and
the Galois correspondence, or from the finite linear-disjointness
statement just proved.

Fix \(\ell\). The geometric Kummer functional
\(\lambda:T_\ell A\to\mathbb Z_\ell(1)\) of the pinned map
is invariant under \(G_L\); its action factors through (F.19.3).
Since \(H\) fixes all roots of unity, \(\lambda\) factors
through \(W=(V_\ell A)_H\). Finite-quotient detection shows,
as in Claim~\ref{clm:FFcore}, that the projected point-Kummer
cocycles span \(W\). Choose finitely many spanning points and
put their finite fields of definition into a single finite
extension \(E/L\). Restriction to \(G_E\) keeps their spans,
since their projected cocycles are homomorphisms into a rational
\(\ell\)-adic vector space and this subgroup has finite index.

Their coordinates and actual multiplicative values are finitely
many elements of \(E\). Descend them by enlarging \(C_0\)
to a finitely generated \(C_1\subset C\) and taking a
regular finitely generated model field \(M_1/C_1\) inside
\(E\), containing \(F_0\), such that
\(E=M_1C\). In particular \(M_1/\mathbb Q\) is finitely
 generated. The same argument gives
\[
 \operatorname{im}(G_E\to G_{M_1})
       =H_1:=G_{M_1\bar C_1}\triangleleft G_{M_1}.         \tag{F.19.4}
\]
Under restriction \(r:G_{M_1}\to G_{F_0}\), the image
\(r(H_1)\) is the open subgroup
\(\operatorname{im}(G_E\to H)\) of \(H\).
In particular \(H_1\) acts trivially on \(W\).

Each chosen Albanese point and its multiplicative value is
\(M_1\)-rational. Its Kummer cocycle is therefore defined on
\(G_{M_1}\). Equality of the two evaluated cocycles on
\(G_E\) descends to \(H_1\), since \(G_E\to H_1\) is
surjective and both cocycles are restrictions of cocycles over
\(M_1\). For \(g\in G_{M_1}\), \(h\in H_1\), the
cocycle conjugation identity, projected to \(W\), now gives
\[
 \lambda\bigl(r(g)c_x(h)\bigr)
             =\chi_\ell(g)\lambda\bigl(c_x(h)\bigr).
\]
The error term is of the form
\((1-r(ghg^{-1}))c_x(g)\), which vanishes in \(W\)
because \(H_1\) is normal and its image lies in \(H\).
The finitely many images span \(W\), proving
\(G_{M_1}\)-equivariance of \(\lambda\).
The integral Tate-Hom statement of
\citet[Section~1.4]{FF_Jossen2014}, now over the genuinely
finitely generated field \(M_1/\mathbb Q\), puts this
functional in the completed normalized-unit lattice.

The rest of Claim~\ref{clm:FFcore} transfers without a new
field hypothesis: norms to punctured lines and transverse curves
make boundary coefficients integral; the integer divisor matrix
makes unit exponents integral; base-point evaluation removes the
constant cohomology class. Equation (F.19.1), high-degree closed
points and pullbacks of \(T-b\) remove the divisible pointwise
scalars. In the constant case choose a proper finite extension of
\(L\), which Hilbertianity supplies. This proves the precise
point-function condition CHW~4.18(5), hence \(\mathcal R(L)\).

Finally let \(L/F\) be as in the statement, without assuming
\(F\) algebraically closed. In an algebraic closure of \(L\)
choose an algebraic closure \(\bar F\) of \(F\). The field
\(L\bar F\) is finitely generated of the same positive
transcendence degree over \(\bar F\), so reconstructs by the
case just proved. Since \(L\bar F/L\) is algebraic,
Claim~\ref{clm:FFdescent} gives \(\mathcal R(L)\).
\end{proof}

\begin{claim}{F.20}{Known}\label{clm:FFmixedcentralizer}
The necessary centralizer condition of candidate F3(a) holds for
every characteristic-zero field admitting a discrete valuation
of residue characteristic \(p>0\). It also holds for every
field finitely generated and transcendental over any subfield.
These statements do not by themselves prove reconstruction.
\end{claim}
\begin{proof}
Hoshi--Tsujimura Corollary~4.1 says that
\(\operatorname{Aut}(E)\to\operatorname{Out}(G_E)\) is
injective if \(E\) embeds in the fraction field of a noetherian
local domain of mixed characteristic. Corollary~4.5(3) makes
\(G_E\) slim under the same hypothesis
\citep{FF_HoshiTsujimura2023}. A mixed-characteristic discrete
valuation ring is such a domain. Every finite extension of its
fraction field has a prolongation with value group contained
in \(\frac1e\mathbb Z\), hence again has a discrete
mixed-characteristic valuation. Thus both conclusions hold
for every finite extension \(E/k\).

An automorphism of \(\bar k\) centralizing \(G_E\) preserves
its fixed field \(E\); its restriction is in the kernel of the
displayed outer-action map. Injectivity makes that restriction
the identity, leaving an element of \(Z(G_E)=1\). This proves
the desired centralizer statement. For a finitely generated
transcendental field use instead Hoshi--Tsujimura Corollary~1.6
and Corollary~4.5(1); that hypothesis is also preserved by finite
extensions. It gives an independent necessary-condition check
on Claim~\ref{clm:FFfunction}.
\end{proof}

\subsection{The sufficient classes in R9, with their quantifiers}

\begin{claim}{F.21}{Proved}\label{clm:FFclasses}
The following are sufficient conditions for \(\mathcal R(k_0)\),
always in characteristic zero.
\begin{enumerate}
\item There is a finitely generated \(D/\mathbb Q\) with
\(D\subset k_0\), \(k_0/D\) algebraic, and proper normal
closure \(N\ne\bar D\) inside \(\bar D\).
\item For some finitely generated \(D/\mathbb Q\), the field
\(k_0\) is a subfield of a finitely generated extension of
\(D^{\mathcal C}\). Here \(\mathcal C\) is a proper
isomorphism-invariant class of finite groups, containing all
cyclic groups and closed under subgroups, quotients and
extensions; \(D^{\mathcal C}\) is the maximal Galois
pro-\(\mathcal C\) extension. Equivalently, for some finitely
generated \(D_0\subset k_0\), the extension \(k_0/D_0\)
is algebraic and its normal-closure group is virtually
pro-\(\mathcal C\). This includes AP-F.37 and AP-F.43
by taking \(\mathcal C\) to be the solvable groups, and
AP-F.44 in general.
\item The field \(k_0\) has positive transcendence degree
over \(\mathbb Q\) and is a subfield of a finitely generated
extension of \(K_0^{\mathcal C}\), where
\(K_0/\overline{\mathbb Q}\) is finitely generated of
positive transcendence degree. Equivalently, it is algebraic
over some finitely generated \(D_0/\mathbb Q\) of positive
transcendence degree whose normal closure \(N/D_0\) has
\[
 \operatorname{Gal}
       (N\overline{\mathbb Q}/D_0\overline{\mathbb Q})
       \quad\hbox{virtually pro-}\mathcal C.
\]
This is AP-F.46. In particular every positive-transcendence
subfield of a finitely generated extension of
\(\overline{\mathbb Q}\) reconstructs, as in AP-F.45.
\item In the setting of AP-F.49, let \(D/\mathbb Q\)
be finitely generated of positive transcendence degree, and
let \(k/D\) be algebraic Galois with
\(\overline{\mathbb Q}\subset k\). If \(k\) is
Hilbertian, every finitely generated \(L/k\) reconstructs.
If instead
\(\operatorname{Gal}(k/D\overline{\mathbb Q})\) is
small, then also every positive-transcendence subfield
\(k_0\subset L\) reconstructs. Small means finitely many
open subgroups of each prescribed finite index.
\end{enumerate}
In (2) and (3), taking the union over \(\mathcal C\) is
equivalent to requiring that an open subgroup of the respective
normal-closure group have no finite quotient containing some
fixed \(A_m\), \(m\ge5\), as a subquotient.
These are sufficient classes and equivalences between their
descriptions, not a necessary classification of all reconstruction
fields. In particular (1) does not assert that arbitrary subfields
of a reconstruction field reconstruct.
\end{claim}
\begin{proof}
For (1), Claim~\ref{clm:FFnormal} gives \(\mathcal R(N)\),
and Claim~\ref{clm:FFdescent} descends along the algebraic
extension \(N/k_0\).

We give the subfield argument for (2), since subfield passage
cannot be inferred just from (1). Put \(k=D^{\mathcal C}\),
and choose a finitely generated \(F/\mathbb Q\) inside
\(L\), containing \(D\), with \(L=Fk\). Given
\(k_0\subset L\), choose a finite transcendence basis \(T\)
of \(k_0/\mathbb Q\), enlarge \(F\) to contain it, and
put \(A=\mathbb Q(T)\). The relative algebraic closure
\(R\) of \(A\) in \(F\) is finite over \(A\)
\citep[Lemma~9.26.11]{FF_StacksFields}, and \(F/R\) is
regular. Let \(N/R\) be the normal closure of \(Rk_0/R\).
Regularity identifies its conjugates over \(R\) with those
after scalar extension to \(F\). Since \(L/F\) is Galois
pro-\(\mathcal C\), all these conjugates lie in \(L\),
and restriction gives
\[
 \operatorname{Gal}(N/R)
   \simeq\operatorname{Gal}(FN/F),
\]
a quotient of \(\operatorname{Gal}(L/F)\). Thus it is
pro-\(\mathcal C\). This normal extension is proper.
Indeed a proper hereditary extension-closed class containing
the cyclic groups omits \(A_n\) for all sufficiently large
\(n\): if it contained arbitrarily large alternating groups,
the embeddings \(G\hookrightarrow S_{|G|}\hookrightarrow
A_{|G|+2}\) would put every finite group in the class.
Every finitely generated characteristic-zero field is
Hilbertian, and the Morse-polynomial construction supplies
\(S_n\)-extensions for every \(n\)
\citep[Theorems~3.4.1, 4.4.5]{FF_Serre2007}; for positive
transcendence degree use also
\citep[Theorem~1.2]{FF_BarySorokerParan2013}.
Consequently its absolute Galois group cannot be pro-
\(\mathcal C\). Apply (1) over \(R\), and then descend
from \(N\) to \(k_0\), since \(R/A\) and \(k_0/A\)
are algebraic.

For the stated equivalence, take the finite Galois closure
\(R'/A\) of \(R/A\). The normal closure of \(k_0/A\)
is contained in the compositum of the finitely many
\(A\)-conjugates of \(NR'\). Over \(R'\) this is
pro-\(\mathcal C\): its group is a closed subgroup of a
finite product of pro-\(\mathcal C\) groups. Thus its
group over \(A\) is virtually pro-\(\mathcal C\).
Conversely an open pro-\(\mathcal C\) subgroup of a normal
closure fixes a finite extension \(D_1/D_0\); the whole
normal closure is contained in \(D_1^{\mathcal C}\).
This supplies exactly the embedding condition in (2).

For (3), choose \(D/\mathbb Q\) finitely generated with
\(K_0=D\overline{\mathbb Q}\). The field
\(k=K_0^{\mathcal C}\) is normal over \(D\): every
\(D\)-automorphism preserves \(K_0\) and permutes its
Galois \(\mathcal C\)-extensions. We can choose \(F\)
as above with \(L=Fk\) normal over \(F\), while
\(L/F\overline{\mathbb Q}\) is pro-\(\mathcal C\).
The same \(A,R,N\) argument now gives
\[
 \operatorname{Gal}
   (N\overline{\mathbb Q}/R\overline{\mathbb Q})
 \quad\hbox{pro-}\mathcal C.
\]
Here regularity survives algebraic extension of constants,
so \(F\overline{\mathbb Q}/R\overline{\mathbb Q}\)
is regular and the indicated group is the corresponding
quotient of \(\operatorname{Gal}
(L/F\overline{\mathbb Q})\). The field
\(R\overline{\mathbb Q}\) has positive transcendence
degree over \(\overline{\mathbb Q}\); it is Hilbertian
and has the same arbitrarily large symmetric-group quotients.
Therefore \(N\overline{\mathbb Q}\ne\bar R\).
This is a proper normal extension of the finitely generated
field \(R/\mathbb Q\); (1) and algebraic descent apply.
The equivalence follows by the same finite normal-closure
argument, now after adjoining \(\overline{\mathbb Q}\).
Conversely, a finite extension of
\(D_0\overline{\mathbb Q}\) fixed by the specified open
subgroup is a finitely generated positive-transcendence
\(\overline{\mathbb Q}\)-field and supplies the required
geometric pro-\(\mathcal C\) envelope. The special case
with no geometric algebraic enlargement gives AP-F.45;
it also follows directly from Claim~\ref{clm:FFfunction}
and algebraic descent.

For (4), its Hilbertian assertion is exactly
Claim~\ref{clm:FFcore}. In the small case the original
Hilbertian proof uses
\citet[Theorem~4.4]{FF_BarySoroker2008} over the Hilbertian
field \(D\overline{\mathbb Q}\). We can also derive
the subfield conclusion directly from (1).
In the notation just used, \(k\cap F\overline{\mathbb Q}\)
is finite over \(D\overline{\mathbb Q}\), because it
is algebraic inside a finitely generated extension.
Thus \(\operatorname{Gal}(L/F\overline{\mathbb Q})\)
is an open subgroup of the given small group, and is small.
Its quotient
\(\operatorname{Gal}(N\overline{\mathbb Q}/
R\overline{\mathbb Q})\) is small as well. It cannot
be the entire absolute Galois group of
\(R\overline{\mathbb Q}\). To check this without a
smallness assertion for arbitrary closed subgroups, choose
a separating transcendence basis containing \(x\).
Outside finitely many \(a\in\overline{\mathbb Q}\),
the extension of the corresponding rational-function field
is unramified at \(x-a\), so a prime over this divisor
has \(\operatorname{ord}(x-a)=1\). The distinct
\(x-a\) therefore give infinitely many distinct square
classes, hence infinitely many quadratic extensions.
This absolute Galois group is not small. Once again
\(N\overline{\mathbb Q}/R\) is proper normal, and (1)
and algebraic descent prove the assertion.

Finally, groups with no \(A_m\) subquotient form a proper
class closed under subgroups, quotients and extensions and
containing all cyclic groups. For extension closure, a simple
nonabelian subquotient in an extension must be a subquotient
of either its normal kernel or its quotient: intersect the
chosen subgroup with the kernel and use simplicity.
Conversely the alternating-group exclusion for any proper
\(\mathcal C\) was proved above. This proves the final
equivalence and retains all the finite-extension and
positive-transcendence restrictions.
\end{proof}

\subsection{The negative classes and the two required ghost tests}

\begin{claim}{F.22}{Proved}\label{clm:FFghosts}
The negative statements AP-F.39--F.42 and AP-F.47--F.48
have the following precise content.
\begin{enumerate}
\item Every characteristic-zero field with finite absolute
Galois group fails \(\mathcal R\); in particular all
algebraically closed and real closed fields fail.
\item Let \(F\) have characteristic zero and equip
\(F(t_1,\ldots,t_r)\), \(r\ge1\), with a monomial
valuation whose positive parameter values are rationally
independent. Its henselization \(H\), and every algebraic
extension of \(H\), fails \(\mathcal R\). This includes
every finite-transcendence-degree field over \(\mathbb Q\)
with a nontrivial henselian valuation of residue
characteristic zero satisfying the Abhyankar equality.
\item For any characteristic-zero \(F\), put
\[
 P(F)=\bigcup_{n\ge1}F((t^{1/n})),\qquad
 A(F)=\{a\in P(F):a\text{ is algebraic over }F(t)\},
\]
with compatible roots throughout. The fields \(F((t))\),
\(P(F)\), \(A(F)\), and all their algebraic extensions
fail \(\mathcal R\). There are topological isomorphisms
\[
 (G_{P(F)},\chi_{P(F)})\simeq(G_F,\chi_F),\qquad
 (G_{A(F)},\chi_{A(F)})\simeq(G_F,\chi_F),
\]
with the full cyclotomic characters, and
\(\lvert A(F)\rvert=\max(\lvert F\rvert,\aleph_0)\).
\end{enumerate}
The explicit obstruction in (2) and (3) is the lift of
\(t_1\mapsto t_1+t_1^2\), respectively
\(t\mapsto t+t^2\), centralizing the full absolute Galois
group. It gives an actual pinned relative automorphism of
the full \(\mathbb G_m\) topos. These ghosts also violate
the genus-zero Isom condition in candidate F3(b), already
on the four-punctured lines with parameters \(t\) and
\(t+t^2\); and they cannot satisfy coherent scalar descent
in candidate F3(c). Thus all three candidates correctly
exclude these particular ghosts, although candidate (b)
is false by its opposite-direction failure in
Claim~\ref{clm:F.7}.
\end{claim}
\begin{proof}
An algebraically closed characteristic-zero field has a
nonidentity automorphism. Choose a nontrivial automorphism
of \(\overline{\mathbb Q}\), fix a transcendence basis,
and extend to the algebraic closure of the resulting
rational-function field. The resulting embedding is onto:
its image is algebraically closed and the whole field is
algebraic over it. Its Galois centralizer condition fails,
so Claim~\ref{clm:F.6} applies. Reconstruction ascends
finite extensions by slicing at their finite \(\acute etale\)
constant objects, as in that claim. Hence a finite absolute
Galois group would force the finite algebraic closure to
reconstruct, a contradiction. In the real closed case one
can also use its nontrivial central element in the group
of order two; triviality of the field's own automorphism
group does not remove that relative structural-cell ghost.

For (2), write \(\Gamma_0=\sum_i\mathbb Z\gamma_i\)
for the monomial value group. Distinct exponent vectors
have distinct values, so least monomials prove both
multiplicativity and residue field \(F\). Henselization
is immediate \citep[Corollary~3.3.26]{FF_ADH2017}, giving
\(v(H^\times)=\Gamma_0\), \(\kappa(H)=F\).
Inside an algebraic closure choose \(\bar F\) and
compatible roots \(u_{i,n}=t_i^{1/n}\). For every finite
\(F'/F\) in \(\bar F\), residue independence of an
\(F\)-basis of \(F'\) proves
\[
 [HF':H]=[F':F],\quad
 v((HF')^\times)=\Gamma_0,\quad \kappa(HF')=F'.
\]
Indeed, dividing a nonzero linear combination by one of
its least-valued coefficients gives a nonzero residue.
The \(n^r\) monomials
\(u_{1,n}^{a_1}\cdots u_{r,n}^{a_r}\),
\(0\le a_i<n\), have distinct value classes modulo
\(\Gamma_0\). They are independent over \(HF'\) and
form a basis of
\(E(F',n)=HF'(u_{1,n},\ldots,u_{r,n})\), whose degree,
value group and residue field are therefore
\[
 [E(F',n):H]=[F':F]n^r,\quad
 v(E(F',n)^\times)=\tfrac1n\Gamma_0,\quad
 \kappa(E(F',n))=F'.
\]
The directed union \(\Omega=\bigcup_{F',n}E(F',n)\)
is henselian with divisible value group and algebraically
closed residue field. Since the residue characteristic
is zero, \citet[Corollary~3.3.22]{FF_ADH2017} makes
\(\Omega\) algebraically closed. This constructs the
entire algebraic closure and gives
\[
 G_H\simeq\widehat{\mathbb Z}(1)^r\rtimes G_F.
                                                        \tag{F.22.1}
\]
The section acts on algebraic constants and fixes all
chosen parameter roots; the kernel multiplies the latter
by compatible roots of unity. Thus the action is the
diagonal cyclotomic action.

The substitution \(t_1\mapsto t_1+t_1^2\), fixing the
other variables and \(F\), preserves the monomial
valuation: each monomial acquires a unit of residue one
and the least value remains unique. It extends uniquely
to an embedding \(\sigma:H\to H\). Hensel's lemma
gives a unique positive-valued root \(r_1\) of
\(r_1+r_1^2=t_1\), with \(v(r_1)=\gamma_1\) and
\(r_1/t_1\) of residue one. Substitution by \(r_1\)
is likewise valuation preserving. The two induced
embeddings of \(H\) are inverse, first on the rational
function field and then by the uniqueness in henselization
universality. Thus \(\sigma\) is an automorphism.

For each \(n\), let \(\beta_n\in1+\mathfrak m_H\)
be the unique root of \(X^n=1+t_1\). These exist because
\(n\) is a unit in the residue field, and uniqueness
gives \(\beta_{mn}^{m}=\beta_n\). On \(E(F',n)\)
define
\[
 \tau|_H=\sigma,\quad \tau|_{F'}=1,\quad
 \tau(u_{1,n})=u_{1,n}\beta_n,\quad
 \tau(u_{i,n})=u_{i,n}\ (i>1).
\]
The displayed basis proves this is a well-defined
automorphism: the images satisfy the transformed equations,
and the image contains the base and all original generators.
Compatibility gives \(\tau\in\operatorname{Aut}(\Omega)\).
Every \(g\in G_H\) preserves \(\bar F\), multiplies
\(u_{i,n}\) by a root of unity, and fixes \(\beta_n\).
Consequently \(g\tau=\tau g\) on the field generators
of every \(E(F',n)\), hence on \(\Omega\). This is a
full-group centralizer calculation, not a claim about
the tame quotient alone.

Claim~\ref{clm:F.6} supplies the pinned relative
\(\mathbb G_m\)-topos automorphism. Its rational-point
function is \(a\mapsto\sigma^{-1}(a)\). A
\(H\)-morphism of \(\mathbb G_m\) is \(cT^m\),
and a topos automorphism forces \(m=\pm1\). The values
at \(1\) and \(2\) force \(c=1,m=1\), while the
value at \(t_1\) is \(r_1\ne t_1\). Thus the map
is not algebraic. Claim~\ref{clm:FFdescent}, used
contrapositively, excludes every algebraic extension.
Equivalently, the same \(\tau\) normalizes the fixed field
of every closed subgroup of \(G_H\), because it centralizes
that subgroup. It therefore restricts to every intermediate
algebraic extension and still moves \(t_1\).

For the intrinsic Abhyankar assertion, a henselian
equicharacteristic-zero valued field has a coefficient
subfield \(F\subset\mathcal O_v\) mapping isomorphically
to its residue field \citep[Proposition~3.3.8]{FF_ADH2017}.
Concretely, lift a residue transcendence basis and enlarge
the resulting coefficient field maximally; any missing
algebraic residue element has a separable minimal polynomial,
whose simple root lifts by Hensel's lemma, contradicting
maximality. Choose elements \(t_1,\ldots,t_r\) whose
positive values form a rational basis of the value group.
Their values imply algebraic independence over \(F\)
and the monomial valuation above. The henselization
embeds in the given henselian field. The Abhyankar equality
is exactly what makes the given field algebraic over
\(F(t_1,\ldots,t_r)^h\), so (2) follows. With
transcendence defect this last algebraicity conclusion
is unavailable, as Claim~\ref{clm:FFremaining} records.

For (3) the exact bounded-coefficient Puiseux closure is
\[
 \Omega=\bigcup_{\substack{F'/F\ \mathrm{finite}\\n\ge1}}
                       F'((t^{1/n})),\qquad F'\subset\bar F.
                                                        \tag{F.22.2}
\]
It is an algebraic closure of \(F((t))\)
\citep[Proposition, p.~314]{FF_vanDenDriesRibenboim1984}.
The finite coefficient restriction matters. A series with
arbitrary coefficients in \(\bar F\) need not be
algebraic over \(F((t))\). Expanding the coefficients
in a finite \(F\)-basis of \(F'\) proves
\(\Omega=P(F)\bar F\). Also \(P(F)\cap\bar F=F\):
an element algebraic over \(F\) cannot have nonzero
valuation, by comparing the unique least term in its
minimal polynomial; reduction then supplies a root in
\(F\) of that irreducible polynomial. Coefficientwise
action therefore identifies \(G_{P(F)}\) with \(G_F\).
It and its inverse are continuous, since every element
of \(\Omega\) uses one finite coefficient extension.
Roots of unity are constants, proving cyclotomic compatibility.

Let \(\Omega_0\) be the elements of \(\Omega\)
algebraic over \(F(t)\). It is algebraically closed and
is an algebraic closure of \(A(F)\). We verify
\[
                         \Omega_0=A(F)\bar F.             \tag{F.22.3}
\]
Take \(a\in\Omega_0\cap F'((v))\),
\(v=t^{1/n}\), enlarging \(F'/F\) to be finite
Galois. If \((e_i)\) and \((e_i^*)\) are trace-dual
\(F\)-bases of \(F'\), write
\[
 a=\sum_i e_i a_i,\qquad
 a_i=\operatorname{Tr}_{F'((v))/F((v))}(e_i^*a).
\]
Every coefficientwise conjugate of \(e_i^*a\) is
algebraic over \(F(t)\), so \(a_i\in A(F)\).
This proves (F.22.3); its converse inclusion is immediate.
Restriction now gives the asserted topological and
cyclotomic identification \(G_{A(F)}\simeq G_F\).
For a finite \(E/F\), the same coefficient expansion
gives \(P(F)E=P(E)\), and the same trace argument gives
\(A(F)E=A(E)\). Hence all the corresponding open-subgroup
and finite-extension data agree as well.

Substitution \(t\mapsto t+t^2\) acts on every finite
Puiseux level by
\[
 t^{1/n}\longmapsto t^{1/n}(1+t)^{1/n},
\]
using the rational binomial series of constant term one.
Its inverse uses
\(r(t)=(-1+\sqrt{1+4t})/2\) and the corresponding
constant-term-one root of \(r(t)/t\). These substitutions
preserve (F.22.2), are compatible as the denominator and
the coefficient field increase, and fix \(\bar F\).
They centralize coefficientwise \(G_F\). They also
preserve \(A(F)\) and \(\Omega_0\): one direction is
algebraic over \(F(t+t^2)\subset F(t)\), the other over
\(F(r(t))\), which is algebraic over \(F(t)\).
For \(F((t))\), an element of its Galois group acts on
the generators of (F.22.2) by coefficient action and
root multipliers. The same generator calculation as in
(2) proves centralization there too. Thus the displayed
substitutions give actual nonalgebraic pinned relative
maps on all three full topoi, and algebraic descent
excludes all their algebraic extensions. Finally \(A(F)\)
is algebraic over and contains \(F(t)\), which proves
its cardinality assertion. Taking \(F=\mathbb Q\)
compares the reconstructing field \(\mathbb Q\) with
the nonreconstructing countable field \(A(\mathbb Q)\)
having the same full cyclotomic Galois data.

To test the candidates on these ghosts, put
\(U_t=\mathbb P^1\setminus\{0,1,\infty,t\}\)
and \(U_{t+t^2}=\mathbb P^1\setminus
\{0,1,\infty,t+t^2\}\). The centralizing lift makes
the semilinear isomorphism a relative fundamental-group
isomorphism over the identity of \(G_k\), with only
geometric inner ambiguity. A scheme isomorphism would
require \(t+t^2\) in
\[
 \left\{t,t^{-1},1-t,(1-t)^{-1},
                    (t-1)/t,t/(t-1)\right\}.
\]
Every such equality is a nonzero polynomial relation on
the transcendental element \(t\); none holds. Thus
candidate (b) excludes these fields even at four punctures.
Candidate (a) fails by the centralizing lift itself.
Candidate (c), with actual coherent scalar descent,
would contradict Claim~\ref{clm:FFc}, since its descended
map reconstructs over the finitely generated field.
The program \texttt{F\_puiseux\_checks.py} verifies both
inverse substitutions modulo \(t^{13}\), the compatible
roots for denominators dividing 12, and the six exact
cross-ratio inequalities. The infinite identities and
the topos maps have been proved above.
\end{proof}

\subsection{A countable mixed-characteristic test of transcendence degree zero}

\begin{claim}{F.23}{Proved}\label{clm:FFstricttest}
Fix a prime \(p\) and a prolongation of its valuation to
\(\overline{\mathbb Q}\), and let
\(K_p=\operatorname{Frac}(\mathbb Z_{(p)}^{sh})\) as
in (F.18.3). It is countable, algebraic over \(\mathbb Q\),
henselian discretely valued, with value group \(\mathbb Z\)
and residue field \(\overline{\mathbb F}_p\).
It satisfies candidate F3(a)'s centralizer condition over
every finite extension and candidate F3(b)'s Isom condition
for all hyperbolic curves over every finite extension.

For every number field \(D\subset K_p\), its normal
closure over \(D\) is \(\overline{\mathbb Q}\).
It is not sub-\(q\)-adic for any prime \(q\), and is
not a finite extension of a proper normal algebraic extension
of a number field. Consequently it lies in none of the
sufficient field classes stated in Claims~\ref{clm:F.4},
\ref{clm:FFnormal}, \ref{clm:FFfunction}, and
\ref{clm:FFclasses}. Its centralizer condition excludes
all the negative ghost classes in
Claim~\ref{clm:FFghosts}.
Thus it is a specified field on neither earlier list with the
smallest possible cardinality and transcendence degree
in characteristic zero. We claim no least field under
inclusion. Claim~\ref{clm:FFstrictreconstruct} subsequently
proves its reconstruction by a different argument.
\end{claim}
\begin{proof}
The strict henselization is a strictly henselian discrete
valuation ring with maximal ideal generated by \(p\)
and residue field \(\overline{\mathbb F}_p\)
\citep[Lemmas~10.155.2, 15.46.11]{FF_StacksHensel}.
It is a filtered colimit of pointed \(\acute etale\)
\(\mathbb Z_{(p)}\)-algebras. Their generic fields are
finite separable extensions of \(\mathbb Q\), so its
fraction field is algebraic over \(\mathbb Q\), and
therefore countable. The generator \(p\) gives the
specified normalization of its value group. The roots of
unity of every order prime to \(p\) lie in \(K_p\),
because their residue roots lift uniquely from finite
extensions of \(\mathbb F_p\). Their union is a
nontrivial divisible subgroup of \(K_p^\times\).
Sub-\(q\)-adic fields are Kummer-faithful
\citep[Remark~3.9(2)]{FF_CS2026}; hence \(K_p\) is not
sub-\(q\)-adic for any \(q\).

Claim~\ref{clm:FFmixedcentralizer} applies to this mixed
characteristic discrete valuation and to every finite
extension, proving candidate (a)'s condition.
The field embeds in
\(\operatorname{Frac}W(\overline{\mathbb F}_p)\),
the completion of the maximal unramified extension of
\(\mathbb Q_p\): the strict henselization universal
property gives a local map into the strictly henselian
Witt ring, and its kernel is zero because \(p\) remains
nonzero. Thus it is generalized sub-\(p\)-adic
in \citet[Definition~4.11]{FF_Mochizuki2003}.
Every finite extension remains a subfield of a finite
extension of that complete field, hence is still generalized
sub-\(p\)-adic. Theorem~4.12 of that source, with all
geometric primes retained, proves the relative hyperbolic
Isom assertion, modulo geometric inner conjugation.

We prove the two-valuation fact needed for the normal
closure, rather than use it without its comparison.
A field with two distinct rank-one henselian valuations
is separably closed. Indeed the valuation rings are
incomparable. There is an \(x\) with
\(v_1(x)>0>v_2(x)\): choose an element integral for the
first valuation and negative for the second, and if its
first value is zero multiply a sufficiently large power
by an element of positive first value. Put
\(e_n=x^n/(1+x^n)\). Then \(e_n\to0\) in \(v_1\)
and \(e_n\to1\) in \(v_2\). Thus
\(a+(b-a)e_n\) approximates any prescribed pair \(a,b\)
simultaneously; this proves the required two-valuation
approximation coefficient by coefficient.

If there were a monic irreducible separable polynomial
\(f\) of degree \(d>1\), irreducibility would hold on
a sufficiently small coefficient neighborhood for \(v_1\).
Here are the details of this often-used assertion. Choose
a root \(\alpha\), and choose \(\delta\) in the base
field of value larger than all the finitely many values
\(v_1(\alpha-\alpha')\), with \(\alpha'\ne\alpha\)
a conjugate, and sufficiently large to control the higher
Taylor coefficients. For a sufficiently close polynomial
\(g\), the polynomial
\[
 \frac{g(\alpha+\delta Y)}{\delta g'(\alpha)}
\]
has integral coefficients, constant term in the maximal
ideal, linear coefficient one and all higher coefficients
in the maximal ideal. Hensel's lemma in the finite
henselian extension generated by \(\alpha\) supplies
a root \(\beta\) of \(g\) with
\(v_1(\beta-\alpha)>v_1(\delta)\).
If an automorphism over the base fixes \(\beta\),
uniqueness of the extended henselian valuation implies
that it cannot move \(\alpha\): otherwise both terms
in \(\alpha-\alpha'=(\alpha-\beta)+(\beta-\alpha')\)
would have value strictly greater than its left-hand side.
Thus the field generated by \(\alpha\) is contained
in that generated by \(\beta\). Their degrees force
\(g\) to be irreducible. This proves the asserted open
irreducibility neighborhood.

On the other hand, choose a monic polynomial with \(d\)
distinct roots in the base field. A sufficiently small
coefficient neighborhood for \(v_2\) consists of split
polynomials, by Hensel's lemma at the distinct simple
roots, with disjoint root neighborhoods. Simultaneous
approximation produces a polynomial in both neighborhoods,
an impossibility. This proves the two-valuation fact.

Now fix a number field \(D\subset K_p\), let
\(\widetilde v\) be the chosen valuation on
\(\overline{\mathbb Q}\), and write
\(D_{\widetilde v}\subset G_D\) for its decomposition
subgroup. It is a proper subgroup. To see this concretely,
choose a rational prime \(q\ne p\) unramified in
\(D\), and an integer \(a\) satisfying
\[
 a\equiv1\pmod{p^3},\qquad a\equiv q\pmod{q^2}.
\]
Its \(q\)-valuations in \(D\) are one, so it is not
a square in \(D\); it is a square in \(\mathbb Q_p\)
by Hensel's lemma (also for \(p=2\)). Thus the chosen
local decomposition group fixes a root of this nontrivial
quadratic extension, and is proper in \(G_D\).

Let \(N/D\) be the normal closure of \(K_p/D\).
Since \(K_p\) is henselian, every element of \(G_{K_p}\)
preserves \(\widetilde v\); consequently
\[
 G_N\subset G_{K_p}\subset D_{\widetilde v}
                       \subsetneq G_D.
\]
The normality of \(G_N\) makes the distinct cosets of
\(D_{\widetilde v}\) give distinct conjugate valuations
on \(N\). Each is henselian, since \(N\) is algebraic
over the corresponding conjugate of \(K_p\). They have
rank one because their value groups are subgroups of
\(\mathbb Q\). The two-valuation fact makes \(N\)
algebraically closed, proving \(N=\overline{\mathbb Q}\).

The same proof gives a slightly stronger exclusion needed
for Claim~\ref{clm:FFnormal}. Choose a conjugate
\(gK_p\) for which the prolongation differs. Their
compositum has the two distinct henselian valuations,
so is \(\overline{\mathbb Q}\), and therefore
\[
                   G_{K_p}\cap gG_{K_p}g^{-1}=1.          \tag{F.23.1}
\]
If \(K_p\) were finite over a normal algebraic extension
\(k/D\), both groups on the left would be open in the
normal group \(G_k\). Their intersection would be open,
forcing \(G_{K_p}\) to be finite. This is impossible:
the valuation of a root of \(X^n-p\) is \(1/n\),
so this polynomial is irreducible over \(K_p\) for every
\(n\), and its finite extension degrees are unbounded.
If \(K_p\) were finitely generated over such a \(k\),
it would be finite because it is algebraic over
\(\mathbb Q\). This excludes that enlarged positive
class too. The arithmetic formation classes are excluded
by the full normal closure and their arbitrarily large
symmetric-group obstruction; the geometric and function-field
classes require positive transcendence degree.

Finally, every negative field family in
Claim~\ref{clm:FFghosts} has a nontrivial centralizer over
it or a finite extension, whereas \(K_p\) satisfies
the opposite condition. It therefore lies on neither
of those earlier lists. Countability and transcendence
degree zero are the least possible numerical values
for an infinite characteristic-zero field; this says
nothing about a least field under inclusion. The failure
of Kummer-faithfulness also explains the exact point at
which Carlson--Stix Proposition~3.18 cannot be used to
upgrade the known Isom assertion to full topos reconstruction.
It is a failure of that attempted proof, not a counterexample
to reconstruction or to centralizer sufficiency. The
\(p\)-primary argument below proves reconstruction
without Kummer-faithfulness.
\end{proof}

\subsection{What point detection retains over the strict henselian test}

\begin{claim}{F.24}{Proved}\label{clm:FFstrictdetection}
Fix a prime \(p\), and let \(K_p\) be the field of
Claim~\ref{clm:FFstricttest}.
\begin{enumerate}
\item For every finite extension \(L/K_p\),
\[
 \bigcap_{n\geq1}(L^\times)^{p^n}=\mu_{p'}\subset K_p^\times,
                                                        \tag{F.24.1}
\]
where \(\mu_{p'}\) denotes all roots of unity of order
prime to \(p\).
\item There is a connected nontrivial finite \'etale Galois
cover of an elliptic curve over \(K_p\) whose pullback
at every closed point is completely split. Thus the
finite-quotient detection statement used in
Claim~\ref{clm:FFcore} is false over \(K_p\).
\item If \(X/K_p\) is of finite type and a pinned topos
morphism \(X_{\mathrm{\acute et}}\to
\mathbb G_{m,K_p,\mathrm{\acute et}}\) has constant
underlying closed-point function with value
\(c\in\mathbb G_m(K_p)\), it factors through the
closed subtopos \(c_{\mathrm{\acute et}}\).
For an integral finite-type curve \(X\), every
nonconstant underlying continuous point function has
finite fibres over closed target points.
\item Every field embedding \(K_p\hookrightarrow K_p\)
is an automorphism.
\end{enumerate}
Consequently neither the invisible prime-to-\(p\)
quotients in (2), a residue-tube point function, nor
a proper self-embedding supplies a reconstruction
counterexample.
\end{claim}

\begin{proof}
The unique valuation on a finite extension \(L\) is
henselian, discrete, and has residue field
\(\overline{\mathbb F}_p\). Normalize it to have value
group \(\mathbb Z\), and write \(e=v(p)>0\).
Hensel's lemma identifies \(\mu_{p'}\) with the
nonzero residue elements; these roots already lie
in \(K_p\). Suppose \(a\) belongs to the intersection
in (F.24.1). Its valuation is divisible by every
\(p^n\), hence is zero. Let \(\zeta\in\mu_{p'}\)
be its residue lift. For each \(n\), choose
\(a_n\in L^\times\) with \(a_n^{p^n}=a\), and
let \(\zeta_n\) be its prime-to-\(p\) residue lift.
Uniqueness gives \(\zeta_n^{p^n}=\zeta\).
Thus \(b=a/\zeta=(a_n/\zeta_n)^{p^n}\), with
\(a_n/\zeta_n\in1+\mathfrak m_L\).
For \(s=v(u)>0\), the binomial formula gives
\[
 v((1+u)^p-1)\geq\min\{ps,e+s\}\geq s+1.
\]
Iteration yields \(v(b-1)\geq n+1\) for every
\(n\), so \(b=1\). Conversely, raising to \(p\)
is an automorphism of \(\mu_{p'}\), proving (1).

Choose an elliptic scheme \(\mathcal E\) over
\(\mathbb Z_{(p)}\): for odd \(p\) one may use
\(y^2=x^3-x\), whose discriminant is \(64\),
and for \(p=2\) use \(y^2+y=x^3\), whose
discriminant is \(-27\). Take \(\ell=2\) in
the first case and \(\ell=3\) in the second.
Over the strict henselization, \(\mathcal E[\ell]\)
is the constant group \((\mathbb Z/\ell)^2\).
On the generic fibre \(E/K_p\),
\([\ell]:E\to E\) is therefore a connected
finite \'etale Galois cover with that nontrivial
deck group. For a closed point with residue field
\(L\), properness extends the point to
\(\mathcal E(\mathcal O_L)\). Pulling back
\([\ell]\) gives a finite \'etale scheme over
the strictly henselian ring \(\mathcal O_L\),
so the pullback is a disjoint union of sections.
Using the origin to split the arithmetic fundamental
group, the resulting quotient
\[
 \pi_1(E)\longrightarrow (\mathbb Z/\ell)^2
\]
is nonzero and kills every point-section image
\(G_L\to\pi_1(E)\), including after every finite
base extension. This proves the precise failure in (2).

For (3), the inverse image of
\(\mathbb G_m\setminus\{c\}\) is an open subset
of \(X\) with no closed points. A finite-type scheme
over a field is Jacobson, so this open is empty.
The open--closed decomposition of the \'etale topos
then factors the morphism through the complementary
closed subtopos. If \(X\) is an integral curve,
the inverse image of any closed target point is
a closed subset of \(X\). Unless it is all of
\(X\), it is zero-dimensional and finite. This
also explains why a point function that collapses
an infinite residue tube cannot be the underlying
function of a nonconstant topos morphism.

Finally let \(\iota:K_p\hookrightarrow K_p\).
Since \(K_p/\mathbb Q\) is algebraic, extend
\(\iota\) to \(g\in G_{\mathbb Q}\), so
\(gK_p\subset K_p\). Transport the chosen
henselian valuation to \(gK_p\), and extend it
to the algebraic extension \(K_p/gK_p\).
It is again henselian of rank one. If it differs
from the original valuation of \(K_p\), the
two-valuation argument in
Claim~\ref{clm:FFstricttest} makes \(K_p\)
algebraically closed, a contradiction. They must
therefore coincide. Their unique prolongations
to \(\overline{\mathbb Q}\) coincide as well,
so \(g\) lies in the decomposition group of
the chosen \(p\)-adic valuation. This group
normalizes its inertia subgroup, whose fixed
field is exactly \(K_p\). Hence \(gK_p=K_p\),
and \(\iota\) is surjective.
\end{proof}

\begin{claim}{F.25}{Proved}\label{clm:FFpdetector}
Let \(K/K_p\) be finite, let \(A/K\) be a semiabelian
variety, and put \(V_pA=T_pA\otimes_{\mathbb Z_p}\mathbb Q_p\).
Every nonzero continuous \(G_K\)-equivariant linear map
\[
                  \lambda:V_pA\longrightarrow\mathbb Q_p(1)
\]
is detected by a point: for some finite \(L/K\) and
\(a\in A(L)\),
\[
             \lambda_*\kappa_A(a)\ne0
                 \quad\hbox{in }H^1(G_L,\mathbb Q_p(1)). \tag{F.25.1}
\]
If \(j:X\to A\) is a morphism based at \(x_0\in X(K)\),
with \(j(x_0)=0\), from a geometrically
integral finite-type \(K\)-scheme and \(j(X_{\bar K})\)
generates \(A_{\bar K}\), the detecting points may be
chosen in \(j(X(L))\). Finitely many such points,
over one finite extension, jointly detect every member of
\(\operatorname{Hom}_{G_K}(V_pA,\mathbb Q_p(1))\).
\end{claim}

\begin{proof}
We first explain the comparisons with complete fields.
A finite extension of \(K_p\) is an excellent henselian
discretely valued field of characteristic zero. Completion
induces an equivalence on finite separable extensions.
Indeed, approximate the coefficients of a separable primitive
polynomial over the completion by coefficients in the field;
the Krasner comparison proved in
Claim~\ref{clm:FFstricttest} identifies the resulting completed
extensions. For full faithfulness, a henselian field of
characteristic zero is relatively algebraically closed in its
completion: approximate a putative algebraic element closely
enough for the same simple-root Hensel argument to put it in
the field. A map between completed finite extensions therefore
sends their algebraic elements back into the corresponding
finite extensions. This proves the equivalence and identifies
the absolute Galois groups. Smooth points over the original
field are dense in the points over the completion: use an
\'etale coordinate chart and Hensel's lemma, then the density
of the field in its completion.

Suppose first that \(A=B\) is abelian. After a finite extension
it has semistable reduction. The exact source is
\citet[IX, Theorem~3.6]{FF_SGA7I1972}: over a regular connected
noetherian scheme of dimension one, an abelian generic fibre
acquires semistable reduction after a finite Galois extension.
Apply it to the valuation ring, and then complete. Write
\(R=\mathcal O_{\widehat L}\), with residue field
\(\overline{\mathbb F}_p\). The finite part of the N\'eron
model supplies a \(p\)-divisible group \(G/R\) such that
\[
 0\longrightarrow V_pG\longrightarrow V_pB
       \longrightarrow W\longrightarrow0,             \tag{F.25.2}
\]
where \(W\) is unramified and hence trivial over this
strictly henselian base. Moreover, the connected part
\(G^0\) is the \(p\)-divisible group of the formal
completion of the connected N\'eron model at the identity.
These are precisely
IX (5.1.4)--(5.1.5), (5.5.8), and Proposition~5.6 of
\cite{FF_SGA7I1972}; Proposition~5.6 includes the residue
prime, not just primes different from it.

The \(p\)-adic cyclotomic character on \(G_{\widehat L}\)
has infinite image. For example, a primitive
\(p^n\)-th root of unity has
\(v(\zeta_{p^n}-1)=v(p)/(p^{n-1}(p-1))\), which is
incompatible with a fixed discrete value group as
\(n\to\infty\). Thus
\(\operatorname{Hom}_{G_{\widehat L}}(W,\mathbb Q_p(1))=0\),
and \(\lambda\) is nonzero on \(V_pG\).
Multiply \(\lambda\) by a power of \(p\) to make this
restriction integral. Tate's generic-fibre theorem
\citep[Theorem~4 and Corollary~1]{FF_Tate1967} extends it
uniquely to a morphism
\[
                       f:G\longrightarrow\mu_{p^\infty}
\]
over \(R\). The stated theorem applies to an integrally
closed noetherian domain with characteristic-zero fraction
field, so this complete DVR satisfies its hypotheses.
The restriction to \(G^0\) is nonzero. Otherwise \(f\)
would factor through the \'etale quotient, which is constant
over \(R\), whereas its generic Tate module has no nonzero
map to \(\mathbb Z_p(1)\).

The equivalence of connected \(p\)-divisible groups with
divisible formal Lie groups
\citep[Section~2.2, Proposition~1]{FF_Tate1967} therefore
gives a nonzero formal group homomorphism
\(\widehat f:\widehat B\to\widehat{\mathbb G}_m\).
In characteristic zero, formal logarithms turn a formal
group homomorphism into its linear differential. Hence
\(d\widehat f\ne0\). On a sufficiently small ball,
\(\log\widehat f\) is a nonzero continuous linear form
in logarithmic coordinates. There is a point of this ball
with nonzero image, and density supplies such a point
\(a\in B(L)\). Its image \(u=\widehat f(a)\) is a
principal unit in \(\widehat L\) with \(\log u\ne0\).
Its rational \(p\)-Kummer class is nonzero: if a power
\(u^{p^s}\) were infinitely \(p\)-divisible, the binomial
valuation proof of (F.24.1), which applies equally to
the complete DVR, would give \(u^{p^s}=1\), contradicting
its logarithm.

For clarity, the Kummer comparison here is a comparison of
actual cocycles. Choose a compatible sequence of division
points of \(a\) in the formal group over the algebraic
closure. These points exist because multiplication by
\(p\) on a divisible formal group is an isogeny. Their
successive Galois differences give the Kummer cocycle
in \(T_pG\), whose image in \(T_pB\) is the usual
point-Kummer cocycle. Applying \(\widehat f\) gives
the compatible roots of \(u\), hence its multiplicative
Kummer cocycle. Thus the nonzero class just obtained is
\(p^s\lambda_*\kappa_B(a)\) for the fixed scaling power.
The completion comparison proves (F.25.1) over \(L\).

For a general semiabelian variety, pass to a finite
extension splitting its torus in
\(1\to T\to A\to B\to0\). If \(\lambda|_{V_pT}\ne0\),
one coordinate of \(T\simeq\mathbb G_m^r\), evaluated
at a sufficiently small nontrivial principal unit, detects
it by (F.24.1). If this restriction is zero, \(\lambda\)
factors through a nonzero functional of \(V_pB\).
The abelian case gives a detecting \(b\in B(L)\).
Hilbert's theorem 90 for the split torus gives
\(A(L)\twoheadrightarrow B(L)\), so a lift detects
\(\lambda\) by functoriality of Kummer cocycles.

Now assume that \(j(X_{\bar K})\) generates \(A_{\bar K}\).
Every point of \(A(\bar K)\) is a finite sum of signed
points of \(j(X(\bar K))\). To justify this, the dimensions
of the images of successive signed sum maps eventually
stabilize; their closure is the algebraic subgroup
generated by the image, hence is \(A\). A dominant sum
map has constructible image containing a nonempty open
\(U\). Since \(U\cap(a+U)\ne\varnothing\), every
\(a\) is the difference of two such sums. All finitely
many points involved are defined over a finite extension.
If every \(j(x)\) annihilated \(\lambda\), additivity
of the Kummer map would therefore annihilate every
point, contrary to the preceding proof.

Finally the space of functionals is finite-dimensional
over \(\mathbb Q_p\). Its intersections with the
kernels of point evaluations decrease to zero, so a
finite subcollection already has zero intersection.
Put these points over one finite compositum. Restriction
on \(H^1(-,\mathbb Q_p(1))\) to a subgroup of finite
index is injective, because restriction followed by
corestriction is multiplication by that index. Thus
the detecting property persists over this compositum.
\end{proof}

\subsection{Reconstruction over strictly henselian algebraic fields}

\begin{claim}{F.26}{Proved}\label{clm:FFstrictreconstruct}
For every prime \(p\) and every finite extension
\(K/\operatorname{Frac}(\mathbb Z_{(p)}^{sh})\),
the field \(K\) satisfies \(\mathcal R(K)\).
Consequently every subfield of such a \(K\) reconstructs.
\end{claim}

\begin{proof}
Throughout, all topoi and their structural isomorphisms are
relative to the stated field. Finite scalar extension of a
pinned map is formed by the corresponding finite \'etale
slice. We may make finite extensions during the proof:
if a unit on \(X_L\) realizes the original point function
over a finite extension \(L/K\), pass to the finite normal
closure of \(L/K\). Every conjugate unit has the same
geometric point function. Equality on geometric points of
the reduced scheme makes the units equal, so the unit
descends to \(K\). This reduction does not assume
reconstruction over any of the fields in advance.

Let \(X/K\) be regular, quasiprojective and connected, with
\(x_0\in X(K)\), and let
\[
 \varphi:X_{\mathrm{\acute et}}\longrightarrow
                    \mathbb G_{m,K,\mathrm{\acute et}}
\]
be pinned. Regularity and the rational point make \(X\)
geometrically integral; in dimension zero it is \(\operatorname{Spec}K\),
where Claim~\ref{clm:F.1} applies directly. We henceforth
assume positive dimension. Its rational points are Zariski
 dense. Indeed, an \'etale coordinate chart at \(x_0\)
and Hensel's lemma lift every point of a sufficiently
small valuation polydisc in \(K^n\) to \(X(K)\).
The image of a proper closed subset under the quasi-finite
coordinate map has dimension less than \(n\). A product
of infinite valuation balls cannot lie in a proper
algebraic subset of \(\mathbb A^n_K\), by induction
on the variables of a nonzero polynomial. Thus every
nonempty open of \(X\) has a rational point. This
also holds after each finite extension and permits the
same argument on all inverse-image opens below. Write \(b(x)\) for its actual
value at a finite-extension point, and \(b_0=b(x_0)\).
Its pulled-back \(p\)-Kummer class is
\[
                   \xi=\varphi^*\kappa_p(T)
                          \in H^1(X,\mathbb Z_p(1)).
\]
In particular, evaluating \(\xi\) at \(x\in X(L)\)
gives \(\kappa_p(b(x))\), for every finite \(L/K\).

\emph{1. Integral boundary residues.}
We first prove a calculation requiring only a compatible
\(p\)-Kummer class with actual values at all finite-extension
points, rather than an entire topos map. On a punctured
rational line, after splitting its punctures over a finite
extension, normalization at a rational point writes the
class as
\[
 \kappa_p(c)+\sum_{i=1}^r a_i
       \kappa_p\left(\frac{T-r_i}{t_0-r_i}\right),
                  \qquad a_i\in\mathbb Z_p.             \tag{F.26.1}
\]
This follows from the Kummer sequence and the vanishing of
the Picard group of a localization of \(K[T]\).
Let \(\pi\) be a uniformizer of this finite extension,
and evaluate at \(r_i+\pi^N\) and \(r_i+\pi^{N+1}\),
where \(N\) is larger than the valuations of all nonzero
root differences. Applying the valuation map on
\(H^1(K,\mathbb Z_p(1))\), the difference between these
two evaluations is \(a_i\): the other factors have
constant valuation. Both actual values have ordinary
integer valuation. Thus \(a_i\in\mathbb Z\). The residue
at infinity, the negative sum of the finite residues,
is integral as well.

For a smooth curve and a chosen rational boundary point
\(q\), Riemann--Roch gives a function whose only pole
is \(q\). Remove its branch values and the images of
all boundary points. The resulting map to a punctured
line is finite \'etale. Corestriction of the class has
actual values: at each point of the line these are the
products of the field norms of the actual values in its
fibre. The residue of a norm is the sum of the residue
degrees times the source residues. It has no factor
equal to the ramification index. Since \(q\) is the
only point above infinity, that residue is exactly
the residue at \(q\). Hence it is integral by (F.26.1).
This argument is valid after making a boundary point
rational by a finite extension.

Choose a smooth projective compactification \(\bar X\)
with normal-crossing boundary, available in characteristic
zero. At a smooth point of a chosen boundary component,
away from the other components, take \'etale coordinates
in which that component is the first coordinate divisor.
Fix the other coordinates and take the curve component
through the point. Its local intersection multiplicity
with the chosen divisor is one. Normalize its closure
and apply the preceding curve argument. Thus each residue
\(m_i\) of \(\xi\) at a geometric boundary component
\(D_i\) belongs to \(\mathbb Z\). After another finite
extension all components and their residue data are defined
over the base. Put \(D=\sum_i m_iD_i\).

\emph{2. Forcing arithmetic equivariance over a finite local field.}
Let \(j:X\to A\) be the generalized Albanese map based
at \(x_0\), so \(j(x_0)=0\). The geometric part of
\(\xi-\kappa_p(b_0)\) is a continuous integral functional
\[
                  \lambda:T_pA\longrightarrow\mathbb Z_p(1)
\]
equivariant for \(H_0=G_K\). The Albanese comparison and
its finite obstruction are exactly those in
\citet[Proposition~4.1 and Corollary~4.2]{FF_SpiessSzamuely2003};
the finite obstruction disappears on taking the inverse
limit, since the Tate module of a finite group is zero.
The map \(j(X_{\bar K})\) generates \(A_{\bar K}\).

Choose a number field \(D_0\subset K\) containing all
the finite displayed geometric data and \(b_0\), and
such that \(D_0K_p=K\). Such a field exists because
\(K/K_p\) is finite. The chosen prolongation of the
\(p\)-adic valuation gives a finite extension
\(M_0=(D_0)_v\) of \(\mathbb Q_p\). The field \(K\)
is the fraction field of the corresponding strict
henselization of \(D_0\): each of these two strict
henselian fields contains the other by the universal
property. The completion/Krasner comparison therefore
identifies
\[
                  H_0=G_K=I_{M_0}\triangleleft G_{M_0}.
                                                               \tag{F.26.2}
\]
These are compatible identifications for the Tate modules
of the varieties defined over \(D_0\); their torsion
points and all relevant point-Kummer cocycles are algebraic
over the number field.

Apply Claim~\ref{clm:FFpdetector} to the finite-dimensional
space
\[
                  S=\operatorname{Hom}_{H_0}
                       (V_pA,\mathbb Q_p(1)).
\]
Choose finitely many points \(x_i\), over one finite
extension \(E/K\), whose Kummer evaluations jointly
detect \(S\). Include their actual values \(b(x_i)\)
in the finite data, and choose a number field
\(D_1\subset E\) containing \(D_0\), all these data,
and a generator for \(E/K_p\). Then
\(E=D_1K_p\), and with \(M_1=(D_1)_v\) we have
\[
        H_1=G_E=I_{M_1}\triangleleft G_{M_1},
                  \qquad G_{M_1}\subset G_{M_0}.         \tag{F.26.3}
\]
Restriction to \(H_1\) preserves joint detection,
by restriction--corestriction over \(\mathbb Q_p\).

For \(g\in G_{M_1}\), let
\(g\lambda=\chi_p(g)\lambda\rho_A(g)^{-1}\).
The normality in (F.26.2) shows that
\(\delta_g=g\lambda-\lambda\) still belongs to \(S\).
Put
\[
 c_i=\kappa_A(j(x_i))\in H^1(G_{M_1},V_pA),\qquad
 d_i=\kappa_p(b(x_i)/b_0)\in H^1(G_{M_1},\mathbb Q_p(1)).
\]
The actual point-value equality is
\[
                \lambda_*\operatorname{res}_{H_1}(c_i)
                        =\operatorname{res}_{H_1}(d_i). \tag{F.26.4}
\]
An \(H^1\)-class restricted from a group to a normal
subgroup is invariant under that group's conjugation
action, including its action on the coefficient module.
Apply this elementary cocycle identity to both \(c_i\)
and \(d_i\) in (F.26.4). It gives
\[
               (\delta_g)_*\operatorname{res}_{H_1}(c_i)=0
                                      \quad\hbox{for all }i.
\]
Joint detection proves \(\delta_g=0\). Thus \(\lambda\)
is \(G_{M_1}\)-equivariant. This argument uses cohomology
classes, so the conjugation coboundaries cause no problem;
it does not require a trivial cyclotomic character and
does not assert that arbitrary topos maps descend to
\(M_1\).

\emph{3. A power of the class is the class of a unit.}
We give the integral Picard calculation to avoid an
incorrect local Tate-Hom assertion for general 1-motives.
Let \(P=\operatorname{Pic}^0_{\bar X}\), and let
\(B_X\subset\bigoplus_i\mathbb ZD_i\) be the group
of boundary divisors algebraically equivalent to zero.
The Picard 1-motive is \([B_X\to P]\), dual to
\([0\to A]\). Its realization gives
\[
 0\longrightarrow T_pP\longrightarrow
 H^1(X_{\bar K},\mathbb Z_p(1))\longrightarrow
 B_X\otimes\mathbb Z_p\longrightarrow0.                \tag{F.26.5}
\]
Here the last map sends the geometric class to its
boundary residues. The construction, including the
finite obstruction before inverse limit, is given in
the proof of \citet[Proposition~4.1]{FF_SpiessSzamuely2003}.
The field has characteristic zero, so its prime-to-characteristic
statement includes the present prime \(p\).

By Step~1, the residue vector \(D\) is integral, while
(F.26.5) puts it in \(B_X\otimes\mathbb Z_p\).
Consequently its class in the finitely generated group
\((\bigoplus_i\mathbb ZD_i)/B_X\) is torsion of order
prime to \(p\). Choose an integer \(N>0\), prime
to \(p\), such that \(ND\in B_X\).
Step~2 says that the geometric class is
\(G_{M_1}\)-invariant. The existence of its invariant
multiple lifting \(ND\) in (F.26.5) means
\[
                 [ND]\in\bigcap_n p^nP(M_1).           \tag{F.26.6}
\]
One can check this directly in the 1-motive, without
any Hom theorem. At level \(p^n\), a lift of \(ND\)
is represented by a \(p^n\)-division point of
\([ND]\in P\). The fixed residue vector allows us
to choose its lattice component equal to \(ND\).
Galois invariance then makes that division point
\(M_1\)-rational: a change of representative would
change the lattice component by \(p^n b\), and
\(p^n b=0\) implies \(b=0\) in the torsion-free
lattice. This holds for every \(n\).

For an abelian variety over a finite extension \(M_1\)
of \(\mathbb Q_p\), the compact group \(P(M_1)\)
contains an open subgroup isomorphic, under its formal
logarithm, to a finite free \(\mathbb Z_p\)-module.
Its infinitely \(p\)-divisible subgroup is therefore
exactly its finite prime-to-\(p\) torsion subgroup.
For completeness, modulo the open subgroup there is
only a finite group. A fixed power of \(p\) kills
its \(p\)-primary part; after multiplying by its
prime-to-\(p\) exponent, an infinitely \(p\)-divisible
element lies in every \(p^n\) multiple of the free
\(\mathbb Z_p\)-module and is zero. The converse for
prime-to-\(p\) torsion is immediate.
Thus (F.26.6) is killed by an integer \(d_0>0\)
prime to \(p\), and \(dD\) is principal, where
\(d=Nd_0\).

All divisor and Picard data were defined over the
number field \(D_0\). Their equality after extension
to \(M_1\) implies the equality over \(D_0\), since
extension of scalars is injective on points of a scheme.
The principal divisor has a rational function over
\(D_0\), uniquely after fixing its value at \(x_0\):
the one-dimensional space of functions with that divisor
descends, or equivalently one applies Hilbert 90 to its
scalar ambiguity. Choose
\(h\in\mathcal O(X)^\times\) with
\(\operatorname{div}_{\bar X}(h)=dD\) and
\(h(x_0)=b_0^d\).

The difference between \(d\lambda\) and the geometric
Kummer class of \(h\) has zero residues. It belongs,
by (F.26.5), to \((T_pP)^{G_{M_1}}\), which is zero:
the torsion subgroup of \(P(M_1)\) is finite, by
the same compact Lie-group argument. Hochschild--Serre
and evaluation at \(x_0\) then give the integral identity
\[
                         d\xi=\kappa_p(h).              \tag{F.26.7}
\]
In particular, for every point over a finite extension,
\[
                         b(x)^d=\zeta_xh(x),
                         \qquad\zeta_x\in\mu_{p'}\subset K_p.
                                                               \tag{F.26.8}
\]
This uses the integral kernel (F.24.1), not merely
the kernel after tensoring cohomology with \(\mathbb Q_p\).

\emph{4. Curves with large-degree points and a resultant.}
We record explicitly two elementary tools. Every nonempty
open of a geometrically integral curve over a finite
extension of \(K_p\) has closed points of unbounded
degree. Choose a nonconstant map from its proper
normalization to \(\mathbb P^1\). The polynomials
\(T^n-\pi\) give closed points of degree \(n\) on
the line, by Eisenstein's criterion. For arbitrarily
large \(n\) their support avoids the finite image
of the omitted points; any point in the fibre has
degree at least \(n\).

Second, if \(V\subset K(X)\) is a finite-dimensional
space of rational functions regular on a nonempty open
\(U\), after a finite extension there is a smooth
geometrically integral curve mapping into \(U\) on
which restriction is injective on \(V\). Choose finitely
many geometric points of \(U\) whose evaluations are
injective on \(V\); this is possible by successively
testing a nonzero function in the remaining kernel.
Put them over a finite extension. A geometrically
irreducible curve through them exists by
\citet[Theorem~1.8, Corollary~1.9]{FF_CharlesPoonen2016},
whose hypotheses allow an arbitrary base field and
a finite set of closed points. Normalize its reduced
structure. In characteristic zero the result is a
smooth curve. Evaluation at its points above the
chosen set proves the asserted injectivity.

Suppose now that the geometric class \(\lambda\)
is nonzero. Then \(h\) in (F.26.7) is nonconstant.
The actual point function is nonconstant as well:
if all values were equal to a base-field constant,
(F.26.8) would put all values of \(h\) in the
base field. A nonconstant regular function cannot
have this property, by restricting to a curve and
using closed points of arbitrarily large degree
on its image.
Thus \(U=\varphi^{-1}(\mathbb G_m\setminus\{1\})\)
is a nonempty open of \(X\). Apply Steps~1--3 on
\(U\) to the pullback of \(T-1\), using a
rational point of that open, which exists by the preceding
density argument.
They give a positive integer \(e\), a rational
function \(g\) that is a unit on \(U\), and
\[
                     (b(x)-1)^e=\eta_xg(x),
                           \qquad\eta_x\in\mu_{p'}.
                                                               \tag{F.26.9}
\]
Consider the finite-dimensional space generated by
\(h^ig^j\), for \(0\leq i\leq e\) and
\(0\leq j\leq d\). Choose a curve as in the preceding
paragraph, so restriction to that curve is injective
on this space. Write
\[
              B=e\deg\operatorname{div}_\infty(h|_C)
                   +d\deg\operatorname{div}_\infty(g|_C).
\]
Choose a closed point \(x\in C\) of degree greater
than \(B\). The rational function
\[
 R=\operatorname{Res}_Z
       \bigl(Z^d-\zeta_xh,\ (Z-1)^e-\eta_xg\bigr)
                                                               \tag{F.26.10}
\]
vanishes at \(x\), by (F.26.8)--(F.26.9), and
its pole divisor on \(C\) has degree at most \(B\).
It must vanish identically on \(C\). All its terms
belong to the chosen finite-dimensional space, with
coefficients in the base field: here it matters that
\(\zeta_x,\eta_x\in K_p\). Injectivity of restriction
therefore gives \(R=0\) in \(K(X)\).

There is a common root \(z\) in an algebraic closure
of \(K(X)\), satisfying
\(z^d=\zeta_xh\) and \((z-1)^e=\eta_xg\).
It is nonconstant because \(h\) is nonconstant.
Every conjugate \(z'\) has the form
\(z'=\alpha z\), \(z'-1=\beta(z-1)\), for
roots of unity \(\alpha,\beta\). Consequently
\((\alpha-\beta)z=1-\beta\). As \(z\) is
transcendental over the algebraic closure of the
constant field, this forces \(\alpha=\beta=1\).
Characteristic zero makes the algebraic extension
separable, so \(z\in K(X)\). Its power is a unit
on the normal scheme \(X\), hence \(z\) is a
unit on \(X\).

The geometric group \(H^1(X_{\bar K},\mathbb Z_p(1))\)
is torsion-free. From \(d\lambda=d\kappa_p(z)\)
we obtain \(\lambda=\kappa_p(z)\) geometrically.
Normalize \(z\) to have value \(b_0\) at \(x_0\).
Hochschild--Serre then yields a unit, now denoted \(h\),
with
\[
                             \xi=\kappa_p(h).            \tag{F.26.11}
\]
This conclusion is proved whenever the geometric class
is nonzero, without assuming in advance that the
class of \(T-1\) is nonzero. The normalized unit
also descends through every finite extension made in
this construction, before its actual point function
has been identified. For a Galois conjugate, the ratio
of the two units has zero geometric \(p\)-Kummer class,
so every coefficient of its divisor on \(\bar X\)
is divisible by every \(p^n\), and hence zero.
Properness makes this ratio constant; its value at
\(x_0\) is one. Galois descent and then evaluation
at \(x_0\) in Hochschild--Serre establish
(F.26.11) over the original base field.

\emph{5. The zero geometric class forces a constant point function.}
If \(\lambda=0\), normalization gives
\(\xi=\kappa_p(b_0)\). By (F.24.1), every
actual value is in \(b_0\mu_{p'}\subset K\),
including values at points over arbitrary finite
extensions. Choose a proper finite extension
\(K'/K\) and \(\beta\in K'\setminus K\).
Since no value equals \(\beta\), the pullback of
\(T-\beta\) is defined on all of \(X_{K'}\).
If its geometric class were nonzero, Step~4 would
give a nonconstant unit \(g\) with
\(b(x)-\beta=\eta_xg(x)\),
\(\eta_x\in\mu_{p'}\subset K\). All values of
\(g\) would then lie in \(K'\), contradicting
nonconstancy by the large-degree-point argument.
Thus this geometric class is zero too. Normalization
at \(x_0\) gives
\[
              \frac{b(x)}{b_0}\in\mu_{p'},\qquad
              \frac{b(x)-\beta}{b_0-\beta}\in\mu_{p'}.
\]
The two scalar coefficients belong to \(K\),
and \(1,\beta\) are linearly independent over
\(K\). These equations force \(b(x)=b_0\) at
every geometric point.

\emph{6. Removing the remaining pointwise scalars.}
We now have a unit \(h\), normalized by
\(h(x_0)=b_0\), realizing \(\xi\) whenever
the point function is nonconstant, and a constant
function otherwise. On a smooth curve with
nonconstant \(h\), (F.24.1) gives
\(b(x)=c_xh(x)\), \(c_x\in\mu_{p'}\subset K\).
The point function of \(T-1\) on its inverse-image
open is nonconstant: otherwise the values of the
nonconstant unit \(h\) on that open would all be
in the base field, since \(c_x\in\mu_{p'}\subset K\).
The same argument applies to every \(T-a\). Thus its
class is nonzero by
Step~5. It has a unit representative \(h_1\)
and \(b(x)-1=d_xh_1(x)\), with \(d_x\in\mu_{p'}\).
Choose a point of degree greater than both a common
pole-degree bound for \(c_xh-d_xh_1-1\) and the
degree of \(h:\bar C\to\mathbb P^1\), so that
\(h(x)\notin K\). The
same divisor argument gives an identity
\[
                           h_1=Ah+B,
                              \qquad A,B\in K^\times.
\]
At every point with \(h(x)\notin K\), comparison
of the coefficients of \(1,h(x)\) forces
\(c_x=c=-A/B\).

For each \(a\in K^\times\), repeat the argument
with \(T-a\). On points with \(h(x)\notin K\),
the known value \(b(x)=ch(x)\) shows that its
unit representative is proportional to \(ch-a\).
To justify the identity rather than merely pointwise
equality, choose a point of degree above both the common
pole bound and the degree of \(h\), hence with
\(h(x)\notin K\); it exists because
\(h\) is nonconstant. At a point with
\(h(x)\in K\), set \(a=ch(x)\). If
\(b(x)\ne a\), this point belongs to the open
where \(T-a\) pulls back to a unit, but that
unit is proportional to \(ch-a\) and vanishes
there, a contradiction. Hence \(b=ch\) at
every point. Normalization at \(x_0\) gives
\(c=1\).

For higher-dimensional \(X\), pass to a finite
extension and join \(x_0\) to an arbitrary geometric
point by the normalization of a geometrically
irreducible curve through them, using
\citep[Corollary~1.9]{FF_CharlesPoonen2016} again.
The two points lift after another finite extension.
The restricted class is the Kummer class of
\(h|_C\). If this is constant, Step~5 applies;
if nonconstant, the curve argument applies. Its
normalized unit agrees with \(h|_C\): a ratio
with zero \(p\)-Kummer class has every divisor
coefficient divisible by all \(p^n\), hence is
constant, and its value at the lifted \(x_0\)
is one. Thus \(b(x)=h(x)\) at every geometric
point of \(X\).

\emph{7. Reconstruction and descent.}
The finite-extension descent in the first paragraph
puts the resulting unit over the original field.
We have proved exactly the point-function condition
for every regular connected affine source with a
rational point in
\citet[Theorem~4.18(5)]{FF_CHW2024}.
That theorem gives \(\mathcal R(K)\) for all
finite-type schemes and all pinned relative topos
morphisms, with absolute weak normalization as
in the question. Finally \(K\) is algebraic over
each of its subfields, so
Claim~\ref{clm:FFdescent} gives the asserted
subfield consequence.
\end{proof}

\begin{claim}{F.27}{Proved}\label{clm:FFdiscretealgebraic}
Let \(F/\mathbb Q\) be algebraic. If \(F\) admits a
nontrivial discrete valuation of residue characteristic
\(p>0\), then \(\mathcal R(F)\) holds. More precisely,
for any prolongation to \(\overline{\mathbb Q}\), the
compositum \(FK_p\) is finite over the strict
henselization field \(K_p\) attached to that prolongation.
Thus the fields in this statement are exactly the
subfields of finite extensions of some \(K_p\).
\end{claim}
\begin{proof}
Embed \(F\) in \(\overline{\mathbb Q}\) and fix a
prolongation of its valuation. Normalize it by
\(v(p)=1\). Since \(v(F^\times)\) is discrete and
contains \(\mathbb Z\), it is \(e^{-1}\mathbb Z\)
for an integer \(e\geq1\). The residue field is
algebraic over \(\mathbb F_p\): every element is
algebraic over \(\mathbb Q\), and restriction to
the number field it generates has finite residue field.

The fraction field \(F^{sh}\) of the chosen strict
henselization of the valuation ring embeds in
\(\overline{\mathbb Q}\). Henselization is immediate,
and strict unramified extensions preserve the value
group. Equivalently, the valuation ring of \(F\)
is a DVR, its strict henselization is a DVR with the
same uniformizer, and its residue field is
\(\overline{\mathbb F}_p\); see
\citep[Corollary~3.3.26]{FF_ADH2017} and
\citep[Tags 04GP, 0AP3]{FF_StacksHensel}.
The chosen local map \(\mathbb Z_{(p)}\to
\mathcal O_{F^{sh}}\) extends to the chosen strict
henselization by its universal property. Consequently
\[
 K_p\subset FK_p\subset F^{sh},
       \qquad v((FK_p)^\times)\subset e^{-1}\mathbb Z.
                                                        \tag{F.27.1}
\]
If \(L/K_p\) is finite inside \(FK_p\), then
\[
 [L:K_p]=e(L/K_p)\leq e.                                \tag{F.27.2}
\]
Here the equality has no defect term: the
henselian-completion equivalence proved in
Claim~\ref{clm:FFpdetector} preserves the degree;
the resulting finite extension of the complete DVR
has algebraically closed residue field. Its integral
closure is finite free, and comparison of lengths
modulo the base uniformizer gives degree equal to
ramification index times residue degree, the latter
being one. The inequality follows from (F.27.1).

The compositum \(FK_p\) is the directed union of
these finite extensions. Choose one of maximal
degree, which exists because their degrees are
positive integers bounded by \(e\). Its compositum
with any other member still satisfies (F.27.2),
so maximality forces that other member to be
contained in it. Hence \(FK_p/K_p\) is finite.
Claim~\ref{clm:FFstrictreconstruct} proves
\(\mathcal R(FK_p)\), and the algebraic descent
of Claim~\ref{clm:FFdescent} proves
\(\mathcal R(F)\).

Conversely a subfield of a finite extension of
\(K_p\) inherits a discrete mixed-characteristic
valuation: its value group is a subgroup of a
cyclic group and is nonzero because it contains
\(v(p)>0\). This proves the final characterization.
\end{proof}

\subsection{Bounded wild ramification}

\begin{claim}{F.28}{Proved}\label{clm:FFtamereconstruct}
Fix the prolongation of the \(p\)-adic valuation used
to define \(K_p\), and set
\[
 T_p=K_p\bigl(p^{1/m}:(m,p)=1\bigr)
        \subset\overline{\mathbb Q},                    \tag{F.28.1}
\]
with compatible roots. Every finite extension \(K/T_p\),
and every subfield of such a \(K\), satisfies
\(\mathcal R\). In particular the positive conclusion
includes a henselian field with nondiscrete value group
\(\mathbb Z_{(p)}\).
\end{claim}
\begin{proof}
We give all modifications of Claim~\ref{clm:FFstrictreconstruct};
the scheme-theoretic criterion and its resultant argument
will be applied exactly as proved there.

\emph{1. Field structure and the exact multiplicative kernel.}
The field \(T_p\) is the maximal tame extension of
\(K_p\). Indeed a finite tame extension of a strictly
henselian DVR is generated by a root of its uniformizer
of order prime to the residue characteristic. One can
see this directly: if \(\pi\) is a uniformizer upstairs
and \(p=u\pi^m\), the unit \(u\) has an \(m\)-th
root by Hensel's lemma, so the extension contains an
\(m\)-th root of \(p\), and the degrees agree.
See also \citep[Lemmas~15.116.2 and 15.116.7, Tags
09EV and 0EXW]{FF_StacksRamification}.
Its Galois group over \(K_p\) is
\(\prod_{\ell\ne p}\mathbb Z_\ell(1)\).

Every finite \(K/T_p\) can be written \(L_0T_p\)
with \(L_0/K_p\) finite. Write
\([L_0:K_p]=e=p^a m\), with \((m,p)=1\).
The residue field is algebraically closed, so
\([L_0:K_p]=e(L_0/K_p)\). If
\(p=u\pi^e\) for a uniformizer \(\pi\) of \(L_0\),
extracting an \(m\)-th root of the unit \(u\)
shows that \(L_0\) contains an \(m\)-th root of
\(p\). Consequently \(L_0\cap T_p\) has degree
exactly \(m\): its degree is prime to \(p\) and
divides \(e\), and the displayed root gives the
opposite lower bound. The Galois degree formula gives
\[
 [K:T_p]=p^a,
 \quad v(K^\times)=p^{-a}\mathbb Z_{(p)},
 \quad\overline K=\overline{\mathbb F}_p.              \tag{F.28.2}
\]
For the value group, adjoining roots of \(p\) of
orders prime to \(p\) gives all of
\(e^{-1}\mathbb Z_{(p)}=p^{-a}\mathbb Z_{(p)}\),
while each finite step over \(L_0\) has degree and
ramification index prime to \(p\) and introduces
no further \(p\)-denominator.
Moreover \(K\) is the maximal tame extension of
\(L_0\). All its finite steps over \(L_0\) are tame.
Conversely its value group is divisible by every
integer prime to \(p\), and its residue field is
algebraically closed. Given such an integer \(n\),
any element \(x\in K^\times\) has an \(n\)-th root:
choose \(y\) with \(nv(y)=v(x)\) and extract an
\(n\)-th root of the unit \(x/y^n\) by Hensel.
This contains every tame uniformizer extension of
\(L_0\).

We also have the exact formula
\[
       \bigcap_{n\geq1}(K^\times)^{p^n}=\mu_{p'}.
                                                        \tag{F.28.3}
\]
For a given element \(b\) in the left side, enlarge
\(L_0\) inside \(K\) to contain \(b\). The extension
\(K/L_0\) remains Galois pro-prime-to-\(p\), being
obtained by base change from \(T_p/K_p\). If
\(c^{p^n}=b\) with \(c\in K\), put
\(r=[L_0(c):L_0]\), so \((r,p)=1\). The norm
identity gives \(N(c)^{p^n}=b^r\). Choose integers
\(u,v\) with \(ur+vp^n=1\); then
\(b=(N(c)^u b^v)^{p^n}\). Thus \(b\) belongs to
the \(p\)-divisible kernel of \(L_0^\times\),
which is \(\mu_{p'}\) by Claim~\ref{clm:FFstrictdetection}.
The reverse inclusion is immediate. This proof
avoids using a discrete valuation estimate on the
dense value group in (F.28.2).

\emph{2. Functional detection survives.}
Let \(A/K\) be semiabelian and let
\(0\ne\lambda\in\operatorname{Hom}_{G_K}
(V_pA,\mathbb Q_p(1))\). Descend \(A\) to a
finite extension \(L_0/K_p\) inside \(K\),
as in Step~1. The finite-dimensional vector space
\[
 S=\operatorname{Hom}_{G_K}(V_pA,\mathbb Q_p(1))
\]
is stable under \(G_{L_0}\), and its action factors
through the pro-prime-to-\(p\) group
\(G_{L_0}/G_K\). Its image in
\(\operatorname{GL}(S,\mathbb Q_p)\) is finite:
a compact linear group preserves a \(\mathbb Z_p\)-lattice;
the kernel of reduction modulo \(p\) (modulo \(4\)
if \(p=2\)) is pro-\(p\), whose intersection with
a pro-prime-to-\(p\) group is trivial. Hence after
a finite extension \(L_1/L_0\) inside \(K\),
\(\lambda\) is \(G_{L_1}\)-equivariant.
Claim~\ref{clm:FFpdetector} gives a finite extension
\(M/L_1\) and \(x\in A(M)\) with
\(\lambda_*\kappa_A(x)\ne0\).

This class stays nonzero over \(MK\). Indeed
\(G_{MK}\) is normal in \(G_M\) with
pro-prime-to-\(p\) quotient. Its invariant subspace
in \(\mathbb Q_p(1)\) is zero: a finite extension
of \(K\) still has bounded \(p\)-denominators in
its value group by (F.28.2), whereas the valuations
of \(\zeta_{p^n}-1\) have unbounded
\(p\)-denominators. Thus the cyclotomic character
has infinite image there. Inflation--restriction
therefore makes
\[
 H^1(G_M,\mathbb Q_p(1))\longrightarrow
 H^1(G_{MK},\mathbb Q_p(1))
\]
injective. Alternatively continuous cohomology of
a pro-prime-to-\(p\) group with a finite-dimensional
\(p\)-adic representation vanishes in degree one,
by its finite image and averaging. The signed-sum
and finite-dimensionality arguments of
Claim~\ref{clm:FFpdetector} now give the same finite
joint detector on a pointed generating variety.

\emph{3. Integral residues become rational with prime-to-\(p\)
denominators.}
Normalize the valuation of \(K\) to have image
\(\mathbb Z_{(p)}\), choosing \(\pi\in K\)
of value one. The punctured-line calculation in
Step~1 of Claim~\ref{clm:FFstrictreconstruct}
uses the rational points \(r_i+\pi^N\) and
\(r_i+\pi^{N+1}\). For sufficiently large \(N\),
all other factor valuations agree at these two
points. The difference of the valuations of their
actual values is now in \(\mathbb Z_{(p)}\),
while its image in the \(p\)-adic completion is
exactly the coefficient \(a_i\in\mathbb Z_p\).
Thus \(a_i\in\mathbb Z_{(p)}\).
The finite-map norm argument for curves and the
transverse-curve argument for higher dimensions
are unchanged; every boundary coefficient is
in \(\mathbb Z_{(p)}\). There are finitely many,
so multiplying the class and its boundary by
one integer \(N\) prime to \(p\) makes all
coefficients integral. The normalized actual
values of this multiplied class are \(b(x)^N\).

\emph{4. Arithmetic equivariance and the unit power.}
Descend all source data to a number field
\(D_0\subset K\) with \(D_0T_p=K\), and let
\(M_0\) be its completion at the chosen place.
Step~1 identifies \(K\) with the maximal tame
extension of the chosen henselization of \(D_0\), so
\[
            G_K=P_{M_0}\triangleleft G_{M_0},          \tag{F.28.4}
\]
where \(P\) denotes wild inertia. After selecting
the finite jointly detecting points and their
actual values, enlarge to a number field \(D_1\)
with \(D_1T_p=E\), as in Step~2 of
Claim~\ref{clm:FFstrictreconstruct}. Then
\(G_E=P_{M_1}\triangleleft G_{M_1}\), where
\(M_1=(D_1)_v\). The same conjugation identity
for the point and value cocycles, followed by
the detector in Step~2 above, makes the geometric
functional \(G_{M_1}\)-equivariant. No topos
map is asserted to descend to \(M_1\).

Apply the Picard 1-motive argument of Step~3 of
Claim~\ref{clm:FFstrictreconstruct} to the
multiplied boundary \(ND\). Its integral lattice
coordinate, fixed at every finite Tate level,
gives rational division points over the finite
\(p\)-adic field \(M_1\). The compact group
\(\operatorname{Pic}^0(M_1)\) has only prime-to-\(p\)
torsion in its infinitely \(p\)-divisible subgroup,
and its Tate invariants vanish. Consequently, for
some integer \(d>0\) prime to \(p\),
\[
              d\xi=\kappa_p(h),\qquad
              b(x)^d=\zeta_x h(x),\quad
              \zeta_x\in\mu_{p'}\subset K.            \tag{F.28.5}
\]
Here \(d\) includes the initial multiplier \(N\).
The last equality is integral and follows from
(F.28.3).

\emph{5. The remaining reconstruction steps.}
The only numerical input to Steps~4--6 of
Claim~\ref{clm:FFstrictreconstruct}, beyond
(F.28.3)--(F.28.5), is the existence of closed
points of arbitrarily large degree on a smooth
curve open. The polynomials
\(Z^{p^n}-\pi\) now supply base points of degree
\(p^n\): adjoining a root enlarges
\(\mathbb Z_{(p)}\) to
\(p^{-n}\mathbb Z_{(p)}\), of index \(p^n\),
so the polynomial has its full degree. A finite
map from the proper curve to \(\mathbb P^1\)
and avoidance of the finitely many omitted images
then supply the required points upstairs.
All the other inputs hold as before: \(K\) is
henselian of rank one, so smooth pointed varieties
have Zariski-dense rational points; it contains
\(\mu_{p'}\); every finite extension has the same
properties by (F.28.2); and Charles--Poonen's
curve-through-points theorem applies in
characteristic zero. Thus the two-function
resultant kills the power \(d\), the proper
finite extension and \(T-\beta\) argument settles
zero geometric functionals, and the \(T-a\)
argument removes the pointwise roots of unity.
Their bounds and the finite-extension descent are
exactly those proved in
Claim~\ref{clm:FFstrictreconstruct}. This proves
its point-function conclusion here. CHW
Theorem~4.18(5) gives \(\mathcal R(K)\), and
algebraic descent gives the assertion for subfields.
\end{proof}

\begin{claim}{F.29}{Proved}\label{clm:FFboundedwild}
Let \(F/\mathbb Q\) be algebraic, and choose a
\(p\)-adic place, normalized by \(v(p)=1\).
Suppose that for some integer \(a\geq0\),
\[
                 v(F^\times)\subset
                       p^{-a}\mathbb Z_{(p)}.          \tag{F.29.1}
\]
Then \(FT_p/T_p\) is finite, and \(\mathcal R(F)\)
holds. Equivalently this condition says that the
\(p\)-parts of the ramification indices of all
finite number subfields, at the chosen place,
are bounded. These fields are exactly the subfields
of finite extensions of \(T_p\).
\end{claim}
\begin{proof}
Take the chosen henselization and strict residue
extension \(F^{sh}\). Henselization is immediate,
and the unramified extensions used in strict
henselization preserve its value group. The
chosen \(K_p\) is contained in \(F^{sh}\).
After adjoining roots of \(p\) of all orders
prime to \(p\), we obtain an extension containing
\(FT_p\) whose value group is still contained
in \(p^{-a}\mathbb Z_{(p)}\). To check the last
assertion even for this nondiscrete valuation,
note that its residue field is algebraically
closed and its units have all roots of order
prime to \(p\) by Hensel. For \((m,p)=1\),
the Kummer class of \(p\) modulo \(m\)-th powers
therefore has exactly the order of \(v(p)\)
in the value group modulo \(m\). Adjoining an
\(m\)-th root adds only the corresponding
prime-to-\(p\) denominators to that value group.

For every finite \(L/T_p\) inside \(FT_p\),
(F.28.2) says
\([L:T_p]=[v(L^\times):\mathbb Z_{(p)}]\).
Thus \([L:T_p]\leq p^a\). The maximal-degree
argument of Claim~\ref{clm:FFdiscretealgebraic}
shows that \(FT_p/T_p\) is finite.
Claim~\ref{clm:FFtamereconstruct} and algebraic
descent prove \(\mathcal R(F)\). Conversely,
(F.28.2) shows that every subfield of a finite
extension of \(T_p\) satisfies (F.29.1).
\end{proof}

\begin{claim}{F.30}{Proved}\label{clm:FFpadicsubfields}
Let \(F\subset\overline{\mathbb Q}_p\) be a field
of finite transcendence degree over \(\mathbb Q\).
Suppose that, for the induced valuation normalized
by \(v(p)=1\),
\[
                  v(F^\times)\subset
                         p^{-a}\mathbb Z_{(p)}
                    \quad\text{for some }a\geq0.       \tag{F.30.1}
\]
Then \(\mathcal R(F)\) holds. In particular this
applies when the induced valuation on \(F\) is
discrete. No henselianity assumption is required.
\end{claim}
\begin{proof}
Choose a finitely generated field
\(D\subset F\) containing a transcendence basis
of \(F/\mathbb Q\). Its closure inside
\(\overline{\mathbb Q}_p\) is a finite extension
\(M/\mathbb Q_p\): it is generated topologically
by \(\mathbb Q_p\) and the finitely many
\(\mathbb Q_p\)-algebraic generators of \(D\).
Let \(K_D\) be the chosen strict henselization
field of the valuation ring of \(D\), inside
its algebraic closure in
\(\overline{\mathbb Q}_p\), and let
\(T_D\) be the maximal tame extension of \(K_D\)
there. Completion and the chosen prolongation
identify
\[
               G_{K_D}=I_M,\qquad G_{T_D}=P_M.         \tag{F.30.2}
\]
Indeed the henselization of \(D\) is dense in
\(M\) and has the same finite separable extensions
by the Krasner--Hensel comparison proved in
Claim~\ref{clm:FFpdetector}. Strict henselization
then fixes inertia, and adjoining the prime-to-\(p\)
roots of a uniformizer fixes wild inertia.

We verify that Claims~\ref{clm:FFpdetector},
\ref{clm:FFstrictreconstruct} and
\ref{clm:FFtamereconstruct} apply to finite
extensions of \(K_D\) and \(T_D\), respectively.
The first proof uses only a henselian discrete
valuation of characteristic zero, algebraically
closed residue field of characteristic \(p\),
its completion, semistable finite-part
\(p\)-divisible groups, and Tate's theorem over
the completed DVR. These hypotheses hold here.
The exact multiplicative kernel follows from
the same discrete binomial estimate, with
\(\mu_{p'}\) supplied by strict henselization.

In the arithmetic-equivariance step of
Claim~\ref{clm:FFstrictreconstruct}, replace
its number fields \(D_i\) by finite extensions
of the fixed \(D\) containing the finitely
many model coefficients, detector coordinates,
and actual point values. Such fields suffice
because all the fields under consideration are
algebraic over \(D\). Their completions
\(M_i\) are finite \(p\)-adic fields, and
(F.30.2) gives exactly the required normal
inertia subgroups. The finite-point conjugation
argument is therefore unchanged. The Picard
1-motive step uses compactness only for
\(\operatorname{Pic}^0(M_i)\), not for the
uncompleted model field. Its resulting torsion
or zero equality for a \(D_i\)-rational Picard
point descends from \(M_i\) to \(D_i\), since
field extension is faithfully flat. A principal
boundary divisor is represented by a normalized
function over \(D_i\) by the same Hilbert~90
argument. Thus the unit power, resultant,
zero-functional argument, scalar removal and
finite-extension descent all apply.

For \(T_D\), the norm proof of the multiplicative
kernel and finite-image descent of the functional
space use only the pro-prime-to-\(p\) quotient
of inertia by wild inertia. The valuation and
tame-root calculations in
Claim~\ref{clm:FFtamereconstruct} likewise use
only the DVR \(K_D\). The argument gives
\(\mathcal R(K)\) for every finite extension
\(K/T_D\).

Finally \(F/D\) is algebraic. The bounded
\(p\)-denominator argument of
Claim~\ref{clm:FFboundedwild}, applied with
\(K_D\) in place of \(K_p\), shows that
\(FT_D/T_D\) is finite. The normalization of
its value group may include the fixed
ramification index of \(M/\mathbb Q_p\);
this changes the finite bound by a fixed factor
and does not affect the argument. Therefore
\(\mathcal R(FT_D)\) holds. The extension
\(FT_D/F\) is algebraic because \(D\subset F\)
and \(T_D/D\) is algebraic. Algebraic descent
proves \(\mathcal R(F)\).
\end{proof}

\subsection{The next cyclotomic test and its undetected functionals}

\begin{claim}{F.31}{Proved}\label{clm:FFcyclotomictest}
For every prime \(p\), put
\[
                   U_p=T_p(\mu_{p^\infty})
                       \subset\overline{\mathbb Q}.     \tag{F.31.1}
\]
This countable, transcendence-degree-zero field has
henselian value group \(\mathbb Q\), residue field
\(\overline{\mathbb F}_p\), and contains every root
of unity. For every finite extension \(E/U_p\),
\[
           C_{\operatorname{Aut}(\overline{\mathbb Q})}
                         (G_E)=1.                       \tag{F.31.2}
\]
It lies on none of the sufficient field lists
preceding this claim, including the bounded-wild class,
and none of the negative centralizer-ghost lists.
Thus \(U_5\) was a specified next test with the
smallest possible cardinality and transcendence degree
in characteristic zero; no inclusion-minimality is
claimed. Claim~\ref{clm:FFcyclotomicreconstruct}
subsequently proves reconstruction for every finite
extension of \(U_p\).
\end{claim}
\begin{proof}
Henselianity and the residue statement follow by
algebraic extension from \(K_p\). The tame
uniformizer roots supply all prime-to-\(p\)
denominators in the value group, while
\(v(\zeta_{p^n}-1)=1/(p^{n-1}(p-1))\)
supplies arbitrarily large \(p\)-denominators.
Their sums generate \(\mathbb Q\).

We first record the group input precisely.
For a finite \(p\)-adic field \(M\), its wild
inertia \(P_M\) is free pro-\(p\) of countably
infinite rank by
\citet[Theorem~2(ii), p.~463]{FF_Iwasawa1955}.
Closed subgroups of a free pro-\(p\) group are
free pro-\(p\); in a finite-rank free pro-\(p\)
group the centralizer of a nontrivial element
is procyclic. The latter statement and the
subgroup property are explicitly used and
proved/cited in
\citet[Theorem~5.1, proof, Case~I]{FF_KochloukovaZalesskii2011}.
Consequently, if \(H\subset P_M\) is nonabelian,
then \(C_{P_M}(H)=1\). Indeed if a nontrivial
\(c\) commuted with noncommuting
\(h_1,h_2\in H\), the closed subgroup
\(\langle c,h_1,h_2\rangle\) would be a
finite-rank free pro-\(p\) group in which both
\(h_i\) belong to the procyclic centralizer
of \(c\), a contradiction.

For
\(H=\ker(\chi_p:P_M\to\mathbb Z_p^\times)\),
nonabelianity is explicit. Choose four free
basis elements and map them to
\(1+e_{12},1+e_{23},1+e_{34},1+e_{45}\)
in the upper unitriangular group
\(\operatorname{UT}_5(\mathbb F_p)\).
Their first and second pair commutators map
to \(1+e_{13}\) and \(1+e_{35}\), whose
commutator is \(1+e_{15}\ne1\).
Both pair commutators belong to \(H\), since
\(\chi_p\) is abelian. This works also for
\(p=2\).

Now let \(E/U_p\) be finite. Choose a finite
extension \(L_0/K_p\) inside \(E\) such that
\(E=L_0U_p\), and then a number field model
with completion \(M\) for \(L_0\), as in
Claim~\ref{clm:FFstrictreconstruct}. We have
\[
 G_{L_0T_p}=P_M,
 \qquad G_E=\ker(\chi_p|_{P_M})=:H,
 \qquad H\triangleleft G_M.                            \tag{F.31.3}
\]
In particular \(E\) is not algebraically closed.
Let \(C=C_{G_M}(H)\). Since \(H\) is normal,
\(C\) is normal in \(G_M\); the preceding
paragraph gives \(C\cap P_M=1\).
As \(P_M\) is normal too,
\([C,P_M]\subset C\cap P_M=1\).
But \(L_0T_p/T_p\) is finite, and
Claim~\ref{clm:FFtamereconstruct}, followed by
the necessary centralizer condition of
Claim~\ref{clm:F.6}, gives
\(C_{G_{\mathbb Q}}(P_M)=1\).
Therefore \(C=1\).

Suppose finally that
\(\sigma\in\operatorname{Aut}(\overline{\mathbb Q})
=G_{\mathbb Q}\) centralizes \(G_E\).
It preserves \(E\). Its transported rank-one
henselian valuation must equal the chosen one:
otherwise the two-valuation argument of
Claim~\ref{clm:FFstricttest} would make \(E\)
algebraically closed. Uniqueness of prolongation
from a henselian field then puts \(\sigma\)
in the chosen decomposition group
\(G_{\mathbb Q_p}\subset G_{\mathbb Q}\).
Since \(G_M\) is open, some positive power
\(\sigma^r\) belongs to \(G_M\), hence to
\(C=1\). The group \(G_{\mathbb Q_p}\) has
no torsion: its wild inertia is torsion-free
free pro-\(p\), and its tame quotient is an
extension of \(\widehat{\mathbb Z}\) by
\(\prod_{\ell\ne p}\mathbb Z_\ell\), both
torsion-free. Thus \(\sigma=1\), proving
(F.31.2).

For completeness the list exclusions do not
rely only on failure of a sufficient hypothesis.
The two-valuation normal-closure proof of
Claim~\ref{clm:FFstricttest} applies to every
finite extension of \(U_p\). Thus its normal
closure over any number field of definition
is \(\overline{\mathbb Q}\), and a distinct
valuation-conjugate has compositum
\(\overline{\mathbb Q}\). Since its absolute
Galois group is infinite by (F.31.3), the same
open-subgroup argument excludes being finite
over a proper normal algebraic field. It also
cannot be contained in a proper normal algebraic
extension of a number field: first adjoin that
number field to \(U_p\), a finite extension,
and apply the normal-closure assertion.
Positive-transcendence-degree function-field
classes are excluded by transcendence degree
zero. The group \(\mu_\infty\subset U_p^\times\)
is nontrivial divisible, excluding the
Kummer-faithful sub-\(q\)-adic classes. For
any valuation of residue characteristic \(q\),
its elements \(\zeta_{q^n}-1\) have unbounded
\(q\)-denominators, excluding the bounded-wild
classes at every prime, under any field embedding.
Finally (F.31.2) excludes all negative
centralizer-ghost families. The proof of
reconstruction in Claim~\ref{clm:FFcyclotomicreconstruct}
requires the separate treatment of invisible
functionals developed next.
\end{proof}

\begin{claim}{F.32}{Proved}\label{clm:FFcyclotomicdetector}
The functional-detection assertion of
Claim~\ref{clm:FFpdetector} fails over the
explicit field \(U_5\). More precisely there
is an elliptic curve \(A/U_5\) and a nonzero
\(G_{U_5}\)-equivariant functional
\[
               V_5A\longrightarrow\mathbb Q_5(1)
\]
whose push-forward of every point Kummer class,
over every finite extension of \(U_5\), is zero.
This is not a counterexample to
\(\mathcal R(U_5)\).
\end{claim}
\begin{proof}
Take
\(A:y^2=x^3-x\), whose discriminant \(64\)
is a unit at \(5\), and choose
\(i\in K_5\) with \(i^2=-1\) and
\(i\equiv2\pmod5\). The automorphism
\([i](x,y)=(-x,iy)\) defines the usual
\(\mathbb Z[i]\)-action. Put
\(\pi=2+i\) and \(\bar\pi=2-i\).
Then \([\pi]\) is an isogeny of degree
\(N(2+i)=5\), and its tangent map is
multiplication by \(2+i\), a unit at \(5\).
Hence \([\pi^n]\) is a finite etale isogeny
of the good elliptic model over
\(\mathcal O_{K_5}\), of degree \(5^n\).
Its kernel is constant over this strictly
henselian ring. The idempotents of
\(\mathbb Z[i]\otimes\mathbb Z_5
\simeq\mathbb Z_5\times\mathbb Z_5\)
split \(T_5A\); the \(\pi\)-primary factor
is the constant etale Tate module
\(T_\pi A\simeq\mathbb Z_5\).
Projection to this factor is nonzero.

For any finite \(L/U_5\), properness extends
\(a\in A(L)\) to the valuation ring
\(\mathcal O_L\). This ring is strictly
henselian, even though it is not discrete:
it is henselian and its residue field remains
\(\overline{\mathbb F}_5\). Pulling back
\([\pi^n]\) along the section gives a finite
etale scheme over \(\mathcal O_L\), hence
it splits. Thus \(a\) has a \(\pi^n\)-division
point in \(A(L)\) for every \(n\). Applying
the same splitting successively to each chosen
division point gives a compatible division
sequence in \(A(L)\). Its full Tate-module
Kummer cocycle is therefore zero, so the
Kummer class projected to \(T_\pi A\) is zero.
Since \(5=\pi\bar\pi\), the other factor
acts invertibly on the \(\pi\)-primary part,
so this is exactly the projection of the
ordinary \(5\)-adic point Kummer class.
But \(U_5\) contains all \(5\)-power roots
of unity. Its cyclotomic representation is
therefore trivial, and any chosen nonzero
identification
\(T_\pi A\otimes\mathbb Q_5\simeq
\mathbb Q_5(1)\) is \(G_{U_5}\)-equivariant.
This is the required undetected functional.

No pinned topos map realizing that functional
has been constructed. In particular a pinned
map with constant closed-point function must
factor through that closed point by the
Jacobson argument of
Claim~\ref{clm:FFstrictdetection}; such a map
cannot realize this nonzero geometric
functional. Thus the detector counterexample
does not decide reconstruction or refute
centralizer sufficiency.
\end{proof}

\subsection{Cyclotomic reconstruction after removing invisible functionals}

\begin{claim}{F.33}{Proved}\label{clm:FFcyclotomiccohomology}
Let \(K/U_p\) be finite. Then the integral and
rational \(p\)-Kummer kernels of \(K^\times\)
are both \(\mu_\infty\). If \(M/\mathbb Q_p\)
is finite and
\(H=\ker(\chi_p|_{P_M})\), restriction
\[
 H^1(G_M,\mathbb Q_p(1))\longrightarrow
                  H^1(H,\mathbb Q_p(1))
                                                        \tag{F.33.1}
\]
is injective. The same is true with inertia
\(I_M\) in place of \(G_M\).
\end{claim}
\begin{proof}
Write \(K=L_0U_p\) with \(L_0/K_p\) finite,
and enlarge \(L_0\) to contain a specified
\(a\in K^\times\). Both \(T_p/K_p\) and
\(K_p(\mu_{p^\infty})/K_p\) are abelian, so
\(K/L_0\) is Galois abelian. Its cyclotomic
character has infinite image.
If \(a\in\bigcap_n(K^\times)^{p^n}\), choose
a compatible sequence of roots in \(K\).
Such a sequence exists because \(K\) contains
all roots of unity and each finite root fibre
is nonempty; the transition maps can equally
be made surjective by adjusting by roots of
unity. Its integral Kummer cocycle factors
through \(\Gamma=\operatorname{Gal}(K/L_0)\).
For a \(\mathbb Q_p(1)\)-valued cocycle \(c\)
of an abelian group, commutativity gives
\[
 (\chi_p(s)-1)c(t)=(\chi_p(t)-1)c(s).
                                                        \tag{F.33.2}
\]
Choose \(s\) with \(\chi_p(s)\ne1\). Equation
(F.33.2) makes \(c\) a coboundary over
\(\mathbb Q_p(1)\). Thus \(\kappa_p(a)\)
is rationally zero already over \(L_0\), so
some \(p\)-power of \(a\) belongs to the
integral kernel \(\mu_{p'}\) of
Claim~\ref{clm:FFstrictdetection}. Hence
\(a\in\mu_\infty\). The converse is immediate,
and the rational kernel is the same because
a power of an element that is a root of unity
can only come from a root of unity.

For (F.33.1), \(P_M/H\) is central in
\(G_M/H\): the cyclotomic character embeds
it in \(\mathbb Z_p^\times\), and conjugation
preserves that character. Choose a central
\(s\in P_M/H\) with \(\chi_p(s)\ne1\).
The same cocycle computation (F.33.2), now
using only centrality of \(s\), shows that
\(H^1(G_M/H,\mathbb Q_p(1))=0\).
Inflation--restriction gives injectivity,
because \(H\) acts trivially on the coefficient.
The identical argument applies to \(I_M/H\).
No triviality of the tame action is assumed.
\end{proof}


\begin{claim}{F.34}{Proved}\label{clm:FFcyclotomicsplitting}
Let \(A\) be a semiabelian variety over a finite extension
\(M/\mathbb Q_p\), put
\(H=\ker(\chi_p|_{P_M})\), and set
\[
 S=\operatorname{Hom}_H(V_pA,\mathbb Q_p(1)).
                                                        \tag{F.34.1}
\]
After a finite extension of the finite local model, the
inertia action on this space has a canonical decomposition
\[
 S=S_0\oplus S_1,\qquad
 I_M|_{S_0}=1,\quad I_M|_{S_1}=\chi_p.
                                                        \tag{F.34.2}
\]
Both summands are stable under the full local Galois group;
its unramified action on either summand need not be trivial.
For an algebraic model over a finite extension \(K/U_p\),
the summand \(S_1\) vanishes on the rational Kummer classes
of all points over all finite extensions of \(K\).
The summand \(S_0\) is jointly detected by finitely many
such point evaluations. If \(j:X\to A\) is pointed and
its geometric image generates \(A\), those evaluations
may all be taken at points \(j(x)\).
\end{claim}
\begin{proof}
We use the convention in which \(\mathbb Q_p(1)\)
has Hodge--Tate weight \(1\). The Hodge--Tate comparison
for a smooth variety with a smooth normal-crossing
compactification applies to a semiabelian variety.
Precisely, \citet[Proposition in Section~3.4 and
Theorem~3.6, especially proof~(d)]{FF_Beilinson2012}
gives the filtered, Galois-compatible de Rham
comparison for this nonproper variety; taking
the associated graded gives the Hodge--Tate
decomposition.
Here \(H^1_{\mathrm{\acute et}}(A_{\overline M},\mathbb Q_p)
=\operatorname{Hom}(V_pA,\mathbb Q_p)\); dualizing
its logarithmic Hodge filtration gives the stated
sign convention for the Tate module.
It makes \(V_pA\) Hodge--Tate with weights \(0,1\).
One can also read the weights from the torus and abelian
quotient after the comparison theorem supplies the
decomposition: the torus has weight \(1\), and the
abelian quotient has weights \(0,1\).
Consequently \(\operatorname{Hom}(V_pA,\mathbb Q_p(1))\)
and its \(G_M\)-stable subrepresentation \(S\) are
Hodge--Tate with weights \(0,1\).
Here closure under subobjects, duals, and tensor
products, with exactness on the weight spaces, is
\citet[Theorem~4.2]{FF_Fontaine2004}.

The quotient \(P_M/H\) is central in \(G_M/H\), as in
Claim~\ref{clm:FFcyclotomiccohomology}. The tame quotient
\(I_M/P_M\) is procyclic of pro-prime-to-\(p\) order.
A central extension with procyclic quotient is abelian,
so \(I_M/H\) is abelian. Its pro-prime-to-\(p\)
subgroup has finite image in \(\operatorname{GL}(S)\):
a compact linear group has an open pro-\(p\) subgroup,
whose intersection with a pro-prime-to-\(p\) image is
trivial. Kill this finite image and the finite torsion
in the cyclotomic character by a finite extension.
On a sufficiently small open subgroup of inertia,
the action is therefore
\[
           \rho_S(g)=\exp\bigl(\log\chi_p(g)\,B\bigr)
              \quad\text{for some }B\in\operatorname{End}_{\mathbb Q_p}(S).
                                                        \tag{F.34.3}
\]
The same formula on \(S\oplus\mathbb Q_p(1)\)
shows that the inertia Lie algebra is the
one-dimensional space generated by \(\operatorname{diag}(B,1)\).
By \citet[Theorem~3.1, with \(d=0\)]{FF_Ohkubo2014},
the Sen operator belongs to the scalar extension of
this inertia Lie algebra to \(\mathbb C_p\).
Its component on \(\mathbb Q_p(1)\) is \(1\), so
it is exactly \(\operatorname{diag}(B,1)\).
The Hodge--Tate decomposition makes that operator
semisimple, with eigenvalues \(0,1\) on \(S\).
This is the precise criterion in
\citet[Section~2.7, printed p.~45]{FF_Fontaine2004},
with the same sign convention.
Thus \(B\) is semisimple over \(\mathbb Q_p\),
with those eigenvalues. Its two eigenspaces give
(F.34.2), after shrinking inertia once more if necessary.
Conjugation by \(G_M\) fixes the normalized generator
\(B\), since it fixes the cyclotomic character on
the central group \(P_M/H\). Hence the two eigenspaces
are \(G_M\)-stable. This argument makes no assertion
that the unramified factors have finite image.

All finite changes of local model used here can be
realized by finite extensions of the chosen algebraic
model, by the henselization/completion equivalence.
For fixed \(S\), the finite inertia images can even
be killed inside the fixed field of \(H\), since
\(H\) acts trivially on \(S\oplus\mathbb Q_p(1)\).
Enlarging the finite model inside \(K\) in this way
does not change \(K\). Alternatively one may make
a finite extension of \(K\), as allowed throughout
the point-function argument.

We prove the asserted invisibility before detection.
For \(\lambda\in S_1\), the formula for the natural
action on Hom,
\[
        g\lambda=\chi_p(g)\lambda\rho_A(g)^{-1},
\]
and (F.34.2) say that, after choosing a basis of
the target, the underlying linear map
\(\lambda':V_pA\to\mathbb Q_p\) is inertia-equivariant
for the trivial target action. After splitting the
torus by a further finite extension, it kills its
Tate module: that module is a sum of copies of
\(\mathbb Q_p(1)\), and inertia has infinite
cyclotomic image. Thus \(\lambda'\) factors through
the Tate module of the abelian quotient \(B_A\).

Every point in question is algebraic over \(\mathbb Q\).
It is therefore defined over a finite extension \(L/K_p\),
after adjoining its finitely many coordinates to a
finite field of definition. Make a further finite
extension so that \(B_A\) has semistable reduction.
The finite-part construction used in
Claim~\ref{clm:FFpdetector} supplies a connected
\(p\)-divisible group \(\mathcal G^0\) over the
complete strict valuation ring, whose formal group
is the formal completion of the identity of the
connected N\'eron model. Restrict \(\lambda'\) to
\(V_p\mathcal G^0\) and multiply by a power of \(p\)
to make it integral. The resulting map
\(T_p\mathcal G^0\to\mathbb Z_p\) is equivariant
over the completion of \(L\). By
\citet[Theorem~4 and Corollary~1]{FF_Tate1967}, it
extends uniquely to a map of \(p\)-divisible groups
\[
                  \mathcal G^0\longrightarrow\mathbb Q_p/\mathbb Z_p.
\]
This map is zero: every finite connected source
has zero map to a constant \'etale finite group.
Thus \(\lambda'\) kills \(V_p\mathcal G^0\).

For any \(a\in B_A(L)\), some nonzero integer
multiple of \(a\) lies in that formal completion.
Indeed, the component group is finite, and every
point of the connected reduction over
\(\overline{\mathbb F}_p\) has finite order:
it and the variety descend to a finite field.
Choose compatible successive \(p\)-division points
inside the formal group. Their differences under
Galois lie in \(T_p\mathcal G^0\), so applying
\(\lambda'\) kills the full compatible Kummer
cocycle, not just its classes at separate finite
levels. Division by the integer multiple is
legitimate with rational coefficients. Restriction
to the group of any finite extension of \(K\)
therefore gives zero. This proves invisibility
of \(S_1\).

For \(\lambda\in S_0\setminus\{0\}\), inertia
equivariance permits application of
Claim~\ref{clm:FFpdetector} over a finite extension
of \(K_p\). It gives a point over a further
finite extension whose image under \(\lambda\)
has nonzero rational Kummer class. The restriction
from this inertia group to its
\(\ker(\chi_p|_P)\) subgroup is injective by
Claim~\ref{clm:FFcyclotomiccohomology}, so the
class remains nonzero after adjoining \(U_p\).
The pointed-generation assertion follows from
the corresponding assertion in
Claim~\ref{clm:FFpdetector}. Finally, a finite
family of evaluations suffices on the
finite-dimensional space \(S_0\): successively
choose an evaluation nonzero on the remaining
common kernel until that kernel is zero.
\end{proof}

\begin{claim}{F.35}{Proved}\label{clm:FFcyclotomicreconstruct}
For every prime \(p\), every finite extension of
\[
 U_p=\operatorname{Frac}(\mathbb Z_{(p)}^{sh})
       \bigl(p^{1/m}:(m,p)=1,\ \mu_{p^\infty}\bigr)
                                                        \tag{F.35.1}
\]
satisfies \'etale reconstruction. Every subfield of
such an extension also satisfies \'etale reconstruction.
In particular, the failure of the functional detector
in Claim~\ref{clm:FFcyclotomicdetector} does not prevent
reconstruction of \(U_5\).
\end{claim}
\begin{proof}
We prove the actual-point-function condition of
\citet[Theorem~4.18(5)]{FF_CHW2024}. Let \(K/U_p\)
be finite, let \(X/K\) be regular, connected and
quasiprojective with \(x_0\in X(K)\), and let
\(\varphi:X_{\mathrm{\acute et}}\to
\mathbb G_{m,K,\mathrm{\acute et}}\) be pinned.
Use the notation \(b(x),b_0,\xi,j:X\to A\) of
Claim~\ref{clm:FFstrictreconstruct}; in particular
\(j(x_0)=0\). As there, \(X\) is geometrically
integral and its rational points are Zariski dense.
The density proof uses only Hensel's lemma and
infinite valuation balls, so applies to the
nondiscrete henselian field \(K\).
The zero-dimensional case is already covered by
Claim~\ref{clm:F.1}.

We may pass to finite extensions of \(K\) during
this proof: a unit realizing the original geometric
point function over such an extension descends by
Galois invariance and equality of functions on a
reduced scheme, exactly as at the beginning of
Claim~\ref{clm:FFstrictreconstruct}. This does not
presuppose reconstruction over the larger field.
All finite fields considered below remain finite
extensions of \(U_p\).

\emph{1. The part detected by actual point values.}
Choose a number field \(D_0\subset K\) containing
all the finite model data and \(b_0\), with
\(K=D_0U_p\), and set \(M_0=(D_0)_v\).
Then
\[
 G_K=H_0=\ker(\chi_p|_{P_{M_0}})
                  \triangleleft G_{M_0}.
                                                        \tag{F.35.2}
\]
The geometric part of \(\xi-\kappa_p(b_0)\)
is an integral functional \(\lambda\) in the
space \(S\) of Claim~\ref{clm:FFcyclotomicsplitting}.
After the finite changes allowed there, write
\(\lambda=\lambda_0+\lambda_1\), with
\(\lambda_i\in S_i\). For every finite-extension
point, the actual-value identity becomes
\[
  (\lambda_0)_*\kappa_A(j(x))
                =\kappa_p(b(x)/b_0)
               \quad\text{over its field of definition}.
                                                        \tag{F.35.3}
\]
Indeed, \(\lambda_1\) kills that entire rational
Kummer class by Claim~\ref{clm:FFcyclotomicsplitting}.

Choose finitely many point evaluations jointly
detecting \(S_0\), over one finite extension \(E/K\).
Include the points, their actual values, and a
generator of \(E/U_p\) in a number field \(D_1\).
Then \(E=D_1U_p\), \(M_1=(D_1)_v\), and
\(H_1=\ker(\chi_p|_{P_{M_1}})\) is normal in
\(G_{M_1}\). Joint detection survives further
finite restriction, by restriction--corestriction
with rational coefficients.
For \(g\in G_{M_1}\), the difference
\(g\lambda_0-\lambda_0\) lies in \(S_0\):
both \(S\) and its weight-zero summand are
\(G_{M_0}\)-stable. The restricted point and
value cocycles in (F.35.3) are invariant under
the conjugation action of \(G_{M_1}\), since
their points and actual values are defined over
\(M_1\). The same cocycle computation as
(F.26.4) therefore makes every chosen evaluation
of \(g\lambda_0-\lambda_0\) zero. Joint detection
gives
\[
                \lambda_0\in
             \operatorname{Hom}_{G_{M_1}}(V_pA,\mathbb Q_p(1)).
                                                        \tag{F.35.4}
\]
No equivariance of \(\lambda_1\) under the full
local Galois group is needed.

There is a normalized rational global class
\(\xi_0\) over \(X_{M_1}\) with geometric part
\(\lambda_0\), and with value zero at \(x_0\).
Here is an explicit construction, so no integrality
of the projection is silently assumed. The point
\(x_0\) splits the arithmetic fundamental-group
sequence. On
\(\pi_1(X_{\overline M_1})\rtimes G_{M_1}\),
define the cocycle by applying \(\lambda_0\)
to the geometric abelian pro-\(p\) component
and assigning zero to the section \(G_{M_1}\).
Equivariance (F.35.4) is exactly the cocycle
identity. Equivalently, pull the universal rational
Kummer class on \(A\) back by \(j\), and apply
\(\lambda_0\). Its evaluations satisfy (F.35.3).
The same construction defines the compatible
class over the algebraic field \(E\).
We henceforth replace \(K\) by \(E\), as permitted.

\emph{2. Boundary residues are rational.}
The summand \(\lambda_1\) is zero on the torus
of the generalized Albanese, as proved in
Claim~\ref{clm:FFcyclotomicsplitting}. Thus the
boundary residues of \(\xi_0\) are the residues
of the original integral class \(\xi\), and
belong to \(\mathbb Z_p\). We now prove that
they also belong to \(\mathbb Q\).

On a punctured rational line over a finite model,
write the normalized class in the form (F.26.1),
with residues \(a_i\in\mathbb Z_p\). Choose
a uniformizer \(\pi\) of the finite local model
in its dense number subfield. Evaluate at
\(r_i+\pi^N\) and \(r_i+\pi^{N+1}\), for \(N\)
larger than the valuations of all nonzero root
differences. Include these points and their
actual values in a further finite number field,
with completion \(M_2\). Over
\(H_2=\ker(\chi_p|_{P_{M_2}})\), equality
(F.35.3) identifies each evaluation with the
Kummer class of its actual value. Both classes
come from \(G_{M_2}\); the restriction injection
of Claim~\ref{clm:FFcyclotomiccohomology} therefore
gives that equality already over \(G_{M_2}\).
Apply the ordinary normalized local valuation
to the difference of the two equalities. The
other puncture factors have constant valuation,
so
\[
        a_i v_p(\pi)
           =v_p\bigl(b(r_i+\pi^{N+1})\bigr)
                -v_p\bigl(b(r_i+\pi^N)\bigr)
                    \in\mathbb Q.
                                                        \tag{F.35.5}
\]
Consequently \(a_i\in\mathbb Q\cap\mathbb Z_p
=\mathbb Z_{(p)}\).
This calculation uses the finite local field,
not a valuation map on Kummer cohomology over
\(K\); the latter would lose this information
because \(v(K^\times)=\mathbb Q\).

For a smooth curve, use the finite \'etale map
to a punctured line with a chosen boundary
point as its only pole, exactly as in Step~1
of Claim~\ref{clm:FFstrictreconstruct}.
Corestriction has actual point values given by
products of field norms, and the residue at
infinity is the chosen residue. The rational
class is defined over a finite local field,
and its norm values are still algebraic over
\(\mathbb Q\), so the preceding comparison
applies to that class without change.
Transverse curves through the smooth locus
of each boundary component, with intersection
multiplicity one, give the higher-dimensional
assertion. Therefore the full residue divisor
\(D\) of \(\xi_0\) has coefficients in
\(\mathbb Z_{(p)}\).

\emph{3. An algebraic unit up to powers and roots of unity.}
Apply the Picard 1-motive calculation (F.26.5)
to the rational invariant geometric class
\(\lambda_0\). Choose a positive integer \(q\)
such that \(q\lambda_0\) is integral and
\(qD\in B_X\). Such an integer exists:
first clear the finitely many denominators
of \(D\); then use
\[
 (B_X\otimes\mathbb Q_p)\cap
       \Bigl(\bigoplus_i\mathbb QD_i\Bigr)
                         = B_X\otimes\mathbb Q
\]
and clear its rational lattice denominators;
finally multiply by a power of \(p\) to make
the entire geometric class integral.
An invariant integral lift with lattice
coordinate exactly \(qD\) gives
\([qD]\in\bigcap_n p^nP(M_1)\), by the
finite-level division-point argument in
Step~3 of Claim~\ref{clm:FFstrictreconstruct}.
That intersection is finite prime-to-\(p\)
torsion. Hence a positive multiple \(dD\)
is principal. The principal divisor descends
to the number-field model and has a function
\(h\in\mathcal O(X)^\times\), normalized by
\(h(x_0)=b_0^d\).

The zero-residue difference between
\(d\lambda_0\) and the geometric Kummer class
of \(h\) lies in \((V_pP)^{G_{M_1}}=0\).
Normalization at \(x_0\) therefore gives the
rational equality
\[
                   d\xi_0=\kappa_p(h/b_0^d).
                                                        \tag{F.35.6}
\]
Evaluate and use (F.35.3). The rational
Kummer kernel in every finite extension of
\(K\) is \(\mu_\infty\), by
Claim~\ref{clm:FFcyclotomiccohomology}. Thus
\[
                b(x)^d=\zeta_xh(x),
                         \qquad \zeta_x\in\mu_\infty\subset K.
                                                        \tag{F.35.7}
\]
We have proved this pointwise relation; we
do not identify the original class \(\xi\)
with the Kummer class of \(h\). Its invisible
summand may be nonzero, as Claim~\ref{clm:FFcyclotomicdetector}
demonstrates.

\emph{4. Closed points of arbitrarily large degree.}
Every finite extension \(L/K\) has infinite
absolute Galois group. Indeed,
Claim~\ref{clm:FFcyclotomictest} identifies it
with a nonabelian subgroup of the free wild
inertia group. An infinite profinite group
has open subgroups of arbitrarily large
index; otherwise its finite quotients would
have bounded order, forcing the group to
be finite. Finite separable extensions of
\(L\) therefore have unbounded degrees.
By the primitive-element theorem their
generators give closed points on
\(\mathbb A^1_L\) of unbounded degrees.
A finite nonconstant map from the proper
normalization of a geometrically integral
curve to \(\mathbb P^1_L\) then gives
closed points of unbounded degrees in
every nonempty open of that curve: choose
base points of larger degree than every
point in the finite image of its omitted
locus. Points above them have at least
the same degree. This replaces the
Eisenstein-polynomial construction, which
would not be valid for the divisible
value group of \(K\).

\emph{5. Recovering the actual point function.}
The remaining algebraic steps of
Claim~\ref{clm:FFstrictreconstruct}, Steps~4--7,
now apply, with \(\lambda_0\) in place of
\(\lambda\) and \(\mu_\infty\) in place of
\(\mu_{p'}\). We spell out why their exact
inputs have all been retained.

If \(\lambda_0\ne0\), the unit \(h\) in
(F.35.6) is nonconstant. Relation (F.35.7)
and the existence of large-degree points
imply that the actual point function is
nonconstant. On the nonempty inverse-image
open of \(T-1\ne0\), apply Steps~1--3 to
that target function. This gives
\((b(x)-1)^e=\eta_xg(x)\), with
\(\eta_x\in\mu_\infty\subset K\).
Use a smooth geometrically integral curve
on which restriction is injective on the
finite span of \(h^ig^j\), as supplied by
\citet[Theorem~1.8 and Corollary~1.9]{FF_CharlesPoonen2016}.
For a closed point of degree greater than
\(e\deg\operatorname{div}_\infty h+
d\deg\operatorname{div}_\infty g\), the
resultant
\[
        \operatorname{Res}_Z
          (Z^d-\zeta_xh,\ (Z-1)^e-\eta_xg)
\]
vanishes identically, by its uniform pole
bound. Its coefficients lie in \(K\),
since all roots of unity do. Restriction
injectivity makes it vanish in \(K(X)\).
Its common root \(z\) is rational: a
conjugate satisfies \(z'=\alpha z\) and
\(z'-1=\beta(z-1)\), with \(\alpha,\beta\)
constant roots of unity, and nonconstancy
forces \(\alpha=\beta=1\). It is a unit
because \(z^d=\zeta_xh\).
Equation (F.35.6) then gives
\(\lambda_0=\kappa_p(z)\) geometrically
with rational coefficients. Rational
coefficients suffice for this division
by \(d\), and a rationally zero unit
Kummer class still has zero divisor.
Normalize \(z(x_0)=b_0\). Evaluation at
\(x_0\) kills the constant cohomology
difference, so (F.35.3) becomes
\(\kappa_p(b(x))=\kappa_p(z(x))\)
at every finite-extension point. Hence
\(b(x)/z(x)\in\mu_\infty\subset K\).

If \(\lambda_0=0\), (F.35.3) puts all
values \(b(x)/b_0\) in \(\mu_\infty\subset K\).
Choose \(\beta\) in a proper finite
extension of \(K\), and apply the same
construction to \(T-\beta\). Such an
extension exists by Step~4. If its
weight-zero class is nonzero, the preceding
resultant gives a nonconstant unit, whereas
all values of that unit lie in the fixed
finite base extension containing \(\beta\):
indeed \(b(x)\in K\), and its root-of-unity
ratio also belongs to \(K\). This holds for
all geometric points, including points defined
over arbitrarily larger finite fields: every
\(b(x)-\beta\) still lies in \(K(\beta)\).
The
large-degree-point argument in Step~5
of Claim~\ref{clm:FFstrictreconstruct}
excludes this. If its weight-zero class
is zero, comparing the two root-of-unity
ratios and using the linear independence
of \(1,\beta\) over \(K\) makes
\(b(x)=b_0\). Thus the zero case has a
constant actual point function.

Finally, for a nonconstant weight-zero
class, (F.35.3) and its realizing unit
give the same constant-scalar ambiguity
as in Step~6 of
Claim~\ref{clm:FFstrictreconstruct}.
Apply the construction also to \(T-1\),
and use a closed point of degree larger
than the relevant pole bound and the
degree of the nonconstant unit as a map
to \(\mathbb P^1\). Linear independence
at that point forces the two units and
the actual values to have one fixed
scalar relation. Applying \(T-a\) deals
with the remaining base-field-value
points. Each nonconstant point function
remains nonconstant on its inverse-image
open: otherwise the inverse image of a
closed target point would contain a dense
open of the integral source and hence
be the whole source. In higher dimension,
the same argument is made on a smooth
geometrically integral curve through the
base point and the point being tested,
using the cited Charles--Poonen theorem.
These are precisely the algebraic
scalar-removal steps already proved there;
they require only (F.35.3), all roots of
unity in the base, and Step~4.

The weight decomposition used throughout
is compatible with restriction to opens,
finite scalar extensions, and pullback
along the curves: it is the eigenspace
decomposition of the normalized inertia
Lie generator. In particular a geometric
Kummer functional of an algebraic unit
belongs to weight zero, and comparisons
of such units in this last step preserve
the normalized class \(\lambda_0\).
We have therefore obtained a unit whose
actual geometric point function is exactly
that of \(\varphi\). Finite descent returns
it to the original \(K\). The point-function
criterion cited at the start proves
\(\mathcal R(K)\).
Every subfield \(F\subset K\) has
\(K/F\) algebraic, since \(K/\mathbb Q\)
is algebraic. Algebraic descent, Claim~\ref{clm:FFdescent},
proves \(\mathcal R(F)\).
\end{proof}

\subsection{Maximal abelian extensions of strict henselization fields}

\begin{claim}{F.36}{Proved}\label{clm:FFabelianHT}
Let \(M/\mathbb Q_p\) be finite, and let \(S\)
be a Hodge--Tate representation of \(G_M\) whose
inertia image is abelian. The inertia representation
on \(S\) is semisimple. In particular there is a
canonical decomposition
\[
 S=S^{I_M}\oplus
       \sum_{g\in I_M}(g-1)S.                         \tag{F.36.1}
\]
Its projection onto the first summand commutes
with \(G_M\). For every inertia-stable subspace
\(N\), that projection sends \(N\) to
\(N\cap S^{I_M}\), and taking inertia invariants
is exact on \(0\to N\to S\to S/N\to0\).
All assertions hold after any finite extension
of \(M\).
\end{claim}
\begin{proof}
Write \(\mathfrak g\) for the Lie algebra of
the compact inertia image. It is abelian.
Let \(B\subset\operatorname{End}_{\mathbb Q_p}(S)\)
be the unital commutative algebra generated by
\(\mathfrak g\). It is a finite-dimensional
commutative algebra over the characteristic-zero
field \(\mathbb Q_p\).

We recall the elementary algebra fact used
here. Such an algebra has a unique maximal
\'etale subalgebra \(B_{\rm ss}\), and
\[
           B=B_{\rm ss}\oplus\operatorname{rad}(B)
                         \quad\text{as vector spaces}.
                                                        \tag{F.36.2}
\]
One proof uses the additive and multiplicative
Jordan decompositions of commuting matrices:
their semisimple parts belong to the algebra,
are polynomials in the original matrices, and
add and multiply as the original matrices do
modulo the nilpotent radical. These parts
give the unique section of the finite separable
algebra \(B/\operatorname{rad}(B)\). Equivalently,
in each local Artin factor, its separable residue
field has a unique lift through the nilpotent
ideal. After extension to \(\mathbb C_p\),
each factor has the form
\(\mathbb C_p\oplus\mathfrak m\), with
\(\mathfrak m\) nilpotent. An element acting
semisimply on the faithful module has zero
\(\mathfrak m\)-component. Thus the semisimple
elements of \(B\otimes\mathbb C_p\) are exactly
\(B_{\rm ss}\otimes\mathbb C_p\).

The Sen operator \(\Theta\) of \(S\) is
semisimple, because \(S\) is Hodge--Tate
\citep[Section~2.7]{FF_Fontaine2004}.
By \citet[Theorem~3.1, \(d=0\)]{FF_Ohkubo2014},
\(\Theta\in\mathfrak g\otimes\mathbb C_p\),
and \(\mathfrak g\) is the smallest
\(\mathbb Q_p\)-subspace of the endomorphism
space with this property. The preceding
algebra fact puts \(\Theta\) in
\(B_{\rm ss}\otimes\mathbb C_p\). Minimality
therefore gives \(\mathfrak g\subset B_{\rm ss}\).
In particular every element of \(\mathfrak g\)
acts semisimply, and these elements commute.

On a sufficiently small open subgroup of
the compact inertia image, matrix logarithm
has values in \(\mathfrak g\), and matrix
exponential is its inverse. Hence every
element of this open subgroup is semisimple.
An arbitrary inertia matrix has a positive
integer power in this subgroup. A matrix
over a characteristic-zero field whose
positive power is invertible and semisimple
is itself semisimple: if that power has
distinct eigenvalues \(\alpha_j\), the
matrix is annihilated by the square-free
polynomial \(\prod_j(T^n-\alpha_j)\).
All inertia matrices commute, so they are
simultaneously diagonalizable over a finite
extension of \(\mathbb Q_p\). Their
representation is semisimple over
\(\mathbb Q_p\) as well.

In this simultaneous decomposition,
\(S^{I_M}\) is the trivial-character part,
and the sum in (F.36.1) is exactly the
sum of the nontrivial-character parts.
This description is independent of a
coefficient extension. It also proves
the assertion about every stable \(N\)
and exactness of invariants. Conjugation
by \(G_M\) preserves inertia and its
trivial-character part, as well as the
explicit second summand in (F.36.1).
The projection therefore commutes with
\(G_M\). Restricting to any open local
Galois subgroup retains both the
Hodge--Tate and abelian-image hypotheses.
\end{proof}

\begin{claim}{F.37}{Proved}\label{clm:FFmaxabelianreconstruct}
Let \(K_p=\operatorname{Frac}(\mathbb Z_{(p)}^{sh})\)
with the chosen embedding in \(\overline{\mathbb Q}\),
and let \(K_p^{\rm ab}\) be its maximal
abelian extension in that algebraic closure.
Every finite extension of \(K_p^{\rm ab}\),
and every subfield of such an extension,
satisfies \'etale reconstruction.
The extension \(K_p^{\rm ab}/U_p\) is infinite,
so this is a proper enlargement of
Claim~\ref{clm:FFcyclotomicreconstruct} at the
level of the specified henselian fields.
\end{claim}
\begin{proof}
Put \(I=I_{\mathbb Q_p}=G_{K_p}\), and
let \(J=\overline{[I,I]}\) be its closed
commutator subgroup. Thus
\(G_{K_p^{\rm ab}}=J\).
The tame quotient of inertia is abelian,
so \(J\subset P_{\mathbb Q_p}\).
It contains the commutator subgroup of
the free pro-\(p\) group
\(P_{\mathbb Q_p}\), and is infinite:
two free generators have a nonidentity
commutator, of infinite order by
torsion-freeness. Every open subgroup
of \(J\) is consequently infinite.
This uses the same verified free wild
inertia theorem as Claim~\ref{clm:FFcyclotomictest}.

Both \(T_p/K_p\) and the cyclotomic
extension \(K_p(\mu_{p^\infty})/K_p\)
are abelian. Hence \(U_p\subset K_p^{\rm ab}\).
In particular these fields contain
\(\mu_\infty\), have value group
\(\mathbb Q\), and have residue field
\(\overline{\mathbb F}_p\).

\emph{1. The enlargement is infinite.}
Set \(L=K_p(\zeta_p)\), including
\(L=K_2\) when \(p=2\), and let
\(\pi=\zeta_p-1\). Then \(\pi\) is
a uniformizer of \(L\), and
\(\Delta=\operatorname{Gal}(L/K_p)\)
has order \(p-1\). For the principal
unit filtration \(U^i=1+\pi^i\mathcal O_L\),
\[
                    U^1/U^2\simeq
                 \overline{\mathbb F}_p
                \quad\text{as }\mathbb F_p\text{-spaces}.
                                                        \tag{F.37.1}
\]
The \(\Delta\)-action on this quotient
is the mod-\(p\) cyclotomic character
\(\bar\chi_p\): its residue action is
trivial, and
\(\sigma(\pi)\equiv\bar\chi_p(\sigma)\pi
\pmod{\pi^2}\).
If \(u^p\in U^1\), then \(u\) is a
unit of residue one, because the residue
field has characteristic \(p\). Writing
\(u=1+\pi c\), the binomial expansion
and \(v_L(p)=p-1\) show that
\(u^p\in U^2\). Therefore (F.37.1)
is a quotient of a subspace of
\(L^\times/(L^\times)^p\).
Since \(|\Delta|\) is prime to \(p\),
taking its character parts is exact.
The \(\bar\chi_p\)-part of this Kummer
space is infinite-dimensional over
\(\mathbb F_p\).

Kummer theory and the twist by
\(\mu_p^{-1}\) identify that part
with \(H^1(G_L,\mathbb F_p)^\Delta\).
Inflation--restriction for \(\Delta\)
identifies the latter with
\(H^1(I,\mathbb F_p)\), because
the positive-degree cohomology of this
prime-to-\(p\) finite group is zero.
Thus \(I^{\rm ab}\) has an
infinite-dimensional elementary abelian
pro-\(p\) quotient. By contrast
\(\operatorname{Gal}(U_p/K_p)\)
is a closed subgroup of a product
of a pro-prime-to-\(p\) procyclic
group and a subgroup of
\(\mathbb Z_p^\times\); its pro-\(p\)
Sylow embeds in the second factor.
Its elementary abelian pro-\(p\)
quotient is finite-dimensional.
A finite kernel could increase that
dimension by only a finite amount.
Hence \(K_p^{\rm ab}/U_p\) cannot be finite.

\emph{2. Exact cohomological inputs for all finite extensions.}
Let \(K/K_p^{\rm ab}\) be finite.
It has the form \(L_0K_p^{\rm ab}\)
for a finite extension \(L_0/K_p\).
Every specified element of \(K\)
can be included in a larger such
\(L_0\). The extension \(K/L_0\)
is Galois abelian and has infinite
cyclotomic image. The abelian cocycle
calculation (F.33.2), together with
the exact Kummer kernel over \(L_0\),
therefore proves
\[
 \ker\bigl(K^\times\to H^1(G_K,\mathbb Z_p(1))\bigr)
 =\ker\bigl(K^\times\to H^1(G_K,\mathbb Q_p(1))\bigr)
 =\mu_\infty.
                                                        \tag{F.37.2}
\]
Compatible roots may be chosen because
all roots of unity are in \(K\).
This is the proof of
Claim~\ref{clm:FFcyclotomiccohomology}
with another abelian extension in
place of \(U_p\); all its hypotheses
have just been checked.

Choose a number field \(D_0\subset K\)
containing any prescribed finite data
and generators of \(K/K_p^{\rm ab}\),
and put \(M_0=(D_0)_v\). Inside the
fixed local Galois group the absolute
group of \(K\) is
\[
            H_0=J\cap G_{M_0}
                         \triangleleft G_{M_0}.
                                                        \tag{F.37.3}
\]
The quotient \(I_{M_0}/H_0\) is
abelian, as a subgroup of \(I/J\),
and its cyclotomic image is infinite.
On this abelian quotient, (F.33.2)
gives \(H^1(I_{M_0}/H_0,\mathbb Q_p(1))=0\).
Its coefficient invariants are zero.
Inflation--restriction for inertia
in \(G_{M_0}/H_0\) then also gives
\(H^1(G_{M_0}/H_0,\mathbb Q_p(1))=0\).
Consequently both restriction maps
\[
 H^1(I_{M_0},\mathbb Q_p(1))\longrightarrow H^1(H_0,\mathbb Q_p(1)),
 \qquad
 H^1(G_{M_0},\mathbb Q_p(1))\longrightarrow H^1(H_0,\mathbb Q_p(1))
                                                        \tag{F.37.4}
\]
are injective. Every finite extension
\(E/K\) equals \(D_1K_p^{\rm ab}\)
for a finite number field \(D_1/D_0\),
and has \(H_1=J\cap G_{M_1}\) with
\(M_1=(D_1)_v\). The same proof gives
(F.37.2)--(F.37.4) at every such stage.

\emph{3. The invisible subspace.}
Use the regular pointed source \(X\),
its generating generalized Albanese
\(j:X\to A\), the actual point function
\(b\), and its geometric functional
\(\lambda\) as in
Claim~\ref{clm:FFcyclotomicreconstruct}.
Include the data in \(D_0\), and set
\[
                   S=\operatorname{Hom}_{H_0}
                                (V_pA,\mathbb Q_p(1)).
                                                        \tag{F.37.5}
\]
It is a \(G_{M_0}\)-stable Hodge--Tate
space, by the exact comparison and
subobject results verified in
Claim~\ref{clm:FFcyclotomicsplitting}.
Its inertia action factors through
the abelian quotient in (F.37.3).
Claim~\ref{clm:FFabelianHT} thus applies.

Let \(N\subset S\) consist of those
functionals whose pushforward of
\(\kappa_A(a)\) is zero for every
\(a\in A(L)\) over every finite
extension \(L/K\). This is a linear
subspace. It is \(G_{M_0}\)-stable:
conjugation sends such an \(L\) to
another finite extension of \(K\),
and sends the point and its Kummer
cocycle to their conjugates. This
uses normality of \(H_0\), and
includes the coefficient action.

The same \(N\) is the common kernel
of all evaluations at \(j(x)\).
To see the nontrivial inclusion, recall
that finitely many signed sums of
\(j(X_{\bar K})\) dominate \(A_{\bar K}\).
Every point of \(A(\bar K)\) is a
difference of points in this dominant
image, since two nonempty opens related
by a translation intersect. Thus it
is a finite signed sum of \(j(x_i)\),
after a finite field extension.
A nonzero rational Kummer evaluation
cannot become zero under a finite
restriction, by restriction--corestriction.
Consequently vanishing on all \(j(x)\)
implies vanishing on every \(a\).
Finite-dimensionality of \(S/N\)
now supplies finitely many such
evaluations with common kernel
exactly \(N\), over one finite
extension \(E/K\).

For every finite local model \(M_1\)
obtained as in Step~2, we have
\[
                        N\cap S^{I_{M_1}}=0.
                                                        \tag{F.37.6}
\]
Indeed, a nonzero inertia-invariant
functional is detected over a finite
extension of the corresponding strict
henselization by
Claim~\ref{clm:FFpdetector}. Its
nonzero class survives adjoining
\(K_p^{\rm ab}\) by (F.37.4),
so it cannot belong to \(N\).

\emph{4. The canonical invariant lift.}
Include the finitely many detecting
points and their actual values in
\(D_1\subset E\), so that
\(E=D_1K_p^{\rm ab}\) and
\(H_1=J\cap G_{M_1}\) is normal
in \(G_{M_1}\). The actual-value
cohomology identity, followed by
conjugation, gives
\[
                  g\lambda-\lambda\in N
                         \quad(g\in G_{M_1}).
                                                        \tag{F.37.7}
\]
Here the selected point and value
cocycles are defined over \(M_1\),
and their restricted classes are
conjugation-invariant. Their common
evaluation kernel is exactly \(N\).
No descent of the topos map is
being assumed.

Take the canonical projection
\(\Pi:S\to S^{I_{M_1}}\) from
Claim~\ref{clm:FFabelianHT}, and put
\(\lambda_0=\Pi\lambda\).
Equation (F.37.7) makes the image
of \(\lambda\) in \(S/N\)
inertia-invariant. Exactness of
invariants and (F.37.6) imply
\[
               \lambda-\lambda_0\in N,
           \qquad \lambda_0\in
             \operatorname{Hom}_{G_{M_1}}(V_pA,\mathbb Q_p(1)).
                                                        \tag{F.37.8}
\]
For the second assertion, \(\Pi\)
commutes with \(G_{M_1}\), and
\(\Pi N=N\cap S^{I_{M_1}}=0\).
Thus applying \(\Pi\) to
(F.37.7) kills every difference.
This is the unique inertia-invariant
lift of the class modulo \(N\).
It is not asserted to have been
\(I_{M_0}\)-invariant before the
finite actual-value descent.

After splitting the torus of \(A\)
over a further finite model, restriction
of a functional to that torus has
inertia-trivial target
\(\operatorname{Hom}(V_pT,\mathbb Q_p(1))\).
Every element of \(N\) has zero
restriction: an inertia-equivariant
map to a trivial representation
factors through \(\Pi\), and
\(\Pi N=0\). The original integral
boundary residues of \(\lambda\)
are therefore unchanged by replacement
with \(\lambda_0\).

\emph{5. Reconstruction from the retained point values.}
Construct the normalized rational
class \(\xi_0\) from \(\lambda_0\)
using the section at \(x_0\),
exactly as in Step~1 of
Claim~\ref{clm:FFcyclotomicreconstruct}.
Since \(\lambda-\lambda_0\in N\),
its evaluation at every finite-extension
point is \(\kappa_p(b(x)/b_0)\).
Its residues are the integral residues
of the original class, by Step~4.
All the subsequent hypotheses in
that reconstruction proof are now
available, as follows.

First, (F.37.4) lifts an equality
of actual evaluations from \(H_2\)
to the finite local model \(G_{M_2}\).
The two nearby punctured-line values
therefore force each boundary residue
into \(\mathbb Q\cap\mathbb Z_p\),
using the ordinary finite-local-field
valuation. Norm transfer and transverse
curves give the same assertion in
every dimension. Second, the invariant
rational Picard 1-motive lift can
be multiplied to have integral Tate
coordinates and an integral lattice
coordinate. The compact local Picard
group then makes a further multiple
of that boundary divisor principal,
and its invariant zero-residue Tate
part is zero. Thus a unit \(h\)
and integer \(d>0\) satisfy
\[
                    b(x)^d=\zeta_xh(x),
                         \qquad\zeta_x\in\mu_\infty\subset K,
                                                        \tag{F.37.9}
\]
by (F.37.2). This is an equality
of point values, not an identification
of the original integral topos class
with a unit class.

Third, every finite extension of
\(K\) has infinite absolute Galois
group, since its group is open in
\(J\). Hence the primitive-element
and finite-map argument in Step~4
of Claim~\ref{clm:FFcyclotomicreconstruct}
supplies closed points of arbitrarily
large degree on every nonempty
curve open. The \(T,T-1\) resultant
applied to (F.37.9) therefore makes
the retained nonzero functional the
geometric Kummer functional of a unit.
If it is zero, every actual value
lies in \(b_0\mu_\infty\subset K\),
and the \(T-\beta\) argument forces
the point function to be constant.
The \(T-a\) comparisons remove the
remaining root-of-unity scalars.

These comparisons are compatible
with the canonical projections just
used. To verify this without making
\(N\) explicit, a pullback on
Albanese Tate modules carries an
inertia-invariant realizing functional
to an inertia-invariant one, and
carries an invisible functional to
an invisible one: point Kummer
classes pull back by actual points.
Uniqueness of the invariant lift
in (F.37.8) gives the required
compatibility, after a common finite
model extension. The same applies
to restrictions to nonempty opens
and to the curves used for testing.
Thus every step of the actual
point-function and scalar-removal
proof is preserved.

It follows that the pinned map's
geometric point function is realized
by a unit over a finite extension.
As in Claim~\ref{clm:FFstrictreconstruct},
equality on geometric points makes
that unit Galois-invariant, so it
descends to the original \(K\).
The point-function criterion
\citep[Theorem~4.18(5)]{FF_CHW2024}
proves \(\mathcal R(K)\).
Since \(K\) is algebraic over
\(\mathbb Q\), it is algebraic
over each of its subfields.
Claim~\ref{clm:FFdescent} proves
the final subfield assertion.
\end{proof}


\begin{claim}{F.38}{Proved}\label{clm:FFnilpotentinputs}
The following two assertions hold.
\begin{enumerate}
\item Let \(\Gamma\) be a profinite group
nilpotent of some fixed finite class, and let
\(\chi:\Gamma\to\mathbb Q_p^\times\) be a
nontrivial continuous character. Then
\[
                 H^1(\Gamma,\mathbb Q_p(\chi))=0.
                                                        \tag{F.38.1}
\]
\item Let \(S\) be a Hodge--Tate representation
of the Galois group of a finite extension
\(M/\mathbb Q_p\). If its inertia Lie algebra
is nilpotent, that Lie algebra is abelian.
After a finite extension of \(M\), the
inertia representation is semisimple and
has the canonical invariant projection
of Claim~\ref{clm:FFabelianHT}.
\end{enumerate}
\end{claim}
\begin{proof}
For the first assertion, induct on the
nilpotency class and put \(Z=Z(\Gamma)\), a closed central subgroup. If
\(\chi|_Z\ne1\), choose \(z\in Z\)
with \(\chi(z)\ne1\). For every cocycle
\(c\), centrality gives
\[
             (\chi(z)-1)c(g)=(\chi(g)-1)c(z).
\]
Thus \(c\) is a coboundary, proving
(F.38.1). If \(\chi|_Z=1\), it descends
to a nontrivial character of \(\Gamma/Z\).
Inflation--restriction gives
\[
 0\longrightarrow H^1(\Gamma/Z,\mathbb Q_p(\chi))
 \longrightarrow H^1(\Gamma,\mathbb Q_p(\chi))
 \longrightarrow H^1(Z,\mathbb Q_p(\chi))^{\Gamma/Z}.
\]
The first group is zero by induction,
since \(\Gamma/Z\) has smaller
nilpotency class. The last group is
zero too: it is the invariant part
of \(\operatorname{Hom}_{\rm cont}(Z,\mathbb Q_p)
\otimes\mathbb Q_p(\chi)\), on which
\(\Gamma\) acts only through the
nontrivial scalar \(\chi\), because
\(Z\) is central. This proves the
first assertion, including the abelian
base case.

For the second, write \(\mathfrak g\)
for the inertia Lie algebra and
\(\Theta\in\mathfrak g\otimes\mathbb C_p\)
for the Sen operator. It is semisimple
by Hodge--Tate comparison
\citep[Section~2.7]{FF_Fontaine2004}.
Hence \(\operatorname{ad}\Theta\)
is semisimple on
\(\operatorname{End}(S)\otimes\mathbb C_p\),
and on its invariant subspace
\(\mathfrak g\otimes\mathbb C_p\).
It is nilpotent on that subspace
because \(\mathfrak g\) is nilpotent.
It must therefore be zero. Thus
\(\Theta\in Z(\mathfrak g)\otimes\mathbb C_p\).
The minimal-subspace assertion of
\citet[Theorem~3.1, \(d=0\)]{FF_Ohkubo2014}
implies
\(\mathfrak g\subset Z(\mathfrak g)\).
So \(\mathfrak g\) is abelian.

Matrix exponential identifies a
sufficiently small open subgroup
of the compact inertia image with
a lattice in this abelian Lie algebra.
That open subgroup is abelian.
A finite extension of \(M\) makes
its inertia image lie in this subgroup:
choose an open subgroup of \(G_M\)
whose intersection with inertia is
inside the given open inertia subgroup.
Claim~\ref{clm:FFabelianHT} applies
to the restricted Hodge--Tate
representation and proves the
remaining assertions. This concerns
a finite extension, not an infinite
passage that might destroy the
finite local-field hypotheses.
\end{proof}

\begin{claim}{F.39}{Proved}\label{clm:FFnilpotentreconstruct}
Let \(I=G_{K_p}=I_{\mathbb Q_p}\),
and define its closed lower central
series by
\(\gamma_1(I)=I\) and
\(\gamma_{r+1}(I)=\overline{[I,\gamma_r(I)]}\).
For each integer \(c\ge1\), put
\[
              K_{p,c}=
                \overline{\mathbb Q}^{\,\gamma_{c+1}(I)}.
                                                        \tag{F.39.1}
\]
Every finite extension of \(K_{p,c}\),
and every subfield of such an extension,
satisfies \'etale reconstruction.
Here \(K_{p,1}=K_p^{\rm ab}\), and
every extension \(K_{p,c+1}/K_{p,c}\)
is infinite. The nilpotency class
\(c\) is fixed and finite in the
assertion; no conclusion about the
union of these fields or a general
solvable closure is included.
\end{claim}
\begin{proof}
Write \(J_c=\gamma_{c+1}(I)\).
It is characteristic in inertia,
hence normal in \(G_{\mathbb Q_p}\).
The quotient \(I/J_c\) is nilpotent
of class at most \(c\).
The field in (F.39.1) is exactly
the compositum of all finite Galois
extensions of \(K_p\) with Galois
group of nilpotency class at most
\(c\): the lower-central identities
pass to finite quotients, and the
intersection of the kernels of the
finite quotients of \(I/J_c\) is
trivial.

We first verify that these stages
are proper and form a strictly
increasing infinite hierarchy.
The free wild inertia group \(P\)
has \(H^2(P,\mathbb F_p)=0\).
The quotient \(I/P\) is pro-prime-to-\(p\),
whose positive-degree cohomology
on \(p\)-primary discrete modules
vanishes by averaging on finite
quotients. Hochschild--Serre gives
\[
                   H^2(I,\mathbb F_p)=0.
                                                        \tag{F.39.2}
\]
Also \(H^1(I,\mathbb F_p)\) is
infinite-dimensional by the principal-unit
calculation in Step~1 of
Claim~\ref{clm:FFmaxabelianreconstruct}.

Every finite \(p\)-group \(E\)
is a quotient of \(I\). Here are
the embedding-problem details.
Its Frattini quotient \(E/\Phi(E)\)
is a finite-dimensional \(\mathbb F_p\)-space,
so the preceding infinite-dimensional
\(H^1\) supplies a surjection from
\(I\) to that quotient. Choose a
central chief series inside
\(\Phi(E)\): each successive
factor has order \(p\) and is
central in the corresponding quotient
of \(E\). Such a series exists
because a nontrivial normal subgroup
of a finite \(p\)-group meets
its center nontrivially, and one
continues in the quotient. Each
successive central extension by
\(\mathbb F_p\) has zero pullback
obstruction by (F.39.2), so the
map from \(I\) lifts through all
of them. Its final image is all
of \(E\): a proper subgroup lies
in a maximal subgroup, and every
maximal subgroup contains
\(\Phi(E)\), contrary to surjectivity
on the Frattini quotient.

For each \(n\), apply this to
\(E=\operatorname{UT}_{c+2}(\mathbb F_p)^n\).
Each unitriangular factor has
nilpotency class \(c+1\), and
its \((c+1)\)-st lower-central
subgroup has order \(p\).
The nested commutator of the
successive elementary transvections
is \(1+E_{1,c+2}\), proving this
also in small characteristic.
It follows that
\(\gamma_{c+1}(I)/\gamma_{c+2}(I)\)
has quotients \((\mathbb Z/p)^n\)
for arbitrarily large \(n\).
It is infinite, proving the
asserted infinite stage extensions.

We will also need that \(J_c\)
and all its open subgroups are
infinite. Independently of the
preceding quotient argument,
\(J_c\) contains \(\gamma_{c+1}(P)\).
Map \(c+1\) free generators of
\(P\) to the successive elementary
transvections just used. Their
nested commutator is nonidentity,
so \(\gamma_{c+1}(P)\ne1\).
Torsion-freeness of \(P\) makes
it infinite. Every open subgroup
of \(J_c\) is infinite as well.

Now fix a finite extension
\(K/K_{p,c}\). It contains
\(K_p^{\rm ab}\), hence
\(U_p\) and \(\mu_\infty\).
Write \(K=L_0K_{p,c}\), with
\(L_0/K_p\) finite, enlarging
\(L_0\) to contain any specified
finite data. The Galois group
\(\operatorname{Gal}(K/L_0)\)
is a closed subgroup of the
nilpotent group \(I/J_c\),
so is nilpotent of class at most
\(c\), with infinite cyclotomic
image. Claim~\ref{clm:FFnilpotentinputs}(1)
replaces the abelian cocycle identity
in the proof of (F.37.2).
It proves the same exact integral
and rational Kummer kernels
\(\mu_\infty\), because an
infinitely \(p\)-divisible element
has its compatible Kummer cocycle
factoring through this nilpotent
quotient. Rational triviality there
implies that a \(p\)-power of
the element belongs to the known
integral kernel over \(L_0\),
and hence the element is a
root of unity.

Choose a finite number model
\(D_0\subset K\), put
\(M_0=(D_0)_v\), and identify
\[
        G_K=H_0=J_c\cap G_{M_0}
                         \triangleleft G_{M_0}.
                                                        \tag{F.39.3}
\]
The quotient \(I_{M_0}/H_0\)
is nilpotent of class at most
\(c\) and has nontrivial
cyclotomic character. Assertion~(1)
again gives zero \(H^1\) of this
quotient with coefficients
\(\mathbb Q_p(1)\), while its
coefficient invariants are zero.
The two inflation--restriction
arguments in (F.37.4) prove
injectivity from both inertia
and the full local Galois group
to \(H_0\). The same statements
hold at every finite field and
finite local-model extension.

For a pointed generalized Albanese,
form the finite-dimensional space
\(S=\operatorname{Hom}_{H_0}(V_pA,\mathbb Q_p(1))\).
It is Hodge--Tate and its inertia
image is nilpotent of class at
most \(c\); its Lie algebra
is therefore nilpotent.
By Claim~\ref{clm:FFnilpotentinputs}(2),
a finite local extension makes
this inertia representation abelian
and semisimple. Realize that
extension on a finite algebraic
model, possibly passing to a
finite extension of \(K\).
Retain the original space \(S\)
under this restriction; it is
unnecessary to replace it by
a larger Hom space. The finite
extension is permissible in the
point-function descent argument.

The all-point-invisible subspace
\(N\subset S\) of Step~3 in
Claim~\ref{clm:FFmaxabelianreconstruct}
has the same definition and the
same normality and generation
properties. It is unchanged by
testing only finite extensions
of the enlarged base: any finite
extension of the original base
can be combined with that base,
and rational restriction--corestriction
preserves a nonzero evaluation.
The inertia detector of
Claim~\ref{clm:FFpdetector}, followed
by the just-proved restriction
injection, again gives
\(N\cap S^{I_{M_1}}=0\) at
every further finite local stage.
Actual point/value descent and
the canonical semisimple projection
therefore give the unique invariant
lift \(\lambda_0\) modulo \(N\),
exactly as in (F.37.7)--(F.37.8).
The projection remains compatible
with torus restrictions, open
restrictions and curve pullback;
in particular it retains the
original boundary residues.

All remaining inputs of the
reconstruction proof in Step~5
of Claim~\ref{clm:FFmaxabelianreconstruct}
have now been established: the
exact Kummer kernel, finite-local
restriction injection, unchanged
integral boundary residues, and
the infinite groups of every
finite extension needed for
large-degree points. Thus the
normalized rational class yields
a power of an algebraic unit,
the \(T,T-1\) resultant removes
that power, and \(T-\beta,T-a\)
remove the constant and root-of-unity
ambiguities in the actual point
function. Finite descent returns
the resulting unit to the
original \(K\), and
\citep[Theorem~4.18(5)]{FF_CHW2024}
gives \(\mathcal R(K)\).
Algebraic descent proves the
subfield assertion, since all
these fields are algebraic
over \(\mathbb Q\).
The proof uses a fixed bound
\(c\) in the cohomological
central-series induction and
the inertia Lie argument;
neither argument has been
applied to an unbounded union.
\end{proof}


\begin{claim}{F.40}{Proved}\label{clm:FFmetabeliantest}
Let \(I=G_{K_p}\), let \(I'=\overline{[I,I]}\),
let \(I''=\overline{[I',I']}\), and put
\[
                  B_p=\overline{\mathbb Q}^{\,I''}.
                                                        \tag{F.40.1}
\]
This maximal metabelian extension of \(K_p\)
is a proper countable algebraic subfield of
\(\overline{\mathbb Q}\), with henselian
value group \(\mathbb Q\), residue field
\(\overline{\mathbb F}_p\), and
\(\mu_\infty\subset B_p\).
For every finite extension \(E/B_p\),
\[
             C_{\operatorname{Aut}(\overline{\mathbb Q})}(G_E)=1.
                                                        \tag{F.40.2}
\]
It belongs to none of the sufficient field
families established in this section through
Claim~\ref{clm:FFnilpotentreconstruct}, including
the subfields of finite extensions of a
fixed bounded-nilpotency stage at any prime,
and to none of the negative centralizer-ghost
families. Moreover,
\[
 K_p^{{\rm ab},\times}\subset
                  \bigcap_{n\ge1}(B_p^\times)^{p^n}.
                                                        \tag{F.40.3}
\]
Thus the exact \(\mu_\infty\) Kummer-kernel
input used in Claims~\ref{clm:FFcyclotomicreconstruct}--
\ref{clm:FFnilpotentreconstruct} fails here.
No reconstruction failure is asserted.
\end{claim}
\begin{proof}
The field \(B_p\) contains
\(K_p^{\rm ab}\), hence \(U_p\).
The valuation, residue and roots-of-unity
assertions follow. It is countable and
of transcendence degree zero because
it is algebraic over \(\mathbb Q\).

We verify the needed nonabelian groups
explicitly. For any finite number model
\(M/\mathbb Q_p\), wild inertia
\(P_M\) is free pro-\(p\) of
countably infinite rank. Choose eight
free generators and map them to
\(1+E_{i,i+1}\), \(1\le i\le8\),
in \(\operatorname{UT}_9(\mathbb F_p)\).
The words
\[
 u=[[x_1,x_2],[x_3,x_4]],\qquad
 v=[[x_5,x_6],[x_7,x_8]]
\]
belong to \(P_M''\), and map to
\(1+E_{1,5}\) and \(1+E_{5,9}\).
Their commutator maps to
\(1+E_{1,9}\ne1\). Thus
\(P_M''\) is nonabelian for every
prime, including \(p=2\).
In particular \(I''\ne1\), proving
that \(B_p\ne\overline{\mathbb Q}\).

For a finite extension \(E/B_p\),
choose a number field \(D\subset E\)
with \(E=DB_p\), and write
\(M=D_v\). Its absolute group is
\[
                 H=G_E=I''\cap G_M
                             \triangleleft G_M.
                                                        \tag{F.40.4}
\]
It lies in \(P_M\) and contains
\(P_M''\), so is nonabelian.
Its centralizer in \(P_M\) is
trivial. Indeed, if \(a\ne1\)
centralized \(H\), the finitely
generated closed free pro-\(p\)
subgroup generated by \(a,u,v\)
would put two noncommuting elements
in the procyclic centralizer of
\(a\), contrary to the free-group
centralizer theorem used in
Claim~\ref{clm:FFcyclotomictest}.

The subgroup \(C_{G_M}(H)\) is
normal in \(G_M\), by (F.40.4).
Its commutator with \(P_M\) lies
in \(C_{G_M}(H)\cap P_M=1\),
so it centralizes \(P_M\).
The fixed field of \(P_M\) is
a finite extension of \(T_p\).
It reconstructs by
Claim~\ref{clm:FFtamereconstruct};
the necessary condition in
Claim~\ref{clm:F.6} therefore makes
the centralizer of \(P_M\) in
\(\operatorname{Aut}(\overline{\mathbb Q})\)
trivial. Thus \(C_{G_M}(H)=1\).

Any global automorphism centralizing
\(H\) preserves its fixed field
\(E\). The rank-one henselian
valuation of \(E\) is unique:
two distinct such valuations would
make it algebraically closed, by
the two-valuation argument proved
in Claim~\ref{clm:FFstricttest},
whereas \(H\) is nonabelian.
Its unique prolongation to
\(\overline{\mathbb Q}\) is
therefore preserved. The global
centralizing automorphism belongs
to \(G_{\mathbb Q_p}\); a positive
power belongs to the open subgroup
\(G_M\) and is then trivial.
Torsion-freeness of \(G_{\mathbb Q_p}\),
verified in Claim~\ref{clm:FFcyclotomictest},
makes the automorphism itself
trivial. This proves (F.40.2).

We next exclude the new nilpotent
families, including abstract field
embeddings. Let \(s=1\) for
odd \(p\), and \(s=2\) for
\(p=2\). For arbitrarily large
\(n\), the finite \(p\)-group
\[
 Q_n=\langle a,b:
       a^{p^n}=b^{p^n}=1,\quad
       bab^{-1}=a^{1+p^s}\rangle
                                                        \tag{F.40.5}
\]
is a semidirect product of two
cyclic groups of order \(p^n\).
The displayed action has order
dividing \(p^n\), so the
presentation is consistent.
It is metabelian, and for every
fixed \(c\),
\[
           \gamma_{c+1}(Q_n)
                  =\langle a^{p^{sc}}\rangle,
            \qquad
           |\gamma_{c+1}(Q_n)|=p^{n-sc}
                   \quad(n>sc).
                                                        \tag{F.40.6}
\]
Every finite \(p\)-group is a
quotient of \(I\), by the
Frattini lifting proof in
Claim~\ref{clm:FFnilpotentreconstruct}.
Since \(Q_n''=1\), these quotients
factor through \(I/I''\).
Consequently the image of
\(\gamma_{c+1}(I)\) in
\(I/I''\) is infinite, by
(F.40.6).
If \(B_p\) were contained in
a finite extension \(L/K_{p,c}\),
the subgroup \(G_L\) would be
open in \(\gamma_{c+1}(I)\)
and contained in \(I''\).
That would make this image
finite, a contradiction.

Suppose instead that an abstract
embedding of \(B_p\) had image
in a finite extension \(L/K_{q,c}\).
Extend it to an automorphism
\(g\) of \(\overline{\mathbb Q}\).
The transported henselian valuation
of \(gB_p\) extends to a henselian
valuation on the algebraic extension
\(L\). Its residue characteristic
is \(p\); the given valuation
of \(L\) has residue characteristic
\(q\). Since \(L\) is not
algebraically closed, uniqueness
of a rank-one henselian valuation
forces equality of the valuations
and \(p=q\). Their unique
prolongations to the algebraic
closure then give \(g\in G_{\mathbb Q_p}\).
Conjugation by this group preserves
both \(I''\) and the closed
lower central series of \(I\).
Thus \(gB_p=B_p\), reducing to
the already excluded inclusion.

For completeness, the preceding
positive families are excluded
by their already verified criteria.
The field contains \(\mu_\infty\),
which is a nontrivial divisible
multiplicative subgroup, so it
cannot be a Kummer-faithful
sub-\(q\)-adic field for any
\(q\). The values of
\(\zeta_{q^m}-1\) have unbounded
\(q\)-denominators at every
residue-characteristic-\(q\)
place, excluding the bounded-wild
families at all primes.
Its transcendence degree zero
excludes exactly the families with an explicit
positive-transcendence hypothesis in
Claim~\ref{clm:FFclasses}(3)--(4) and the
finitely generated transcendental fields
of Claim~\ref{clm:FFfunction}.
We do not assert arbitrary subfield
closure for a proper normal envelope
over a transcendental finitely generated field.

The proper-normal-envelope families
are excluded by the same complete
two-valuation argument as in
Claim~\ref{clm:FFstricttest}.
Explicitly, for any number field
\(D\) in a finite extension
\(E/B_p\), choose
\(g\in G_D\) moving the selected
valuation. Such a \(g\) exists
because the decomposition group
of a finite place of a number
field is proper; the local-square,
global-nonsquare CRT construction
in that claim proves this directly.
The fields \(E\) and \(gE\)
then give two distinct rank-one
henselian valuations on their
compositum. They remain distinct
there because otherwise their
unique prolongations to the
algebraic closure would coincide.
Hence \(EgE=\overline{\mathbb Q}\).
The normal closure of \(E/D\)
is therefore \(\overline{\mathbb Q}\).
Together with the precise normal-closure
equivalence in Claim~\ref{clm:FFclasses}(2),
this excludes the formation subfield
families as stated there. It also
excludes containment in a
proper normal algebraic extension
of any number field, after
adjoining that number field if
necessary. It also excludes being
finite over such a proper normal
field: two conjugate copies of
\(G_E\) would then be open in
its normal Galois group but have
trivial intersection, forcing
the infinite group \(G_E\) to
be finite. Finally (F.40.2)
excludes all the negative
centralizer-ghost families.

To prove (F.40.3), take
\(a\in K_p^{{\rm ab},\times}\).
Over \(K_p^{\rm ab}\), adjoin
all compatible \(p\)-power
roots of every conjugate of
\(a\) under \(I\). There are
only finitely many such conjugates,
since \(a\) is algebraic over
\(K_p\). All roots of unity
are already in the base, so
this is an abelian pro-\(p\)
extension of \(K_p^{\rm ab}\),
stable under \(I\). Its Galois
group over \(K_p\) is an
extension of an abelian group
by an abelian group and hence
metabelian. It lies in \(B_p\).
In particular \(a\) has
compatible \(p\)-power roots
in \(B_p\), proving (F.40.3).
The element \(a=p\) already
shows that the Kummer kernel
is larger than \(\mu_\infty\).
\end{proof}

The field \(B_p\), for example \(B_2\),
is consequently a specified remaining
test with the smallest possible
cardinality and transcendence degree
in characteristic zero. An
inclusion-minimal field is not
claimed. Whether \(\mathcal R(B_p)\)
holds remains \textbf{Open} in
this work. The failed exact
Kummer-kernel step identifies a
precise boundary of the preceding
proofs; it is not a topos map or
a refutation of centralizer
sufficiency.


\begin{claim}{F.41}{Proved}\label{clm:FFmetabelianbarrier}
The invariant-lift step of
Claim~\ref{clm:FFmaxabelianreconstruct}
fails for an actual elliptic curve over
the metabelian field \(B_p\), even
after every finite change of local
arithmetic model. More precisely, put
\(H=I''=G_{B_p}\), and take
\[
                  A:\quad y^2+xy=x^3-p.
                                                        \tag{F.41.1}
\]
There is a finite local field
\(M_0/\mathbb Q_p\) of degree at most
two, whose chosen global quadratic
model lies in \(K_p^{\rm ab}\), such
that the following holds. The space
\[
 S=\operatorname{Hom}_{H}
          (V_pA,\mathbb Q_p(1))
                                                        \tag{F.41.2}
\]
is two-dimensional. Let \(N\) be the
subspace of functionals whose pushforward
of the point Kummer class is zero for
every \(a\in A(E)\) and every finite
extension \(E/B_p\). Then for every
finite extension \(M/M_0\),
\[
       N\simeq\mathbb Q_p(1),\qquad
       S/N\simeq\mathbb Q_p,\qquad
       S^{I_M}=0,\qquad
       (S/N)^{I_M}=S/N\ne0.
                                                        \tag{F.41.3}
\]
The first two isomorphisms describe
the actual local Galois representations.
Consequently the visible invariant
class has no inertia-invariant lift.
This is a counterexample to a proposed
replacement lemma, not a counterexample
to \(\mathcal R(B_p)\).
\end{claim}
\begin{proof}
For the displayed integral Weierstrass
equation one computes
\[
             c_4=1,\qquad
             \Delta=p(1-432p),\qquad
             v_p(j)=-1.
                                                        \tag{F.41.4}
\]
In particular \(j\ne0,1728\).
Tate's convergent inverse to the
\(j\)-series gives \(q\in\mathbb Q_p\)
with \(v_p(q)=1\) and a Tate curve
\(E_q\) with this \(j\)-invariant.
Two characteristic-zero elliptic curves
with this \(j\)-invariant become
isomorphic over a quadratic extension:
their geometric automorphism group
is \(\{1,-1\}\), so the isomorphism
torsor is quadratic. Thus
\(A_{M_0}\simeq E_q\) for a field
\(M_0\) of degree at most two.
A square class in \(\mathbb Q_p^\times\)
has a rational representative, by
openness of squares and density of
\(\mathbb Q\). Hence \(M_0\) is
the completion of a chosen quadratic
number field \(D_0\). Since its
inertia character is abelian,
\(D_0\subset K_p^{\rm ab}\).
One could also use the split node
with tangent cone \(y(y+x)\) in
(F.41.1); allowing the quadratic
extension avoids needing a separate
split-reduction criterion.

The exact primary uniformization
input is Tate's Theorem~1, together
with the description of finite-order
points on pages~18--19 of the author
text \citep{FF_Tate1995}. It gives
the compatible local Galois sequences
\[
 \begin{gathered}
   \varphi:M^\times/q^{\mathbb Z}
                         \xrightarrow{\sim}A(M),\qquad
   0\longrightarrow\mathbb Z_p(1)
    \longrightarrow T_pA
    \longrightarrow\mathbb Z_p\longrightarrow0,\\
   0\longrightarrow\mathbb Q_p(1)
    \longrightarrow V_pA
    \longrightarrow\mathbb Q_p\longrightarrow0.
 \end{gathered}
                                                        \tag{F.41.5}
\]
The extension class of the last
sequence is the local Kummer class
\(\kappa_M(q)\). This use of \(q\)
is entirely local: we do not assert
that \(q\) is algebraic over
\(\mathbb Q\).

Every element of \(I'\) lies in
\(G_{M_0}\), since the quadratic
character factors through an abelian
quotient. Its cyclotomic character
is one. In a basis adapted to
(F.41.5), the image of \(I'\)
therefore consists of unipotent
upper-triangular matrices. They
commute, so \(H=I''\) acts trivially
on \(V_pA\). It also acts trivially
on \(\mathbb Q_p(1)\). Thus (F.41.2)
is the full two-dimensional Hom space.
Dualizing and twisting (F.41.5) gives
\[
 0\longrightarrow\mathbb Q_p(1)
  \longrightarrow S
  \xrightarrow{\ r\ }\mathbb Q_p
  \longrightarrow0,
                                                        \tag{F.41.6}
\]
where \(r\) is restriction to the
torus Tate line. Its extension class
is \(\kappa_M(q)\), up to sign.

For any finite \(M/M_0\), valuation
shows that \(\kappa_M(q)\) is nonzero
in \(H^1(G_M,\mathbb Q_p(1))\):
the valuation component in the
rationalized Kummer completion is
\(v_M(q)\ne0\). Restriction to
\(I_M\) is injective, by
inflation--restriction and
\(\mathbb Q_p(1)^{I_M}=0\).
Consequently (F.41.6) does not split
on \(I_M\). Its connecting map
from \(\mathbb Q_p\) is nonzero,
and the left term has no inertia
invariants. Therefore
\[
                         S^{I_M}=0.
                                                        \tag{F.41.7}
\]

We now determine the invisible space,
with all finite extensions included.
If \(r(\lambda)=0\), take any
finite \(E/B_p\) and \(a\in A(E)\).
Choose a number field \(D\subset E\)
containing \(D_0\) and the coordinates
of \(a\), and put \(M=D_v\).
Then \(G_E\subset G_M\), using
the chosen henselian valuation.
By (F.41.5), write
\(a=\varphi(u)\), \(u\in M^\times\).
Choose roots \(u_n\) recursively
with \(u_{n+1}^p=u_n\) and
\(u_n^{p^n}=u\). The points
\(\varphi(u_n)\) are coherent
\(p\)-power divisions of \(a\).
For \(g\in G_M\), their Kummer
cocycle is
\[
       g\varphi(u_n)-\varphi(u_n)
           =\varphi\bigl(g(u_n)/u_n\bigr),
       \qquad g(u_n)/u_n\in\mu_{p^n}.
                                                        \tag{F.41.8}
\]
It lies in the torus subgroup at
every level, compatibly. Hence
\(\lambda\) kills the full Tate-module
Kummer class, not merely its separate
finite-level images. Restricting to
\(G_E\) proves
\(\ker(r)\subset N\).

Conversely let \(r(\lambda)\ne0\).
We construct an actual point of
\(A(B_p)\) which detects it.
The group \(H\) is itself a closed
subgroup of wild inertia, because
\(I/P\) is abelian. It is nontrivial
by Claim~\ref{clm:FFmetabeliantest}.
Closed subgroups of a free pro-\(p\)
group are free pro-\(p\), so \(H\)
has a nontrivial continuous quotient
\(\mathbb F_p\). This uses freeness
of \(H\) itself, not merely the
presence of a free subgroup.
Kummer theory and \(\mu_p\subset B_p\)
give \(a\in B_p^\times\) which is
not a \(p\)-th power.
Since \(\mu_{p^\infty}\subset B_p\),
its compatible Kummer class is a
homomorphism
\[
       \kappa_H(a):H\longrightarrow\mathbb Z_p(1)
                                                        \tag{F.41.9}
\]
whose reduction modulo \(p\) is
nonzero. It remains nonzero after
rationalization.

We may normalize \(a\) to an
arbitrarily small principal unit
without losing this property.
All compatible roots of \(p\)
lie in \(B_p\): prime-to-\(p\)
roots lie in its tame subfield,
and the radical tower over
\(K_p^{\rm ab}\) is metabelian
over \(K_p\), as in (F.40.3).
For any \(s\in\mathbb Q\), a
chosen \(p^s\) has zero rational
Kummer class on \(H\), since an
integer multiple of that class
is a multiple of \(\kappa_H(p)=0\).
Divide \(a\) by \(p^{v_p(a)}\),
and divide its residue by the
matching prime-to-\(p\) root of
unity in \(K_p\). The result
\(u_0\in B_p\) has positive
valuation of \(u_0-1\) and the
same nonzero rational Kummer class.
Raising it to a sufficiently large
\(p\)-power produces
\(u=u_0^{p^m}\) as close to one
as desired, still with nonzero
rational Kummer class.

Choose a number field
\(D\subset B_p\) containing \(D_0\)
and \(u\), and put \(M=D_v\).
The auxiliary analytic point
\(a_0=\varphi(u)\in A(M)\) lies
arbitrarily close to the identity.
Equations (F.41.8)--(F.41.9) give
\[
        \lambda_*\kappa_H(a_0)
              =r(\lambda)\kappa_H(u)\ne0.
                                                        \tag{F.41.10}
\]
Here \(a_0\) need not have globally
algebraic coordinates, and is only
an intermediate local point.

After multiplying \(\lambda\) by
a nonzero integer, it is integral
on \(T_pA\). Choose \(n\) such
that the integral class in
(F.41.10) has nonzero image modulo
\(p^n\). The subgroup
\([p^n]A(M)\) is open: in a
characteristic-zero formal coordinate
at the identity, multiplication by
\(p^n\) has nonzero derivative and
is an analytic isomorphism between
sufficiently small neighborhoods.
Thus the finite-level Kummer class
is constant on an open neighborhood
of \(a_0\).

An etale coordinate at the rational
identity, defined over \(D\), shows
that \(A(D^h)\) is dense in the
small identity neighborhood of
\(A(M)\). Indeed choose its coordinate
in the dense subfield \(D\); Hensel's
lemma supplies the inverse point
over the chosen henselization
\(D^h\). Since \(B_p\) is henselian
and contains \(D\), this chosen
\(D^h\) lies in \(B_p\). We can
therefore choose
\(a'\in A(D^h)\subset A(B_p)\)
with \(a'-a_0\in[p^n]A(M)\).
The nonzero finite-level evaluation
in (F.41.10) persists for \(a'\).
On \(H\), all the representations
are trivial, so a nonzero integral
Kummer homomorphism into the
torsion-free module \(\mathbb Z_p(1)\)
remains nonzero rationally.
Thus \(\lambda\notin N\).

We have proved \(N=\ker(r)\).
Equations (F.41.6)--(F.41.7)
now give all of (F.41.3).
The accompanying standard-library program
checks the invariant polynomial identity
(F.41.4) and the three unitriangular
commutator identities of
Claim~\ref{clm:FFmetabeliantest} over
\(\mathbb Z\), hence at every prime.
It does not compute a global Tate
parameter or certify the infinite arguments.
\end{proof}

The most direct ghost attempt from
this example would assign to a
point its multiplicative Tate
uniformization parameter. Two
requirements remain unproved.
First the analytic parameter of
a globally algebraic point need
not be known to lie in the global
field \(B_p\); the uniformization
theorem only places it in a finite
\(p\)-adic completion. Second, a
homomorphism to the geometric
Tate module is not a map of full
etale topoi. Such a map must also
give compatible maps after removing
every finite set of target points.
A chosen analytic fundamental
domain does not supply these
Zariski-open restrictions. We have
constructed neither the required
actual point function nor this
compatible family of morphisms.
The constant-point construction is
already excluded by
Claim~\ref{clm:FFstrictdetection}.


\begin{claim}{F.42}{Proved}\label{clm:FFmetabeliankernel}
Put \(F_p=K_p^{\rm ab}\), and define
\[
 W_p=\{a\in B_p^\times:
                  a^m\in F_p^\times
                  \text{ for some integer }m\ge1\}.
                                                        \tag{F.42.1}
\]
For every finite extension \(E/B_p\),
let \(\kappa_E\) denote the rational
\(p\)-primary Kummer map. Then
\[
 \ker\bigl(\kappa_E:E^\times
             \longrightarrow H^1(G_E,\mathbb Q_p(1))\bigr)
       =\bigcap_{n\ge1}(E^\times)^{p^n}=W_p.
                                                        \tag{F.42.2}
\]
In particular no new element in a
finite extension enlarges this kernel.
The subgroup recovers its abelian
constant field by the equation
\[
                     W_p\cap(1-W_p)=F_p\setminus\{0,1\}.
                                                        \tag{F.42.3}
\]
There are two precise value-detection
consequences.
\begin{enumerate}
\item If \(a,b\in E\setminus\{0,1\}\)
have the same pair
\((\kappa_E(a),\kappa_E(1-a))\),
then \(a=b\) or both \(a,b\in B_p\).
Thus the pair is injective outside
\(B_p\), with no collision between
such a value and a value in \(B_p\).
\item If \(a,b\in E\) satisfy
\[
                  \kappa_E(a-c)=\kappa_E(b-c)
                  \quad\text{for every }c\in\mathbb Q
                  \text{ with }a,b\ne c,
                                                        \tag{F.42.4}
\]
then \(a=b\) or both \(a,b\in F_p\).
Conversely any two values in \(F_p\)
satisfy (F.42.4). Thus all rational
target translates detect every value
outside \(F_p\) exactly.
\end{enumerate}
\end{claim}
\begin{proof}
Let \(a\in E^\times\) have zero
rational Kummer class. Choose a finite
Galois number field \(D/\mathbb Q\)
containing \(a\) and finitely many
generators for \(E/B_p\). We may
replace \(E\) temporarily by
\(DB_p\), since restriction preserves
the zero class. Its local completion
\(M=D_v\) is a finite Galois extension
of \(\mathbb Q_p\). With the fixed
embedding into the local algebraic
closure, put
\[
 \begin{gathered}
 I=G_{K_p},\qquad J=I\cap G_M,\qquad
 H=I''\cap J=G_{DB_p},\\
 \Gamma=J/H,\qquad C=(I'\cap J)/H.
 \end{gathered}
                                                        \tag{F.42.5}
\]
The Galois choice of \(M\) makes
\(J\) and \(H\) normal in \(I\),
so every conjugation used below acts
on this same model. The group \(C\)
is abelian, the quotient \(\Gamma/C\)
is abelian, and the cyclotomic
character \(\chi\) is trivial on
\(C\) and has infinite image on
\(\Gamma/C\). The latter follows
because \(J\) is open in \(I\).

The class of \(a\) on \(J\)
inflates from \(H^1(\Gamma,\mathbb Q_p(1))\),
since its restriction to \(H\) is
zero. The abelian cocycle identity
used in Claim~\ref{clm:FFcyclotomiccohomology}
gives
\[
                 H^1(\Gamma/C,\mathbb Q_p(1))=0.
\]
Inflation--restriction therefore makes
the map
\[
 H^1(\Gamma,\mathbb Q_p(1))
       \longrightarrow H^1(C,\mathbb Q_p(1))
                                                        \tag{F.42.6}
\]
injective. Every \(g\in I'\)
centralizes \(C\): the commutator
of \(g\) with an element of
\(I'\cap J\) lies both in
\(I''\) and in \(J\), hence
in \(H\). Also \(\chi(g)=1\).
Thus \(g\) acts trivially on the
right side of (F.42.6), and hence
on its left side. Functoriality of
Kummer classes gives
\[
                         \kappa_J(g(a)/a)=0
                         \qquad(g\in I').
                                                        \tag{F.42.7}
\]

All elements \(g(a)/a\) lie in
the same finite field \(M\).
Restriction
\(H^1(G_M,\mathbb Q_p(1))
\longrightarrow H^1(J,\mathbb Q_p(1))\)
is injective: \(J=I_M\) and
\(\mathbb Q_p(1)^J=0\), so
inflation--restriction has zero
kernel. The rational Kummer kernel
of \(M^\times\) is precisely its
finite group of roots of unity.
Indeed valuation detects the
uniformizer component, and on a
sufficiently small principal-unit
group the \(p\)-adic logarithm
identifies the remaining rationalized
Kummer completion with the additive
local field. It follows from
(F.42.7) that
\[
                            g(a)/a\in\mu(M)
                            \qquad(g\in I').
                                                        \tag{F.42.8}
\]
Choose a single integer \(m\)
annihilating the finite group
\(\mu(M)\). Then \(a^m\) is
fixed by every \(g\in I'\).
Since \(a\) is globally algebraic,
\(a^m\in\overline{\mathbb Q}^{I'}=F_p\).
Moreover \(\mu_m\subset F_p\),
so \(F_p(a)/F_p\) is cyclic.
The field \(B_p\) is the maximal
abelian extension of \(F_p\), with
groups \(G_{F_p}=I'\) and
\(G_{B_p}=I''\). Consequently
\(a\in B_p\), and hence
\(a\in W_p\).

Conversely suppose \(a^m\in F_p^\times\).
For every \(n\), a root \(a_n\)
with \(a_n^{p^n}=a\) satisfies
\(a_n^{mp^n}=a^m\). Because
\(F_p\) contains all roots of
unity, the extension
\(F_p(a_n)/F_p\) is cyclic and
therefore lies in \(B_p\).
Choosing these roots recursively
gives compatible divisions of \(a\)
in \(B_p\), so its integral and
rational Kummer classes vanish
over every \(E\).
Finally \(G_E\) acts trivially on
\(\mathbb Z_p(1)\), making the
integral Kummer class a homomorphism
to a torsion-free module. Its
rationalization vanishes exactly
when it vanishes integrally. This
is equivalent to divisibility by
every \(p^n\); compatible choices
exist since all \(p\)-power roots
of unity are already in \(E\).
This proves (F.42.2).

For (F.42.3), the inclusion of
\(F_p\setminus\{0,1\}\) is immediate.
Conversely let \(x\in W_p\) and
\(1-x\in W_p\). For every
\(g\in I'\), there are roots
of unity \(\zeta,\eta\in F_p\)
such that
\[
                    g(x)=\zeta x,\qquad
                    g(1-x)=\eta(1-x).
\]
If \(g(x)\ne x\), their sum gives
\((\zeta-\eta)x=1-\eta\).
The case \(\zeta=\eta\) would
force both to equal one, contrary
to the choice of \(g\). Otherwise
the equality puts \(x\) in \(F_p\),
also contradicting \(g(x)\ne x\).
Thus every \(g\in I'\) fixes
\(x\), and \(x\in F_p\).

For the first value-detection assertion,
equality of the two classes puts
\[
                 \alpha=a/b\in W_p,\qquad
                 \beta=(1-a)/(1-b)\in W_p.
\]
If \(\alpha=\beta\), summing the
two defining equalities gives
\(\alpha=1\), so \(a=b\).
Otherwise
\(b=(1-\beta)/(\alpha-\beta)\)
and \(a=\alpha b\) both belong
to \(B_p\).

For the second assertion assume
\(a\ne b\). Choose distinct
rational \(c_0,c_1\) avoiding
\(a,b\). The two ratios
\((a-c_i)/(b-c_i)\in W_p\subset B_p\)
are distinct; if equal, subtraction
would make both one and give
\(a=b\). Solving the two linear
equations therefore puts
\(a,b\in B_p\).
Choose a finite Galois number field
\(D\) containing them. For every
\(g\in I'\), consider
\[
 R_g(T)=\frac{(g(a)-T)(b-T)}{(g(b)-T)(a-T)}
                                 \in D(T).
                                                        \tag{F.42.9}
\]
For each rational \(c\) away
from its finitely many exceptional
values, (F.42.4) puts
\(\alpha_c=(a-c)/(b-c)\) in
\(W_p\). A power of \(\alpha_c\)
is fixed by \(g\), so
\(R_g(c)=g(\alpha_c)/\alpha_c\)
is a root of unity. It belongs
to the same number field \(D\).
There are only finitely many roots
of unity in \(D\), so one value
occurs for infinitely many \(c\).
A rational function with infinitely
many occurrences of one value is
constant. Its value at infinity
is one, hence \(R_g(T)=1\).
Comparing coefficients gives
\[
              g(a)-a=g(b)-b=\delta,\qquad
              \delta(b-a)=0.
\]
Since \(a\ne b\), both are
fixed by \(g\). This holds
for every \(g\in I'\), proving
\(a,b\in F_p\). The converse
follows directly from (F.42.2).
\end{proof}

This supplies exact value detection
from the collection of all rational
target translates, despite the
failure of the single-functional
invariant lift in
Claim~\ref{clm:FFmetabelianbarrier}.
It does not yet prove reconstruction:
one must still derive sufficiently
coherent arithmetic control of those
translated Kummer maps from an
actual pinned topos map, and handle
the values in \(F_p\) which the
displayed family does not distinguish.
Nor does it make an analytic Tate
parameter globally algebraic. The
local version of the same argument
has the local fixed field
\(\overline{\mathbb Q}_p^{I'}\)
as its constants; that field already
contains \(\mathbb Q_p\) and the
local parameter \(q\) in
Claim~\ref{clm:FFmetabelianbarrier}.


\begin{claim}{F.43}{Proved}\label{clm:FFmetabelianpointdetection}
Let \(E/B_p\) be finite and let
\(a,b\in E\). There is a proper
finite extension \(L/E\) and
\(\beta\in L\setminus E\) such
that equality of the two pairs
\[
 \begin{split}
 &\bigl(\kappa_L(a-\beta),
                  \kappa_L(1-a+\beta)\bigr),\\
 &\bigl(\kappa_L(b-\beta),
                  \kappa_L(1-b+\beta)\bigr)
 \end{split}
                                                        \tag{F.43.1}
\]
implies \(a=b\). In fact every
proper finite extension and every
\(\beta\in L\setminus E\) work.

Consequently the following criterion
detects the actual point function of
a pinned map. Let \(X/B_p\) be
of finite type and let
\(f,g:X_{\acute et}\to\mathbb A^1_{B_p,\acute et}\)
be pinned morphisms over \(BG_{B_p}\).
For every finite \(L/B_p\) and
\(\beta\in L\), set
\[
 Y_\beta=\mathbb A^1_L\setminus\{\beta,\beta+1\},
 \qquad
 U_\beta=f_L^{-1}(Y_\beta)\cap g_L^{-1}(Y_\beta)
                         \subset X_L.
                                                        \tag{F.43.2}
\]
Suppose the pullbacks by \(f_L\)
and \(g_L\) of the two Kummer
classes of \(T-\beta\) and
\(\beta+1-T\) agree after restriction
to \(U_\beta\). Then the two
actual geometric point functions
are equal, including their values
in \(K_p^{\rm ab}\).
This is a uniqueness criterion for
point functions; it neither constructs
the comparison morphism nor asserts
that equality of point functions
alone identifies arbitrary topos maps.
\end{claim}
\begin{proof}
Every finite extension of \(B_p\)
has a nontrivial absolute Galois
group, by the free-group argument
in Claim~\ref{clm:FFmetabeliantest}.
Thus a proper finite extension
\(L/E\) exists. If
\(\beta\in L\setminus E\), both
\(a-\beta\) and \(b-\beta\)
lie outside \(E\), hence outside
\(B_p\); neither equals zero or
one. The first value-detection
assertion of
Claim~\ref{clm:FFmetabeliankernel},
applied over \(L\), proves the
first conclusion.

For the criterion, take any
\(x\in X(\overline{\mathbb Q})\).
Choose a finite \(E/B_p\) over
which \(x\) is defined. Pinning
and the retained relative structure
give actual values
\(a=f(x),b=g(x)\in E\), by
\citep[Proposition~3.6]{FF_CHW2024}.
Choose \(L/E\) and \(\beta\)
as above. The point \(x_L\)
lies in both inverse-image opens
in (F.43.2), since neither \(a\)
nor \(b\) is \(\beta\) or
\(\beta+1\). Evaluating the two
assumed cohomology equalities at
\(x_L\) gives (F.43.1), so
\(f(x)=g(x)\). This argument
works at every geometric point
and makes no exception for values
in \(K_p^{\rm ab}\).
\end{proof}

The criterion uses genuine operations
available for pinned morphisms:
finite scalar extension, inverse images
of Zariski opens, and pullback and
evaluation of Kummer classes. Its
remaining hypothesis is substantial.
The visible nonsplit functional in
Claim~\ref{clm:FFmetabelianbarrier}
does not by itself produce compatible
equalities for these translated pairs,
or a scheme morphism with which to
compare them. Deriving those equalities
is the unresolved existence step.


\begin{claim}{F.44}{Proved}\label{clm:FFmetabelianaffine}
Let \(E/B_p\) be finite. Every
pinned automorphism of
\(\mathbb A^1_{E,\acute et}\) over
\(BG_E\), with its relative
structural isomorphism retained,
has actual geometric point function
\[
                          x\longmapsto ux+v
                          \qquad
                          (u\in E^\times,\ v\in E).
                                                        \tag{F.44.1}
\]
The two coefficients are unique.
This is an existence theorem for
the actual point function in the
automorphism case. It does not
assert that the entire topos
automorphism equals the one induced
by (F.44.1), and does not classify
arbitrary pinned morphisms.
\end{claim}
\begin{proof}
Write \(f:\overline{\mathbb Q}
\to\overline{\mathbb Q}\) for
the actual point function. It is
a bijection, because the inverse
topos automorphism supplies its
inverse. It is \(G_E\)-equivariant,
and for every finite \(L/E\)
it restricts to a bijection of
\(L\), by pinned relative point
evaluation \citep[Proposition~3.6]{FF_CHW2024}.
In particular \(f(0),f(1)\in E\)
are distinct. Postcompose with
the affine scheme automorphism
\[
                   T\longmapsto
                   \frac{T-f(0)}{f(1)-f(0)}.
\]
It suffices to prove that the
new point function is the identity.
Henceforth
\[
                            f(0)=0,\qquad f(1)=1.
                                                        \tag{F.44.2}
\]

\emph{The cohomology of each puncture.}
Let \(a\in\overline{\mathbb Q}\),
and choose a finite extension
\(L/E\) containing \(a\). Then
\(f(a)\in L\). The scalar-extended
topos automorphism restricts to an
equivalence
\[
 (\mathbb A^1_L\setminus\{a\})_{\acute et}
       \longrightarrow
 (\mathbb A^1_L\setminus\{f(a)\})_{\acute et}.
                                                        \tag{F.44.3}
\]
Indeed its bijective geometric point
function takes precisely \(a\)
to \(f(a)\), so the inverse-image
open is precisely the displayed
source. Geometric continuous
\(p\)-adic cohomology in degree
one, with coefficients
\(\mathbb Z_p(1)\), is free of
rank one, with basis
\(\kappa(T-a)\). The equivalence
(F.44.3) therefore acts on that
integral geometric line by a unit
\(\alpha_a\in\mathbb Z_p^\times\).
The unit is unchanged by further
finite scalar extension; equivalently
one may compute after the algebraic
scalar extension to
\(\overline{\mathbb Q}\), obtained
as the inverse limit of finite
etale base extensions.

The Hochschild--Serre sequence and
the global lift \(\kappa(T-a)\)
show that the pullback of
\(\kappa(T-f(a))\), with rational
coefficients, is
\[
                         \alpha_a\kappa(T-a)+\eta_{a,L},
           \qquad \eta_{a,L}\in H^1(G_L,\mathbb Q_p(1)).
                                                        \tag{F.44.4}
\]
The relative structural isomorphism
identifies the coefficient system
and its base pullback here. Evaluating
at any \(x\in L\setminus\{a\}\)
gives
\[
       \kappa_L(f(x)-f(a))
          =\alpha_a\kappa_L(x-a)+\eta_{a,L}.
                                                        \tag{F.44.5}
\]
These equations always compare points
and classes over one common finite
base \(L\).

\emph{Recovering the affine lines.}
Put \(F=K_p^{\rm ab}\). For distinct
\(a,b\) and \(x\notin\{a,b\}\)
in \(\overline{\mathbb Q}\), choose
one finite \(L/E\) containing all
three. Subtracting the evaluations
at \(x\) and \(b\) in (F.44.5)
cancels the arithmetic term and gives
\[
 \kappa_L\!\left(\frac{f(x)-f(a)}{f(b)-f(a)}\right)
     =\alpha_a
          \kappa_L\!\left(\frac{x-a}{b-a}\right).
                                                        \tag{F.44.6}
\]
The corresponding comparison using
the puncture \(b\) gives
\[
 \kappa_L\!\left(\frac{f(b)-f(x)}{f(b)-f(a)}\right)
     =\alpha_b
          \kappa_L\!\left(\frac{b-x}{b-a}\right).
                                                        \tag{F.44.7}
\]
The coefficients are nonzero. By
Claim~\ref{clm:FFmetabeliankernel},
both classes on the right vanish
exactly when \((x-a)/(b-a)\in F\).
Both on the left vanish exactly
when \((f(x)-f(a))/(f(b)-f(a))\in F\).
Including the endpoints, and using
the inverse automorphism for the
reverse inclusion, we obtain
\[
                 f\bigl(a+(b-a)F\bigr)
                     =f(a)+(f(b)-f(a))F.
                                                        \tag{F.44.8}
\]
Thus \(f\) preserves all affine
\(F\)-lines, in both directions.

\emph{An elementary affine-geometry argument.}
We give the argument also in this
infinite-dimensional setting. Regard
\(\overline{\mathbb Q}\) as a vector
space over \(F\). Its dimension
is at least two. A bijection fixing
zero and preserving lines maps any
plane spanned by independent vectors
\(x,y\) onto the plane spanned by
\(f(x),f(y)\). To see containment,
each point off the two axes is on
a line joining a point of \(Fx\)
to a point of \(Fy\); for example
\(sx+ty\) is on the line joining
\(2sx\) and \(2ty\). Use the
inverse bijection for equality.
Noncollinearity is preserved as well.
Within such a plane, parallel lines
are characterized by disjointness,
so parallelism is preserved.
The parallelogram intersections give
\[
                      f(x+y)=f(x)+f(y)
                   \quad\text{if }x,y\text{ are independent.}
                                                        \tag{F.44.9}
\]
For dependent vectors, introduce a
vector \(z\) outside their line.
First apply (F.44.9) to
\(z=x+(-x+z)\) and to \(-x+z\)
to obtain \(f(-x)=-f(x)\).
If \(x+y\ne0\), compute
\(f(x+y+z)\) both as
\(f(x+y)+f(z)\) and as
\(f(x)+f(y+z)\). This gives
additivity also for dependent vectors;
the zero cases follow immediately.

For each nonzero vector \(x\),
line preservation defines a bijection
\(\sigma_x:F\to F\) by
\(f(cx)=\sigma_x(c)f(x)\).
For independent \(x,y\), compare
the two additive expansions of
\(f(c(x+y))\) to see
\(\sigma_x=\sigma_y=\sigma_{x+y}\).
For dependent nonzero vectors,
use a third independent vector.
Thus a single \(\sigma\) works
for every vector. Additivity of
\(f\) makes \(\sigma\) additive,
and comparing \(f(cd x)\) in
two ways makes it multiplicative.
It is bijective and takes one
to one, hence is an automorphism
of \(F\). In view of (F.44.2),
\[
      f(x+y)=f(x)+f(y),\qquad
      f(cx)=\sigma(c)f(x)\ (c\in F),\qquad
      f|_F=\sigma.
                                                        \tag{F.44.10}
\]
In particular \(f(F)=F\) and
\(f\) fixes \(\mathbb Q\).
Only two-dimensional spans were
used, so there is no finite-dimension
assumption hidden in this argument.

\emph{Recovering multiplication.}
In (F.44.5) take \(a=0\).
Evaluating at \(x=1\) and using
(F.44.2) gives \(\eta_{0,L}=0\).
Consequently, on every common finite
extension \(L/E\),
\[
                            \kappa_L(f(x))
                                   =\alpha_0\kappa_L(x).
                                                        \tag{F.44.11}
\]
Let \(x,y\ne0\), and choose
one such \(L\) containing both.
It also contains all their images
under \(f\). For each rational
\(c\) with \(x+c\ne0\),
multiplicativity of Kummer classes
and (F.44.11) yield
\[
       \kappa_L(f((x+c)y))
                     =\kappa_L(f(x+c))+\kappa_L(f(y)).
\]
Using (F.44.10), this becomes
\[
       \kappa_L\!\left(\frac{f(xy)}{f(y)}+c\right)
                     =\kappa_L(f(x)+c).
                                                        \tag{F.44.12}
\]
The equalities hold for all rational
translates at which both sides
are defined. The final assertion
of Claim~\ref{clm:FFmetabeliankernel}
therefore gives either
\[
                        \frac{f(xy)}{f(y)}=f(x)
\]
or both displayed values lie in
\(F\). If \(x\notin F\),
then \(f(x)\notin F\), so
only the first alternative is
possible. If \(x\in F\),
semilinearity already gives
\(f(xy)=\sigma(x)f(y)=f(x)f(y)\).
Zero factors are harmless. Hence
\(f\) is multiplicative everywhere.

The normalized point function is
therefore a field automorphism of
\(\overline{\mathbb Q}\).
Its \(G_E\)-equivariance says exactly
that it centralizes \(G_E\).
Claim~\ref{clm:FFmetabeliantest}
applies to every finite \(E/B_p\)
and makes this centralizer trivial.
Thus the normalized \(f\) is
the identity. Undoing the affine
normalization gives (F.44.1), with
\(u=f_{
m original}(1)-f_{
m original}(0)\)
and \(v=f_{
m original}(0)\).
\end{proof}

The unresolved automorphism kernel
now has a precise form: after
composition with a unique affine
scheme automorphism, an additional
pinned automorphism, if one exists,
must fix every geometric point.
The proof above does not identify
its action on the entire etale
site, and does not provide the
missing existence theorem for
noninvertible pinned maps.


\subsection{Faithfulness of generic restriction on the pointwise kernel}

\begin{claim}{F.45}{Proved}\label{clm:FFgenerickernel}
Let \(k\) be a characteristic-zero field, let \(X/k\)
be a smooth geometrically integral curve, and put
\(F=k(X)\). A pinned automorphism
\(\varphi:X_{\mathrm{\acute et}}\to
X_{\mathrm{\acute et}}\) over \(BG_k\) induces a generic
outer automorphism
\[
 \alpha_\varphi\in
 \operatorname{Aut}_{G_k}(G_F)/
              \operatorname{Inn}(G_{\bar k(X)}).
                                                        \tag{F.45.1}
\]
If its actual function on \(X(\bar k)\) is the identity
and (F.45.1) is trivial, then \(\varphi\) is isomorphic
to the identity over \(BG_k\).

Consequently, for every finite \(E/B_p\), the remaining
pointwise kernel in Claim~\ref{clm:FFmetabelianaffine}
injects into the relative birational outer automorphism
group of \(G_{E(T)}\). Its image preserves the geometric
conjugacy class of every finite-point inertia subgroup,
and its induced map on the corresponding residue Galois
group is the identity. In particular, an inner generic
action cannot hide a nontrivial automorphism of the full
etale topos.
\end{claim}

\begin{proof}
Write \(j:\operatorname{Spec}F\to X\), and choose a
separable closure of \(F\) containing \(\bar k\).
We retain all structural isomorphisms over \(BG_k\).

\emph{The generic action.}
An automorphism preserves the unique generic point and
therefore permutes the nonempty open subtopoi. The filtered
category of quasicompact separated etale objects over
these opens, with objects and arrows identified after
further shrinking, is the category of finite etale
\(F\)-schemes. Indeed, every finite separable extension
of \(F\), and every morphism between finitely many such
extensions, spreads out over a nonempty open of \(X\).
Conversely, two morphisms that agree generically agree
after shrinking. Thus \(\varphi\) induces an
autoequivalence of finite continuous \(G_F\)-sets.
The specified relative structure fixes the subcategory
of sets inflated from \(G_k\), with its specified
comparison. Galois-category reconstruction gives exactly
(F.45.1), with inner ambiguity only in the geometric
kernel.

Assume from now on that the geometric point function is
the identity and (F.45.1) is trivial. Choose the resulting
natural identity for the generic inverse-image functor,
over the category of \(G_k\)-sets. It extends from finite
continuous \(G_F\)-sets to all continuous \(G_F\)-sets:
each such set is a union of its finite orbits.

\emph{The maximal etale extension of a generic object.}
For a finite continuous \(G_F\)-set \(S\), the sheaf
\(j_*S\) is represented by the etale locus of the finite
normalization of \(X\) in the finite etale \(F\)-algebra
corresponding to \(S\). The normalization is finite
because \(X\) is a finite-type regular curve over a
field. To check the assertion, an \(F\)-map from the
generic fibre of an etale normal \(X\)-scheme extends
uniquely to its normalization map. At every point in its
image the extension of discrete valuation rings has
ramification index one and separable residue extension,
so the image lies in the etale locus. This gives precisely
the defining sections of \(j_*S\).

Every quasicompact separated etale \(X\)-scheme \(U\)
with generic fibre \(S\) is an open subobject
\[
                           U\ \lhook\joinrel\longrightarrow j_*S.
                                                        \tag{F.45.2}
\]
For example, this follows by applying the same
normalization argument to each component: a separated
quasi-finite etale map into a normal curve is an open
subscheme of its finite normalization. All components of
\(U\) dominate \(X\), since etale maps are open.
The complement in the etale normalization locus consists
of finitely many closed points.

For a geometric point \(\bar x\) above a closed point
\(x\), choose a prolongation of its valuation to the
separable closure of \(F\). Let \(D_x\) and \(I_x\)
be the corresponding decomposition and inertia groups.
Then
\[
               (j_*S)_{\bar x}=S^{I_x},\qquad
               D_x/I_x=G_{\kappa(x)}.
                                                        \tag{F.45.3}
\]
The stalk \(U_{\bar x}\) in (F.45.2) is a
\(D_x\)-stable subset of \(S^{I_x}\). This description
also follows directly by base change to the strict
henselization at \(\bar x\): its unramified sheets are
exactly the \(I_x\)-fixed generic sheets.

\emph{The relevant normalizer.}
In this setting
\[
                            N_{G_F}(I_x)=D_x.
                                                        \tag{F.45.4}
\]
Indeed, the fixed field \((F^{\rm sep})^{I_x}\) is the
fraction field of the strict henselization at the chosen
place. It has a henselian discrete rank-one valuation,
separably closed residue field of characteristic zero,
and value group \(\mathbb Z\). It is not separably
closed: a uniformizer has no square root in it. An element
normalizing \(I_x\) transports this valuation to another
rank-one henselian valuation on the same field. If the
two valuations were distinct, they would be incomparable,
and the two-valuation argument proved in
Claim~\ref{clm:FFstricttest} would make the field separably
closed. They therefore coincide. Its unique prolongation
to \(F^{\rm sep}\) is then preserved as well, so the
normalizing element belongs to \(D_x\). The reverse
inclusion is the usual normality of inertia in the
decomposition group.

\emph{The specialization comparison is forced.}
The chosen generic identity and the adjunction
\(j^*\dashv j_*\), together with invertibility of
\(\varphi^*\), give natural identifications
\[
                             \varphi^*j_*S\simeq j_*S.
                                                        \tag{F.45.5}
\]
Actual pointwise identity gives an isomorphism
\(\bar x^*\varphi^*\simeq\bar x^*\), by
\citet[Proposition~3.6]{FF_CHW2024}. Combining it with
(F.45.5) yields a natural automorphism of the functor
\[
       \{\text{finite continuous }G_F\text{-sets}\}
          \longrightarrow\{\text{sets}\},\qquad
                            S\longmapsto S^{I_x}.
                                                        \tag{F.45.6}
\]
There is no arbitrary specialization permutation here.
The functor (F.45.6) is prorepresented by the profinite
\(G_F\)-set \(G_F/I_x\). Concretely, write it as the
inverse system of finite \(G_F\)-sets
\(G_F/(I_xN)\), where \(N\) ranges over open normal
subgroups. A continuous map from this system to a finite
\(G_F\)-set is exactly an \(I_x\)-fixed element.
Pro-Yoneda therefore identifies the natural automorphisms
of (F.45.6) with
\[
                      N_{G_F}(I_x)/I_x=D_x/I_x.
                                                        \tag{F.45.7}
\]
Equivalently, the comparison acts on every \(S^{I_x}\)
by one element of \(D_x/I_x\). Hence it preserves
every \(D_x\)-stable subset, in particular
\(U_{\bar x}\). If the relative point comparisons are
chosen explicitly over \(BG_k\), their action on sets
inflated from \(G_k\) is the identity; this even makes
the element in (F.45.7) trivial, since
\(G_{\kappa(x)}\to G_k\) is injective. Preservation of
the subsets is all that is needed.

It follows from (F.45.5)--(F.45.7) that \(\varphi^*\)
fixes the subobject (F.45.2) inside \(j_*S\): its
generic stalk is fixed by construction and all its closed
geometric stalks are fixed by the preceding paragraph.
Geometric points detect equality of subobjects. This
produces an isomorphism \(\varphi^*U\simeq U\) for
every quasicompact separated etale \(U\). These
isomorphisms are natural. A morphism between two such
objects is determined by its generic restriction, because
the source is reduced with all components dominant and
the target is separated. Both ways around the proposed
naturality square agree generically and hence everywhere.

The quasicompact separated etale objects generate the
etale topos, and automorphisms preserve this subcategory;
see \citet[Theorem~A.8, Corollaries~A.9 and~A.11]
{FF_CHW2024}. Thus the natural identity extends to the
whole inverse-image functor. It respects the specified
structure over \(BG_k\): for a finite \(G_k\)-set the
corresponding constant etale cover is already its own
maximal extension \(j_*S\), and the generic identity
was chosen with precisely that relative compatibility.
This proves the first assertion.

For the final assertion take \(X=\mathbb A^1_E\).
The same specialization calculation, without assuming
the generic action trivial, identifies the prorepresented
functors attached to \(I_x\) and its transform. Thus
the generic action preserves the inertia conjugacy class
at each fixed geometric point. The relative point
comparison makes the conjugator geometric and makes the
map on the residue Galois group the identity. The first
assertion proves injectivity on the pointwise kernel.
\end{proof}

\paragraph{Exact birational scope at the metabelian field.}
For every finite \(E/B_p\), Koenigsmann's
\citet[Theorem~4.1, Corollaries~4.3 and~4.5]
{FF_Koenigsmann2003} apply: \(E\) is neither separably
closed nor real closed, and has a rank-one mixed
characteristic \((0,p)\) henselian valuation with
\(p\) dividing the order of \(G_E\).
They characterize the geometric decomposition and inertia
groups of the complete curve and the fundamental-group
quotients of its opens. Therefore the generic action in
Claim~\ref{clm:FFgenerickernel}, already fixing every
finite geometric place, fixes the remaining place at
infinity as well. Corollary~4.5 deduces a birational
Isom conclusion only when a separate fundamental-group
Isom conclusion is available for an open. It does not
prove that the pointwise birational stabilizer here is
trivial.

The more recent relative and semi-absolute theorems of
\citet[Theorems~A and~E, Definition~1.1,
Proposition~2.6]{FF_Tsujimura2026} retain an abelian-variety
Kummer-faithfulness hypothesis. For \(\Sigma=\{p\}\)
that hypothesis fails here: the elliptic curve in
Claim~\ref{clm:FFmetabelianbarrier} has trivial
\(G_{B_p}\)-action on its full \(p\)-adic Tate module,
so its nonzero rational \(p\)-power torsion points have
compatible \(p\)-divisions in \(B_p\). They are not
prime-to-\(p\) torsion. In fact every nonempty
\(\Sigma\) is excluded. If \(\ell\in\Sigma\)
and \(\ell\ne p\), the pro-\(p\) image of
\(G_{B_p}\) on the Tate module of any positive-dimensional
abelian variety injects into \(\operatorname{GL}_n(\mathbb F_\ell)\):
its intersection with the pro-\(\ell\) principal
congruence kernel is trivial. The image is therefore
finite. After a finite base extension all
\(\ell\)-power torsion is rational. A nonzero point
of this divisible \(\ell\)-primary subgroup is
\(\Sigma\)-divisible, since multiplication by each
other prime is bijective on that subgroup, but is not
\(\Sigma'\)-torsion. For \(\Sigma=\{p\}\) use
the preceding Tate example. This is a failure of that
source's hypothesis; the positive fields already proved
above show why it is not a reconstruction counterexample.
No Isom theorem for \(B_p\),
or vanishing of the remaining stabilizer, is inferred
from either source.

The commensurator theorem of
\citet[Theorem~6.8]{FF_HoshiMochizukiTsujimura2025}
also assumes TKND-AVKF. Its more general
Theorem~5.15 assumes QAA; Definition~5.12(ii)--(iii)
retains conditions on proper-cover cohomology and
function fields on all covers. These are not supplied
by the genus-zero conclusions below.


\subsection{Cyclic rational covers determine every inertia exponent}

\begin{claim}{F.46}{Proved}\label{clm:FFallprimekernel}
Let \(k\) be a characteristic-zero field and \(X/k\)
a smooth geometrically integral curve with function field
\(F\). The following data determine an automorphism of
\(X_{\mathrm{\acute et}}\) over \(BG_k\): a continuous
automorphism \(\alpha\) of \(G_F\) over \(G_k\),
and a \(G_k\)-equivariant bijection
\(\theta:X(\bar k)\to X(\bar k)\), such that
\(\alpha\) sends the geometric conjugacy class of
\(I_x\) to that of \(I_{\theta(x)}\) for every
\(x\in X(\bar k)\). The resulting actual point
function is \(\theta\), and the resulting generic
action is \(\alpha\), with the usual geometric inner
ambiguity.

Now let \(E/B_p\) be finite, and let \(\varphi\)
be a pointwise-trivial pinned automorphism of
\(\mathbb A^1_{E,\mathrm{\acute et}}\) over \(BG_E\).
For every finite extension \(L/E\) and every positive
integer \(m\), its generic action fixes
\[
                    H^1(G_{L(T)},\mu_m)
\]
pointwise. Equivalently, it acts trivially on the maximal
abelian profinite quotient of \(G_{L(T)}\), after
every finite constant base change. All geometric inertia
exponents are one, at every prime. Neither assertion
identifies the remaining nonabelian generic action with
an inner automorphism.
\end{claim}

\begin{proof}
\emph{Reconstructing the generating site.}
We first give a converse to the construction in
Claim~\ref{clm:FFgenerickernel}. Fix prolongations and
geometric residue embeddings at the closed points of
\(X\). The category of quasicompact separated etale
\(X\)-schemes has the following description:
\[
 \begin{gathered}
 \left(S,\ (B_x)_{x\in|X|_{\rm cl}}\right),
 \qquad S\text{ a finite continuous }G_F\text{-set},\\
 B_x\subset S^{I_x}\text{ a }D_x\text{-stable subset}
       \qquad(x\in|X|_{\rm cl}).
 \end{gathered}
                                                        \tag{F.46.1}
\]
where \(B_x=S^{I_x}\) for all but finitely many closed
\(x\). Its morphisms are the \(G_F\)-maps on generic
sets that take every \(B_x\) into its counterpart.
Indeed, form the finite normalization belonging to
\(S\), take its etale locus \(j_*S\), and remove
the finitely many closed points specified by the omitted
\(D_x\)-orbits. This reconstructs the object. Conversely,
Claim~\ref{clm:FFgenerickernel} identifies every object
as precisely such an open. A generic map extends to the
finite normalizations, and its image lies in the indicated
open exactly when the conditions on the \(B_x\) hold.
This proves the asserted description of arrows as well.
Multiple generic components are allowed in \(S\).
Nonseparated objects are not needed: this separated
subcategory is the generating site by
\citet[Theorem~A.8 and Corollary~A.11]{FF_CHW2024}.

Let \(N=G_{\bar k(X)}\). For a chosen geometric
\(x\), choose \(g_x\in N\) such that
\[
                  \alpha(I_x)=g_x I_{\theta(x)}g_x^{-1}.
                                                        \tag{F.46.2}
\]
On generic sets use pullback by \(\alpha\), so the new
action on \(S\) is \(g\cdot_{\rm new}s=\alpha(g)s\).
Replace the local subset at \(x\) by
\[
                                g_x B_{\theta(x)}.
                                                        \tag{F.46.3}
\]
This is a subset of \((\alpha^*S)^{I_x}\). It is
\(D_x\)-stable because (F.45.4) and (F.46.2) imply
\(\alpha(D_x)=g_xD_{\theta(x)}g_x^{-1}\).
Another choice of \(g_x\) differs on the right by an
element of \(D_{\theta(x)}\), again by (F.45.4), and
therefore gives the same subset. Changes of valuation
prolongations give the usual identical description up
to canonical generic-set isomorphism. The
\(G_k\)-equivariance of \(\theta\) makes the rule
well-defined on closed-point orbits. Only finitely many
full subsets are changed, since \(\theta\) bijects
those orbits.

The rule preserves arrows and has an inverse obtained
from \(\alpha^{-1},\theta^{-1}\). Finite limits are
computed on the generic and closed geometric stalks, as
are jointly surjective etale covering families. Thus it
is an equivalence of the generating sites and induces
an equivalence of their topoi. On sets inflated from
\(G_k\), the generic action is the specified identity
and \(g_x\in N\) acts trivially. Hence the equivalence
has the required structure over \(BG_k\).

For completeness, the residue-field identification is
also determined. Since \(\theta\) is an equivariant
bijection, the stabilizers of \(x\) and \(\theta(x)\)
in \(G_k\) coincide. Their embedded finite residue
fields inside \(\bar k\) are consequently the same.
The isomorphism on \(D_x/I_x\) has the identity
projection to this stabilizer, since both \(\alpha\)
and the conjugation in (F.46.2) are relative over
\(G_k\). The stalk identifications in (F.46.3)
therefore identify the geometric point functor at \(x\)
with the target point functor at \(\theta(x)\), with
its specified base comparison. By
\citet[Proposition~3.6]{FF_CHW2024}, the actual point
function is exactly \(\theta\). The generic action
is the prescribed one by construction. This proves the
first assertion.

\emph{The cyclic rational cover.}
Fix now \(L/E\) finite and \(a\in L\), and base
change \(\varphi\) to \(L\). Its point function is
still the identity. On
\(V=\mathbb A^1_L\setminus\{a\}\), geometric first
cohomology with \(\mu_m\)-coefficients is generated
by \(\kappa_m(T-a)\). The equivalence therefore gives
\[
 \varphi^*\kappa_m(T-a)
       =u_m\kappa_m(T-a)+\eta_m,
 \quad u_m\in(\mathbb Z/m\mathbb Z)^\times,
 \quad\eta_m\in H^1(G_L,\mu_m).
                                                        \tag{F.46.4}
\]
The displayed geometric generator is already lifted by
the indicated global unit, so the remaining term is
constant, by inflation--restriction. Evaluation at
\(T=a+1\) gives \(\eta_m=0\).

The field \(L\) contains \(\mu_m\). Thus the
extension
\[
              L(T)\subset L(U),\qquad T=U^m+a,
                                                        \tag{F.46.5}
\]
is cyclic Galois of degree \(m\). Its absolute Galois
subgroup is the kernel of the homomorphism
\(\kappa_m(T-a):G_{L(T)}\to\mu_m\).
By (F.46.4), the generic automorphism \(\alpha\)
preserves this subgroup. Its restriction
\(\alpha_U\) is an automorphism of \(G_{L(U)}\)
over \(G_L\).

The exact local theory in
\citet[Theorem~4.1 and Corollaries~4.3--4.5]
{FF_Koenigsmann2003} applies to \(L\), and after every
finite constant extension, both for \(\mathbb P^1_T\)
and \(\mathbb P^1_U\). It gives the associated
\(G_L\)-equivariant permutation \(\theta_U\) of
geometric places. More explicitly, first make a chosen
geometric place rational over a finite constant field
\(L'/L\). Corollary~4.3 characterizes its arithmetic
decomposition conjugacy class. Any conjugating element
can be multiplied by the inverse of an element of the
target decomposition group with the same projection to
\(G_{L'}\), because that projection is surjective.
The resulting conjugator is geometric. Applying this
argument over all finite constants, uniqueness of the
recovered valuation classes makes the point permutations
compatible with scalar extension and with the
\(G_L\)-action. This is the geometric conjugacy
transport required in (F.46.2).
This permutation lies above the
identity permutation on \(\mathbb P^1_T\). To see
the needed compatibility explicitly, choose a place of
the common separable closure. Its decomposition and
inertia groups for \(L(U)\) are the intersections of
the corresponding groups for \(L(T)\) with
\(G_{L(U)}\). The subgroup is preserved by
\(\alpha\), so these intersections and their place
restrictions are transported compatibly. Uniqueness in
the local characterization identifies the transported
places with those of \(\theta_U\). In particular,
\[
                           \theta_U(u)^m=u^m
                  \qquad(u\in\bar L).
                                                        \tag{F.46.6}
\]
The unique place over \(T=\infty\) is fixed, as is
the unique place \(U=0\) over \(T=a\). Therefore
\(\alpha_U,\theta_U\) satisfy the first assertion
for the affine line \(\mathbb A^1_U\). They define
a pinned automorphism of its full etale topos over
\(BG_L\), rather than merely of its generic finite
cover category.

Claim~\ref{clm:FFmetabelianaffine} now applies to this
lifted automorphism. Its actual point function is affine.
Since it fixes zero and satisfies (F.46.6), it has the
form
\[
                       \theta_U(u)=c u,
                       \qquad c\in\mu_m(L).
                                                        \tag{F.46.7}
\]
On the other hand, (F.46.4) says that \(\alpha\)
acts on the cyclic deck quotient of (F.46.5) by the
power \(u_m\). Conjugation and transport of inertia
classes are compatible. Indexing the deck rotations by
their action on the geometric points of the cover,
this gives
\[
                    \theta_U(\zeta u)
                          =\zeta^{u_m}\theta_U(u)
                    \qquad(\zeta\in\mu_m).
                                                        \tag{F.46.8}
\]
Combining
(F.46.7) and (F.46.8) gives
\(\zeta^{u_m}=\zeta\) for every \(\zeta\in\mu_m\).
Hence \(u_m=1\) in \(\mathbb Z/m\mathbb Z\).
We have proved, for every \(a\) and \(m\),
\[
                  \varphi^*\kappa_m(T-a)=\kappa_m(T-a).
                                                        \tag{F.46.9}
\]

\emph{Descent of an arbitrary rational Kummer class.}
For \(r\in L(T)^\times\), take its nonvanishing
open \(V_r\subset\mathbb A^1_L\) and a finite
Galois constant extension \(L'/L\) splitting its
divisor. Over \(L'\), it is a constant times a product
of integer powers of the linear units in (F.46.9).
The relative structural isomorphism fixes the constant
classes. Thus
\(\delta=\varphi^*\kappa_m(r)-\kappa_m(r)\)
restricts to zero on \((V_r)_{L'}\).
Inflation--restriction for this purely constant cover
puts \(\delta\) in the image of
\(H^1(\operatorname{Gal}(L'/L),\mu_m)\), hence in
the constant subgroup from \(H^1(G_L,\mu_m)\).
Choose \(t_0\in V_r(L)\). Actual pointwise identity
gives \(t_0^*\delta=0\), and evaluation at \(t_0\)
is a left inverse on constant classes. Therefore
\(\delta=0\) before the scalar extension as well.

Every element of \(H^1(G_{L(T)},\mu_m)\) is a
Kummer class. Moreover \(\mu_m\subset L\), so
these cohomology groups are exactly the continuous
characters with values in \(\mu_m\). Their equality
for every \(m\) is equivalent to equality on the full
abelian profinite quotient. The exponents on finite
geometric inertia are detected by (F.46.9), and the
single peripheral relation on the projective line gives
the same exponent at infinity. All of them are one.
\end{proof}

This improves the direct \(p\)-primary calculation:
the other primes are controlled by the rational cyclic
covers and the already-proved affine point-function
rigidity, even though point Kummer classes at those
primes vanish over \(B_p\). A nonabelian generic
automorphism acting trivially on all these abelian
quotients has not been excluded.


The distinction from inner triviality is concrete even
for an abstract four-punctured line. Let \(P\) be the
free pro-\(p\) group on \(x,y,z\), with fourth
peripheral generator \((xyz)^{-1}\). The automorphism
\[
 x\longmapsto(xy)x(xy)^{-1},\qquad
 y\longmapsto(xy)y(xy)^{-1},\qquad
 z\longmapsto z
                                                        \tag{F.46.10}
\]
has inverse given by conjugating the first two generators
by \((xy)^{-1}\). It fixes \(xy\), \(xyz\), every
peripheral conjugacy class, and abelianization. It is not
inner. In the torsion-free class-two pro-\(p\) quotient,
the independent central commutators are
\([x,y],[x,z],[y,z]\). Conjugation by an element with
abelian coordinates \((a,b,c)\) could agree with
(F.46.10) on \(x,y\) only if \(a=b=1,c=0\).
Its commutator with \(z\) would then be
\([x,z][y,z]\ne1\), a contradiction. This quotient
exists also for \(p=2\): use coordinates in
\(\mathbb Z_p^3\times\mathbb Z_p^3\), with the
central product correction
\((a_i a'_j)_{i<j}\), so no division by two is involved.
This is an abstract obstruction to an abelianization-only
argument. No compatibility of (F.46.10) with the actual
\(G_E\)-action, with every finite cover, or with the
full affine-line etale topos is asserted.

The class-two calculation was executed modulo 2, 3, 4, 5 and 8. The order-64 model was checked on all 4096 products. The program and saved output record these finite checks:

\path{agent_F/F_metabelian_checks.py}.


\subsection{The pointwise kernel fixes the genus-zero part of the site}

\begin{claim}{F.47}{Proved}\label{clm:FFgenuszerokernel}
Let \(E/B_p\) be finite and let \(\varphi\) be a
pointwise-trivial pinned automorphism of
\(\mathbb A^1_{E,\mathrm{\acute et}}\) over \(BG_E\).
Its inverse-image functor is naturally isomorphic to the
identity on the full subcategory of quasicompact separated
etale objects whose smooth projective geometric components
have genus zero. The same assertion holds after every
finite constant extension.

This constraint is strictly stronger than peripheral and
abelianization control for an abstract free profinite
four-puncture group. Specifically, the word automorphism
\[
 x\longmapsto(xy)x(xy)^{-1},\qquad
 y\longmapsto(xy)y(xy)^{-1},\qquad z\longmapsto z
                                                        \tag{F.47.1}
\]
of the free \emph{profinite} group on \(x,y,z\) moves
a degree-three genus-zero cover, despite fixing every
peripheral conjugacy class and abelianization. Hence it
cannot be the geometric action of such a \(\varphi\)
on a four-punctured open. This is a statement about the
full profinite word automorphism (F.47.1); it is not an
assertion about every full automorphism whose pro-\(p\)
image equals the corresponding pro-\(p\) word automorphism.
\end{claim}

\begin{proof}
We first take a geometrically connected etale object
\(Y\to\mathbb A^1_E\) with projective completion
\(C\) of genus zero, and write
\(Y'=\varphi^*Y\). Pullback to the slice topoi gives
an equivalence \(Y'_{\mathrm{\acute et}}\to
Y_{\mathrm{\acute et}}\) over the specified base.
It induces an isomorphism of the generic Galois groups
over \(G_E\). The constant fields remain \(E\),
because the relative generic action preserves the
projection to \(G_E\). Let \(C'\) denote the smooth
projective completion of \(Y'\).

Koenigsmann's local theory identifies all geometric
inertia groups on \(C\) and \(C'\), not merely those
at the points of the two displayed opens. Killing their
closed normal span gives the etale fundamental group of
the complete curve, compatibly over \(G_E\). Hence
the geometric fundamental groups of \(C\) and \(C'\)
are isomorphic. In characteristic zero the geometric
abelianization has rank \(2g\), by
\citet[Expos\'e~XIII, Corollary~2.12]{FF_SGA1}.
It follows that \(C'\) also has genus zero.

The same local theory supplies a compatible
\(G_E\)-equivariant point bijection
\[
                         h:C'(\bar E)\longrightarrow C(\bar E).
                                                        \tag{F.47.2}
\]
Over the displayed opens it is the actual point function
of the slice equivalence. Indeed, the specialization
functors \(S\mapsto S^{I_x}\) used in
Claim~\ref{clm:FFgenerickernel} identify the inertia
class of every actual point with precisely its transported
class. At a rational point the arithmetic conjugator can
be made geometric using the surjection of its decomposition
group onto the constant Galois group, as explained in
Claim~\ref{clm:FFallprimekernel}.

Choose a finite Galois extension \(L/E\) over which
\(C\) has a rational point \(c\). Then
\(c'=h^{-1}(c)\) is \(L\)-rational as well.
A smooth projective genus-zero curve with a rational
point is \(\mathbb P^1\). Choose coordinates on
\(C_L\) and \(C'_L\) taking \(c\) and \(c'\)
to infinity. The generic group isomorphism, together
with (F.47.2), transports the inertia classes of the
two remaining affine lines. The site construction in
Claim~\ref{clm:FFallprimekernel}, applied after these
coordinate identifications, gives a pinned automorphism
of the affine-line etale topos whose actual point function
is \(h\) in these coordinates.

Claim~\ref{clm:FFmetabelianaffine} makes this function
affine. It therefore extends to an \(L\)-isomorphism
\(C'_L\to C_L\), realizing all of (F.47.2).
It carries \(Y'_L\) onto \(Y_L\), because on these
opens the actual point function is the bijection of the
slice equivalence. This isomorphism descends to \(E\):
the point bijection is \(G_E\)-equivariant, so every
Galois conjugate of the \(L\)-isomorphism has the same
geometric point function and is therefore equal to it.
This is ordinary finite Galois descent of scheme maps.
We obtain a unique \(E\)-isomorphism
\[
                                  h_Y:Y'\xrightarrow{\sim}Y
                                                        \tag{F.47.3}
\]
with the required actual point function.

It is an isomorphism over \(\mathbb A^1_E\).
The slice equivalence lies over \(\varphi\), whose
base point function is the identity. Thus both composites
from \(Y'\) to \(\mathbb A^1_E\) have the same
geometric point function, and the corresponding scheme
maps are equal. For any arrow \(Y\to Z\) between
two objects under consideration, the same pointwise
compatibility gives
\(h_Z\circ\varphi^*(Y\to Z)=(Y\to Z)\circ h_Y\).
Equality follows because the sources are reduced and
geometric points are dense. This proves naturality.

For a connected object with a finite constant field
extension, apply the same argument over that extension;
the relative structure identifies the extension on both
sides. For a disconnected object, use its finitely many
connected components. Equivalently, one can split all
components and choose their rational points over one
finite Galois constant extension, then descend the unique
pointwise isomorphisms. This proves the full first
assertion and its compatibility with further constants.

\emph{A finite quotient detecting the abstract word automorphism.}
Write \(w=(xyz)^{-1}\). In \(S_3\), assign
\[
           (x,y,z,w)\longmapsto
           \big((12),(23),(12),(13)\big).
                                                        \tag{F.47.4}
\]
The four permutations have product one and generate
\(S_3\). By the characteristic-zero presentation of the
geometric fundamental group of a four-punctured line
in \citet[Expos\'e~XIII, Corollary~2.12]{FF_SGA1},
this gives a connected degree-three etale cover over
\(\bar E\). Each of the four branch permutations
is a transposition. Its smooth completion therefore has
\[
                         2g-2=-2\cdot3+4=-2,
\]
so has genus zero. The cover and its map descend to a
finite constant extension of \(E\), by finite
presentation.

The automorphism (F.47.1) fixes \(xy\) and \(z\),
and therefore fixes \(w\). It conjugates the first
two permutations in (F.47.4) by their nontrivial
three-cycle product, leaving the last two unchanged.
The original and transformed degree-three permutation
representations are not isomorphic. A simultaneous
conjugator would have to centralize both \((12)\) and
\((13)\), so would be the identity; but the first
permutation has changed. Thus (F.47.1) moves this
genus-zero cover and contradicts the first assertion
if proposed as the action of \(\varphi\).

This finite calculation was executed in
\texttt{F\_metabelian\_checks.py}; its output verifies
the product, transitivity, Riemann--Hurwitz integer, and
the absence of any of the six possible simultaneous
conjugators. Since \(S_3\) is not a \(p\)-group,
the test concerns the stated free profinite word action
and does not identify all possible lifts of a pro-\(p\)
action. No arithmetic-compatible ghost is constructed.
\end{proof}

The genus-zero subcategory is not a generating family
for the whole curve etale site. For example, a nonempty
etale object with projective geometric genus at least
one cannot be covered by genus-zero etale curves: a
dominant map between their projective completions would
contradict Riemann--Hurwitz. Nor has any cofinality of
Galois closures of genus-zero covers been established.
The action on arbitrary positive-genus covers remains the missing
part of the birational kernel problem.


\subsection{Genus-zero generation of abelian covers and elliptic Galois covers}

\begin{claim}{F.48}{Proved}\label{clm:FFlatteskernel}
Let \(E/B_p\) be finite, and let \(\varphi\) be a pinned
automorphism of \((\mathbb A^1_E)_{\acute et}\) over \(BG_E\)
whose actual geometric point function is the identity.
The natural scheme realizations of
Claim~\ref{clm:FFgenuszerokernel} extend to the following
objects, after every finite extension of constants:
finite fibre products of the genus-zero objects, their
open subobjects, and finite etale quotients by finite
groups. In particular they extend to every subcover of
the Galois closure of a connected finite etale
genus-zero cover of an open of \(\mathbb A^1\).

They consequently extend to every etale restriction of
a finite abelian cover of \(\mathbb P^1_E\), of
arbitrary genus, and to every finite Galois cover
\[
                  \pi:C\longrightarrow\mathbb P^1_E
       \quad\hbox{with }g(C)=1.
                                                        \tag{F.48.1}
\]
Here ``extend'' means that the induced slice point
function is realized by a unique scheme isomorphism
\(\varphi^*V\to V\) over the chosen base open, and
these isomorphisms are natural in every morphism between
the stated objects. The statement includes a quadratic
cover with four geometric branch points. It also
includes every degree-three cover with genus-one
source, without a Galois assumption.

Here \(C\) is smooth and geometrically integral.
Thus, after identifying the inverse image of any
connected cyclic torsor of the stated kind with the
original one, its possible fibrewise deck rotation is
a single global deck transformation. Point-dependent
rotations on cyclic covers of arbitrary genus are
excluded. No generation or cofinality
claim for all positive-genus covers is asserted.
\end{claim}

\begin{proof}
All constructions below may be made after removing
finitely many points of the base. The resulting objects
are quasicompact separated etale objects, and their
inverse images are represented by objects of the same
kind. We keep the structural comparisons over the
constant field throughout.

\emph{Closure under the stated constructions.}
Say that an etale object \(V\to X\), for an open
\(X\subseteq\mathbb A^1_E\), is realized if its
actual slice point function is induced by an
\(X\)-isomorphism
\[
                       a_V:\varphi^*V\xrightarrow{\sim}V.
                                                        \tag{F.48.2}
\]
Such an isomorphism is unique: two maps between these
reduced finite-type schemes, with separated target, that
agree on all geometric points agree as scheme maps.
Moreover, for any arrow \(V\to W\) between realized
objects, the two maps \(\varphi^*V\to W\) formed
using (F.48.2) agree on geometric points by the
naturality of point evaluation. They therefore agree
as scheme maps. Realizations are consequently natural
on every full subcategory on which they exist.

If \(V_1,\ldots,V_r\) are realized, their fibre
product is realized by the product of the maps
\(a_{V_i}\). Indeed inverse image preserves finite
limits, and a geometric point of the product is
determined by its projections. If \(W\subseteq V\)
is an open subobject of a realized object, the image of
\(\varphi^*W\) under \(a_V\) is precisely \(W\):
the two opens have the same geometric points. Restriction
of \(a_V\) gives \(a_W\). This includes open-and-closed
connected components. These conclusions use the actual
point-function comparison, rather than just invariance
of the isomorphism class of \(V\).

Suppose that a finite group \(H\) acts on a realized
finite etale object \(Z\to X\). Naturality for every
action arrow identifies the transported action with the
given action on \(Z\). Inverse image preserves its
coequalizer. The finite etale quotient therefore has
an induced isomorphism
\[
             \varphi^*(Z/H)\ \xrightarrow{\sim}\ Z/H.
                                                        \tag{F.48.3}
\]
It realizes the slice point function, because the finite
etale surjection \(Z\to Z/H\) is surjective on
geometric points and that point function is natural.

For clarity, the Galois-closure application can be made
entirely inside these constructions. If \(V\to X\)
is connected finite etale of degree \(d\), take in
\(V^d_X\) the open-and-closed subscheme of ordered
\(d\)-tuples of pairwise distinct geometric sheets.
It is the frame torsor of \(V\). A connected component
\(Z\) is its Galois closure, with group the monodromy
image in \(S_d\). Once \(V\) is realized, its frame
torsor and each specified component are realized by
the preceding paragraph. Each finite subcover of
\(Z\) is a quotient \(Z/H\), so (F.48.3) applies.
This proves the first assertion using
Claim~\ref{clm:FFgenuszerokernel}.

\emph{All abelian covers of the rational base.}
After a finite extension of constants, every connected
cyclic cover of degree \(m\) is a Kummer cover.
Indeed all \(m\)-th roots of unity are available,
and Kummer theory applies to its cyclic function-field
extension. After a further finite extension that splits
the divisor and extracts a constant root, write it as
\[
       y^m=f(T)=c^m\prod_{i=1}^r(T-a_i)^{e_i},
           \qquad c\ne0,\quad e_i\in\mathbb Z.
                                                        \tag{F.48.4}
\]
Remove the finitely many \(a_i\), the original affine
infinity, and any additional omitted base points.
Each cover \(u_i^m=T-a_i\) of the resulting open
has genus-zero completion and is realized by
Claim~\ref{clm:FFgenuszerokernel}. Their finite fibre
product is therefore realized, as are its connected
components and finite group quotients.

More explicitly, the classes of the distinct
\(T-a_i\) are independent in
\(\bar E(T)^\times/(\bar E(T)^\times)^m\), as is
seen from their respective valuations at \(a_i\).
Their compositum is thus a connected Galois cover
with group \(\mu_m^r\). The map
\[
       (u_1,\ldots,u_r)\longmapsto
                     y=c\prod_i u_i^{e_i}
                                                        \tag{F.48.5}
\]
identifies (F.48.4) with its quotient by the kernel
of the character
\((\zeta_i)\mapsto\prod_i\zeta_i^{e_i}\).
The character is surjective when the original cover
is geometrically connected of degree \(m\). If
constant splitting produces several components, apply
this argument to each component and descend their
union. This proves the cyclic assertion. A finite
abelian Galois extension is a compositum of cyclic
subextensions, by the elementary structure theorem
for its finite abelian Galois group. Its cover is a
connected component of the corresponding finite fibre
product, so every abelian cover is realized as well.

\emph{The elliptic Latt\`es construction.}
We first give the quadratic case explicitly. After a
finite Galois extension \(L/E\), a branch point of
\(\pi\) has an \(L\)-rational point above it, which
we choose as the origin of the elliptic curve \(A=C_L\).
The involution of the double cover is then \([-1]\).
Enlarge \(L\) so that all points of \(A[3]\) are
rational. Write
\[
                 \pi:A\longrightarrow A/\{\pm1\}
                                      =\mathbb P^1_L.
\]
Multiplication by three commutes with \([-1]\), and
therefore induces a rational morphism of degree nine
\[
          f_3:\mathbb P^1_L\longrightarrow\mathbb P^1_L,
                 \qquad f_3\circ\pi=\pi\circ[3].
                                                        \tag{F.48.6}
\]
Let \(B\) be the branch-value set of \(\pi\), and
put \(X=\mathbb P^1_L\setminus(B\cup\{\infty\})\),
with \(\infty\) interpreted in the originally chosen
affine base coordinate. The maps
\[
 \begin{split}
  Z&=(\pi\circ[3])^{-1}(X)\longrightarrow X,\\
  V&=f_3^{-1}(X)\longrightarrow X
 \end{split}
                                                        \tag{F.48.7}
\]
are finite etale: \([3]\) is etale, and the branch
values of the composite lie in \(B\). The second
cover has genus-zero completion.

The first cover in (F.48.7) is Galois, with its
eighteen deck transformations
\[
            P\longmapsto \epsilon P+t,
             \qquad t\in A[3],\quad\epsilon\in\{1,-1\}.
                                                        \tag{F.48.8}
\]
They are distinct and their number equals the degree
\(2\cdot9\) of the map. Its group is
\(G=A[3]\rtimes\{\pm1\}\). The intermediate
cover \(V\) is the quotient by
\(H=\{\pm1\}\). The core of \(H\) in \(G\)
is trivial: conjugation of \([-1]\) by translation
by \(t\ne0\) gives \(P\mapsto-P+2t\), which
does not belong to \(H\). Thus \(Z\) is the
Galois closure of the genus-zero cover \(V\to X\).
The first part of the proof realizes \(Z\).
Its quotient by the normal translation subgroup
\(A[3]\) is the original cover
\[
               \pi^{-1}(X)\longrightarrow X,
                                                        \tag{F.48.9}
\]
with quotient map supplied by \([3]\). It is therefore
realized as well. The finite group and its faithful
degree-nine permutation action are also checked in
\path{agent_F/F_metabelian_checks.py}.

\emph{The other cyclic genus-one covers.}
The same argument does not require the deck group to
have order two. Over a finite extension, write a
generator of the cyclic deck group as
\(P\mapsto h(P)+t\) on an elliptic model of \(C\).
If \(h=1\), the group consists of translations and
its quotient has genus one, contrary to (F.48.1).
Thus \(h-1\) is a nonzero endomorphism of the elliptic
curve, hence a surjective isogeny. The generator has
a fixed point; choose it as origin, after a further
finite extension. The deck group \(H\) now consists
of group automorphisms of \(A\), and
\(A/H=\mathbb P^1\).

Choose an integer \(n\ge3\) such that
\[
                 n^2>\deg(h-1)
                  \quad\text{for every }1\ne h\in H,
                                                        \tag{F.48.10}
\]
and enlarge the finite constant field to contain
\(A[n]\). The map \([n]\) commutes with \(H\)
and induces a degree-\(n^2\) map \(f_n\) on
\(A/H=\mathbb P^1\). The composite
\(\pi\circ[n]\) is Galois with group
\(A[n]\rtimes H\). The core of \(H\) is trivial.
Indeed, an element \(h\) in that core must remain
in \(H\) after conjugation by every translation
\(t\in A[n]\), so \((h-1)t=0\) for all such
\(t\). If \(h\ne1\), this would put \(n^2\)
distinct points in the kernel of an isogeny of degree
strictly less than \(n^2\), contradicting (F.48.10).
After removing the branch values and the affine
infinity, the preceding Galois-closure argument applies
unchanged. Quotienting by \(A[n]\) recovers
\(\pi\).

\emph{Every Galois genus-one cover.}
Let now \(\pi:A\to\mathbb P^1\) be a Galois
cover with finite deck group \(G\), and genus one
source. Work after a finite extension so that the
origins, deck transformations and torsion used below
are defined. The translations in \(G\) form a normal
subgroup \(T\). Its quotient
\[
                  \lambda:A\longrightarrow B=A/T
\]
is an isogeny of degree \(d=|T|\). The group
\(G/T\) is cyclic: the linear-part map embeds it in
the finite group of origin-preserving automorphisms
of an elliptic curve, which embeds by its action on
an invariant differential in the multiplicative
constant field. This action on differentials is
faithful in characteristic zero. The residual action
on \(B\) cannot consist only of translations,
since \(B/(G/T)=\mathbb P^1\). Its cyclic generator
therefore has a fixed point. Choose that point as
origin of \(B\), and a preimage as origin of \(A\),
so that \(\lambda\) remains a homomorphism. Put
\(H=G/T\), now acting on \(B\) by group
automorphisms, and let
\(\pi_B:B\to B/H=\mathbb P^1\).

Choose \(N\) divisible by \(d\), large enough
that \(N^2>\deg(h-1)\) for all \(1\ne h\in H\).
The Latt\`es argument shows that
\(\pi_B\circ[N]_B\), on the branch complement,
is the Galois closure of a genus-zero cover and is
therefore realized. The dual isogeny
\(\lambda^\vee:B\to A\) satisfies
\(\lambda\circ\lambda^\vee=[d]_B\). Thus
\[
 \pi\circ\lambda^\vee\circ[N/d]_B
                       =\pi_B\circ[N]_B.
                                                        \tag{F.48.11}
\]
The original cover \(\pi\) is a subcover of this
realized finite Galois cover. It is realized by the
quotient part of the proof. This establishes the
full genus-one Galois assertion. The construction alone
does not treat every non-Galois genus-one map; the
following argument also settles all degree-three maps.

\emph{Degree-three genus-one covers, without a Galois assumption.}
The same generated category also contains every degree-three
cover with genus-one source. We give the finite monodromy
argument explicitly. It is the genus-one case of the
classical trigonal construction described in
\citet[Section~9.1, Remark~9.13]{FF_RecillasRodriguez2003};
the finite verification below makes this consequence
independent of a convention about the genus in a
moduli-space statement.

Work over \(\overline{\mathbb Q}\). A connected
degree-three cover has geometric monodromy \(C_3\)
or \(S_3\). The first case is already abelian.
In the second case let \(a\) be the number of branch
transpositions and \(b\) the number of branch
three-cycles. Riemann--Hurwitz and the product-one
relation give
\[
          a+2b=6,\qquad a>0\text{ even},\qquad
          (a,b)\in\{(6,0),(4,1),(2,2)\}.
                                                        \tag{F.48.12}
\]
Let \(q:S_4\twoheadrightarrow S_3\) be the action
on the three partitions of four letters into two
unordered pairs. Its kernel is the Klein four group.
Every branch tuple in (F.48.12) that generates
\(S_3\) and has product one has a lift through
\(q\) that generates \(S_4\), has product one,
and preserves each permutation index. Here the index
of a permutation on \(d\) letters is \(d\) minus
its number of cycles.

This last assertion is a finite lemma verified by
complete enumeration in the supplied program. The
five nonidentity elements of \(S_3\) give
\(5^4+5^5+5^6=19375\) tuples of the possible
lengths. Filtering by index sum six, product one,
and generation of \(S_3\) leaves respectively
72, 270, and 240 tuples. A transposition has two
index-preserving lifts through \(q\), and a
three-cycle has four. For every one of the 582
retained tuples the program enumerates all 64
possible index-preserving lifts and finds exactly
twelve with product one and generated group
\(S_4\). It verifies the quotient map by all
\(24^2\) group products and saves an explicit
lift for every retained tuple. Thus the finite
assertion checks all cases, not just selected
ramification patterns. The program and certificates
are supplied at the following paths.

\path{agent_F/F_metabelian_checks.py}

\path{agent_F/F_recillas_certificates.json}

Apply the characteristic-zero Riemann existence
theorem to a chosen lift at the actual branch
points of the original cover. Its degree-four
source has
\(2g-2=-2\cdot4+6=-2\), hence genus zero.
Its Galois closure has group \(S_4\). The quotient
through \(q\) has exactly the original
\(S_3\) branch tuple, so the original cubic
cover is the subcover associated with the inverse
image of a point stabilizer in \(S_3\).
It therefore belongs to the category already
realized above.

Here descent from the complex construction retains the
covering diagram, not merely its branch data. Let
\(X=\mathbb P^1_{\overline{\mathbb Q}}\setminus B\),
where \(B\) is the actual finite branch set, enlarged
by the original affine infinity if necessary. The
Riemann existence theorem identifies its complex finite
etale covers with finite topological covering spaces
\citep[Expose~XII, Theoreme~5.1]{FF_SGA1}.
Base change from \(\overline{\mathbb Q}\) to
\(\mathbb C\) is an equivalence on finite etale
covering categories: apply
\citet[Expose~XIII, Proposition~4.6(a)]{FF_SGA1}
with \(k=\overline{\mathbb Q}\), the smooth curve
\(X\), and \(Y=\operatorname{Spec}\mathbb C\).
The required resolution hypothesis holds for curves,
and in characteristic zero the prime-to-characteristic
quotient is the entire fundamental group. Since
\(\pi_1(Y)=1\), the resulting fundamental-group map
is an isomorphism. This equivalence is fully faithful,
so the \(S_4\)-action, quotient diagram and isomorphism
with the original cubic cover all descend to
\(\overline{\mathbb Q}\). They use finitely many
algebraic coefficients and morphisms, and therefore
descend to one finite extension \(E'/E\).
Uniqueness of the scheme realization with the specified
actual point function then descends that realization
to \(E\), as detailed below.

\emph{Descent and the fibre rotations.}
All the constructions just used occur after one finite
Galois extension \(L/E\). The resulting realization
of the original etale object is unique from its
geometric point function. That function is
\(G_E\)-equivariant, so the realization is fixed by
\(\operatorname{Gal}(L/E)\) and descends to \(E\).
If the originally specified etale restriction omits
further points of \(C\), use the open-subobject part
of the proof. Naturality follows from the uniqueness
argument at the beginning.

In any connected cyclic torsor case, compare the descended
realization with any chosen isomorphism of the
transported torsor with the original one. Their
composite is an automorphism of the connected cyclic
cover over the base, hence a single global deck
transformation. Accordingly the putative roots of unity
\(c_x\) are constant. This conclusion is about
the slice point function on this object; it does not
identify the induced equivalence on all its further
positive-genus covers.
\end{proof}


The generated subcategory in Claim~\ref{clm:FFlatteskernel}
is in fact not cofinal among all geometric covers.
By the proved fixed-genus composition-factor theorem,
there is a finite set \(\mathcal S_0\) of
nonalternating nonabelian simple groups such that every
simple composition factor of a genus-zero monodromy
group is cyclic, alternating, or belongs to
\(\mathcal S_0\)
\citep[Theorem~1.1, with characteristic zero and genus
zero]{FF_Guralnick2003}.
Every function field obtained by the operations in
Claim~\ref{clm:FFlatteskernel} is contained in a finite
compositum of Galois closures of genus-zero function-field
extensions. The Galois group of that compositum is a
subdirect product of the individual monodromy groups,
so it has no additional simple composition factors.
For completeness, this last assertion follows by induction
on the number of factors: the kernel of projection to
one factor projects normally into each remaining factor;
apply the induction hypothesis to that kernel and then
the composition series of the resulting extension.

Choose a nonalternating nonabelian finite simple group
\(S\notin\mathcal S_0\). Such groups exist, for
example among the infinitely many simple groups of
Lie type. The characteristic-zero fundamental-group
presentation and Riemann existence theorem give an
\(S\)-Galois cover of an open of
\(\mathbb P^1_{\overline{\mathbb Q}}\)
\citep[Expos\'es~XII, Th\'eor\`eme~5.1, and~XIII,
Corollaire~2.12]{FF_SGA1}.
It descends over a finite extension of the constants.
This cover cannot be dominated by an object of the
stated generated subcategory, since then \(S\)
would be a composition factor of one of the preceding
subdirect products. This is a scope restriction on the
closure argument. Noncofinality does not prove that the
remaining automorphism kernel is nontrivial.

\subsection{Vanishing of the Kummer shear on elliptic Galois slices}

\begin{claim}{F.49}{Proved}\label{clm:FFellipticshear}
Let \(E/B_p\) be finite, let \(\varphi\) be a pinned
relative automorphism of \((\mathbb A^1_E)_{\acute et}\)
with identity actual geometric point function, and let
\[
                  \pi:C\longrightarrow\mathbb P^1_E
                                                        \tag{F.49.1}
\]
be a finite Galois cover with smooth geometrically
integral genus-one source. On a branch-complement open,
use the unique actual scheme realization of
Claim~\ref{clm:FFlatteskernel} to regard its generic
slice action as an outer automorphism of \(G_{E(C)}\).
For every finite constant extension \(L/E\), every
\(m\ge1\), and every \(f\in L(C)^\times\), this
outer action satisfies
\[
                     \varphi_C^*\kappa_m(f)=\kappa_m(f).
                                                        \tag{F.49.2}
\]
Equivalently it is the identity on the full abelian
profinite quotient of \(G_{L(C)}\), after every such
constant extension. This includes arbitrary translated
rational functions on the elliptic slice.

Thus every cyclic generic cover of this slice is
preserved as a covering object. Neither its actual
point-function realization nor identity on the entire
elliptic slice topos is asserted by this claim.
\end{claim}

\begin{proof}
We use integral Kummer characters throughout. Since
\(\mu_\infty\subset B_p\), the coefficient sheaves
\(\mathbb Z_\ell(1)\) are constant on every field
under discussion. The relative structural comparison
fixes these coefficients and all characters inflated
from the constant field. Compatible roots of a
nonzero rational function define its character
\(\kappa_{\ell^\infty}(f)\); reduction modulo
\(\ell^r\) is \(\kappa_{\ell^r}(f)\).
Inner changes in the generic action do not affect
these characters.

\emph{Completed-point compatibility on the elliptic slices.}
The geometric local theory used in
Claims~\ref{clm:FFgenerickernel} and
\ref{clm:FFgenuszerokernel} supplies a point
bijection on every smooth projective covering curve,
compatible with finite covering maps. For a
realized genus-one object we must check that this
bijection also agrees at the omitted boundary points
with the extension of its actual scheme realization.
Agreement on the branch complement alone does not
justify that assertion.

Use the dominating Latt\`es cover constructed in
Claim~\ref{clm:FFlatteskernel}. Thus, after finite
constants, the original genus-one Galois cover is a
quotient of
\(Z=B\to\mathbb P^1\) with deck group
\(B[N]\rtimes H\), where \(B\) is elliptic,
\(H\ne1\) consists of origin-preserving group
automorphisms, and \(B/H\) has genus zero. Choose
\(N\) divisible by the degree of the translation
quotient as in that construction, and increase it
so that
\[
              N^2>\sum_{1\ne h\in H}\deg(1-h).
                                                        \tag{F.49.3}
\]
For every \(t\in B[N]\), put
\(H_t=\tau_tH\tau_{-t}\). The quotient
\(Z/H_t\) has genus zero. These quotients jointly
separate all geometric points of \(Z\). Indeed,
if two distinct points \(z,z'\) have the same
image in the quotient by \(H_t\), then
\[
                  z'=h(z)+(1-h)t
\]
for some \(h\in H\setminus\{1\}\). For each
fixed \(h\), the nonzero isogeny \(1-h\) has
at most \(\deg(1-h)\) solutions \(t\) to this
equation, even before restricting to \(B[N]\).
The union of these exceptional sets has size less
than \(N^2\), proving separation.

Claim~\ref{clm:FFgenuszerokernel} matches the
completed generic point bijection with the scheme
isomorphism on each \(Z/H_t\). On the branch
complement, naturality in
Claim~\ref{clm:FFlatteskernel} identifies the
quotient arrows with the scheme realization on
\(Z\); these scheme diagrams extend to their
smooth projective completions. The generic point
bijections are also compatible with every quotient
arrow. Joint separation therefore makes the
completed point bijection on \(Z\) equal to the
completion of its scheme realization. The finite
surjection from \(Z\) to the original genus-one
cover then gives the same equality there, by
surjectivity on geometric points. The construction
works for every genus-one Galois slice used below.
No assertion about separation of branches by the
projections of a normalized fibre product is used.

\emph{The possible error is geometrically unramified.}
For the moment work geometrically and choose an
origin on \(\bar C\), writing it as an elliptic
curve \(A\). Each finite collection of points,
functions and maps used below is defined after a
finite extension of constants, where the actual
scheme realizations in
Claim~\ref{clm:FFlatteskernel} are available and
natural. Their geometric cohomology groups are
unchanged by these extensions. The preceding
argument identifies their completed geometric
point functions as well.

Every geometric inertia exponent of the generic
base action is one, by
Claim~\ref{clm:FFallprimekernel}. The inertia at a
point of a finite covering curve is a finite-index
subgroup of the corresponding base inertia. After
the geometric point identification supplied by the
actual realization, the restricted exponent is
therefore also one. This holds at every point of
the smooth completion, including points over
infinity. Conjugator choices do not change the
conclusion: geometric curve inertia is procyclic,
and conjugation by its decomposition group acts
through the cyclotomic character, which is trivial
on the constant field here.

It follows that the geometric difference
\[
 b_1(f):=\varphi_C^*\kappa_{\ell^\infty}(f)
                         -\kappa_{\ell^\infty}(f)
       \ \in H^1(A,\mathbb Z_\ell(1))
                                                        \tag{F.49.4}
\]
is unramified at all points. The indicated group
means geometric cohomology of the proper curve.
Indeed the residue of a Kummer character on point
inertia is the valuation of its function, and all
these residues of the difference vanish. The
resulting unramified character factors through the
proper curve fundamental group. It is an integral
class, not just a rational cohomology class.

\emph{Multiplication covers and their translations.}
For every positive integer \(n\), the map
\(\pi\circ[n]:A\to\mathbb P^1\), after finite
constants, is again a Galois cover with genus-one
source. To check this without assuming that all
deck transformations fix the origin, write each
deck transformation of \(\pi\) as
\(P\mapsto h(P)+t\), where \(h\) is a group
automorphism. It has exactly \(n^2\) lifts through
\([n]\), namely
\(P\mapsto h(P)+s\) with \(ns=t\). These give
\(n^2\deg(\pi)\) distinct automorphisms of the
composite, equal to its degree. All are defined over
one finite constant extension. In particular the
deck transformations include all translations
\(\tau_t\), \(t\in A[n]\).

Denote this covering object by \(A_n\), retaining
its map to the original rational base, and denote
its induced generic action by \(\varphi_n^*\).
Claim~\ref{clm:FFlatteskernel} realizes \(A_n\),
its deck arrows, and the map
\([n]:A_n\to A_1\) naturally on their branch
complements. Hence on cohomology
\[
 \varphi_n^*[n]^*=[n]^*\varphi_1^*,\qquad
 \varphi_n^*\tau_t^*=\tau_t^*\varphi_n^*.
                                                        \tag{F.49.5}
\]
These equations are not an assumed extension to
unrelated rational subfields: all their arrows are
over the original base. The possible inner
ambiguities in the generic group comparisons are
irrelevant on \(H^1\).
Define \(b_n(h)\) as in (F.49.4). The inertia
argument makes every such class lie in
\(H^1(A,\mathbb Z_\ell(1))\).
Translations act as the identity on this group,
and \([n]^*\) acts as multiplication by \(n\).
One may see both statements from
\(A(\mathbb C)=\mathbb C/\Lambda\): translations
are homotopic to the identity and multiplication
acts by \(n\) on the lattice. The algebraic and
profinite identifications follow from the
characteristic-zero comparison theorem
\citep[Expose~XII, Theoreme~5.1]{FF_SGA1}.
The group is free of rank two over \(\mathbb Z_\ell\).

\emph{An integral averaging identity.}
Suppose first that the principal divisor
\[
                       D=\operatorname{div}(f)
                              =\sum_i m_i[P_i]
\]
has some coefficient equal to \(1\) or \(-1\).
Choose \(Q_i\) with \(nQ_i=P_i\). Since \(D\)
is principal, \(\sum_i m_i=0\) and
\(\sum_i m_iP_i=0\) in the elliptic group.
Consequently \(s=\sum_i m_iQ_i\) belongs to
\(A[n]\). Adjust the chosen \(Q_i\) at a
coefficient \(\pm1\) by the appropriate element
of \(A[n]\) so that this sum becomes zero.
The divisor
\[
                         D_0=\sum_i m_i[Q_i]
\]
now has degree zero and zero sum in the elliptic
group, so it is principal, say
\(D_0=\operatorname{div}(h)\).
All these choices use only finitely many algebraic
constants. Divisors then give
\[
          f\circ[n]=c\prod_{t\in A[n]}\tau_t^*h
                  \qquad(c\in\overline{\mathbb Q}^{\times}).
                                                        \tag{F.49.6}
\]
A constant has zero geometric Kummer character.
Apply the difference operator on \(A_n\) to
(F.49.6), and use (F.49.5). It gives the equality
in integral geometric cohomology
\[
           n\,b_1(f)
              =[n]^*b_1(f)
              =\sum_{t\in A[n]}\tau_t^*b_n(h)
              =n^2b_n(h).
                                                        \tag{F.49.7}
\]
Torsion-freeness permits cancellation of \(n\).
Taking \(n=\ell^r\) for every \(r\ge1\) shows
that \(b_1(f)\) belongs to
\(\ell^r H^1(A,\mathbb Z_\ell(1))\) for every
\(r\), and hence is zero.

This proves the assertion for every rational
function, not just the ones with a coefficient
\(\pm1\). Given \(f\), choose a function \(q\)
with a simple zero outside the support of
\(\operatorname{div}(f)\), and with no cancellation
there. For example, take a degree-two function on
the elliptic curve and subtract a constant avoiding
its finitely many branch values and the images of
that support. Both \(q\) and \(fq\) have such
a simple zero. The proved case and additivity give
\(b_1(f)=b_1(fq)-b_1(q)=0\).
Thus the geometric restriction of every error in
(F.49.4) vanishes, for every \(\ell\).

\emph{Removing the arithmetic error by the norm.}
Return to an arbitrary finite \(L/E\) and
\(f\in L(C)^\times\). Put
\(F=L(T)\), \(K=L(C)\), and
\(d=[K:F]\). The full integral difference
\[
       \delta=\varphi_C^*\kappa_{\ell^\infty}(f)
                         -\kappa_{\ell^\infty}(f)
          \in\operatorname{Hom}_{\mathrm{cts}}
                     (G_K,\mathbb Z_\ell(1))
                                                        \tag{F.49.8}
\]
has zero restriction to the geometric kernel by
the preceding argument. It therefore factors as
an inflated character
\(\gamma:G_L\to\mathbb Z_\ell(1)\).
This factorization is at the full compatible
integral level; no conclusion is drawn merely
from separate finite-level classes.

Corestriction commutes with the induced generic
automorphisms, because the actual slice realization
identifies the covering map with \(\pi\). The
Kummer norm identity and
Claim~\ref{clm:FFallprimekernel} on the rational
base imply
\[
 \operatorname{cor}_{K/F}(\delta)
       =\varphi^*\kappa_{\ell^\infty}(N_{K/F}f)
                   -\kappa_{\ell^\infty}(N_{K/F}f)=0.
                                                        \tag{F.49.9}
\]
Since \(\delta\) is the restriction of the
character inflated from \(\gamma\), the
restriction--corestriction identity identifies the
left side with \(d\gamma\), inflated to \(G_F\).
Inflation is injective, and
\(\operatorname{Hom}_{\mathrm{cts}}
(G_L,\mathbb Z_\ell(1))\) is torsion-free.
Therefore \(\gamma=0\), and \(\delta=0\).
This norm argument needs no rational point on
\(C_L\), so it also handles genus-one torsors
before choosing an origin.

Reducing the equality modulo every \(\ell^r\)
and taking the product over primes proves
(F.49.2). Every finite Kummer character comes
from some \(f\) by Kummer theory, and
\(\mu_m\subset L\), so these equalities are
exactly identity on the full abelian profinite
quotient. This conclusion determines cyclic
covering objects; it does not supply the further
actual point comparisons required to identify
the nonabelian generic action or the whole slice
topos.
\end{proof}


% Root and geometric_routes independently approved the complete proof.
% Integrated as F.50; retained separately as the reviewed source.
\subsection{Integral norms and a higher-genus slice}

\begin{claim}{F.50}{Proved: Kummer rigidity for an explicit genus-two slice}
\label{clm:FFhighergenusnorm}
Retain the pointwise pinned automorphism hypotheses of
Claims~\ref{clm:FFellipticshear} and the preceding
genus-zero and Lattes results, over a finite extension
\(E/B_p\). Let \(C\) be the smooth projective curve
\[
 C:\quad y^2=(u^2-1)(u^2-2)(u^2-3),\qquad T=u^2.
                                                        \tag{F.50.1}
\]
This is a genus-two Galois cover of \(\mathbb P^1_T\)
with group \(V_4\). Its transported actual point
function, including every completed point, is realized
by a unique scheme isomorphism over \(\mathbb P^1_T\).
Using that isomorphism to identify its generic action,
the automorphism fixes
\[
                    H^1(G_{L(C)},\mu_m)
                                                        \tag{F.50.2}
\]
for every finite extension \(L/E\) and every positive
integer \(m\).

More generally, the same Kummer conclusion holds for a
realized smooth proper slice \(C\) whenever it has
coherently realized connected finite etale covers
\(D_s\to C\), with completed-point compatibility,
whose pushforward images in integral geometric first
homology are contained in
\(\ell^sH_1(C,\mathbb Z_\ell)\), for arbitrarily
large \(s\) and each prime \(\ell\).
\end{claim}

The completed-point hypothesis in the general criterion
is substantive. Preserving the isomorphism class of an
Abel cover does not establish it. The specific curve
(F.50.1) meets the hypothesis by its two elliptic
quotients, as proved below. This is a positive-genus
Kummer constraint; it does not identify all further
nonabelian covers or prove reconstruction over \(B_p\).

\begin{proof}
\emph{1. The integral norm criterion.}
Work first over an algebraic closure of the constants.
For every coherently realized curve \(X\), use its
actual-point scheme isomorphism to identify the
transported generic group with \(G_{\overline E(X)}\).
The induced action fixes each geometric place and its
inertia exponent, the latter by restriction of the
base exponent one. Hence for each rational function
\(f\ne0\) its Kummer error is unramified:
\[
 b_X(f):=\alpha_X\kappa_\ell(f)-\kappa_\ell(f)
             \in H^1(X_{\overline E},\mathbb Z_\ell(1)).
                                                        \tag{F.50.3}
\]
The localization sequence proves this integrally:
the residues at all places agree, so their difference
lies in proper first cohomology.

Let \(\rho:D\to C\) be any connected finite etale
cover in the stated coherent family. The norm
\(\overline E(D)^\times\to\overline E(C)^\times\)
is surjective. Here is the dimension proof needed for
that assertion. The function field of a curve over an
algebraically closed field is \(C_1\): for a
homogeneous form of degree \(d\) in \(N>d\) variables,
choose the variables in \(H^0(C,\mathcal O_C(mA))\),
where \(A\) is a fixed positive divisor. After clearing
the fixed coefficient poles, the equation is a system
of at most \(dm\deg A+O(1)\) homogeneous equations in
\(Nm\deg A+O(1)\) scalar variables, by Riemann--Roch.
For sufficiently large \(m\) there are fewer equations
than variables, so the projective dimension theorem
gives a nonzero common zero over the algebraically
closed constant field. Apply this to
\[
               N_{\overline E(D)/\overline E(C)}(v)-f z^d,
               \qquad d=[\overline E(D):\overline E(C)].
                                                        \tag{F.50.4}
\]
It has degree \(d\) in \(d+1\) variables. A nonzero
solution has \(z\ne0\), since the norm of a nonzero
field element is nonzero. Dividing by \(z\) supplies
\(h\in\overline E(D)^\times\) with \(N(h)=f\).
All these data descend to a finite constant extension.
This is the usual Tsen argument, included here in the
precise form used.

Functoriality of corestriction and of the coherent
scheme realizations gives
\[
                   b_C(f)=\operatorname{Tr}_\rho b_D(h).
                                                        \tag{F.50.5}
\]
The transfer direction matters: under integral
Poincare duality for the two compact curves,
\(\operatorname{Tr}_\rho:H^1(D,\mathbb Z_\ell(1))
\to H^1(C,\mathbb Z_\ell(1))\) corresponds to
\(\rho_*:H_1(D,\mathbb Z_\ell)\to
H_1(C,\mathbb Z_\ell)\). Thus the hypothesis on
the image of homology makes (F.50.5) divisible by
\(\ell^s\) in the fixed free \(\mathbb Z_\ell\)-module
\(H^1(C,\mathbb Z_\ell(1))\). Varying \(s\) forces
\(b_C(f)=0\).

For the full Abel mod-\(n\) cover of \(C\), the
homology image is exactly \(nH_1(C,\mathbb Z)\):
its fundamental subgroup is the kernel of
\(\pi_1(C)\to H_1(C,\mathbb Z/n)\).
Every element of its homology image is divisible by
\(n\), and the \(n\)-th power of a loop realizing
any integral homology basis element gives the reverse
inclusion. Therefore these Abel covers would satisfy
the transfer requirement if their completed-point
compatibility were known. No such compatibility for
arbitrary curves is being assumed.

\emph{2. The three quotients of the specified curve.}
The six roots of the polynomial in (F.50.1) are
distinct, so its double cover of \(\mathbb P^1_u\)
has genus two. The independent quadratic equations
\(u^2=T\) and \(y^2=(T-1)(T-2)(T-3)\) give
the \(V_4\)-cover over \(T\). Its three quadratic
quotients are
\[
 \begin{gathered}
 \mathbb P^1_u,\qquad A_1:\ y^2=(T-1)(T-2)(T-3),\\
 A_2:\ z^2=T(T-1)(T-2)(T-3),\qquad z=uy.
 \end{gathered}
                                                        \tag{F.50.6}
\]
The latter two are elliptic curves, each with origin
at its branch point \(T=1\). Their maps to
\(\mathbb P^1_T\) are the quotients by \(\pm1\).
The point \(P_0=(u,y)=(1,0)\) on \(C\) maps to
both origins.

The previous genus-zero and Lattes results realize
all three quotient objects naturally, including
their completed-point maps. The three maps from
\(C\) jointly separate every geometric completed
point. For a group-theoretic verification, a
geometric point has cyclic inertia \(D\leq V_4\),
and equality of its images in a quotient by one
of the three order-two subgroups \(H_i\) means
that the relative translating element lies in
\(H_iD\). One has
\(\bigcap_iH_iD=D\): this is immediate if
\(D=1\), and if \(|D|=2\) one of the \(H_i\)
is exactly \(D\). Thus no distinct completed
points have all three images equal. Naturality
on the branch-complement opens, the unique
extension of a curve isomorphism to its smooth
completion, and this separation show that the
realization of \(C\) agrees with its entire
completed-point bijection. This step does not
assume that arbitrary normalized products
separate boundary branches.

\emph{3. Elliptic division covers provide the required
homology divisibility.}
Let \(\rho_i:C\to A_i\) be the two elliptic quotient
maps, based at \(P_0\). They induce homomorphisms
\[
 \varphi:J_C\longrightarrow A_1\times A_2,
 \qquad
 \psi:A_1\times A_2\longrightarrow J_C,
                                                        \tag{F.50.7}
\]
where \(\varphi\) is the product of the pushforwards
and \(\psi\) the sum of the pullbacks. Their composite
is
\[
                              \psi\varphi=[2]_{J_C}.
                                                        \tag{F.50.8}
\]
Indeed if \(\tau(u,y)=(-u,y)\) and
\(\iota(u,y)=(u,-y)\), the two degree-two norm
identities give the endomorphisms \(1+\tau\) and
\(1+\iota\tau\) on \(J_C\). The hyperelliptic
involution \(\iota\) acts as \(-1\), and their
sum is \(2\). In particular \(\varphi\) has
finite kernel killed by two and is an isogeny,
since both dimensions are two.

For \(n\geq1\), take the connected component
containing \((P_0,0,0)\) of
\[
 D_n=C\mathbin{\times_{A_1\times A_2,[n]}}
                                  (A_1\times A_2).
                                                        \tag{F.50.9}
\]
It is a finite etale cover of \(C\). Its two elliptic
factors, considered over \(\mathbb P^1_T\), are
the Galois genus-one Lattes covers
\(A_i\xrightarrow{[n]}A_i\to\mathbb P^1_T\),
after a finite constant extension making the torsion
rational. They and all their displayed arrows have
the coherent completed-point realizations already
proved. Ordinary fiber products and the selected
component now give that same compatibility for
\(D_n\). This time there is no normalization
ambiguity: (F.50.9) is etale over the smooth curve
\(C\), hence smooth, and its geometric points are
exactly the indicated tuples. The projection to
\(C\) and the two elliptic coordinates therefore
separate all its completed points.

Let \(\Lambda=H_1(C(\mathbb C),\mathbb Z)\) and
\(\Lambda_A=H_1((A_1\times A_2)(\mathbb C),\mathbb Z)\),
using any embedding of the finite coefficient field
into \(\mathbb C\). The covering subgroup in
(F.50.9) consists exactly of the loops whose
\(\varphi_*\)-image lies in \(n\Lambda_A\).
Consequently each element \(v\) in the image of
\(H_1(D_n,\mathbb Z)\to\Lambda\) satisfies
\[
                 \varphi_*v\in n\Lambda_A,
                 \qquad 2v=\psi_*\varphi_*v\in n\Lambda.
                                                        \tag{F.50.10}
\]
For \(n=\ell^s\), this puts the image in
\(\ell^s\Lambda\) if \(\ell\ne2\), and in
\(2^{s-1}\Lambda\) if \(\ell=2\). Integral
comparison gives the same assertion for etale
\(\mathbb Z_\ell\)-homology. The loss of one power
at two is uniform and harmless. Step~1 therefore
proves zero geometric Kummer error on \(C\) at
every prime, including two.

\emph{4. Arithmetic constants and finite coefficients.}
Fix a finite \(L/E\) and \(f\in L(C)^\times\).
Since \(\mu_\infty\subset L\), its compatible
\(\ell\)-adic Kummer class is a genuine continuous
character. The difference under the identified
generic automorphism is zero on the geometric
kernel by the preceding steps, hence is the
inflation of a character
\(\gamma:G_L\to\mathbb Z_\ell(1)\).
Corestriction to \(L(T)\) is zero, since rational
Kummer classes are fixed. On the inflated character
it is multiplication by \([L(C):L(T)]=4\).
The group \(\operatorname{Hom}_{\mathrm{cont}}
(G_L,\mathbb Z_\ell(1))\) is torsion-free, so
\(\gamma=0\), including when \(\ell=2\).
Reducing the integral Kummer character gives every
finite Kummer class, and combining the primes
proves (F.50.2).
For the general criterion the identical argument uses
the positive degree \([L(C):L(T)]\) of its finite
slice in place of four; no divisibility assumption
on that degree is needed.
\end{proof}

\paragraph{Exact limitation and the next obstruction.}
The construction works because the Jacobian of this
curve admits an isogeny to elliptic quotients which
are themselves realized covers of the same rational
base. More generally the argument applies whenever
finitely many such quotient maps induce an injection
on rational first homology and their division
fiber products have completed-point compatibility;
a fixed isogeny denominator only changes the bounded
loss in (F.50.10). For a Jacobian with an absolutely
simple factor of dimension greater than one, these
elliptic maps do not provide that injection. The
obstruction cannot be removed merely by taking a
finite cover \(D\to C\): norm composed with
pullback is multiplication by its degree on
\(J_C\), so \(J_C\) remains an isogeny factor
of \(J_D\). If the latter were built entirely
from elliptic factors, the former would be as well.
The
fixed-prime configuration theorem establishes a
completion topology, not the required actual point
maps on these higher-dimensional Abel covers.
Thus the remaining compatibility problem has not
been solved for an arbitrary positive-genus slice.

\bibliographystyle{plainnat}
\bibliography{ANABELIAN_SUCCESSOR_references}
\end{document}
